\documentclass[11pt]{article}
\usepackage[margin=1.1in]{geometry}
\usepackage[T1]{fontenc}
\usepackage{lmodern}
\usepackage{amsmath,amssymb,amsthm}
\usepackage{microtype}
\usepackage{xcolor}
\usepackage{enumitem}
\usepackage{array}
\usepackage{longtable}
\usepackage{hyperref}
\hypersetup{colorlinks=true,linkcolor=blue!50!black,citecolor=blue!50!black,urlcolor=blue!50!black}

\newtheorem{theorem}{Theorem}[section]
\newtheorem{lemma}[theorem]{Lemma}
\newtheorem{corollary}[theorem]{Corollary}
\newtheorem{proposition}[theorem]{Proposition}
\theoremstyle{definition}
\newtheorem{definition}[theorem]{Definition}
\newtheorem{example}[theorem]{Example}  %
\theoremstyle{remark}
\newtheorem{remark}[theorem]{Remark}

\emergencystretch=2em

\DeclareRobustCommand{\CPTODO}[1]{{\color{red}\textbf{[CP-TODO:} #1\textbf{]}}}

\providecommand{\eps}{\varepsilon}
\providecommand{\Z}{\mathbb{Z}}
\providecommand{\Q}{\mathbb{Q}}
\providecommand{\N}{\mathbb{N}}
\providecommand{\R}{\mathbb{R}}
\newcommand{\abs}[1]{\lvert #1\rvert}
\newcommand{\Ent}{\mathrm{H}}                 %
\newcommand{\Hb}{h}                           %
\DeclareMathOperator{\supp}{supp}
\DeclareMathOperator{\rk}{rank}
\newcommand{\bR}{\underline{R}}               %
\newcommand{\mat}[3]{\langle #1,#2,#3\rangle} %
\newcommand{\cp}{\mathbin{.}}                 %

\newcommand{\Lop}{\mathrm{Lop}\text{-}\allowbreak\mathrm{AE}\text{-}\allowbreak\mathrm{SparseTri}}
\newcommand{\Lopc}{\#\Lop}                    %
\newcommand{\Hyp}{\mathsf{H}}                 %
\newcommand{\Hypk}[1]{\mathsf{H}_{#1}}        %
\newcommand{\es}{\varepsilon^*_5}             %
\newcommand{\bk}{b_\kappa}                    %
\newcommand{\fbat}{f}                         %

\newcommand{\Cov}{\mathrm{Cov}}               %
\newcommand{\cB}{c_B}                         %
\newcommand{\ZB}{Z_B}                         %
\newcommand{\lamB}{\lambda_B}                 %
\newcommand{\delB}{\delta_B}                  %
\newcommand{\alB}{\alpha_B}                   %
\newcommand{\RB}{R_B}                         %
\newcommand{\Lam}{\Lambda}                    %
\newcommand{\Lamp}{\Lambda^{+}}               %
\newcommand{\Lampr}{\Lambda_{\mathrm{pr}}}    %
\newcommand{\Lamcert}{\Lambda_{\mathrm{cert}}}%
\newcommand{\cg}{c_{\mathrm g}}               %
\newcommand{\cstar}{c^{*}}                    %
\newcommand{\cbar}{\bar c}                    %
\newcommand{\co}{c_{\mathrm o}}               %
\newcommand{\ccol}{c_{\mathrm{col}}}          %
\newcommand{\Hcond}[1]{\textup{(H#1)}}        %
\newcommand{\Hscond}[1]{\textup{(H#1$_\sigma$)}} %
\newcommand{\Tcond}[1]{\textup{(T#1)}}        %
\newcommand{\Dcond}{\textup{(D1)}}            %
\newcommand{\TPE}{\mathrm{T}_{\mathrm{PE}}}   %
\newcommand{\Tcol}{\mathrm{T}_{\mathrm{col}}} %

\newcommand{\ThmD}{Theorem~D}                               %
\newcommand{\ThmDk}{Theorem~D\texorpdfstring{\ensuremath{_\kappa}}{\_kappa}}     %
\newcommand{\ThmDprod}{Theorem~D\texorpdfstring{\ensuremath{_\otimes}}{\_x}}     %
\newcommand{\ThmDsig}{Theorem~D\texorpdfstring{\ensuremath{_\sigma}}{\_sigma}}   %
\newcommand{\ThmB}{Theorem~B}                               %
\newcommand{\AlgD}{Algorithm~D}                             %
\newcommand{\Qgain}[2]{Q^{#1}_{#2}}           %
\newcommand{\Erange}[1]{E^{#1}}               %
\newcommand{\gbound}[2]{\gamma^{#1}_{#2}}     %
\newcommand{\xmin}{x_{\min}}                  %

\newcommand{\Qs}{\mathcal{Q}}                 %
\newcommand{\slack}{\mathrm{slack}}           %
\newcommand{\Pzero}{P_0}                      %
\newcommand{\Ptt}{P_{23}}                     %
\newcommand{\Pthree}{P_{32}}                  %
\newcommand{\Pff}{P_{45}}                     %
\newcommand{\Sall}{S_{\mathrm{all}}}          %
\newcommand{\Ukept}{\mathcal{U}}              %
\newcommand{\Fstd}{\mathcal{F}^{\mathrm{std}}_q} %
\newcommand{\Fmrg}{\mathcal{F}^{\mathrm{mrg}}_q} %
\newcommand{\Fstdq}[1]{\mathcal{F}^{\mathrm{std}}_{#1}} %
\newcommand{\Fmrgq}[1]{\mathcal{F}^{\mathrm{mrg}}_{#1}} %
\newcommand{\Bp}{B'}                          %
\newcommand{\Oterm}{O}                        %

\newcommand{\om}{\omega}                      %
\newcommand{\alstar}{\alpha^*}                %

\newcommand{\Esd}{\eta_{\mathrm{SD}}}                 %

\newcommand{\Lleaf}{\Lambda^{\mathrm{leaf}}}          %
\newcommand{\Leg}{\mathcal{L}}                        %
\newcommand{\LamT}[1]{\Lambda^{\mathrm{tile}}_{#1}}   %
\newcommand{\tpad}{t_{\mathrm{pad}}}                  %
\newcommand{\Ebat}{\mathcal{E}}                       %

\newcommand{\KL}{\mathrm{KL}}                        %
\newcommand{\Polq}{\mathcal P}                       %
\newcommand{\Tile}{G}                                %
\newcommand{\UG}{\mathcal U_{\mathrm G}}             %
\newcommand{\RG}{\mathcal R_{\mathrm G}}             %
\newcommand{\LamG}{\Lambda_{\mathrm G}}              %
\newcommand{\AG}{A_{\mathrm G}}                      %
\newcommand{\gG}{g_{\mathrm G}}                      %
\newcommand{\Ic}{\mathcal I}                         %
\newcommand{\Jc}{\mathcal J}                         %
\newcommand{\fCW}{f_{\mathrm{CW}}}                   %
\newcommand{\fG}{f_{\mathrm G}}                      %
\newcommand{\fdir}{f_{\mathrm{dir}}}                 %
\newcommand{\Eenc}[1]{\mathcal E^{#1}}               %
\newcommand{\Pprod}{P_{\otimes}}                     %
\newcommand{\Ppr}{P_{\mathrm{pr}}}                   %
\newcommand{\AlgDprod}{Algorithm~D\texorpdfstring{\ensuremath{_\otimes}}{\_x}}

\newcommand{\Tcert}{T_{\mathrm{cert}}}               %
\newcommand{\rdir}{r_{\mathrm{dir}}}                 %
\newcommand{\rDS}{r_{\mathrm{DS}}}                   %
\newcommand{\fcert}{\fbat_{\mathrm{cert}}}           %
\newcommand{\Bt}{\mathcal B}                         %
\newcommand{\gMM}{g_{\mathrm{MM}}}                   %

\newcommand{\Acs}{\mathcal A}                %
\newcommand{\ktyp}{\kappa_{\mathrm{typ}}}    %
\newcommand{\Dpi}{\mathsf D_\pi}             %
\newcommand{\Sr}{\mathcal S}                 %
\newcommand{\Pl}{\mathrm{Pl}}                %
\newcommand{\Hzl}{\mathcal H_{Z|L}}          %
\newcommand{\Hlz}{\mathcal H_{L|Z}}          %
\newcommand{\unif}{\mathbf u}                %
\newcommand{\wt}{\mathrm{wt}}                %
\newcommand{\wtop}{\mathrm{wt}_{\mathrm{top}}}
\newcommand{\wbar}{\overline{\mathrm{wt}}}
\newcommand{\Pre}{\mathrm{Pre}}              %
\newcommand{\PhiT}{\Phi_{\mathrm T}}         %
\newcommand{\PhiG}{\Phi_{\mathrm G}}         %
\newcommand{\cH}{c_{\mathrm H}}              %
\newcommand{\cmpn}{\mathrm{comp}}            %
\newcommand{\typ}{\mathrm{type}}             %
\newcommand{\hgt}{\mathrm{ht}}               %
\newcommand{\Vhyp}[1]{\textup{(V$_{#1}$)}}   %
\newcommand{\bz}{\beta_0}                    %
\newcommand{\norm}[1]{\lVert#1\rVert}

\makeatletter
\newcommand{\cpvdef}[3]{\expandafter\def\csname cpv@#1\endcsname{\ensuremath{#2}}}
\newcommand{\cpv}[1]{\ifcsname cpv@#1\endcsname\csname cpv@#1\endcsname\else\CPTODO{value #1 undefined}\fi}
\makeatother

\newcommand{\RangeVone}{0.17227}%
\cpvdef{range_v1}{0.17227}{PROVED}%
\newcommand{\EpsGammaVoneOpt}{0.005825}%
\cpvdef{epsgamma_v1_opt}{0.005825}{COMPUTED}%
\newcommand{\EpsGammaBest}{0.00612}%
\cpvdef{epsgamma_v2}{0.00612}{CERTIFIED}%
\newcommand{\HypShrTwentySix}{0.0061945}%
\cpvdef{epsgamma_shr26}{0.0061945}{CLAIMED}%
\newcommand{\HintedOMvVone}{0.17227}%
\cpvdef{hinted_omv_v1}{0.17227}{PROVED}%
\newcommand{\HintedOMvRfourteen}{0.18840}%
\cpvdef{hinted_omv_v2}{0.18840}{PROVED}%
\newcommand{\HintedOMvRtwenty}{0.1907801}%
\cpvdef{hinted_omv_R20}{0.1907801}{CERTIFIED}%
\newcommand{\HintedOMvRthirtytwo}{0.204684}%
\cpvdef{hinted_omv_R32}{0.204684}{CERTIFIED}%
\newcommand{\HypHalfRfourteen}{0.17843}%
\cpvdef{hyp_half_v2}{0.17843}{PROVED}%
\newcommand{\OnlineSetDisj}{0.182773}%
\cpvdef{online_setdisj_R77}{0.182773}{CERTIFIED}%
\newcommand{\SetDisjEta}{0.006973}%
\cpvdef{setdisj_eta_v2}{0.006973}{CERTIFIED}%
\newcommand{\SetDisjEtaPruned}{0.007467}%
\cpvdef{setdisj_eta_pruned_v2}{0.007467}{CERTIFIED}%
\newcommand{\ThmDRangePzero}{0.29459}%
\cpvdef{thmD_range_P0}{0.29459}{CERTIFIED}%
\newcommand{\AlphaVXXZ}{0.321334}%
\cpvdef{alpha_vxxz}{0.321334}{PUBLISHED}%
\newcommand{\AlphaNow}{0.5}%
\cpvdef{alpha_preprints}{1/2}{CLAIMED}%
\newcommand{\HintedOMvLitDet}{0.17227}%
\newcommand{\HintedOMvMcOurs}{0.2731}%
\cpvdef{hinted_omv_mc}{0.2731}{PROVED}%
\cpvdef{omega_vxxz}{2.371552}{PUBLISHED}%

\cpvdef{title_range}{0.2945}{CERTIFIED}%
\cpvdef{title_omv}{0.2046}{CERTIFIED}%
\cpvdef{range_avw26}{0.1204}{PROVED}%
\cpvdef{alpha_lg12}{0.30298}{PUBLISHED}%
\cpvdef{alpha_cop97}{0.29462}{PUBLISHED}%

\cpvdef{P0_lam}{\tfrac{7}{19}}{EXACT}%
\cpvdef{P0_u6}{\tfrac{97}{100}}{EXACT}%
\cpvdef{P0_u7}{\tfrac{979}{1000}}{EXACT}%
\cpvdef{P0_d}{456000}{EXACT}%
\cpvdef{alphaP0_lo}{0.2945976}{CERTIFIED}%
\cpvdef{alphaP0_hi}{0.2945977}{CERTIFIED}%
\cpvdef{slackP0_lo}{3.46\cdot10^{-5}}{CERTIFIED}%
\cpvdef{slackP0_hi}{3.47\cdot10^{-5}}{CERTIFIED}%
\cpvdef{cbarP0_hi}{1.077861}{CERTIFIED}%
\cpvdef{tcapP0_lo}{0.3249}{CERTIFIED}%
\cpvdef{tpadP0_lo}{0.6498}{CERTIFIED}%
\cpvdef{lna1P0_lo}{0.72548}{CERTIFIED}%
\cpvdef{HXsixP0_hi}{0.705479}{CERTIFIED}%
\cpvdef{HZsixP0_lo}{0.735621}{CERTIFIED}%

\cpvdef{P23_lam}{\tfrac{197197}{10^6}}{EXACT}%
\cpvdef{P23_u6}{\tfrac{498851}{500000}}{EXACT}%
\cpvdef{P23_u7}{\tfrac{998717}{10^6}}{EXACT}%
\cpvdef{P32_lam}{\tfrac{187849}{10^6}}{EXACT}%
\cpvdef{P32_u6}{\tfrac{999089}{10^6}}{EXACT}%
\cpvdef{P32_u7}{\tfrac{99953}{10^5}}{EXACT}%
\cpvdef{P45_lam}{\tfrac{182991}{10^6}}{EXACT}%
\cpvdef{P45_u6}{\tfrac{999779}{10^6}}{EXACT}%
\cpvdef{P45_u7}{\tfrac{999899}{10^6}}{EXACT}%
\cpvdef{per_entry_P23}{0.194639}{CERTIFIED}%
\cpvdef{D4_Qa}{0.1}{EXACT}%
\cpvdef{D4_Qb}{0.3}{EXACT}%
\cpvdef{D4_Qc}{0.4}{EXACT}%
\cpvdef{D4_Qd}{0.49}{EXACT}%
\cpvdef{D4_Qe}{\tfrac12}{EXACT}%
\cpvdef{D4_Qf}{0.75}{EXACT}%
\cpvdef{D4_Qg}{1}{EXACT}%
\cpvdef{D4_a}{0.174273}{CERTIFIED}%
\cpvdef{D4_b}{0.178422}{CERTIFIED}%
\cpvdef{D4_c}{0.180571}{CERTIFIED}%
\cpvdef{D4_d}{0.182550}{CERTIFIED}%
\cpvdef{D4_e}{0.182773}{CERTIFIED}%
\cpvdef{D4_f}{0.188519}{CERTIFIED}%
\cpvdef{D4_g}{0.194639}{CERTIFIED}%
\cpvdef{D7_tau2}{0.9509}{PROVED}%
\cpvdef{R11_range}{0.178938}{CERTIFIED}%
\cpvdef{R45_range}{0.187823}{CERTIFIED}%
\cpvdef{D14_HLS}{1.79}{PROVED}%

\cpvdef{excess_crude_P0}{0.2946164}{CERTIFIED}%
\cpvdef{excess_crude_t}{\tfrac{39}{100000}}{EXACT}%
\cpvdef{excess_crude_x}{37600}{EXACT}%
\cpvdef{excess_sharp_P0}{0.2954}{CERTIFIED}%
\cpvdef{excess_sharp_x}{870}{EXACT}%
\cpvdef{excess_sharp_gamma}{8.62\cdot10^{-5}}{CERTIFIED}%
\cpvdef{alphaC_computed}{0.294629}{COMPUTED}%
\cpvdef{R27_ceiling}{0.2299}{COMPUTED}%
\cpvdef{R27_alpha_bound}{0.6226}{COMPUTED}%
\cpvdef{A10_C}{0.0070125}{CERTIFIED}%
\cpvdef{A10_strip}{1.0238}{COMPUTED}%
\cpvdef{A11_vxxz_half}{0.307314}{CERTIFIED}%
\cpvdef{A11_vxxz_half_short}{0.3073}{CERTIFIED}%
\cpvdef{A11_tab1_half}{0.308262}{CERTIFIED}%
\cpvdef{A11_tab1_one}{0.340984}{CERTIFIED}%
\cpvdef{A11_tab1_e_half}{0.33}{EXACT}%
\cpvdef{A11_tab1_e_one}{0.35}{EXACT}%
\cpvdef{A11_tab1_g_one}{0.001363}{PUBLISHED}%
\cpvdef{alpha_lgu}{0.31389}{PUBLISHED}%
\cpvdef{A11_lgu_half}{0.303886}{COMPUTED}%
\cpvdef{A11_pre_half}{0.370242}{COMPUTED}%
\cpvdef{A11_pre_one}{0.5}{COMPUTED}%

\cpvdef{G3_gamma_exact}{0.0039983470}{CERTIFIED}%
\cpvdef{G3_gamma}{0.003998}{CERTIFIED}%
\cpvdef{G3_eps}{0.29425}{CERTIFIED}%
\cpvdef{G3_eps_bind}{0.2942510}{CERTIFIED}%
\cpvdef{G3_t}{\tfrac{49}{2500}}{EXACT}%
\cpvdef{G4_epsgamma}{1.17651\cdot10^{-3}}{CERTIFIED}%
\cpvdef{G6_a_G}{4.7905\cdot10^{-3}}{CERTIFIED}%
\cpvdef{G6_a_E}{0.149907}{CERTIFIED}%
\cpvdef{G6_b_G}{4.6996\cdot10^{-3}}{CERTIFIED}%
\cpvdef{G6_b_E}{0.171888}{CERTIFIED}%
\cpvdef{G6_c_G}{4.5841\cdot10^{-3}}{CERTIFIED}%
\cpvdef{G6_c_E}{0.199816}{CERTIFIED}%
\cpvdef{G6_d_G}{4.3723\cdot10^{-3}}{CERTIFIED}%
\cpvdef{G6_d_E}{0.249726}{CERTIFIED}%
\cpvdef{G6_e_G}{4.1991\cdot10^{-3}}{CERTIFIED}%
\cpvdef{G6_e_E}{0.289745}{CERTIFIED}%
\cpvdef{G6_epsgamma}{1.21667\cdot10^{-3}}{CERTIFIED}%
\cpvdef{G10_a_E}{0.1400}{CERTIFIED}%
\cpvdef{G10_a_G}{2.679\cdot10^{-2}}{CERTIFIED}%
\cpvdef{G10_b_E}{0.1500}{CERTIFIED}%
\cpvdef{G10_b_G}{2.160\cdot10^{-2}}{CERTIFIED}%
\cpvdef{G10_c_E}{0.1600}{CERTIFIED}%
\cpvdef{G10_c_G}{1.770\cdot10^{-2}}{CERTIFIED}%
\cpvdef{G10_d_E}{0.1700}{CERTIFIED}%
\cpvdef{G10_d_G}{1.474\cdot10^{-2}}{CERTIFIED}%
\cpvdef{G10_e_E}{0.1720}{CERTIFIED}%
\cpvdef{G10_e_G}{1.423\cdot10^{-2}}{CERTIFIED}%
\cpvdef{G10_f_E}{0.1800}{CERTIFIED}%
\cpvdef{G10_f_G}{1.246\cdot10^{-2}}{CERTIFIED}%
\cpvdef{G10_g_E}{0.1900}{CERTIFIED}%
\cpvdef{G10_g_G}{1.070\cdot10^{-2}}{CERTIFIED}%
\cpvdef{G10_h_E}{0.2000}{CERTIFIED}%
\cpvdef{G10_h_G}{9.314\cdot10^{-3}}{CERTIFIED}%
\cpvdef{G10_i_E}{0.2200}{CERTIFIED}%
\cpvdef{G10_i_G}{7.334\cdot10^{-3}}{CERTIFIED}%
\cpvdef{G10_j_E}{0.2500}{CERTIFIED}%
\cpvdef{G10_j_G}{5.548\cdot10^{-3}}{CERTIFIED}%
\cpvdef{G10_k_E}{0.2700}{CERTIFIED}%
\cpvdef{G10_k_G}{4.799\cdot10^{-3}}{CERTIFIED}%
\cpvdef{G10_l_E}{0.2900}{CERTIFIED}%
\cpvdef{G10_l_G}{4.254\cdot10^{-3}}{CERTIFIED}%
\cpvdef{G10_m_E}{0.29425}{CERTIFIED}%
\cpvdef{G10_m_G}{3.998\cdot10^{-3}}{CERTIFIED}%
\cpvdef{G10_n_E}{0.2945}{CERTIFIED}%
\cpvdef{G10_n_G}{2.749\cdot10^{-3}}{CERTIFIED}%
\cpvdef{G10_o_E}{0.2950}{CERTIFIED}%
\cpvdef{G10_o_G}{1.161\cdot10^{-3}}{CERTIFIED}%
\cpvdef{G10_p_E}{0.2952}{CERTIFIED}%
\cpvdef{G10_p_G}{4.399\cdot10^{-4}}{CERTIFIED}%
\cpvdef{G10_q_E}{0.2954}{CERTIFIED}%
\cpvdef{G10_q_G}{8.62\cdot10^{-5}}{CERTIFIED}%
\cpvdef{G11_cross}{0.134}{COMPUTED}%
\cpvdef{G11_env_lo}{4.0\cdot10^{-3}}{COMPUTED}%
\cpvdef{G11_env_hi}{4.79\cdot10^{-3}}{COMPUTED}%
\cpvdef{G11_end}{0.295465}{COMPUTED}%
\cpvdef{G12_epsgamma}{3.75\cdot10^{-3}}{CERTIFIED}%
\cpvdef{G12_cw_epsgamma}{1.2\cdot10^{-3}}{COMPUTED}%
\cpvdef{G13_cross}{0.17151}{COMPUTED}%
\cpvdef{G13_crude_a_G}{3.2709\cdot10^{-4}}{CERTIFIED}%
\cpvdef{G13_crude_a_E}{0.1999}{CERTIFIED}%
\cpvdef{G13_crude_b_G}{9.910\cdot10^{-4}}{CERTIFIED}%
\cpvdef{G13_crude_b_E}{0.1800}{CERTIFIED}%

\cpvdef{B2_grid}{97}{COUNT}%
\cpvdef{B3_ka}{\tfrac12}{EXACT}%
\cpvdef{B3_fa}{0.295435}{CERTIFIED}%
\cpvdef{B3_kb}{1}{EXACT}%
\cpvdef{B3_fb}{0.299509}{CERTIFIED}%
\cpvdef{B3_kc}{1.1}{EXACT}%
\cpvdef{B3_fc}{0.212445}{CERTIFIED}%
\cpvdef{B3_kd}{1.15}{EXACT}%
\cpvdef{B3_fd}{0.204707}{CERTIFIED}%
\cpvdef{B3_ke}{1.2}{EXACT}%
\cpvdef{B3_fe}{0.200830}{CERTIFIED}%
\cpvdef{B3_kf}{1.25}{EXACT}%
\cpvdef{B3_ff}{0.198539}{CERTIFIED}%
\cpvdef{B3_kg}{1.3}{EXACT}%
\cpvdef{B3_fg}{0.197066}{CERTIFIED}%
\cpvdef{B3_kh}{1.4}{EXACT}%
\cpvdef{B3_fh}{0.195402}{CERTIFIED}%
\cpvdef{B3_ki}{1.5}{EXACT}%
\cpvdef{B3_fi}{0.194637}{CERTIFIED}%
\cpvdef{B3_kj}{1.6}{EXACT}%
\cpvdef{B3_fj}{0.194351}{CERTIFIED}%
\cpvdef{B3_kk}{1.645}{EXACT}%
\cpvdef{B3_floor}{0.194324}{CERTIFIED}%
\cpvdef{B8_gamma_half}{3.998\cdot10^{-3}}{CERTIFIED}%

\cpvdef{C12_cg_q7}{0.710640}{COMPUTED}%
\cpvdef{C13_cg}{0.118440}{COMPUTED}%
\cpvdef{C13_half_alpha}{\tfrac17}{EXACT}%
\cpvdef{C16_range}{0.296579}{CERTIFIED}%
\cpvdef{C16_theta}{0.074144993}{CERTIFIED}%
\cpvdef{C19_window}{2187}{EXACT}%
\cpvdef{C21_cg}{0.6106}{CERTIFIED}%

\cpvdef{E7_a}{(7,\tfrac34,\tfrac{333}{1000})}{CERTIFIED}%
\cpvdef{E7_b}{(7,\tfrac12,\tfrac{2757}{7000})}{CERTIFIED}%
\cpvdef{E7_c}{(8,\tfrac13,\tfrac{179}{400})}{CERTIFIED}%
\cpvdef{E7_d}{(10,\tfrac14,\tfrac{197}{400})}{CERTIFIED}%
\cpvdef{E7_e}{(11,\tfrac15,\tfrac{2853}{5500})}{CERTIFIED}%
\cpvdef{E7_f}{(13,\tfrac17,\tfrac{887}{1625})}{CERTIFIED}%
\cpvdef{E7_g}{(18,\tfrac1{10},\tfrac{10703}{18000})}{CERTIFIED}%
\cpvdef{E9_half_half}{8.1\cdot10^{-5}}{COMPUTED}%
\cpvdef{E9_half_one}{3.7\cdot10^{-4}}{COMPUTED}%
\cpvdef{E9_tenth_half}{4.5\cdot10^{-5}}{COMPUTED}%
\cpvdef{E9_tenth_one}{2.0\cdot10^{-4}}{COMPUTED}%
\cpvdef{E13_top}{0.594611}{CERTIFIED}%
\cpvdef{E15_a}{0.172744}{CERTIFIED}%
\cpvdef{E15_bc}{0.586372}{CERTIFIED}%
\cpvdef{E16_r_one}{\tfrac15}{EXACT}%
\cpvdef{E16_r_two}{\tfrac3{11}}{EXACT}%
\cpvdef{E16_r_out}{0.2775}{PROVED}%
\cpvdef{E16_mu}{0.5286}{PUBLISHED}%
\cpvdef{E18_g}{0.03414}{CERTIFIED}%

\cpvdef{F2_kappa}{2.91}{COMPUTED}%
\cpvdef{F9_kappa}{6.025}{COMPUTED}%
\cpvdef{F9_entry}{6.789}{COMPUTED}%
\cpvdef{F10_R20}{0.190780}{COMPUTED}%
\cpvdef{F17_CD}{0.5438150}{CERTIFIED}%

\cpvdef{P23_slack_lo}{3\cdot10^{-7}}{CERTIFIED}%
\cpvdef{P32_slack_lo}{2\cdot10^{-7}}{CERTIFIED}%
\cpvdef{alphaP23_lo}{0.293802}{CERTIFIED}%
\cpvdef{alphaP23_hi}{0.293803}{CERTIFIED}%
\cpvdef{coP23_lo}{1.057504}{CERTIFIED}%
\cpvdef{coP23_hi}{1.057505}{CERTIFIED}%
\cpvdef{alphaP32_lo}{0.293664}{CERTIFIED}%
\cpvdef{alphaP32_hi}{0.293665}{CERTIFIED}%
\cpvdef{cprimeP32_lo}{0.766575}{CERTIFIED}%
\cpvdef{cprimeP32_hi}{0.766576}{CERTIFIED}%
\cpvdef{D10_pre}{0.187644}{CERTIFIED}%
\cpvdef{D13_c}{0.154529}{CERTIFIED}%
\cpvdef{D14_HLS_exact}{1.795316}{PROVED}%
\cpvdef{osd_avw26}{0.12}{PUBLISHED}%

\cpvdef{F9_entry_lo}{6.788}{CERTIFIED}%
\cpvdef{F18_err_a}{5\cdot10^{-5}}{COMPUTED}%
\cpvdef{F18_err_b}{9\cdot10^{-4}}{COMPUTED}%
\cpvdef{F18_err_c}{0.9\%}{COMPUTED}%
\cpvdef{F18_factor}{4.6}{COMPUTED}%
\cpvdef{F18_beta_a}{0.01}{EXACT}%
\cpvdef{F18_beta_b}{0.1}{EXACT}%
\cpvdef{F23_G6}{0.0012944331520783}{CERTIFIED}%
\cpvdef{F23_G7}{0.0012082016806783}{CERTIFIED}%

\cpvdef{A11_tab1_om_half}{2.000100}{PUBLISHED}%
\cpvdef{A11_tab1_g_half}{0.0001}{PUBLISHED}%
\cpvdef{A11_tab1_om_one}{2.001363}{PUBLISHED}%
\cpvdef{A11_t_half}{\tfrac{1}{10000}}{EXACT}%
\cpvdef{A11_t_tab_half}{\tfrac{1}{250}}{EXACT}%
\cpvdef{A11_t_one}{\tfrac{3}{50}}{EXACT}%
\cpvdef{A11_grid_n}{97}{COUNT}%
\cpvdef{A11_grid_pos}{16}{COUNT}%
\cpvdef{cert_tol}{10^{-12}}{EXACT}%
\cpvdef{cert_tol_e}{10^{-9}}{EXACT}%
\cpvdef{cert_step}{10^{-6}}{EXACT}%

\cpvdef{u7_lamcert_t0}{1.70916\cdot10^{-3}}{CERTIFIED}%
\cpvdef{u7_G3_margin}{0.00401}{CERTIFIED}%
\cpvdef{u7_ubar}{\tfrac{999999}{10^6}}{EXACT}%
\cpvdef{u7_G6_a_lam}{\tfrac{33}{50}}{EXACT}%
\cpvdef{u7_G6_b_lam}{\tfrac{73}{200}}{EXACT}%
\cpvdef{u7_G6_e_lam}{\tfrac{41}{400}}{EXACT}%
\cpvdef{u7_G6_a_t}{\tfrac{23}{1000}}{EXACT}%
\cpvdef{u7_G6_b_t}{\tfrac{1133}{50000}}{EXACT}%
\cpvdef{u7_G6_c_t}{\tfrac{2221}{100000}}{EXACT}%
\cpvdef{u7_G6_d_t}{\tfrac{1067}{50000}}{EXACT}%
\cpvdef{u7_G6_e_t}{\tfrac{2051}{100000}}{EXACT}%
\cpvdef{u7_Pprod_lam}{\tfrac{181}{1000}}{EXACT}%
\cpvdef{u7_Pprod_slack}{1.84\cdot10^{-5}}{CERTIFIED}%
\cpvdef{u7_Pprod_alpha}{0.2935253}{CERTIFIED}%
\cpvdef{u7_G10_a_c}{\tfrac{1263}{200}}{EXACT}\cpvdef{u7_G10_a_t}{\tfrac{7777}{100000}}{EXACT}\cpvdef{u7_G10_a_f}{\tfrac{2347}{2500}}{EXACT}%
\cpvdef{u7_G10_b_c}{\tfrac{3109}{500}}{EXACT}\cpvdef{u7_G10_b_t}{\tfrac{1717}{25000}}{EXACT}\cpvdef{u7_G10_b_f}{\tfrac{2111}{2500}}{EXACT}%
\cpvdef{u7_G10_c_c}{\tfrac{766}{125}}{EXACT}\cpvdef{u7_G10_c_t}{\tfrac{6089}{100000}}{EXACT}\cpvdef{u7_G10_c_f}{\tfrac{379}{500}}{EXACT}%
\cpvdef{u7_G10_d_c}{\tfrac{1209}{200}}{EXACT}\cpvdef{u7_G10_d_t}{\tfrac{5423}{100000}}{EXACT}\cpvdef{u7_G10_d_f}{\tfrac{3391}{5000}}{EXACT}%
\cpvdef{u7_G10_e_c}{\tfrac{603}{100}}{EXACT}\cpvdef{u7_G10_e_t}{\tfrac{2651}{50000}}{EXACT}\cpvdef{u7_G10_e_f}{\tfrac{6627}{10000}}{EXACT}%
\cpvdef{u7_G10_f_c}{\tfrac{5971}{1000}}{EXACT}\cpvdef{u7_G10_f_t}{\tfrac{2427}{50000}}{EXACT}\cpvdef{u7_G10_f_f}{\tfrac{3017}{5000}}{EXACT}%
\cpvdef{u7_G10_g_c}{\tfrac{5903}{1000}}{EXACT}\cpvdef{u7_G10_g_t}{\tfrac{273}{6250}}{EXACT}\cpvdef{u7_G10_g_f}{\tfrac{1067}{2000}}{EXACT}%
\cpvdef{u7_G10_h_c}{\tfrac{5843}{1000}}{EXACT}\cpvdef{u7_G10_h_t}{\tfrac{3953}{100000}}{EXACT}\cpvdef{u7_G10_h_f}{\tfrac{2337}{5000}}{EXACT}%
\cpvdef{u7_G10_i_c}{\tfrac{5741}{1000}}{EXACT}\cpvdef{u7_G10_i_t}{\tfrac{1649}{50000}}{EXACT}\cpvdef{u7_G10_i_f}{\tfrac{1731}{5000}}{EXACT}%
\cpvdef{u7_G10_j_c}{\tfrac{5629}{1000}}{EXACT}\cpvdef{u7_G10_j_t}{\tfrac{1313}{50000}}{EXACT}\cpvdef{u7_G10_j_f}{\tfrac{937}{5000}}{EXACT}%
\cpvdef{u7_G10_k_c}{\tfrac{5573}{1000}}{EXACT}\cpvdef{u7_G10_k_t}{\tfrac{579}{25000}}{EXACT}\cpvdef{u7_G10_k_f}{\tfrac{951}{10000}}{EXACT}%
\cpvdef{u7_G10_l_c}{\tfrac{5529}{1000}}{EXACT}\cpvdef{u7_G10_l_t}{\tfrac{13}{625}}{EXACT}\cpvdef{u7_G10_l_f}{\tfrac{123}{10000}}{EXACT}%
\cpvdef{u7_G10_n_x}{5000}{EXACT}\cpvdef{u7_G10_o_x}{1400}{EXACT}\cpvdef{u7_G10_p_x}{1100}{EXACT}%
\cpvdef{u7_aG_min}{2.196}{PROVED}%
\cpvdef{u7_cmin}{\tfrac{25}{21}}{EXACT}%
\cpvdef{u7_LG_lo}{1.822}{PROVED}%
\cpvdef{u7_aGL}{0.37}{PROVED}%
\cpvdef{u7_Ppr_lam}{\tfrac{21}{100}}{EXACT}\cpvdef{u7_Ppr_u6}{\tfrac{9957}{10000}}{EXACT}\cpvdef{u7_Ppr_u7}{\tfrac{4987}{5000}}{EXACT}%
\cpvdef{u7_cr_a_c}{\tfrac{2519}{500}}{EXACT}\cpvdef{u7_cr_a_t}{\tfrac{143}{100000}}{EXACT}\cpvdef{u7_cr_a_f}{\tfrac{653}{1000}}{EXACT}%
\cpvdef{u7_cr_b_c}{\tfrac{101}{20}}{EXACT}\cpvdef{u7_cr_b_t}{\tfrac{83}{20000}}{EXACT}\cpvdef{u7_cr_b_f}{\tfrac{1749}{2000}}{EXACT}%

\cpvdef{u8_grid_lo}{22}{COUNT}%
\cpvdef{u8_grid_mid}{74}{COUNT}%
\cpvdef{u8_t_lo}{10^{-4}}{EXACT}%
\cpvdef{u8_t_a}{0.25}{EXACT}\cpvdef{u8_t_b}{0.26}{EXACT}\cpvdef{u8_t_z}{0.99}{EXACT}\cpvdef{u8_t_skip}{0.28}{EXACT}%
\cpvdef{u8_t_cap}{0.3249}{EXACT}%
\cpvdef{u8_t1}{0.06}{EXACT}\cpvdef{u8_M1}{1.10841\cdot10^{-2}}{CERTIFIED}%
\cpvdef{u8_t2}{0.1}{EXACT}\cpvdef{u8_M2}{2.56438\cdot10^{-2}}{CERTIFIED}%
\cpvdef{u8_M3}{0.172905}{CERTIFIED}%
\cpvdef{u8_t4}{0.5}{EXACT}\cpvdef{u8_M4}{0.34703}{CERTIFIED}%
\cpvdef{u8_t5}{0.65}{EXACT}\cpvdef{u8_M5}{0.53169}{CERTIFIED}%
\cpvdef{u8_t6}{0.8}{EXACT}\cpvdef{u8_M6}{0.74681}{CERTIFIED}%
\cpvdef{u8_M7}{1.06018}{CERTIFIED}%
\cpvdef{u8_k_a}{0.25}{EXACT}\cpvdef{u8_f_a}{0.294753}{CERTIFIED}%
\cpvdef{u8_k_b}{0.75}{EXACT}\cpvdef{u8_f_b}{0.296915}{CERTIFIED}%
\cpvdef{u8_k_c}{1.02}{EXACT}\cpvdef{u8_f_c}{0.295371}{CERTIFIED}%
\cpvdef{u8_k_d}{1.025}{EXACT}\cpvdef{u8_f_d}{0.288475}{CERTIFIED}%
\cpvdef{u8_k_e}{1.05}{EXACT}\cpvdef{u8_f_e}{0.235760}{CERTIFIED}%
\cpvdef{u8_plateau_end}{1.023803}{CERTIFIED}%
\cpvdef{u8_gD_one}{2.237\cdot10^{-2}}{CERTIFIED}%
\cpvdef{u8_exc_half}{8.37\cdot10^{-4}}{CERTIFIED}%
\cpvdef{u8_exc_one}{4.91\cdot10^{-3}}{CERTIFIED}%
\cpvdef{u8_ccol_P32}{1.7666}{COMPUTED}%
\cpvdef{u8_alpha_P32}{0.2937}{COMPUTED}%

\cpvdef{C19_copy_lo}{0.0574}{CERTIFIED}%
\cpvdef{C21_half_alpha}{0.160667}{EXACT}%
\cpvdef{C13_alpha}{\tfrac27}{EXACT}%
\cpvdef{C11_q5}{6\ln2/7}{PROVED}%
\cpvdef{C19_window_c}{3^4\cdot3^3}{EXACT}%

\cpvdef{E7pt_a}{(7,\tfrac34,\tfrac{999}{1000},\tfrac1{2000})}{EXACT}%
\cpvdef{E7pt_b}{(7,\tfrac12,\tfrac{919}{1000},\tfrac{237}{5000})}{EXACT}%
\cpvdef{E7pt_c}{(8,\tfrac13,\tfrac{179}{200},\tfrac{337}{5000})}{EXACT}%
\cpvdef{E7pt_d}{(10,\tfrac14,\tfrac{197}{200},\tfrac{47}{5000})}{EXACT}%
\cpvdef{E7pt_e}{(11,\tfrac15,\tfrac{951}{1000},\tfrac{83}{2500})}{EXACT}%
\cpvdef{E7pt_f}{(13,\tfrac17,\tfrac{887}{1000},\tfrac{53}{625})}{EXACT}%
\cpvdef{E7pt_g}{(18,\tfrac1{10},\tfrac{973}{1000},\tfrac{49}{2500})}{EXACT}%
\cpvdef{E7E_a}{\tfrac{333}{1000}}{EXACT}%
\cpvdef{E7E_b}{\tfrac{2757}{7000}}{EXACT}%
\cpvdef{E7E_c}{\tfrac{179}{400}}{EXACT}%
\cpvdef{E7E_d}{\tfrac{197}{400}}{EXACT}%
\cpvdef{E7E_e}{\tfrac{2853}{5500}}{EXACT}%
\cpvdef{E7E_f}{\tfrac{887}{1625}}{EXACT}%
\cpvdef{E7E_g}{\tfrac{10703}{18000}}{EXACT}%
\cpvdef{E7Ed_a}{0.333000}{EXACT}%
\cpvdef{E7Ed_b}{0.393857}{EXACT}%
\cpvdef{E7Ed_c}{0.447500}{EXACT}%
\cpvdef{E7Ed_d}{0.492500}{EXACT}%
\cpvdef{E7Ed_e}{0.518727}{EXACT}%
\cpvdef{E7Ed_f}{0.545846}{EXACT}%
\cpvdef{E7Ed_g}{0.594611}{EXACT}%
\cpvdef{E7b_a}{\tfrac{999}{4000}}{EXACT}%
\cpvdef{E7b_b}{\tfrac{2757}{14000}}{EXACT}%
\cpvdef{E7b_c}{\tfrac{179}{1200}}{EXACT}%
\cpvdef{E7b_d}{\tfrac{197}{1600}}{EXACT}%
\cpvdef{E7b_e}{\tfrac{2853}{27500}}{EXACT}%
\cpvdef{E7b_f}{\tfrac{887}{11375}}{EXACT}%
\cpvdef{E7b_g}{\tfrac{10703}{180000}}{EXACT}%
\cpvdef{E7sl_a}{3.76\cdot10^{-3},\ 3.80\cdot10^{-3}}{CERTIFIED}%
\cpvdef{E7sl_b}{9.77\cdot10^{-4},\ 9.83\cdot10^{-4}}{CERTIFIED}%
\cpvdef{E7sl_c}{7.73\cdot10^{-4},\ 8.01\cdot10^{-4}}{CERTIFIED}%
\cpvdef{E7sl_d}{1.03\cdot10^{-3},\ 1.01\cdot10^{-3}}{CERTIFIED}%
\cpvdef{E7sl_e}{6.61\cdot10^{-4},\ 6.52\cdot10^{-4}}{CERTIFIED}%
\cpvdef{E7sl_f}{6.09\cdot10^{-4},\ 5.91\cdot10^{-4}}{CERTIFIED}%
\cpvdef{E7sl_g}{5.09\cdot10^{-4},\ 5.14\cdot10^{-4}}{CERTIFIED}%
\cpvdef{E15bc_1}{(0.10,\,0.91)}{EXACT}%
\cpvdef{E15bc_2}{(0.05,\,0.96)}{EXACT}%
\cpvdef{E15bc_3}{(0.20,\,0.85)}{EXACT}%
\cpvdef{E15bc_4}{(0.30,\,0.75)}{EXACT}%
\cpvdef{E15bc_5}{(0.40,\,0.65)}{EXACT}%
\cpvdef{E15bc_6}{(0.10,\,0.95)}{EXACT}%
\cpvdef{E15s_1}{\tfrac{10}{91}}{EXACT}%
\cpvdef{E15s_2}{\tfrac{5}{96}}{EXACT}%
\cpvdef{E15s_3}{\tfrac{4}{17}}{EXACT}%
\cpvdef{E15s_4}{\tfrac{2}{5}}{EXACT}%
\cpvdef{E15s_5}{\tfrac{8}{13}}{EXACT}%
\cpvdef{E15s_6}{\tfrac{2}{19}}{EXACT}%
\cpvdef{E15lo_1}{0.01}{EXACT}%
\cpvdef{E15lo_2}{0.01}{EXACT}%
\cpvdef{E15lo_3}{0.05}{EXACT}%
\cpvdef{E15lo_4}{0.05}{EXACT}%
\cpvdef{E15lo_5}{0.05}{EXACT}%
\cpvdef{E15lo_6}{0.05}{EXACT}%
\cpvdef{E15r_1}{0.054584}{CERTIFIED}%
\cpvdef{E15r_2}{0.029730}{CERTIFIED}%
\cpvdef{E15r_3}{0.098500}{CERTIFIED}%
\cpvdef{E15r_4}{0.118157}{CERTIFIED}%
\cpvdef{E15r_5}{0.133200}{CERTIFIED}%
\cpvdef{E15r_6}{0.054584}{CERTIFIED}%
\cpvdef{E15q_1}{0.029459}{CERTIFIED}%
\cpvdef{E15q_2}{0.014729}{CERTIFIED}%
\cpvdef{E15q_3}{0.058919}{CERTIFIED}%
\cpvdef{E15q_4}{0.088379}{CERTIFIED}%
\cpvdef{E15q_5}{0.117839}{CERTIFIED}%
\cpvdef{E15v_1}{0.017227}{CERTIFIED}%
\cpvdef{E15v_4}{0.051681}{CERTIFIED}%
\cpvdef{E15v_5}{0.068908}{CERTIFIED}%
\cpvdef{mu_vxxz}{0.527661}{PUBLISHED}%
\cpvdef{E7_slack_min}{5.09\cdot10^{-4}}{CERTIFIED}%
\cpvdef{E16_s}{0.84}{PROVED}%

\cpvdef{hinted_omv_shr26}{0.145}{CLAIMED}%
\cpvdef{hinted_omv_mc_quarter}{\tfrac14}{PROVED}%
\cpvdef{u8_plateau_f}{0.29459778}{CERTIFIED}%
\cpvdef{u8_k_plateau_fl}{1.025}{EXACT}%
\cpvdef{cert_ln_tol}{2^{-200}}{EXACT}%
\cpvdef{abs_chars}{1920}{COUNT}%
\cpvdef{A11_v_a}{\tfrac{153657}{500000}}{EXACT}%
\cpvdef{A11_v_b}{\tfrac{154131}{500000}}{EXACT}%
\cpvdef{A11_v_c}{\tfrac{42623}{125000}}{EXACT}%
\cpvdef{B2_grid_lo_list}{10^{-4},\allowbreak\,3\cdot10^{-4},\allowbreak\,8\cdot10^{-4},\allowbreak\,0.002,\allowbreak\,0.004,\allowbreak\,0.0074,\allowbreak\,0.012,\allowbreak\,0.0196,\allowbreak\,0.028,\allowbreak\,0.036,\allowbreak\,0.048,\allowbreak\,0.06,\allowbreak\,0.07,\allowbreak\,0.085,\allowbreak\,0.1,\allowbreak\,0.12,\allowbreak\,0.145,\allowbreak\,0.17,\allowbreak\,0.2,\allowbreak\,0.24,\allowbreak\,0.28,\allowbreak\,0.3249}{EXACT}%
\cpvdef{R27_r}{0.5011}{COMPUTED}%

\makeatletter
\newcommand{\Vref}[1]{\ifcsname vtworef@#1\endcsname\csname vtworef@#1\endcsname\else\CPTODO{label #1 of the other paper undefined}\fi}
\@namedef{vtworef@c:hyp}{2.1}
\@namedef{vtworef@c:sec:prelim}{2}
\@namedef{vtworef@m:cor:abfk}{15.2}
\@namedef{vtworef@m:cor:exactboxes}{8.1}
\@namedef{vtworef@m:cor:omv}{12.1}
\@namedef{vtworef@m:cor:prunedpoints}{9.4}
\@namedef{vtworef@m:lem:hashes}{4.1}
\@namedef{vtworef@m:lem:identity}{3.1}
\@namedef{vtworef@m:prop:setdisjone}{14.4}
\@namedef{vtworef@m:rem:converse}{10.2}
\@namedef{vtworef@m:rem:exact}{6.2}
\@namedef{vtworef@m:rem:omv512}{12.3}
\@namedef{vtworef@m:sec:barriers}{16}
\@namedef{vtworef@m:sec:fixedq}{8}
\@namedef{vtworef@m:sec:leaves}{6}
\@namedef{vtworef@m:sec:squared}{5}
\@namedef{vtworef@m:sec:tensor}{3}
\@namedef{vtworef@m:thm:correct}{5.1}
\@namedef{vtworef@m:thm:fixedq}{8.2}
\@namedef{vtworef@m:thm:general}{7.4}
\@namedef{vtworef@m:thm:powers}{16.3}
\@namedef{vtworef@m:thm:setdisj}{14.1}
\@namedef{vtworef@t:cor:omv}{17.14}
\makeatother

\title{Prescribed entries of thin matrix products via the laser method}
\author{Ruhollah Majdoddin\thanks{\url{https://github.com/Majdoddin/formal-math}: the scripts that check the computer-checked statements of this paper, with their outputs (Section~\ref{cp:sec:certificates}).}\\ \normalsize\texttt{r.majdodin@gmail.com}}%
\date{}

\begin{document}
\maketitle

\begin{abstract}
Given an $N\times D$ and a $D\times N$ integer matrix with $D\le N^\eps$ and a set of $N^2/D^\kappa$ positions of their product, can the entries at these positions be computed in time $N^2/D^\gamma$ for some $\gamma>0$? Alman and Vassilevska Williams proved this for $\eps<\cpv{range_avw26}$ at $\kappa=\frac12$, and arXiv:2610.17420 for every $\eps<\cpv{range_v1}$, a limit of Schönhage's partial-product tile.

We put a fast rectangular algorithm from the laser method on top and answer the sub-products its terms ask for with a bottom, paying for requests per entry, per column or shared through the support distribution of its terms. From Coppersmith's level-1 analysis and Le Gall's counting theorem alone, shared requests give every $\kappa>0$ and every $\eps\le\cpv{thmD_range_P0}$, with an explicit saving: $\gamma=\cpv{G3_gamma}$ for every $\eps\le\cpv{G3_eps}$ at $\kappa=\frac12$. With the rectangular product of Vassilevska Williams, Xu, Xu and Zhou as bottom we get every $\eps\le\cpv{A11_vxxz_half}$ at $\kappa=\frac12$, and every $\eps$ below the dual exponent $\alpha$ at $\kappa=1$, whatever $\alpha$ is; no algebraic algorithm that works for every $\kappa$ passes $\alpha$.

A transfer theorem shows that a family with one preimage per output has typical support exponent zero exactly when its base decomposition is triangular, and that a family with dual exponent at most $1/2$ reaches it when its typical support exponent is below half of it. All time bounds are deterministic.

With the data structure of arXiv:2610.17420 as bottom, the hinted OMv Conjectures~5.2 and~5.7 of van den Brand, Nanongkai and Saranurak fail for hints of size $n^\tau$, $\tau<\cpv{hinted_omv_R32}$, Conjecture~5.12 fails on larger regions, and Online Set Disjointness has worst-case queries polynomially below $\sqrt D$ after preprocessing a power of $D$ below $m^2$ when $D\le m^\eps$, $\eps<\cpv{online_setdisj_R77}$. Hypothesis~1.9 of Abboud, Bringmann, Fischer and Künnemann fails on a larger region containing the triples their Theorem~1.10 uses for output exponents $r\le\frac15$.
\end{abstract}

\begin{quote}\small\itshape
\textbf{Use of AI.} This paper was written with Anthropic's Claude models, Claude Opus 5.5 and Claude Fable 5.1, through Claude Code, in October and November 2026. Under the author's direction they searched the literature, searched for and developed the proofs, wrote and ran the ancillary scripts, drafted and revised the text, and checked it against the cited sources; the research questions, the choice among the directions proposed and the decision on every statement are the author's. The computations can be re-run from the ancillary files (Section~\ref{cp:sec:certificates}). The author takes full responsibility for the content.
\end{quote}

\tableofcontents

\section{Introduction}\label{cp:sec:intro}

Let $X$ be an $N\times D$ and $Y$ a $D\times N$ matrix with integer entries, $D\le N^\eps$, and let $W$ be a set of at most $N^2/D^\kappa$ positions of the product. How fast can the entries of $XY$ at the positions of $W$ be computed? Computing them one by one takes $\abs WD=N^2D^{1-\kappa}$ steps, and fast rectangular matrix multiplication computes all of $XY$ in $N^{2+o(1)}$ steps when $\eps$ is below the dual exponent $\alpha$ of matrix multiplication. The statement $\Hypk\kappa(\eps,\gamma)$ (Definition~\ref{cp:P1}) asks for more: a deterministic algorithm that takes $O(N^2/D^\gamma)$ steps for some $\gamma>0$, for all $N$ and all $2\le D\le N^\eps$, a saving of a power of $D$ below $N^2$ even when $D$ is a constant. Alman and Vassilevska Williams~\cite{AV26} proved $\Hyp(\eps,\gamma)=\Hypk{1/2}(\eps,\gamma)$ for $\eps<\cpv{range_avw26}$ and derived from it the first polynomial improvements over the textbook algorithms for 3SUM and APSP. The paper~\cite{v2} extended the range to every $\eps<\es=\ln4/(5\ln5)=\cpv{range_v1}\ldots$ with Schönhage's partial-product tile~\cite{Sch81}.

This paper shows that the range is not limited by the tile. We put a fast rectangular algorithm from the laser method on top and answer the sub-products its terms ask for with a simple bottom. From Coppersmith's level-1 analysis and Le Gall's counting theorem alone this gives $\Hypk\kappa(\eps,\gamma)$ for every $\kappa>0$ and every $\eps\le\cpv{thmD_range_P0}$ (Corollary~\ref{cp:A2}), and with the published fast rectangular product of Vassilevska Williams, Xu, Xu and Zhou~\cite{VXXZ24} as the bottom, $\Hyp(\eps,\gamma)$ for every $\eps\le\cpv{A11_vxxz_half}$ (Corollary~\ref{cp:A11}). The same construction, with the data structure of~\cite{v2} as the bottom, gives deterministic refutations of the hinted OMv conjectures of van den Brand, Nanongkai and Saranurak~\cite{BNS19} for hints of size $n^\tau$ with $\tau<\cpv{hinted_omv_R32}$, and Online Set Disjointness for universes of size $D\le m^{\eps}$ with $\eps<\cpv{online_setdisj_R77}$. Every algorithm of this paper for which a running time is claimed is deterministic.

\subsection{Thin products and the tile}\label{cp:sec:intro:v1}

The method of~\cite{AV26} and~\cite{v2} is a data structure: after a preprocessing in $O(N^2/D^{\gamma})$ time and space, every single entry of $XY$ is computed in $O(D^q)$ time, for every query exponent $q>0$ (Theorem~\ref{cp:P3}). It is built from fixed powers of a small identity, the ten-term identity in~\cite{AV26} and Schönhage's five-term tile in~\cite{v2}, and its range $\es$ is a limit of the fixed powers of this tile when the data structure must serve every query exponent at once; at a fixed query exponent the same data structure reaches beyond $\es$~\cite[Theorem~\Vref{m:thm:fixedq}]{v2}. Here we ask what lies beyond the tile. For the range the answer is: the composites of this paper reach every $\eps$ up to the dual exponent of the family of fast rectangular algorithms they use on top, and a little beyond it at a finite split (Remark~\ref{cp:A3}, Corollary~\ref{cp:B8}).

\subsection{The composite method}\label{cp:sec:intro:method}

A composite algorithm has two parts. Its \emph{top} is an exact bilinear algorithm $B$ for the product of a $k\times b$ by a $b\times k$ matrix, $b=k^{\alB}$, taken from a family of border decompositions that the laser method produces (Section~\ref{cp:sec:tops}). It is raised to a tensor power $B^{\otimes n}$ whose size $K=k^n$ grows with $D$, and $X$ and $Y$ are cut into $K\times b^n$ and $b^n\times K$ grids of blocks with inner dimension $D_A=\lceil D/b^n\rceil$. Each term $l$ of $B^{\otimes n}$ asks for entries of a sub-product $X_lY_l$ of two linear combinations of blocks. Its \emph{bottom} answers these requests; three bottoms occur: inner products of length $D_A$ (Theorem~\ref{cp:T6}), the data structure of~\cite{v2} (Theorem~\ref{cp:P3}), and a fast rectangular product (Lemma~\ref{cp:A10:bottom}).

The cost depends on how the requests are paid for. \emph{Per entry}, every wanted entry pays one request for every term that covers its block, $K^{c}$ requests with $c$ the coverage exponent of the top (Proposition~\ref{cp:T8}). \emph{Per column}, a whole column pays once for the union of the covering sets of its outputs (Proposition~\ref{cp:T9}). \emph{Shared}, a term is paid once at each inner position for all wanted entries at that position (Lemma~\ref{cp:T2}). The shared count is governed by the profile of the top (Definition~\ref{cp:T1}), a log-moment of the distribution of the supports of its terms over the outputs. Its slope at $0$, the \emph{typical support exponent} $\cg$ of the family, decides the range: \ThmDk{} (Theorem~\ref{cp:T7}) gives $\Hypk\kappa(\eps,\gamma)$ for every $\kappa>0$ and every $\eps\le\alpha$ when $\cg=0$, where $\alpha$ is the dual exponent \emph{of the top's family}, the limit of $\ln b/\ln k$, and not the dual exponent of matrix multiplication.

\subsection{The range}\label{cp:sec:intro:range}

\paragraph{From the level-1 analysis.} Coppersmith analysed a mixture of the Coppersmith--Winograd tensors $F_6$ and $F_7$ at level~1 and obtained $\alpha>\cpv{alpha_cop97}$~\cite{Cop97}. We replace the standard border decomposition of $F_q$~\cite{CW90} by a merged one with the same coverage at every output and a different distribution of the supports over its terms (Proposition~\ref{cp:L1}, Lemma~\ref{cp:L2}), whose level-1 family has $\cg=0$ (Lemma~\ref{cp:L8}), and we count the laser restriction with Le Gall's Theorem~4.2~\cite{LG12}, which allows the unequal block sizes of these points (Lemma~\ref{cp:L4}). Theorem~\ref{cp:A1} gives, for every admissible rational level-1 point $P$, every $\kappa>0$ and every $\eps\le\alpha(P)$, the statement $\Hypk\kappa(\eps,\gamma)$ with some $\gamma>0$. At the point $\Pzero=(\cpv{P0_lam},\cpv{P0_u6},\cpv{P0_u7})$ we certify $\alpha(\Pzero)>\cpv{alphaP0_lo}$ (Lemma~\ref{cp:L11}), so $\Hypk\kappa(\eps,\gamma)$ holds for every $\kappa>0$ and every $\eps\le\cpv{thmD_range_P0}$ (Corollary~\ref{cp:A2}). The proof uses no result of~\cite{v2}: its bottom is an inner product, and its only cited input is Le Gall's count. The supremum of $\alpha(P)$ over level-1 points is known only numerically, $\cpv{alphaC_computed}$ (computed, Remark~\ref{cp:A4}). At $\kappa=\frac12$ and a finite split the range exceeds $\alpha(\Pzero)$ slightly: every $\eps<\cpv{excess_crude_P0}$ with the profile bound of Lemma~\ref{cp:L8}, and every $\eps\le\cpv{excess_sharp_P0}$ with the sharper bound of Lemma~\ref{cp:G1} (Remark~\ref{cp:A3}). The per-entry composites have a lower ceiling: with the coverage given by the product bound, their range is below $\cpv{R27_ceiling}$ (computed, not certified) for every analysis of powers of $F_q$ with $q\le12$ (Remark~\ref{cp:A5}); sharing the queries is what passes it.

\paragraph{A converse.} In the multiplication model, where an algorithm is a division-free straight-line program over $\Q$ for each $(N,D,W)$ and its cost is its number of non-scalar multiplications, an algorithm that works for every small $\kappa$ forces $\omega(1,\eps,1)=2$, hence $\eps\le\alpha$, even if it saves nothing (Proposition~\ref{cp:A6}). The algorithm of \ThmD{} is of this algebraic form, so $\alpha(P)\le\eps^{\mathrm{time}}_{\mathrm{alg}}\le\eps^{\mathrm{mult}}_{\mathrm{alg}}\le\alpha$ for every admissible $P$ (Corollary~\ref{cp:A7}). Since $\alpha>\AlphaVXXZ>\alpha(\Pzero)$ is known~\cite{VXXZ24}, this says nothing new about $\alpha$; it says that no algebraic method that works for every $\kappa$ passes the true $\alpha$. A lower bound in the multiplication model limits algebraic algorithms only, and at a single $\kappa$ the converse caps nothing at $\kappa=1$ (Remark~\ref{cp:A8}).

\paragraph{Fast rectangular bottoms.} At a fixed $\kappa$ a fast rectangular product as the bottom uses this room. In Theorem~\ref{cp:A10} every term of the top chooses between inner products and a fast product of its whole sub-instance, and the output is tiled so that the costs become powers of $D$. If $\omega(1,e,1)\le2+g$ with $\alpha(\Pzero)<e<1$, then $\Hypk\kappa(\eps,\gamma)$ holds for every $\kappa\in(0,1]$ and every $\eps$ below an explicit function $F(\kappa;e,g)$, and $F(1;e,0)=e$; so $\Hypk1(\eps,\gamma)$ holds with some $\gamma>0$ for every $\eps<\alpha$, whatever the value of $\alpha$ is. With the published bounds of~\cite{VXXZ24} we certify (Corollary~\ref{cp:A11}): every $\eps\le\cpv{A11_vxxz_half}$ at $\kappa=\frac12$ from $\alpha>\AlphaVXXZ$; every $\eps<\AlphaVXXZ$ at $\kappa=1$; and every $\eps\le\cpv{A11_tab1_one}$ at $\kappa=1$ from $\omega(1,\cpv{A11_tab1_e_one},1)\le\cpv{A11_tab1_om_one}$. The last value lies above the published bound on $\alpha$ but does not improve $\alpha$: it uses a bound on $\omega(1,\cpv{A11_tab1_e_one},1)$ above~$2$.

So the range has two tiers: every $\eps\le\cpv{thmD_range_P0}$ at every $\kappa>0$ from the level-1 analysis alone, and every $\eps\le\cpv{A11_vxxz_half}$ at $\kappa=\frac12$ with the published bound of~\cite{VXXZ24} as the bottom. Neither tier is the true dual exponent: the bound $\alpha\ge\cpv{alpha_preprints}$ is claimed in two preprints of October 2026~\cite{AvdBZ26,Uzn26}, no top of this paper reaches it, and at $\kappa=1$ only the fast bottom carries the range to $\alpha$.

\paragraph{Which decompositions give $\cg=0$.} Section~\ref{cp:sec:transfer} answers this for border decompositions with as many terms as output variables. If the supports of such a decomposition are triangular with respect to an integer weight, then every family of restrictions of its powers by variable maps with few preimages per output, rank exponent~$2$, linear copy size and a window on the weights of the preimages of the kept outputs has $\cg=0$, so \ThmDk{} reaches its family's $\alpha$ (Theorem~\ref{cp:C2}, the transfer theorem). Conversely, if the decomposition is not triangular, every partition of every such family has $\cg>0$ for its structural profile (in which a term covers an output when its support meets a preimage of the output), and $\cg>0$ itself when every kept output has one preimage (Theorem~\ref{cp:C6}, Corollary~\ref{cp:C7}); so for the families with one preimage per kept output $\cg=0$ holds exactly for the triangular decompositions. But $\cg=0$ is more than \ThmD{} needs: $\cg<\alpha/2$ suffices (Corollary~\ref{cp:C2:half}), and Example~\ref{cp:C13} is a non-triangular decomposition with $\cg=\cpv{C13_cg}<\alpha/2=\cpv{C13_half_alpha}$, whose family reaches its $\alpha$ all the same. The section ends with a conditional instance (Theorem~\ref{cp:C20}): if, for every $\beta<\AlphaVXXZ$, some feasible point of the program of~\cite{VXXZ24} satisfies their condition~(27) with $\omega'=2$ at $\beta$, which is the numerical claim behind their published bound and is not verified here, then $\Hyp(\eps,\gamma)$ holds for every $\eps<\AlphaVXXZ$.

\subsection{Savings and batches}\label{cp:sec:intro:gamma}

The constant of $\Hypk\kappa$ is very large, because the top is a large fixed power, but the saving $\gamma$ is explicit (Section~\ref{cp:sec:gamma}). At $\Pzero$, $\Hyp(\eps,\gamma)$ holds for every $\eps<\alpha(\Pzero)$ and every $\gamma<\min\{\cpv{G3_gamma_exact},1/\eps-1/\alpha(\Pzero)\}$, in particular $\Hyp(\eps,\cpv{G3_gamma})$ for every $\eps\le\cpv{G3_eps}$ (Theorem~\ref{cp:G3}). Products of the level-1 top with the tile of~\cite{v2} raise $\gamma$ at a fixed $\eps$ between about $\cpv{G11_cross}$ and $\cpv{u7_Pprod_alpha}$, the dual exponent of the level-1 point at which they are taken (computed); Theorem~\ref{cp:G10} certifies a curve from $\gamma<\cpv{G10_a_G}$ for $\eps<\cpv{G10_a_E}$ to $\gamma<\cpv{G10_q_G}$ for $\eps\le\cpv{G10_q_E}$. With the profile bounds of this paper these products do not raise $\eps\gamma$: every product pair has $\eps\gamma$ at most the better of the values of its two pure factors (Proposition~\ref{cp:G9}), which are numerical optima, and the largest certified value, $\cpv{G12_epsgamma}$, is below $\cpv{epsgamma_v1_opt}$, the numerical optimum of the tile alone (Remark~\ref{cp:G12}). Since~\cite{v2} derives the exponents of its applications (3SUM, APSP and others) through $\eps\gamma$, the composites of this paper are a statement about the range, not about those exponents.

For batches (Section~\ref{cp:sec:batch}) a phase~1 sees only $X$ and $Y$, and a phase~2 answers any batch of at most $N^2/D^\kappa$ positions. We certify a lower bound $\fcert(\kappa)$ on the range of this two-cost statement for every $\kappa>0$ (Theorem~\ref{cp:B3}): its largest certified value is $\cpv{B3_fb}$, at $\kappa=1$; every $\eps<\alpha(\Pzero)$ is reached for every $\kappa\le\cpv{u8_plateau_end}$ (Corollary~\ref{cp:B7}); and from $\kappa=\cpv{B3_kk}$ on it is its value at $t=1$, at least $\cpv{B3_floor}$, the per-entry value at $\Pzero$. The single-entry structure of Theorem~\ref{cp:D1}, at the point $\Ptt$, answers a batch entry by entry, so the range of the batch statement is at least $\cpv{per_entry_P23}$ for every $\kappa\ge1$. These are lower bounds from these constructions; for $\kappa>1$ no upper bound is known. Batches of polynomially many full columns are easy by fast multiplication alone (Proposition~\ref{cp:B9}).

\subsection{Consequences}\label{cp:sec:intro:cons}

\paragraph{Hinted OMv.} Single entries and whole columns need only the coverage of the top and the union of the covering sets of a column, which the level-1 tops have uniformly at every output (Lemma~\ref{cp:L10}). With the data structure of~\cite{v2} as the bottom, the per-entry and column composites give deterministic structures with preprocessing $O(N^2/D^\gamma)$ in time and space: single entries in $O(D^{Q'})$ for every $\eps<\cpv{per_entry_P23}$ (Theorem~\ref{cp:D1}, certified at a point $\Ptt$), and whole columns and rows in $O(ND^{Q'})$ for every $\eps<\cpv{hinted_omv_R32}$ (Theorem~\ref{cp:D5}, certified at a point $\Pthree$), with $Q'<1$. By the reductions of~\cite{v2} (Lemma~\ref{cp:P5}), Conjectures~5.2 and~5.7 of~\cite{BNS19} fail, deterministically, for every $\tau<\cpv{hinted_omv_R32}$, and Conjecture~5.12 fails for every $\tau_1<\cpv{per_entry_P23}$ and every $\tau_2$ (Corollary~\ref{cp:D2}), for every $\tau_2<\cpv{hinted_omv_R32}\tau_1$ and for every $\tau_1<\cpv{hinted_omv_R32}\tau_2$ (Corollary~\ref{cp:D6}). Randomized refutations reach further: Monte Carlo algorithms over fields of characteristic two refute Conjectures~5.2 and~5.7 for every $\tau<\cpv{hinted_omv_mc}$~\cite[Corollary~\Vref{t:cor:omv}]{v2}. The distance oracles of Haeupler, Long and Saranurak~\cite{HLS24} cite Conjecture~5.7 for their bound on the update time; this bound now needs new evidence for update times $n^{2-\delta}$ with $\delta\le\cpv{hinted_omv_R32}$, and the hinted matrix--vector hypothesis of van den Brand, Song and Zhou~\cite{BSZ24} for dynamic attention fails for every $\tau<\cpv{hinted_omv_R32}$ (Section~\ref{cp:D14}, Proposition~\ref{cp:D:BSZ}). We give no faster algorithm for either problem.

\paragraph{Set disjointness.} With a query exponent below $\frac12$, the per-entry composite gives Online Set Disjointness: two families of $m$ sets over a universe of size $D\le m^\eps$, $\eps<\cpv{online_setdisj_R77}$, are preprocessed in time and space $O(m^2/D^\gamma)$, after which $\abs{S_I\cap T_J}$ is computed for any pair in $O(D^Q)$ time with $Q<\frac12$, worst case and online (Theorem~\ref{cp:D3}, at $\Ptt$). The conclusion of the bounded-universe lower bound of Goldstein, Lewenstein and Porat~\cite[Theorem~4]{GLP19} fails for universes of size $u=\lceil N^c\rceil$ with $c<\cpv{D13_c}$ (Corollary~\ref{cp:D13}); its hypothesis, the Strong SetDisjointness Conjecture~\cite{GKLP17,GLP19}, is false already by the data structure of~\cite{v2} (Remark~\ref{cp:D:SDfalse}; this is also in~\cite[Theorem~\Vref{m:thm:setdisj}]{v2}), and no composite raises the exponent of that refutation (Proposition~\ref{cp:D8}).

\paragraph{Sparse triangles.} Hypothesis~1.9 of Abboud, Bringmann, Fischer and Künnemann~\cite{ABFK24} states that PS-AE-Triangle$(m^a,m^b,m^c,m)$, all-edges triangle detection in a tripartite graph with parts of sizes $m^a$, $m^b$, $m^c$ and at most $m$ edges between the last two, cannot be solved in time $O(m^{\min\{1+a,\omega(a,b,c)\}-\epsilon})$. It is one statement quantified over all triples and is already refuted at triples in the range of~\cite[Theorem~\Vref{m:thm:general}]{v2} (\cite[Corollary~\Vref{m:cor:abfk}]{v2}); what matters is where it fails. Put $\kappa=(b+c-1)/a$ and $s=\min(b,c)/\max(b,c)$. Through a reduction to transposed batches (Lemma~\ref{cp:E10}), the hypothesis fails at every triple with $\kappa\le1$ and $a/\min(b,c)<\alpha(\Pzero)$, and at $\kappa=1$ up to $\cpv{B3_fb}$ (Proposition~\ref{cp:E11}); with a fast bottom, at every triple with $\kappa=1$ and $a/\min(b,c)<\alpha$, whatever $\alpha$ is, and with the bound of~\cite{VXXZ24} on $\omega(1,\cpv{A11_tab1_e_one},1)$ up to $\cpv{A11_tab1_one}$ (Proposition~\ref{cp:E12}); and with the unequal tops of Section~\ref{cp:sec:unequal}, at every triple with $\kappa\le1$ and $a/\min(b,c)<E(s)$, where the step function $E$ reaches $\cpv{E7E_g}$ (Corollary~\ref{cp:E13}). This region contains the triples at which the proof of~\cite[Theorem~1.10]{ABFK24} uses the hypothesis, for every output exponent $r\in(0,\cpv{E16_r_one}]$, and for every $r\in(0,\cpv{E16_r_two}]$ with the value of their exponent $\sigma(r)$ that our results force (Proposition~\ref{cp:E16}); so these instances of the hypothesis are false, and Theorem~1.10 loses its support for these $r$, while it remains valid as an implication. For $r>\cpv{E16_r_out}$ the triples lie outside the region.

\subsection{Unequal outer sides}\label{cp:sec:intro:unequal}

For an $n_1\times D$ by $D\times n_2$ product with $n_1\le n_2$, \ThmDsig{} (Theorem~\ref{cp:E3}) uses tops for unequal outer sides, and Theorem~\ref{cp:E8} gives a saving $\gamma$ that is uniform in the shape: with $s=\ln n_1/\ln n_2$, every $D\le n_1^\eps$ with $\eps\le E(s)$ is served, where $E(\frac12)=\cpv{E7Ed_b}\ldots$ and $E(s)=\cpv{E7Ed_g}\ldots$ for $s\le\frac1{10}$. These ranges are measured against the smaller side; the unequal top moves $\eps$ from $\alpha(\Pzero)$ towards the input-size ceiling~$1$, not towards $\alpha$ measured against the larger side. In the multiplication model, an every-$\kappa$ algorithm at a shape forces $\eps\le\eps^*(\sigma)=\alstar(\sigma)/\sigma$ (Proposition~\ref{cp:E1}), and $\max(\alpha,E(\sigma))\le\eps^*(\sigma)\le\min\{1,\alpha(1+\sigma)/(2\sigma)\}$ (Proposition~\ref{cp:E2}).

\subsection{The degree profile}\label{cp:sec:intro:profile}

Section~\ref{cp:sec:profile} explains the numbers above through one function of the top, the distribution of the supports of its terms, measured by a log-moment, the leaf-degree profile. A batch of wanted entries saves a value of its Legendre function (Proposition~\ref{cp:F2}); a data structure with a fixed query exponent follows a frontier built from the same function (Proposition~\ref{cp:F7}); and a range is not a value of the Legendre function but is set by the encoding cost at a zero of it (Remark~\ref{cp:F8}). For the tile of~\cite{v2} and for powers of one fixed decomposition the profile is an exact log-moment and the batch statement is an identity; the frontier is a construction whose optimality is not claimed. For the level-1 family we bound the profile only from above, so every level-1 saving printed in this paper is a lower bound on what the construction saves.

\subsection{Comparison with other work}\label{cp:sec:intro:comparison}

Table~\ref{cp:tab:comparison} compares the results with those of~\cite{AV26}, of the first version of~\cite{v2} (12 October 2026), of the public working note of shreshth-s~\cite{Shr26} (16--24 October 2026), and of the second version of~\cite{v2}.

\begin{table}[!htbp]
\centering\footnotesize
\renewcommand{\arraystretch}{1.15}\setlength{\tabcolsep}{3pt}
\begin{tabular}{@{}>{\raggedright}p{0.18\textwidth}>{\raggedright}p{0.08\textwidth}>{\raggedright}p{0.1\textwidth}>{\raggedright}p{0.1\textwidth}>{\raggedright}p{0.16\textwidth}>{\raggedright\arraybackslash}p{0.29\textwidth}@{}}
statement & \cite{AV26} & \cite{v2}, v1 & \cite{Shr26} & \cite{v2}, v2 & this paper\\\hline
$\Hyp(\eps,\gamma)$ for every $\eps<$ & $\cpv{range_avw26}$ & $\cpv{range_v1}$ & -- & $\cpv{hyp_half_v2}$ & $\cpv{thmD_range_P0}$ (every $\kappa$, Cor.~\ref{cp:A2}); $\cpv{excess_sharp_P0}$ (Rem.~\ref{cp:A3}); $\cpv{A11_vxxz_half}$ with the bound of~\cite{VXXZ24} as bottom (Cor.~\ref{cp:A11})\\
largest $\eps\gamma$ & -- & $\cpv{epsgamma_v1_opt}$ (numerical optimum) & $\cpv{epsgamma_shr26}$ (claimed) & $\cpv{epsgamma_v2}$ & $\cpv{G12_epsgamma}$ (Thm.~\ref{cp:G10}); products never above their best factor (Prop.~\ref{cp:G9})\\
BNS19 Conj.~5.2, 5.7 fail for every $\tau<$ & $\cpv{range_avw26}$ det. & $\cpv{hinted_omv_v1}$ det. & $\tau\le\cpv{hinted_omv_shr26}$ det. (claimed) & $\cpv{hinted_omv_R20}$ det.; Monte Carlo $\cpv{hinted_omv_mc_quarter}$, $\cpv{hinted_omv_mc}$ & $\cpv{hinted_omv_R32}$ det. (Thm.~\ref{cp:D5}, at $\Pthree$); per entry $\cpv{per_entry_P23}$ (Thm.~\ref{cp:D1}, at $\Ptt$)\\
Online Set Disjointness for $D\le m^\eps$, $\eps<$ & $\cpv{osd_avw26}$ & $\cpv{range_v1}$ & -- & $\cpv{hyp_half_v2}$ det. & $\cpv{online_setdisj_R77}$ det., worst case, online (Thm.~\ref{cp:D3}, at $\Ptt$)\\
Strong SetDisjointness Conjecture & -- & -- & -- & refuted & refuted (Rem.~\ref{cp:D:SDfalse}, also in~\cite[Theorem~\Vref{m:thm:setdisj}]{v2}); no composite raises the exponent (Prop.~\ref{cp:D8})\\
ABFK Hypothesis~1.9 fails at $a/\min(b,c)<$ & -- & (implied) & -- & $\cpv{range_v1}$ for $\kappa\le1$ & $\alpha(\Pzero)$ for $\kappa\le1$ (Prop.~\ref{cp:E11}); $\alpha$ at $\kappa=1$ (Prop.~\ref{cp:E12}); $E(s)$ up to $\cpv{E7E_g}$ (Cor.~\ref{cp:E13}); ABFK's Theorem~1.10 triples for $r\le\cpv{E16_r_one}$ (Prop.~\ref{cp:E16})\\\hline
\end{tabular}
\caption{Comparison (det.\ = deterministic). The entries of~\cite{Shr26} are claims of a public note that we have not checked in full; those of~\cite{v2} are the values of its current draft. Every entry of this paper is deterministic and in the time model, except the refutation of the Strong SetDisjointness Conjecture, which concerns space and query time only (Remark~\ref{cp:D:SDfalse}).}\label{cp:tab:comparison}
\end{table}

Since the first version of~\cite{v2}, others have built on its tile. The note of shreshth-s~\cite{Shr26} and Essayag and Zabokritskiy~\cite{EZ26} improve the exponents of 3SUM and APSP that~\cite{v2} derives from $\Hyp$, and Karingula and Lovett~\cite{KL26} give a simplified proof of the polynomial improvements of~\cite{AV26}. These works concern $\eps\gamma$, the quantity the applications use, and not the range. The second version of~\cite{v2} adds exact counts that raise the range at a fixed query exponent (\cite[Section~\Vref{m:sec:fixedq}]{v2}), the statements that make $\es$ a limit of the tile (\cite[Section~\Vref{m:sec:barriers}]{v2}), and refutations of the Strong SetDisjointness Conjecture and of Hypothesis~1.9 of~\cite{ABFK24} in the range of~\cite[Theorem~\Vref{m:thm:general}]{v2} (\cite[Theorem~\Vref{m:thm:setdisj} and Corollary~\Vref{m:cor:abfk}]{v2}); together with thin products over fields of characteristic two~\cite{cop} it also gives the randomized refutations of the hinted OMv conjectures quoted above. The laser-method inputs of this paper are Coppersmith's level-1 analysis~\cite{Cop97} with Le Gall's Theorem~4.2~\cite{LG12}, and, as bottoms or for comparison, the bounds of Le Gall and Urrutia~\cite{LGU18} and of~\cite{VXXZ24}; the bound $\alpha\ge\cpv{alpha_preprints}$ claimed in~\cite{AvdBZ26,Uzn26} enters only conditional remarks, which say so.

\subsection{Models, numbers and organisation}\label{cp:sec:intro:conventions}

Every theorem says in which of two models it holds (Section~\ref{cp:sec:models}): running time on a deterministic word RAM with $O(\log N)$-bit words, or the number of non-scalar multiplications of a division-free straight-line program over $\Q$. A time bound of an algebraic algorithm implies the same multiplication bound, not conversely; the converses of Sections~\ref{cp:sec:range}, \ref{cp:sec:transfer} and~\ref{cp:sec:unequal} are multiplication statements and limit algebraic algorithms only.

A number is \emph{certified} when it is obtained in exact rational or interval arithmetic with outward rounding, and \emph{computed} when it is obtained in floating point; only certified numbers appear in statements of the form ``for every $\eps<$'', and every printed threshold is rounded in the safe direction. Section~\ref{cp:sec:certificates} lists, for every certified number, the script that certifies it and the line of its output, and Appendix~\ref{cp:app:certificates} lists the rational points. Each composite number is certified at a named level-1 point: $\Pzero$ for \ThmD{} and the batch curve, $\Ptt$ for single entries and Online Set Disjointness, $\Pthree$ for columns. Two numbers that look alike belong to different points and do not contradict each other: $\cpv{B3_floor}$ is the per-entry value at $\Pzero$, and $\cpv{per_entry_P23}$ the per-entry threshold at $\Ptt$.

Sections~\ref{cp:sec:prelim}--\ref{cp:sec:level1} are the toolkit: the problem and the results of~\cite{v2} and of the laser method that we use, the top algorithms with \ThmD{} and \ThmDk, and the level-1 family. Section~\ref{cp:sec:composites} draws the per-entry and column consequences, Section~\ref{cp:sec:range} the range with shared queries, the converse and the fast bottoms, Section~\ref{cp:sec:gamma} the explicit savings and the product curve, and Section~\ref{cp:sec:batch} the batches. Section~\ref{cp:sec:transfer} is the transfer theorem and its converse, Section~\ref{cp:sec:unequal} treats unequal outer sides and Hypothesis~1.9 of~\cite{ABFK24}, Section~\ref{cp:sec:profile} the degree profile, and Section~\ref{cp:sec:certificates} the certificates. Appendix~\ref{cp:app:vxxz} gives the structure of the degeneration of~\cite{VXXZ24} that Theorem~\ref{cp:C20} uses.

\section{Preliminaries}\label{cp:sec:prelim}

This section fixes the problem, the two models of computation in which the theorems of this paper are stated, and the results of~\cite{v2} and of the laser method that later sections use. Apart from the one-line Lemma~\ref{cp:P4}, nothing in it is new.

\subsection{The problem and the statement \texorpdfstring{$\Hypk\kappa(\eps,\gamma)$}{H\_kappa(eps,gamma)}}\label{cp:sec:problem}

\begin{definition}[$\Lopc$ and $\Hypk\kappa(\eps,\gamma)$]\label{cp:P1}
For integers $N\ge1$ and $D\ge2$, an instance of $\Lopc(N,D)$ consists of matrices $X\in\{0,1\}^{N\times D}$ and $Y\in\{0,1\}^{D\times N}$ and a set $W\subseteq[N]\times[N]$ of positions; the output is the list of the integers $(XY)[a,b]$ for $(a,b)\in W$. This is the counting version of the lopsided all-edges sparse triangle problem of Alman and Vassilevska Williams~\cite{AV26}, in the matrix form of~\cite[Section~\Vref{c:sec:prelim}]{v2}.
For $\eps>0$, $\kappa>0$ and $0<\gamma\le\kappa$, $\Hypk\kappa(\eps,\gamma)$ is the statement that there are a deterministic algorithm and a constant $C$ such that for all integers $N\ge1$ and $D\ge2$ with $D\le N^\eps$ the algorithm solves every instance of $\Lopc(N,D)$ with $\abs W\le N^2/D^\kappa$ in at most $CN^2/D^\gamma$ steps of a word RAM with $O(\log N)$-bit words. For $0<\eps,\gamma\le\frac12$ we write $\Hyp(\eps,\gamma):=\Hypk{1/2}(\eps,\gamma)$; this is~\cite[Definition~\Vref{c:hyp}]{v2}.
\end{definition}

The set $W$ is part of the input: the algorithm sees $W$ together with $X$ and $Y$. The cap $\gamma\le\kappa$ is forced, since the algorithm must output up to $\lfloor N^2/D^\kappa\rfloor$ numbers. Savings are powers of $D$, not of $N$, because $D$ may be a constant. The definition is stated for $0/1$ matrices, as in~\cite{v2}; the algorithms of Section~\ref{cp:sec:tops} work unchanged, with the same bounds, for integer matrices whose entries have absolute value $N^{O(1)}$, since every number they compute is $N^{O(1)}$ (the word-size step in the proof of Theorem~\ref{cp:T6}).

\paragraph{Two models.}\label{cp:sec:models}
Every theorem of this paper says in which of two models it holds. In the \emph{time model} the cost is the number of steps of a deterministic word RAM with $O(\log N)$-bit words, as in Definition~\ref{cp:P1}. In the \emph{multiplication model} an algorithm is, for each $N$, $D$ and $W$, a division-free straight-line program over $\Q$ in the entries of $X$ and $Y$ whose control flow depends on $N$, $D$ and $W$ only and which computes the polynomials $(XY)[a,b]$, $(a,b)\in W$, as polynomials (equivalently, correctly for all integer matrices; correctness on $0/1$ matrices alone is not enough, since two polynomials can agree on $\{0,1\}^n$ without being equal); its cost is its number of non-scalar multiplications. An algorithm of this algebraic form whose running time is $T$ uses at most $T$ non-scalar multiplications, because every arithmetic operation on numbers of $O(1)$ words is at least one step. The converse fails, since a word-RAM algorithm need not be algebraic. So a lower bound in the multiplication model limits every algebraic algorithm also in the time model, and says nothing about word-RAM algorithms that are not of this form.

\begin{definition}[the $\kappa$-batch statement]\label{cp:P2}
Let $\eps>0$ and $\kappa>0$, and put $\bk:=\max\bigl(1,\lfloor N^2/D^\kappa\rfloor\bigr)$. The input consists of $X\in\Z^{N\times D}$ and $Y\in\Z^{D\times N}$ with $2\le D\le N^\eps$ and entries of absolute value $N^{O(1)}$. Phase~1 receives $X$ and $Y$. Phase~2 receives a batch $\mathcal B\subseteq[N]\times[N]$ of positions and outputs $(XY)[i,j]$ for $(i,j)\in\mathcal B$. The \emph{$\kappa$-batch statement at $\eps$} says: there are a deterministic two-phase algorithm and constants $\gamma_{\mathrm{pre}}>0$, $g>0$ and $C$ such that for all $N$ and $D$ as above, phase~1 takes at most $CN^2D^{-\gamma_{\mathrm{pre}}}$ steps of the word RAM, and then phase~2, started from the memory left by phase~1, answers every batch with $\abs{\mathcal B}\le\bk$ in at most
\[
CN^2D^{-g}\quad\text{steps if }\kappa\le1,\qquad C\,\bk\,D^{1-g}\quad\text{steps if }\kappa\ge1 .
\]
We write $\fbat(\kappa)$ for the supremum of the $\eps$ for which some algorithm satisfies the $\kappa$-batch statement at $\eps$.
\end{definition}

The two targets agree at $\kappa=1$ up to rounding ($b_1D\le N^2$). The maximum with $1$ matters only for $\kappa>2/\eps$, where it makes a single entry the smallest batch. Phase~2 may write into memory; how several batches are answered from one phase-1 state is the subject of Lemma~\ref{cp:B4}.

\subsection{The data structure and the tile of~\texorpdfstring{\cite{v2}}{[Maj26a]}}\label{cp:sec:v1}

\begin{theorem}[{\cite[Theorem~\Vref{m:thm:general}]{v2}}]\label{cp:P3}
For every $\eps_A<\es=\ln4/(5\ln5)=\cpv{range_v1}\ldots$ and every $q_A>0$ there is a $\gamma_A>0$ such that, for $X\in\Z^{N\times D}$ and $Y\in\Z^{D\times N}$ with $2\le D\le N^{\eps_A}$ and entries of absolute value $N^{O(1)}$, a deterministic preprocessing in $O(N^2/D^{\gamma_A})$ time and space allows any single entry of $XY$ to be computed deterministically in $O(D^{q_A})$ time. For every $\eps_A<\es$ and $\kappa>0$ there is a $\gamma>0$ such that the entries of $XY$ at any set $W$ of at most $N^2/D^\kappa$ positions can be computed deterministically in $O(N^2/D^\gamma)$ time.
\end{theorem}

The model is time on the word RAM of Definition~\ref{cp:P1}. Sections~\ref{cp:sec:tops}, \ref{cp:sec:composites} and~\ref{cp:sec:batch} use only the first sentence, the data structure, as a \emph{bottom}: its sub-instances are the products of the top algorithm's encodings. The range $\es$ of Theorem~\ref{cp:P3} is that of the tile used in~\cite{v2}, Schönhage's partial-product tensor~\cite{Sch81}; \cite[Theorem~\Vref{m:thm:powers}]{v2} shows that $\es$ is a limit of the fixed powers of this tile when the data structure must serve every query exponent at once, and not a limit of the data structure or of the method: at a fixed query exponent the same data structure reaches beyond $\es$.

Section~\ref{cp:sec:gamma} combines the level-1 top of Section~\ref{cp:sec:level1} with the tile of~\cite{v2}. It uses two results of~\cite{v2}: re-indexing maps with at least $\mu_\infty^2\binom Lm$ good sets exist and are found deterministically~\cite[Lemma~\Vref{m:lem:hashes}]{v2}, and the $2m+1$ squared runs together compute, for every pair $(Q_k,Q_{k'})$ of good sets, every entry of the block product $X_{Q_k}Y_{Q_{k'}}$ as a single output~\cite[Theorem~\Vref{m:thm:correct}]{v2}. The notation of the tile, $\Tile(c)$ and its parameters, is fixed in Section~\ref{cp:sec:prod}. Section~\ref{cp:sec:gamma} needs that identity for block matrices:

\begin{lemma}[block entries]\label{cp:P4}
The identity of~\cite[Theorem~\Vref{m:thm:correct}]{v2} holds when the matrices $X_Q$ and $Y_Q$ are block matrices over any ring, with every product taken left factor first.
\end{lemma}

\begin{proof}
Both sides of the identity are bilinear forms in the entries of the matrices $X_Q$ and $Y_Q$ with integer coefficients that do not depend on these entries: the input arrays $a_s$ and $b_s$ of~\cite[Section~\Vref{m:sec:squared}]{v2} are selections of entries, and the encodings are linear. The identity holds for all integer matrices, so the two coefficient tensors are equal. An equality of coefficient tensors gives the identity in every ring, with each product taken left factor first.
\end{proof}

\subsection{The hinted OMv conjectures}\label{cp:sec:omv}

Van den Brand, Nanongkai and Saranurak~\cite{BNS19} state three hinted variants of Online Matrix--Vector multiplication; their input arrives in phases, all computations are over the Boolean semiring, and $t=n^\tau$, $t_k=n^{\tau_k}$ with $0<\tau,\tau_k<1$. In the \emph{v-hinted Mv} problem~\cite[Definition~5.1]{BNS19}, Phase~1 receives $M\in\{0,1\}^{n\times t}$, Phase~2 receives $V\in\{0,1\}^{t\times n}$, and Phase~3 receives an index $i$ and outputs $MV_{[n],i}$. In the \emph{Mv-hinted Mv} problem, Phase~1 receives $N\in\{0,1\}^{n\times n}$ and $V\in\{0,1\}^{t\times n}$, Phase~2 a list $I\in[n]^t$ of column indices, and Phase~3 an index $j$, after which $N_{[n],I}V_{[t],j}$ is output. In the \emph{uMv-hinted uMv} problem, Phase~1 receives $U\in\{0,1\}^{n\times t_1}$, $N\in\{0,1\}^{n\times n}$ and $V\in\{0,1\}^{t_2\times n}$, Phase~2 a list $I\in[n]^{t_1}$, Phase~3 a list $J\in[n]^{t_2}$, and Phase~4 a pair $(i,j)$, after which the entry $(UN_{I,J}V)_{i,j}$ is output; rows and columns of $N_{I,J}$ may repeat.
\begin{itemize}
\item \emph{Conjecture 5.2 of~\cite{BNS19}}: every algorithm for v-hinted Mv with polynomial Phase~1 needs $\Omega(n^{\omega(1,1,\tau)-\eps})$ operations in Phase~2 or $\Omega(n^{1+\tau-\eps})$ in Phase~3, for every constant $\eps>0$.
\item \emph{Conjecture 5.7}: every algorithm for Mv-hinted Mv with polynomial Phase~1 needs $\Omega(n^{\omega(1,\tau,1)-\eps})$ in Phase~2 or $\Omega(n^{1+\tau-\eps})$ in Phase~3, for every $\eps>0$.
\item \emph{Conjecture 5.12}: every algorithm for uMv-hinted uMv with polynomial Phase~1 needs $\Omega(n^{\omega(1,\tau_1,1)-\eps})$ in Phase~2, or $\Omega(n^{\omega(\tau_2,\tau_1,1)-\eps})$ in Phase~3, or $\Omega(n^{\tau_1+\tau_2-\eps})$ in Phase~4, for every $\eps>0$.
\end{itemize}
Here $\omega(\cdot,\cdot,\cdot)$ is the rectangular exponent of Section~\ref{cp:sec:laser}. By ``Conjecture 5.12'' we always mean the conjecture on uMv-hinted uMv just stated; \cite[Corollary~5.10]{BNS19} cites the Mv-hinted Mv conjecture under the label 5.12. A conjecture of this form \emph{fails} at given exponents if some algorithm violates all of its alternatives for some $\eps>0$.

The refutations of~\cite{v2} use the data structure only through two interfaces, and we state them in that form. Both compute over $\Z$ with $0/1$ inputs and test whether an entry is positive.

\begin{lemma}[the reductions of~{\cite[Corollary~\Vref{m:cor:omv} and Proposition~\Vref{m:rem:omv512}]{v2}}]\label{cp:P5}
Let $\eps\in(0,1)$, and suppose there are $\gamma>0$, $Q'<1$ and a deterministic algorithm that, for all $N\ge1$ and $2\le D\le N^\eps$, preprocesses $X\in\{0,1\}^{N\times D}$ and $Y\in\{0,1\}^{D\times N}$ in $O(N^2/D^\gamma)$ time.
\begin{enumerate}[label=(\alph*)]
\item If afterwards any whole column of $XY$ can be computed in $O(ND^{Q'})$ time, then Conjectures 5.2 and 5.7 of~\cite{BNS19} fail for every $\tau\le\eps$, with Phase~2 in $O(n^{2-\gamma\tau})$ and Phase~3 in $O(n^{1+Q'\tau})$ time, and Conjecture 5.12 fails for every $\tau_1\in(0,1)$ and every $\tau_2\le\eps\tau_1$, with a Phase~2 that only stores $I$, Phase~3 in $O(n^{1+\tau_1-\gamma\tau_2})$ and Phase~4 in $O(n^{\tau_1+Q'\tau_2})$ time.
\item If afterwards any single entry of $XY$ can be computed in $O(D^{Q'})$ time, then the conclusions of (a) hold, and Conjecture 5.12 fails also for every $\tau_1\le\eps$ and every $\tau_2\in(0,1)$, with Phase~2 in $O(n^{2-\gamma\tau_1})$, a Phase~3 that only stores $J$, and Phase~4 in $O(n^{\tau_2+Q'\tau_1})$ time.
\end{enumerate}
All these refutations are deterministic.
\end{lemma}

\begin{proof}
These are the proofs of~\cite[Corollary~\Vref{m:cor:omv} and Proposition~\Vref{m:rem:omv512}]{v2}, which use the data structure of~\cite{v2} only through the two interfaces above. The region $\tau_1<\eps\tau_2$ of~\cite[Corollary~\Vref{m:cor:omv}]{v2} for Conjecture 5.12 lies inside the region of (b)~\cite[Proposition~\Vref{m:rem:omv512}]{v2}; with whole rows it is Corollary~\ref{cp:D7}. (a) For Conjecture 5.2 take $N=n$, $D=t$, $X=M$ and preprocess $X$ and $Y=V$ in Phase~2; for Conjecture 5.7 take $X=N_{[n],I}$, the $n\times t$ matrix formed in Phase~2, and $Y=V$. Then $D=n^\tau\le N^\eps$, Phase~2 costs $O(n^{2-\gamma\tau})$, which is at most $n^{\omega(1,1,\tau)-\gamma\tau}$ and $n^{\omega(1,\tau,1)-\gamma\tau}$ since both exponents are at least $2$, and Phase~3 is one column of $XY$, in $O(n^{1+Q'\tau})$. For Conjecture 5.12 with $\tau_2\le\eps\tau_1$, in Phase~3 the matrices $X:=N_{I,J}$ ($t_1\times t_2$) and $Y:=V$ are known; cut the $n$ columns of $V$ into $n/t_1$ blocks of $t_1$ columns and preprocess the product of $X$ with each block, with $t_1$ in the role of $N$ and $D=t_2\le t_1^\eps$, in total $(n/t_1)\,O(t_1^2/t_2^\gamma)=O(n^{1+\tau_1-\gamma\tau_2})$, below $n^{\omega(\tau_2,\tau_1,1)}\ge n^{1+\tau_1}$. In Phase~4, $(UN_{I,J}V)_{i,j}=1$ if and only if $U_{i,a}=1$ and $(XY)[a,j]>0$ for some $a\le t_1$; these $t_1$ entries form one column of one block product, computed in $O(t_1t_2^{Q'})=O(n^{\tau_1+Q'\tau_2})$, below $n^{\tau_1+\tau_2}$.
(b) A column consists of $N$ entries, so (a) applies. For Conjecture 5.12 with $\tau_1\le\eps$, in Phase~2 the matrices $X:=U$ ($n\times t_1$) and $Y:=N_{I,[n]}$ ($t_1\times n$) are known; preprocess them with $N=n$ and $D=t_1$, in $O(n^{2-\gamma\tau_1})$, below $n^{\omega(1,\tau_1,1)}$. Column $b$ of $UN_{I,J}$ is column $J_b$ of $XY$, so $(UN_{I,J}V)_{i,j}=1$ if and only if $(XY)[i,J_b]>0$ and $V_{b,j}=1$ for some $b\le t_2$: these are $t_2$ single entries, in $O(n^{\tau_2+Q'\tau_1})$, below $n^{\tau_1+\tau_2}$. In each case the polynomial Phase~1 only stores its input, and every phase is below the corresponding alternative of the conjecture by a polynomial factor.
\end{proof}

\subsection{Facts from the laser method}\label{cp:sec:laser}\phantomsection\label{cp:P6}

\paragraph{Tensors and decompositions.} All tensors are over $\Q$, in three sets of variables $x$, $y$, $z$; the $z$-variables are the \emph{outputs}. The matrix product tensor is $\mat abc=\sum_{i\in[a],j\in[b],l\in[c]}x_{ij}y_{jl}z_{li}$, and $s\cp T$ is the direct sum of $s$ copies of $T$ on disjoint variables. A \emph{border decomposition} of $T$ with $r$ terms is an identity $\sum_{l=1}^ru_l(\lambda)\otimes v_l(\lambda)\otimes w_l(\lambda)=\lambda^hT+\lambda^{h+1}T'(\lambda)$ over $\Q[\lambda]$, where $u_l$, $v_l$, $w_l$ are linear forms in the $x$-, $y$- and $z$-variables with coefficients in $\Q[\lambda]$; $h$ is its approximation degree, and its \emph{product degree} is the largest $\lambda$-degree of a coefficient of a product $u_lv_lw_l$, the scalar coefficient of the term included. With $h=0$ and constant coefficients it is an exact decomposition (a bilinear algorithm). The $z$-form of term $l$ is $w_l=(w_l[o])_o$; its \emph{support} is $\supp w_l=\{o:w_l[o]\ne0\}$, for a border decomposition as a polynomial in $\lambda$. The \emph{rank} $R(T)$ and the \emph{border rank} $\bR(T)$ are the least numbers of terms of an exact and of a border decomposition. A \emph{restriction} replaces the variables by linear forms in other variables; a \emph{zero-out} sets some variables to zero, and a \emph{variable map} sends every variable to $0$ or to one variable with coefficient~$1$. A restriction of a border decomposition of $T$ is, term by term, a border decomposition of the restricted tensor, with the same approximation degree and no larger product degree, since restriction is linear. The idea of passing from a border decomposition to an exact one by interpolation in $\lambda$ is Bini's~\cite{Bin80}; Lemma~\ref{cp:T5:lemma4} uses it in the form ``a border decomposition with product degree $e$ and $r$ terms gives an exact one with at most $(e+1)r$ terms''.

\paragraph{Coppersmith and Winograd's tensor.} For an integer $q\ge1$ let
\[
F_q=\sum_{i=1}^q\bigl(x_iy_iz_0+x_iy_0z_i+x_0y_iz_i\bigr)+x_{q+1}y_0z_0+x_0y_{q+1}z_0+x_0y_0z_{q+1}
\]
in the variables $x_0,\dots,x_{q+1}$ and likewise for $y$ and $z$ (this is $F'_q$ in the notation of~\cite[Section~3]{LG12}, and the tensor of~\cite[Section~7]{CW90}). Its variables are grouped into \emph{blocks}: $x_0$ is block~$0$, $x_1,\dots,x_q$ form block~$1$, $x_{q+1}$ is block~$2$, and likewise for $y$ and $z$. Every monomial of $F_q$ has block sum~$2$, so $F_q=\sum T_{IJK}$ over the six block triples $(I,J,K)\in\{0,1,2\}^3$ with $I+J+K=2$, with constituents $T_{011}=\mat11q$, $T_{101}=\mat q11$, $T_{110}=\mat1q1$ and $T_{002}=T_{020}=T_{200}=\mat111$. The \emph{standard decomposition} $\Fstd$ of~\cite[Section~7]{CW90} (in the form of~\cite[Section~3]{LG12}, multiplied by $\lambda^3$) is
\[
\lambda\sum_{i=1}^qA_i-B+(1-q\lambda)C=\lambda^3F_q+O(\lambda^4),
\]
\[
A_i=(x_0+\lambda x_i)(y_0+\lambda y_i)(z_0+\lambda z_i),\quad B=\Bigl(x_0+\lambda^2\sum_{i=1}^qx_i\Bigr)\Bigl(y_0+\lambda^2\sum_{i=1}^qy_i\Bigr)\Bigl(z_0+\lambda^2\sum_{i=1}^qz_i\Bigr),
\]
\[
C=(x_0+\lambda^3x_{q+1})(y_0+\lambda^3y_{q+1})(z_0+\lambda^3z_{q+1}),
\]
a border decomposition of $F_q$ with $q+2$ terms and approximation degree $3$; per coordinate its product degree is $4$ for $A_i$, $6$ for $B$ and $10$ for $C$. On tensor powers $F_q^{\otimes M}$ the blocks of a variable string are the strings of the blocks of its letters, and the constituent of a block-triple string is the tensor product of the constituents of its letters.

\paragraph{Rectangular exponents.} %
$\omega(u,v,w)$ is the infimum of the $\theta$ with $R(\mat{\lceil n^u\rceil}{\lceil n^v\rceil}{\lceil n^w\rceil})=O(n^{\theta})$; it is symmetric under all permutations of its arguments and homogeneous. The dual exponent is $\alpha:=\sup\{a\ge0:\omega(1,a,1)=2\}$, and for unequal outer sides $\alstar(\sigma):=\sup\{\beta\ge0:\omega(\sigma,\beta,1)=1+\sigma\}$. Rank and border rank give the same exponents~\cite{Bin80}. Coppersmith showed $\alpha>\cpv{alpha_cop97}$~\cite{Cop97}, Le Gall $\alpha>\cpv{alpha_lg12}$~\cite{LG12}, and the best published bound is $\alpha\ge\AlphaVXXZ$~\cite{VXXZ24}; the bound $\alpha\ge\cpv{alpha_preprints}$ is claimed in two preprints of October 2026~\cite{AvdBZ26,Uzn26}. No theorem of this paper depends on these bounds except where it says so.

\paragraph{Strassen's bound.} %
Over $\Q$, if a straight-line program computes the bilinear forms of a tensor $T$ with $L$ non-scalar multiplications, then $R(T)\le2L$~\cite{Str73}. This is the only link between the multiplication model and rank that the converses of Sections~\ref{cp:sec:range} and~\ref{cp:sec:unequal} and Proposition~\ref{cp:C17} use.

\paragraph{Le Gall's graph theorem.} %
The count of the laser method that Section~\ref{cp:sec:level1} uses is~\cite[Theorem~4.2]{LG12} (arXiv:1204.1111v2, p.~16), in the setting of~\cite[Section~4.1]{LG12}. Let $\tau$ be a positive integer, $N$ a large integer, $\Lambda=\{(I,J,K)\in[\tau]^N\times[\tau]^N\times[\tau]^N: I_\ell+J_\ell+K_\ell=\tau\text{ for all }\ell\}$, where $[\tau]=\{0,\dots,\tau\}$, and let $f_1,f_2,f_3$ be the three coordinate maps. Let $U\subseteq\Lambda$ be such that there are integers $N_1,N_2,N_3$ with $\abs{\{u\in U:f_i(u)=I\}}\in\{0,N_i\}$ for every $I$ and every $i$, and put $T_i=\abs{f_i(U)}$, so that $\abs U=T_1N_1=T_2N_2=T_3N_3$; as in~\cite[Section~4.1]{LG12}, the $N_i$ and the $T_i$ are strictly increasing functions of $N$. Let $U^*\subseteq U$ satisfy $\abs{f_i(U^*)}=T_i$, with integers $N_i^*$ such that $\abs{\{u\in U:f_i(u)=I\}}=N_i$ if and only if $\abs{\{u\in U^*:f_i(u)=I\}}=N_i^*$. The graph $G$ on $U$ joins two triples that agree in some coordinate map, and a \emph{removal operation} deletes all vertices with a given sequence at a given position; on tensors it is the zero-out of the variables of one block.

\begin{theorem}[{\cite[Theorem~4.2]{LG12}}]\label{cp:P6:LG}
For any constant $\eps>0$, the graph $G$ can be converted, with only removal operations, into an edgeless graph with $\Omega\bigl(T_1N_1^*/(N_1+N_2+N_3)^{1+\eps}\bigr)$ vertices, all of them being in $U^*$.
\end{theorem}

The theorem allows arbitrary $N_1$, $N_2$, $N_3$; its special case $N_1=N_2=N_3$ with $U^*=U$ is the theorem of Coppersmith and Winograd~\cite[Theorem~4.1]{LG12}, which does not cover the points used in Section~\ref{cp:sec:level1}, since there $N_1\ne N_3$ (Remark~\ref{cp:rem:LGcite}). Strings of length $N$ above are the strings of blocks of the tensor power; the parameter $N$ is a large constant, independent of the matrix sizes of the algorithms.

\section{Top algorithms and the composite}\label{cp:sec:tops}

A composite algorithm has a fast rectangular algorithm on top, raised to a tensor power whose length grows with $D$, and a bottom that answers the sub-products the top's terms ask for. This section proves the tools that every later section uses: the profile of a top algorithm and how the queries of a wanted set are shared through it (Lemmas~\ref{cp:T2}--\ref{cp:T4}), how a top algorithm with a given profile is built from a family of border decompositions (Lemma~\ref{cp:T5:lemma4} and Corollary~\ref{cp:T5}), and the composite with inner products as its bottom, \ThmD{} for one top algorithm (Theorem~\ref{cp:T6}) and \ThmDk{} for families (Theorem~\ref{cp:T7}). The section ends with the two composites that use the data structure of Theorem~\ref{cp:P3} as their bottom and pay per entry and per column (Propositions~\ref{cp:T8} and~\ref{cp:T9}). All decompositions are over $\Q$.

\subsection{Top algorithms, profiles and input families}\label{cp:sec:profile-def}

\begin{definition}[top algorithms and profiles]\label{cp:T1}
An \emph{exact top algorithm} for $\mat kbk$ is an exact bilinear algorithm $B$ over $\Q$ for the product of a $k\times b$ by a $b\times k$ matrix: terms $l\in[\RB]$ with linear forms $u_l$, $v_l$ in the entries of the two factors and $z$-forms $w_l=(w_l[o])_{o\in[k]^2}$ such that $(XY)_o=\sum_lw_l[o]\,(u_l\cdot X)(v_l\cdot Y)$ for every output $o$. Terms with zero $z$-form are dropped. For an output $o$ of a decomposition (exact or border) let $\Cov(o)=\{l:o\in\supp w_l\}$, the \emph{covering set} of $o$; $\abs{\Cov(o)}$ is the \emph{coverage} of $o$, and the \emph{maximal coverage} of $B$ is $\cB=\max_o\abs{\Cov(o)}$.

For a partition $(P_c)$ of the terms into classes let $\kappa_c=\max_o\abs{\Cov(o)\cap P_c}$ and define the \emph{profile}
\[
Z(t)=\sum_c\abs{P_c}^{1-t}\kappa_c^t\qquad(t\in[0,1],\ 0^0=1).
\]
$Z(0)$ is the number of terms, $Z(1)=\sum_c\kappa_c$, $Z(t)\le Z(0)$, and $\ln Z$ is convex in $t$. For a top algorithm $B$ with a partition we write
\[
\lamB(t)=\frac{\ln\ZB(t)}{\ln k}-2+2t,\qquad\delB=\lamB(0),\qquad\alB=\frac{\ln b}{\ln k},
\]
so that $\RB=k^{2+\delB}$.

An \emph{input family} is a sequence $(A_M)$, $M$ in an infinite set of integers, where $A_M$ is a border decomposition over $\Q[\lambda]$ of $s_M\cp\mat{a_M}{b_M}{a_M}$ (each copy identified with $\mat{a_M}{b_M}{a_M}$ by a fixed bijection of its variables, which changes no support) with $r_M$ terms of nonzero $z$-form, product degree at most $e_M$, and a partition into $m_M$ classes with profile $Z_M$ and maximal coverage $C_M$. Its hypotheses are:
\begin{itemize}
\item[\Hcond1] $a_M\to\infty$, $b_M\ge2$, and $\alpha_M:=\ln b_M/\ln a_M\to\alpha$;
\item[\Hcond2] $\rho_M:=\ln(r_M/s_M)/\ln a_M\to2$;
\item[\Hcond3] $(\ln(e_M+1)+\ln m_M)/\ln a_M\to0$;
\item[\Hcond4] $\cbar:=\sup_M\ln C_M/\ln a_M<\infty$.
\end{itemize}
We call the limit $\alpha$ in \Hcond1 the \emph{dual exponent of the family}. In this section, and wherever an input family is in play, $\alpha$ means this limit and not the dual exponent of matrix multiplication of Section~\ref{cp:sec:laser}. We put $\cstar:=\limsup_M\ln C_M/\ln a_M\le\cbar$, and define the per-level profile, the profile of the family and its \emph{typical support exponent}
\[
\mu_M(t)=\frac{\ln\bigl(s_M^{t-1}Z_M(t)\bigr)}{\ln a_M}-2+2t,\qquad\Lam(t)=\limsup_M\mu_M(t),\qquad\cg=\liminf_{t\to0+}\frac{\Lam(t)}{t},
\]
and $\Lamp(t)=\max\bigl(\Lam(t),\,t\cstar-2+2t\bigr)$ for $t\in(0,1]$.
\end{definition}

\begin{lemma}[basic bounds]\label{cp:T1:basic}
\begin{enumerate}[label=(\alph*)]
\item A border decomposition of $s\cp\mat aba$ has at least $sa^2$ terms of nonzero $z$-form. In particular every exact top algorithm for $\mat kbk$ has $\RB\ge k^2$, and $\lamB(t)\ge\delB\ge0$ for every $t\in[0,1]$.
\item For an input family, $\mu_M(t)\ge\rho_M-2\ge0$ for every $M$ and $t\in[0,1]$; hence $\Lam\ge0$ and $\cg\ge0$.
\end{enumerate}
\end{lemma}

\begin{proof}
(a) The $z$-flattening $\Phi(\lambda)=\sum_l(u_l\otimes v_l)w_l^{\mathsf T}$ of the decomposition has rank at most the number of terms of nonzero $z$-form over $\Q(\lambda)$, and $\lambda^{-h}\Phi(\lambda)=\Phi_T+\lambda\Phi'(\lambda)$, where $\Phi_T$ is the flattening of $s\cp\mat aba$. Its $sa^2$ outputs are bilinear forms on disjoint sets of monomials, so $\Phi_T$ has rank $sa^2$. A nonzero $sa^2$-minor of $\Phi_T$ is the value at $\lambda=0$ of the same minor of $\Phi_T+\lambda\Phi'(\lambda)$, which is therefore a nonzero polynomial; so the rank over $\Q(\lambda)$ is at least $sa^2$.
For the profile: every term has support at least $1$, so counting the pairs (term of class $c$, output in its support) gives $\sum_o\abs{\Cov(o)\cap P_c}\ge\abs{P_c}$, and the maximum over the $sa^2$ outputs is at least the mean: $\kappa_c\ge\abs{P_c}/(sa^2)$. Hence $Z(t)\ge\sum_c\abs{P_c}(sa^2)^{-t}=r(sa^2)^{-t}$ for $r$ terms. For $B$ ($s=1$, $a=k$) this gives $\lamB(t)\ge\ln\RB/\ln k-2=\delB\ge0$.
(b) With $Z_M(t)\ge r_M(s_Ma_M^2)^{-t}$, $\mu_M(t)\ge\ln(r_M/s_M)/\ln a_M-2=\rho_M-2$, and $r_M\ge s_Ma_M^2$ by (a).
\end{proof}

\begin{remark}\label{cp:T1:rem}
$\cg$ is the slope at $0$ of the limit profile, not the limit of the slopes of the $\mu_M$: if all class-$c$ terms have the same support size and every output the same class coverage, then $\mu_M(t)=(\rho_M-2)+\ln(\mathrm{mean}_lS_l^t)/\ln a_M$ with $S_l=\abs{\supp w_l}$, whose slope at $0$ is the geometric-mean support exponent, and $\cg$ can be larger than the limit of these slopes. Under the hypothesis of the first sentence, Markov's inequality applied to $S^t$ shows that for every $u'>\cg$ all but a fraction $a_M^{-\Omega(1)}$ of the terms have support at most $a_M^{u'+o(1)}$; this is what ``typical support'' refers to.
\end{remark}

\subsection{Shared queries and decoding}\label{cp:sec:shared}

Let $B$ be an exact top algorithm for $\mat kbk$ with a partition. Its $n$-th tensor power $B^{\otimes n}$ has the terms $l=(l_1,\dots,l_n)\in[\RB]^n$ with $w_l=w_{l_1}\otimes\dots\otimes w_{l_n}$; it computes $\mat K{b^n}K$, $K=k^n$, and its outputs are the blocks $o=(o_1,\dots,o_n)\in([k]^2)^n=[K]^2$. Since $w_l[o]=\prod_aw_{l_a}[o_a]$, $o$ lies in $\supp w_l$ if and only if $o_a\in\supp w_{l_a}$ for every $a$. For a wanted set $W$ in a $(KN_{\mathrm{sub}})\times(KN_{\mathrm{sub}})$ output matrix an entry is a triple $(o,r,s)$: a block $o$ and a position $(r,s)\in[N_{\mathrm{sub}}]^2$ inside it. Let $W_{rs}=\{o:(o,r,s)\in W\}$ and $T_{rs}=\bigcup_{o\in W_{rs}}\Cov(o)$.

\begin{lemma}[shared queries]\label{cp:T2}
\begin{enumerate}[label=(\alph*)]
\item For every $W$ and every $t\in[0,1]$,
\[
\sum_{(r,s)}\abs{T_{rs}}\le N_{\mathrm{sub}}^2\sum_{c_1,\dots,c_n}\min\Bigl(\prod_a\abs{P_{c_a}},\,\bar w\prod_a\kappa_{c_a}\Bigr)\le N_{\mathrm{sub}}^2\,\bar w^{\,t}\,\ZB(t)^n,\qquad\bar w=\abs W/N_{\mathrm{sub}}^2.
\]
\item If moreover $\abs W\le N^2/D^\kappa$ and $N\le N':=KN_{\mathrm{sub}}$, then $\sum_{(r,s)}\abs{T_{rs}}\le N'^2D^{-t\kappa}K^{\lamB(t)}$.
\end{enumerate}
\end{lemma}

\begin{proof}
(a) Partition the terms of $B^{\otimes n}$ into the classes $P_{c_1}\times\dots\times P_{c_n}$. The class-$(c_1,\dots,c_n)$ terms covering an output $o$ of the power are the tuples with $l_a\in\Cov(o_a)\cap P_{c_a}$ for every $a$, so there are at most $\prod_a\kappa_{c_a}$ of them. For one position, the class-$(c_1,\dots,c_n)$ part of $T_{rs}$ therefore has at most $\prod_a\abs{P_{c_a}}$ elements, and at most $\abs{W_{rs}}\prod_a\kappa_{c_a}$ by the union bound over the wanted blocks. Summing the minimum of these two bounds over the $N_{\mathrm{sub}}^2$ positions, a sum of minima is at most the minimum of the sums, $N_{\mathrm{sub}}^2\prod_a\abs{P_{c_a}}$ and $\abs W\prod_a\kappa_{c_a}$; this is the first inequality. Finally $\min(A,B)\le A^{1-t}B^t$ for $A,B\ge0$ and $t\in[0,1]$, and the sum over the tuples of $\prod_a\abs{P_{c_a}}^{1-t}\kappa_{c_a}^t$ is $\ZB(t)^n$.
(b) $\bar w=\abs W/N_{\mathrm{sub}}^2\le N^2D^{-\kappa}K^2/N'^2\le K^2D^{-\kappa}$, and $\ZB(t)^n=K^{\lamB(t)+2-2t}$ by the definition of $\lamB$. So the bound of (a) is at most $N_{\mathrm{sub}}^2K^{2t}D^{-t\kappa}K^{\lamB(t)+2-2t}=N'^2D^{-t\kappa}K^{\lamB(t)}$.
\end{proof}

The bound holds for every $W$ of the size, not only on average, and for every $t$; the number of classes, $m^n$ tuples, never enters, because the bound is a sum that factorises.

\begin{lemma}[distinct representatives]\label{cp:T3}
Every exact bilinear algorithm for $\mat kbk$ with $b\ge1$ has an injective map $\psi$ from its outputs to its terms with $\psi(o)\in\Cov(o)$ for every output $o$.
\end{lemma}

\begin{proof}
The $k^2$ output bilinear forms $(XY)_o$ are linearly independent. Since $(XY)_o=\sum_lw_l[o](u_l\cdot X)(v_l\cdot Y)$, a linear relation among the rows $(w_l[o])_l$, $o\in[k]^2$, of the $k^2\times\RB$ matrix $w$ would give the same relation among the $(XY)_o$; so $w$ has rank $k^2$. A nonzero $k^2\times k^2$ minor of $w$ has a nonzero term in its Leibniz expansion, that is, a bijection $o\mapsto\psi(o)$ onto $k^2$ of the columns with every $w_{\psi(o)}[o]\ne0$.
\end{proof}

\begin{lemma}[pruned Kronecker decoding]\label{cp:T4}
Let $\cB$ be the maximal coverage of $B$, fix a position $(r,s)$ with $W_{rs}\ne\emptyset$, and let the values $p_l$, $l\in T_{rs}$, be given. The outputs $y_o=\sum_lw_l[o]\,p_l$, $o\in W_{rs}$, of $B^{\otimes n}$ can be computed with at most $(n+1)(\cB+1)\abs{T_{rs}}$ operations on table entries and at most $n\cB\abs{T_{rs}}$ arithmetic operations, given direct-address tables on a universe of size $\RB^n=K^{2+\delB}$ that are allocated and initialised once, in $O(nK^{2+\delB})$ time, for all positions.
\end{lemma}

\begin{proof}
For $i=0,\dots,n$ consider mixed strings $e=(o_1,\dots,o_i,l_{i+1},\dots,l_n)$ and the values
\[
v_i(e)=\sum_{l_1,\dots,l_i}\ \prod_{a\le i}w_{l_a}[o_a]\;p_{(l_1,\dots,l_n)},
\]
so that $v_0=p$, $v_n(o)=y_o$, and $v_i(o_{\le i},l_{>i})=\sum_{l\in\Cov(o_i)}w_l[o_i]\,v_{i-1}(o_{<i},l,l_{>i})$. Define the needed sets
\[
M_i=\{(o_{\le i},l_{>i}):o\in W_{rs},\ l_a\in\Cov(o_a)\text{ for all }a>i\}.
\]
Then $M_n=W_{rs}$, $M_0=T_{rs}$, and $M_{i-1}=\{(o_{<i},l,l_{>i}):(o_{\le i},l_{>i})\in M_i,\ l\in\Cov(o_i)\}$: $M_{i-1}$ arises from $M_i$ by replacing the letter $o_i$ by each of its at most $\cB$ covering terms and removing duplicates, and the recursion for $v_i$ reads values only at elements of $M_{i-1}$. With $\psi$ of Lemma~\ref{cp:T3}, the map $(o_{\le i},l_{>i})\mapsto(\psi(o_1),\dots,\psi(o_i),l_{>i})$ is injective and maps $M_i$ into $M_0=T_{rs}$, since $\psi(o_a)$ covers $o_a$; so $\abs{M_i}\le\abs{T_{rs}}$ for every $i$. Compute the $M_i$ downwards from $M_n$ and the $v_i$ upwards from $v_0$; stage $i$ costs at most $\cB\abs{M_i}$ insertions and $\cB\abs{M_i}$ multiply-adds.
Coding an output letter $o_a$ by the term $\psi(o_a)$, the mixed strings lie in the universe $[\RB]^n$ ($\RB\ge k^2$). One direct-address table of that size per stage, each cell holding a stamp (the index of the current position) and a value, is allocated once; then every insertion and look-up costs $O(1)$, and no initialisation per position and no sorting are needed. Strings and stamps fit in $O(1)$ words, since $n\log\RB=O(\log D)$ in the applications below.
\end{proof}

B is a fixed finite object in all applications, so $\cB$ and $\psi$ are constants, and only the outer power $B^{\otimes n}$, whose length $n=\Theta(\log D)$ grows, is decoded coordinate by coordinate.

\subsection{From an input family to a top algorithm}\label{cp:sec:lemma4}

\begin{lemma}[top algorithms from border decompositions]\label{cp:T5:lemma4}
Let $A$ be a border decomposition over $\Q[\lambda]$ of $s\cp\mat aba$ with $r$ terms, product degree at most $e$, and a partition into $m$ classes with profile $Z_A$ and maximal coverage $C_A$; put $\mu(t)=\ln(s^{t-1}Z_A(t))/\ln a-2+2t$ and $c=\ln C_A/\ln a$. For every $j\ge1$ there is an exact top algorithm $B$ over $\Q$ for $\mat k{b^j}k$, $k=a^j$, with a partition, such that $\alB=\ln b/\ln a$ and, for every $t\in[0,1]$,
\[
\lamB(t)\le\Bigl(1-\frac1j\Bigr)\Bigl[\max\Bigl(\mu(t),\,tc-2+2t+\frac{\ln m}{\ln a}\Bigr)+\frac{\ln(2(e+1))}{\ln a}\Bigr]+\frac1j\,\frac{\ln((e+1)r)}{\ln a}.
\]
In particular $\delB\le(1-\frac1j)\bigl[\max(\rho-2,\,\ln m/\ln a-2)+\ln(2(e+1))/\ln a\bigr]+\frac1j\ln((e+1)r)/\ln a$ with $\rho=\ln(r/s)/\ln a$.
\end{lemma}

\begin{proof}
\emph{Step 1 (interpolation).} Write the decomposition as $\sum_lu_l(\lambda)\otimes v_l(\lambda)\otimes w_l(\lambda)=\lambda^hT+\lambda^{h+1}T'(\lambda)$ with $T=s\cp\mat aba$. Every coefficient of the left side is a polynomial in $\lambda$ of degree at most $e$, the product degree. For $e+1$ distinct rationals $\lambda_0,\dots,\lambda_e$ the coefficient of $\lambda^h$ of such a polynomial is a fixed rational combination $\sum_p\theta_pf(\lambda_p)$ of its values. Hence $T=\sum_p\sum_l\theta_p\,u_l(\lambda_p)\otimes v_l(\lambda_p)\otimes w_l(\lambda_p)$: an exact decomposition $A'$ over $\Q$ with at most $(e+1)r$ terms $(p,l)$. The $z$-form of $(p,l)$ is $\theta_pw_l(\lambda_p)$, whose support lies in the polynomial support of $w_l$. Give $(p,l)$ the class of $l$. Then $\abs{P'_c}\le(e+1)\abs{P_c}$ (terms with zero $z$-form are dropped) and $\kappa'_c\le(e+1)\kappa_c$, so $Z_{A'}(t)\le(e+1)Z_A(t)$, and the maximal coverage is at most $(e+1)C_A$.

\emph{Step 2 (constructive asymptotic sum inequality with class-homogeneous groups).} Let $T_0=\mat aba$. We build exact algorithms $G_j$ for $s\cp T_0^{\otimes j}$ with partitions, $G_1=A'$. Given $G_j$ (terms $g$, classes $d$, coverages $\kappa^j_d$), write $s\cp T_0^{\otimes(j+1)}=(s\cp T_0)\otimes T_0^{\otimes j}=\sum_m(u_m\otimes v_m\otimes w_m)\otimes T_0^{\otimes j}$ over the terms $m$ of $A'$. The $m$-th summand is one product of $T_0^{\otimes j}$: with $X^{(m)}=(u_m\otimes\mathrm{id})(x)$ and $Y^{(m)}=(v_m\otimes\mathrm{id})(y)$, matrices of shapes $a^j\times b^j$ and $b^j\times a^j$, it contributes $w_m\otimes(X^{(m)}Y^{(m)})$. Assign the terms of $A'$ to groups of $s$ slots so that every group holds terms of a single class: class $c$ fills $\lceil\abs{P'_c}/s\rceil$ groups, the last one possibly with empty slots, which get zero inputs. Compute the $s$ products of each group $\gamma$ at once with $G_j$, slot $\sigma'$ of $G_j$ receiving the product of the term $m(\gamma,\sigma')$ in that slot. The terms of $G_{j+1}$ are the pairs $(\gamma,g)$, and the $z$-form of $(\gamma,g)$ is $\sum_{\sigma'}w_{m(\gamma,\sigma')}\otimes w_g[\sigma',\cdot\,]$ over the filled slots, where $w_g[\sigma',\cdot\,]$ is $g$'s $z$-form on slot $\sigma'$; by bilinearity of block-matrix products (the $X$-blocks always on the left) $G_{j+1}$ is an exact algorithm for $s\cp T_0^{\otimes(j+1)}$. This is Schönhage's construction~\cite{Sch81} in constructive form.
The class of $(\gamma,g)$ is (class of $\gamma$, class of $g$), so $\abs{P^{j+1}_{(c,d)}}=\lceil\abs{P'_c}/s\rceil\,\abs{P^j_d}$. An output $(\sigma,o_1,o')$ of $s\cp T_0^{\otimes(j+1)}$ covered by a class-$(c,d)$ term $(\gamma,g)$ has a slot $\sigma'$ with $w_{m(\gamma,\sigma')}[\sigma,o_1]\ne0$ and $w_g[\sigma',o']\ne0$; mapping $(\gamma,g)$ to $(m(\gamma,\sigma'),g)$ for the least such $\sigma'$ is injective, because $m$ determines $\gamma$, and its image consists of pairs with $m\in P'_c$ covering $(\sigma,o_1)$ and $g\in P^j_d$ covering $(\sigma'(m),o')$. So $\kappa^{j+1}_{(c,d)}\le\kappa'_c\kappa^j_d$, and
\[
Z_{G_{j+1}}(t)\le Y(t)\,Z_{G_j}(t),\qquad Y(t)=\sum_c\lceil\abs{P'_c}/s\rceil^{1-t}{\kappa'_c}^t.
\]
\emph{Step 3 (padding).} $\lceil y\rceil\le y+1$ and $(y+1)^{1-t}\le y^{1-t}+1$ for $t\in[0,1]$, so
\[
Y(t)\le s^{t-1}Z_{A'}(t)+\sum_c{\kappa'_c}^t\le(e+1)s^{t-1}Z_A(t)+m\bigl((e+1)C_A\bigr)^t\le2(e+1)\max\bigl(s^{t-1}Z_A(t),\,mC_A^t\bigr).
\]
\emph{Step 4 (one copy).} Let $B$ be $G_j$ restricted to the variables of slot~$1$ (all other variables set to $0$): an exact algorithm over $\Q$ for $T_0^{\otimes j}=\mat{a^j}{b^j}{a^j}$. It keeps the classes; class sizes do not grow and coverages are maxima over fewer outputs. So $\ZB(t)\le Z_{G_j}(t)\le Y(t)^{j-1}Z_{A'}(t)$, and with $Z_{A'}(t)\le Z_{A'}(0)\le(e+1)r$ and $\ln k=j\ln a$,
\[
\lamB(t)=\frac{\ln\ZB(t)}{j\ln a}-2+2t\le\Bigl(1-\frac1j\Bigr)\Bigl(\frac{\ln Y(t)}{\ln a}-2+2t\Bigr)+\frac1j\Bigl(\frac{\ln((e+1)r)}{\ln a}-2+2t\Bigr).
\]
By Step 3, $\ln Y(t)/\ln a-2+2t\le\ln(2(e+1))/\ln a+\max\bigl(\mu(t),\,tc-2+2t+\ln m/\ln a\bigr)$; dropping $-2+2t\le0$ in the last term gives the claim. $\alB=\ln b^j/\ln a^j$.
\end{proof}

\begin{corollary}[padding lemma]\label{cp:T5}
Let $(A_M)$ satisfy \Hcond1--\Hcond4. For every $t\in(0,1]$ and every $\eta>0$ there are $M$ and $j$ such that the top algorithm $B=B_{M,j}$ of Lemma~\ref{cp:T5:lemma4} applied to $A_M$ satisfies $\alB=\alpha_M$, $\abs{\alB-\alpha}\le\eta$, $\delB\le\eta$ and $\lamB(t)\le\Lamp(t)+\eta$. Moreover:
\begin{enumerate}[label=(\roman*)]
\item $\Lamp(t)=\Lam(t)$ for every $t\le2/(\cstar+2)$;
\item $\Lamp(1)=\Lam(1)=\cstar$;
\item $\Lamp(t)\le t\cstar$ for every $t\in(0,1]$.
\end{enumerate}
\end{corollary}

\begin{proof}
Put $c_M=\ln C_M/\ln a_M$ and $\mathrm{pad}_M(t)=tc_M-2+2t+\ln m_M/\ln a_M$. Lemma~\ref{cp:T5:lemma4} applied to $A_M$ and $j$ gives
\[
\lamB(t)\le\Bigl(1-\frac1j\Bigr)\Bigl[\max\bigl(\mu_M(t),\mathrm{pad}_M(t)\bigr)+\frac{\ln(2(e_M+1))}{\ln a_M}\Bigr]+\frac1j\,\frac{\ln((e_M+1)r_M)}{\ln a_M}.
\]
The limsup of a maximum of two real sequences is the maximum of their limsups. $\limsup_M\mu_M(t)=\Lam(t)$ by definition, and $\limsup_M\mathrm{pad}_M(t)=t\cstar-2+2t$, because $t\ge0$ and $\ln m_M/\ln a_M\to0$ by \Hcond3. Hence $\max(\mu_M(t),\mathrm{pad}_M(t))\le\Lamp(t)+\eta/3$ for all large $M$. Also $\ln(2(e_M+1))/\ln a_M\le\eta/3$ for all large $M$ by \Hcond3 and \Hcond1, and $\abs{\alpha_M-\alpha}\le\eta$ by \Hcond1. For $\delB$ take $t=0$ in the same bound: $\mu_M(0)=\rho_M-2\to0$ by \Hcond2, and $\mathrm{pad}_M(0)=\ln m_M/\ln a_M-2<0$ for large $M$. Fix one $M$ large enough for all these requirements. The bracket is then at most $\Lamp(t)+2\eta/3$, which is nonnegative because $\mu_M\ge0$ (Lemma~\ref{cp:T1:basic}), so the factor $1-1/j$ only lowers it; and for this fixed $M$, $(1/j)\ln((e_M+1)r_M)/\ln a_M\le\eta/3$ for all large $j$. This gives $\lamB(t)\le\Lamp(t)+\eta$, and in the same way $\delB\le\eta$.
(i) If $t\le2/(\cstar+2)$ then $t\cstar-2+2t=t(\cstar+2)-2\le0\le\Lam(t)$.
(ii) $Z_M(1)=\sum_c\kappa_c\ge\sum_c\abs{\Cov(o^*)\cap P_c}=C_M$ for an output $o^*$ of maximal coverage, so $\mu_M(1)\ge c_M$ and $\Lam(1)\ge\cstar$, which is the padding term at $t=1$. Conversely $\kappa_c\le C_M$ gives $Z_M(1)\le m_MC_M$, so $\mu_M(1)\le c_M+\ln m_M/\ln a_M$ and $\Lam(1)\le\cstar$.
(iii) By Hölder's inequality $Z_M(t)\le C_M^t\sum_c\abs{P_c}^{1-t}\le C_M^tm_M^tr_M^{1-t}$, so $\mu_M(t)\le(1-t)(\rho_M-2)+t(c_M+\ln m_M/\ln a_M)$, and the limsup is at most $t\cstar$ by \Hcond2 and \Hcond3; also $t\cstar-2+2t\le t\cstar$.
\end{proof}

The padding term $t\cstar-2+2t$ cannot simply be dropped by a general argument: \Hcond1--\Hcond4 do not imply that it is below $\Lam(t)$ for $t>2/(\cstar+2)$, so the correct quantity for large $t$ is the maximum $\Lamp$. Part (iii) is not used below.

\subsection{\texorpdfstring{\ThmD}{Theorem D}: inner products as the bottom}\label{cp:sec:thmD}

The next theorem is the composite for one fixed top algorithm. Its proof is the second half of the proof of \ThmDk{} below, and it is stated separately because Sections~\ref{cp:sec:range} and~\ref{cp:sec:transfer} apply it to top algorithms that do not come from one input family.

\begin{theorem}[\ThmD]\label{cp:T6}
Let $B$ be an exact top algorithm over $\Q$ for $\mat kbk$ with $2\le b\le k$, $\RB=k^{2+\delB}$ terms and a partition of its terms with profile $\lamB$, and $\alB=\ln b/\ln k$. Let $\kappa>0$, $x>0$ and $t\in(0,1]$ satisfy
\[
\Qgain B\kappa(x,t):=\frac{t\kappa(1+\alB x)-x\lamB(t)-1}{1+\alB x}>0. \tag{D1$^B$}
\]
Then for every $\eps$ with $0<\eps<\Erange B(x):=(1+\alB x)/\bigl(1+x(1+\delB)\bigr)$ and every $\gamma$ with
\[
0<\gamma<\gbound B\kappa(x,t,\eps):=\min\Bigl\{\Qgain B\kappa(x,t),\ \frac1\eps-\frac{1+x(1+\delB)}{1+\alB x},\ \frac{2x}{1+\alB x}\Bigr\},
\]
$\Hypk\kappa(\eps,\gamma)$ holds (time model), by \AlgD{} below with top algorithm $B$.
\end{theorem}

\paragraph{\AlgD.} Fix $B$, $\kappa$, $x$, $t$, $\eps$ and $\gamma$ as in the theorem. Given $N$, $D$ with $2\le D\le N^\eps$, matrices $X$, $Y$ and $W$ with $\abs W\le N^2/D^\kappa$:
\begin{enumerate}
\item Put $n=\lfloor x\ln D/((1+\alB x)\ln k)\rfloor$, $K=k^n$, $B'=b^n$ and $D_A=\lceil D/B'\rceil$. Pad $X$ and $Y$ with zero columns and rows to the inner dimension $D'=B'D_A$, and to $N'=K\lceil N/K\rceil$ rows and columns; put $N_{\mathrm{sub}}=N'/K$, and view $X$ as a $K\times B'$ matrix of blocks of size $N_{\mathrm{sub}}\times D_A$ and $Y$ as a $B'\times K$ matrix of blocks of size $D_A\times N_{\mathrm{sub}}$.
\item Scale $B$ to integer coefficients (common denominators $d_u$, $d_v$, $d_w$ of the forms $u$, $v$, $w$). For every $l\in[\RB]^n$ compute the encodings $X_l=\sum u_l[I,J]X_{I,J}$ and $Y_l=\sum v_l[J,I]Y_{J,I}$ of $B^{\otimes n}$ recursively, one coordinate at a time.
\item Bucket $W$ by inner position: every wanted entry is a triple $(o,r,s)$ with a block $o\in[K]^2$ and a position $(r,s)\in[N_{\mathrm{sub}}]^2$; sort by $(r,s)$ with a counting sort over the $N_{\mathrm{sub}}^2$ positions.
\item For every position $(r,s)$ with $W_{rs}\ne\emptyset$: compute $M_n,\dots,M_0=T_{rs}$ (Lemma~\ref{cp:T4}); for every $l\in T_{rs}$ compute $p_l=(X_lY_l)[r,s]$ as an inner product of length $D_A$ of row $r$ of $X_l$ with column $s$ of $Y_l$; decode $y_o=\sum_lw_l[o]p_l$ for $o\in W_{rs}$ by Lemma~\ref{cp:T4}; output $y_o/(d_ud_vd_w)^n$.
\end{enumerate}

\begin{proof}[Proof of Theorem~\ref{cp:T6}]
\emph{Correctness.} $B^{\otimes n}$ is an exact bilinear algorithm over $\Q$ for $\mat K{B'}K$. Applied to the block matrices of Step~1 it gives, by bilinearity of block-matrix products with the $X$-blocks on the left, that block $o$ of $XY$ equals $\sum_lw_l[o]X_lY_l$; so entry $(r,s)$ of block $o$ is $\sum_{l\in\Cov(o)}w_l[o]\,p_l$ with $\Cov(o)\subseteq T_{rs}$, and Lemma~\ref{cp:T4} computes exactly these sums. With integer coefficients the sum is $(d_ud_vd_w)^n$ times the wanted integer, so the division is exact.

\emph{Sizes.} By the choice of $n$, $K\le D^{x/(1+\alB x)}<kK$, so $K$ is within a factor $k$ of $D^{x/(1+\alB x)}$. Further $B'=K^{\alB}\le D^{\alB x/(1+\alB x)}\le D$, so $D'\le2D$, and since $B'>(D^{x/(1+\alB x)}/k)^{\alB}=D^{\alB x/(1+\alB x)}/b$ we get $D^{1/(1+\alB x)}\le D_A\le2bD^{1/(1+\alB x)}$. Since $\eps<\Erange B(x)<(1+\alB x)/x$, we have $K\le D^{x/(1+\alB x)}\le N^{\eps x/(1+\alB x)}\le N$, so $N\le N'\le2N$. (What the padding needs is $K\le N$ and $B'\le D$; the bound ``$K\le D$'' is false once $x(1-\alB)>1$, and it is not used.)

\emph{Costs.} We bound each step by $O(N^2\log D\cdot D^{-e_i})$ with an exponent $e_i$; all constants depend on $B$, $x$, $\kappa$ and $\eps$ only. By Lemma~\ref{cp:T1:basic}, $\lamB(t)\ge\delB\ge0$.
\begin{itemize}
\item \emph{Queries} (Step~4). By Lemma~\ref{cp:T2}(b), $\sum_{rs}\abs{T_{rs}}\le N'^2D^{-t\kappa}K^{\lamB(t)}$; each inner product costs $D_A$ operations and the decoding $O(n\cB)$ per element of $T_{rs}$ (Lemma~\ref{cp:T4}). With the size bounds this is $O(nN^2D^{-t\kappa+(x\lamB(t)+1)/(1+\alB x)})$, that is $e_1=\Qgain B\kappa(x,t)$.
\item \emph{Encodings} (Step~2). Stage $a$ of the recursion has $\RB^a(kb)^{n-a}$ blocks, each computed with at most $kb$ operations on blocks of $N_{\mathrm{sub}}D_A$ entries; since $\RB\ge k^2\ge kb$, the total is at most $nkb\,\RB^nN_{\mathrm{sub}}D_A=O(nNK^{1+\delB}D_A)$, likewise for the $Y_l$. With $N\ge D^{1/\eps}$ this is $O(nN^2D^{-1/\eps+(1+x(1+\delB))/(1+\alB x)})$: $e_2=1/\eps-(1+x(1+\delB))/(1+\alB x)$.
\item \emph{Tables} (Lemma~\ref{cp:T4}). $O(nK^{2+\delB})=O(nN^2D^{-2/\eps+(2+\delB)x/(1+\alB x)})$; the exponent $e_3=2/\eps-(2+\delB)x/(1+\alB x)$ exceeds $e_2$ by $1/\eps-(x-1)/(1+\alB x)>0$, because $1/\eps>(1+x(1+\delB))/(1+\alB x)$.
\item \emph{Bucketing} (Step~3). $O(N_{\mathrm{sub}}^2+\abs W)$, with $N_{\mathrm{sub}}^2=N'^2/K^2=O(N^2D^{-2x/(1+\alB x)})$ and $\abs W\le N^2D^{-\kappa}$: $e_4=2x/(1+\alB x)$ and $e_5=\kappa$.
\item \emph{Input and output.} Reading and padding the input costs $O(ND)=O(N^2D^{-(1/\eps-1)})$, and $e_6=1/\eps-1\ge e_2$ because $1+x(1+\delB)\ge1+\alB x$; the output has $\abs W$ numbers ($e_5$).
\end{itemize}
\emph{Word size.} After scaling, the entries of $X_l$ and $Y_l$ are at most $(kbH)^n$ in absolute value, $H$ the largest integer coefficient, which is $K^{O(1)}=N^{O(1)}$; the values $p_l$ and $v_i$ of Lemma~\ref{cp:T4} are at most $(\RB H)^{2n}DN^{O(1)}=N^{O(1)}$, also for input entries of absolute value $N^{O(1)}$. So every number fits in $O(1)$ words and every arithmetic operation takes $O(1)$ steps.

\emph{Conclusion.} Since $\lamB(t)\ge0$, $\Qgain B\kappa(x,t)<t\kappa\le\kappa=e_5$. So $\min_ie_i\ge\gbound B\kappa(x,t,\eps)>\gamma$, and with $\gamma<\gamma'<\gbound B\kappa(x,t,\eps)$ the running time is $O(N^2\log D\cdot D^{-\gamma'})\le CN^2D^{-\gamma}$ for all $D\ge2$; also $\gamma<\kappa$, as Definition~\ref{cp:P1} requires. All bounds above hold for every $D\ge2$, including $n=0$. The algorithm is deterministic.
\end{proof}

\begin{remark}[the algebraic form of \AlgD]\label{cp:T6:alg}
For each $N$, $D$ and $W$, \AlgD{} is a division-free straight-line program over $\Q$ in the entries of $X$ and $Y$: its control flow ($n$, $K$, $D_A$, padding, bucketing, the needed sets $M_i$, the order of the operations) depends on $N$, $D$ and $W$ only; its arithmetic consists of integer linear combinations (encodings), products of an $X$-encoding entry with a $Y$-encoding entry (the inner products, the only non-scalar multiplications), linear combinations with the fixed coefficients of $B$ (decoding) and one multiplication by the rational constant $(d_ud_vd_w)^{-n}$. It is correct as a polynomial identity, since $B^{\otimes n}$ is an exact bilinear algorithm over $\Q$, so it computes $XY$ at $W$ for all integer matrices. Its number of non-scalar multiplications is at most its running time. This is the form the converses of Section~\ref{cp:sec:range} need.
\end{remark}

\subsection{\texorpdfstring{\ThmDk}{Theorem D\_kappa}: input families}\label{cp:sec:thmDk}

\begin{theorem}[\ThmDk]\label{cp:T7}
Let $(A_M)$ satisfy \Hcond1--\Hcond4 with $\alpha\in(0,1)$, profile $\Lam$ and $\cstar$ as in Definition~\ref{cp:T1}. Let $\kappa>0$, $x>0$ and $t\in(0,1]$ satisfy
\[
\Qgain{\mathrm D}\kappa(x,t):=\frac{t\kappa(1+\alpha x)-x\Lamp(t)-1}{1+\alpha x}>0.\tag{D1}
\]
Then for every $\eps$ with $0<\eps<\Erange{\mathrm D}(x):=(1+\alpha x)/(1+x)$, $\Hypk\kappa(\eps,\gamma)$ holds (time model) for every
\[
0<\gamma<\gbound{\mathrm D}\kappa(x,t,\eps):=\min\Bigl\{\Qgain{\mathrm D}\kappa(x,t),\ \frac1\eps-\frac{1+x}{1+\alpha x},\ \frac{2x}{1+\alpha x}\Bigr\}.
\]
Moreover:
\begin{enumerate}[label=(\alph*)]
\item If $\cg=0$, then for every $\kappa>0$ and every $0<\eps\le\alpha$, the dual exponent of the family, there is $\gamma>0$ with $\Hypk\kappa(\eps,\gamma)$.
\item For $0<\eps<\alpha$, $\Hypk\kappa(\eps,\gamma)$ holds for every $\gamma$ with
\[
0<\gamma<\gbound{\mathrm D}\kappa(\eps):=\min\Bigl\{\sup_{0<t\le1}\Bigl(t\kappa-\frac{\Lamp(t)}{\alpha}\Bigr),\ \frac1\eps-\frac1\alpha,\ \frac2\alpha\Bigr\}.
\]
\item For every $t\in(0,1]$ with $t\kappa\alpha>\Lamp(t)$ and $t\kappa<1$ put $\xmin(t)=(1-t\kappa)/(t\kappa\alpha-\Lamp(t))$. Then $\Hypk\kappa(\eps,\gamma)$ holds with some $\gamma>0$ for every
\[
\eps<\sup_t\Bigl[\alpha+\frac{1-\alpha}{1+\xmin(t)}\Bigr].
\]
\end{enumerate}
\end{theorem}

For $t\le2/(\cstar+2)$ the condition \Dcond{} reads with $\Lam$ in place of $\Lamp$ (Corollary~\ref{cp:T5}(i)).

\begin{proof}
Fix $\kappa$, $x$, $t$ with \Dcond, $\eps<\Erange{\mathrm D}(x)$ and $\gamma<\gbound{\mathrm D}\kappa(x,t,\eps)$. The quantities $\Qgain B\kappa(x,t)$, $\Erange B(x)$ and $\gbound B\kappa(x,t,\eps)$ of Theorem~\ref{cp:T6} are continuous in $(\alB,\delB,\lamB(t))$ and equal $\Qgain{\mathrm D}\kappa(x,t)$, $\Erange{\mathrm D}(x)$ and $\gbound{\mathrm D}\kappa(x,t,\eps)$ at $(\alpha,0,\Lamp(t))$; they decrease as $\delB$ or $\lamB(t)$ increase. So there is $\eta>0$ such that every top algorithm $B$ with $\abs{\alB-\alpha}\le\eta$, $\delB\le\eta$ and $\lamB(t)\le\Lamp(t)+\eta$ satisfies (D1$^B$), $\eps<\Erange B(x)$ and $\gamma<\gbound B\kappa(x,t,\eps)$. Corollary~\ref{cp:T5} gives such a $B$, for $\mat k{b}k$ with $b=b_M^j\ge2$; since $\alB=\alpha_M$ is close to $\alpha<1$, also $b\le k$. Theorem~\ref{cp:T6} with this $B$ gives $\Hypk\kappa(\eps,\gamma)$. The order of the choices is: $\eps$ and $\kappa$, then $t$ and $x$, then $\gamma$, then $\eta$, then $(M,j)$, that is $B$; the algorithm then runs for all $N$ and $D$, with $n$ chosen from $D$.

(a) Since $\cg=\liminf_{t\to0+}\Lam(t)/t=0$, there are arbitrarily small $t>0$ with $\Lam(t)\le t\kappa\alpha/2$; take such a $t\le2/(\cstar+2)$, so that $\Lamp(t)=\Lam(t)$. Then $t\kappa(1+\alpha x)-x\Lam(t)-1\ge xt\kappa\alpha/2+t\kappa-1>0$ for every $x>2/(t\kappa\alpha)$, so \Dcond{} holds. Since $\alpha<1$, $(1+\alpha x)/(1+x)>\alpha$ for every $x>0$, so every $\eps\le\alpha$ satisfies $\eps<\Erange{\mathrm D}(x)$, and then $\gbound{\mathrm D}\kappa(x,t,\eps)>0$.

(b) Fix $0<\eps<\alpha$ and $\gamma$ below the minimum, and choose $t$ with $t\kappa-\Lamp(t)/\alpha>\gamma$. As $x\to\infty$, $\Qgain{\mathrm D}\kappa(x,t)\to t\kappa-\Lamp(t)/\alpha$; $1/\eps-(1+x)/(1+\alpha x)$ decreases to $1/\eps-1/\alpha>\gamma$; and $2x/(1+\alpha x)$ increases to $2/\alpha>\gamma$. So a large $x$ satisfies all three, and $\eps<\alpha<\Erange{\mathrm D}(x)$. The first term of the minimum is positive for every $\kappa>0$ when $\cg=0$, by the choice of $t$ in (a).

(c) \Dcond{} reads $x(t\kappa\alpha-\Lamp(t))>1-t\kappa$, that is $x>\xmin(t)$, and $\Erange{\mathrm D}(x)=\alpha+(1-\alpha)/(1+x)$ decreases in $x$. For $\eps<\alpha+(1-\alpha)/(1+\xmin(t))$ take $x>\xmin(t)$ close enough to $\xmin(t)$ that $\eps<\Erange{\mathrm D}(x)$; then $\gbound{\mathrm D}\kappa(x,t,\eps)>0$.
\end{proof}

The level-1 family of Section~\ref{cp:sec:level1} has $\cg=0$ (Lemma~\ref{cp:L8}), and (a) applied to it is Theorem~\ref{cp:A1}. The tilt $t$ in (b) and (c) ranges over all of $(0,1]$ by Corollary~\ref{cp:T5}; for $t\le2/(\cstar+2)$, $\Lamp=\Lam$.

\subsection{Composites with the data structure of~\texorpdfstring{\cite{v2}}{[Maj26a]} as bottom}\label{cp:sec:PE}

The next two composites answer the top's sub-products with the data structure of Theorem~\ref{cp:P3}, and pay for the queries per wanted entry or per column instead of sharing them. They need only the coverage of the top algorithm and the union of the covering sets of a column, not its profile. For $\alpha>0$ and $c>0$ consider the hypotheses:
\begin{itemize}
\item[$\TPE(\alpha,c)$:] for every $\eta>0$ there is an exact top algorithm $B$ over $\Q$ for $\mat kbk$ with $k,b\ge2$, at most $k^{2+\eta}$ terms, $\ln b/\ln k\ge\alpha-\eta$, and every output in at most $k^{c+\eta}$ terms;
\item[$\Tcol(\alpha,c')$:] for every $\eta>0$ there is an exact top algorithm $B$ over $\Q$ for $\mat kbk$ with $k,b\ge2$, at most $k^{2+\eta}$ terms and $\ln b/\ln k\ge\alpha-\eta$ such that for every column $\nu$ of its output matrix the set $U_B(\nu)$ of the terms covering some output of that column has at most $k^{1+c'+\eta}$ elements.
\end{itemize}

\begin{proposition}[per-entry composite]\label{cp:T8}
Let $\TPE(\alpha,c)$ hold with $\alpha>\es$, and let $Q\in(0,1]$ with $c>Q\alpha$. For every $\eps<E_Q(\alpha,c):=c\es/(Q\es+c-Q\alpha)$ there are $\gamma>0$, $Q'<Q$ and a deterministic algorithm that, for all $N\ge1$ and $2\le D\le N^\eps$ and integer matrices $X\in\Z^{N\times D}$, $Y\in\Z^{D\times N}$ with entries of absolute value $N^{O(1)}$, preprocesses in time and space $O(N^2/D^\gamma)$, after which any single entry of $XY$ is computed in $O(D^{Q'})$ time (time model).
\end{proposition}

\begin{proof}
(a) \emph{Choice of $x$.} $\varphi(x)=(1+\alpha x)/(x+1/\es)$ is strictly increasing in $x$, since its derivative has the sign of $\alpha/\es-1>0$, and at $x_Q=Q/(c-Q\alpha)$ it equals $E_Q$: $1+\alpha x_Q=c/(c-Q\alpha)$ and $x_Q+1/\es=(Q\es+c-Q\alpha)/(\es(c-Q\alpha))$. So for $\eps<E_Q$ there is $x\in(0,x_Q)$ with $\varphi(x)>\eps$, and since $x(c-Q\alpha)<Q$ is strict, continuity gives $\eta_0>0$, $\eps_A<\es$ and $q_A>0$ with $x\bigl((c+\eta_0)-Q(\alpha-\eta_0)\bigr)<Q-q_A$ and $(1+(\alpha-\eta_0)x)/(x+1/\eps_A)>\eps+\eta_0$.
(b) Theorem~\ref{cp:P3} for $(\eps_A,q_A)$ gives $\gamma_A>0$.
(c) Take $\eta\le\eta_0$ with $x\eta<\gamma_A/2$ and $B$ from $\TPE(\alpha,c)$ with this $\eta$: its rank excess is $\delB\le\eta$, $\alB\ge\alpha-\eta_0$, and its maximal coverage is $\cB\le k^{c+\eta_0}$. $B$ depends on $\eps$ only.
(d) \emph{Algorithm.} For an instance put $n=\lfloor x\ln D/((1+\alB x)\ln k)\rfloor$, $K=k^n$, $B'=b^n$, $D_A=\lceil D/B'\rceil$, pad as in \AlgD{} (by factors at most $2$), encode $X_l$, $Y_l$ for all $l\in[\RB]^n$ recursively after scaling $B$ to integer coefficients, and preprocess every product $P_l=X_lY_l$ with the data structure of Theorem~\ref{cp:P3}; all of them are kept, so space equals time. A query $(i,j)$, with top block $o=(I,J)$ and inner position $(r,s)$, returns the sum of $w_l[o]P_l[r,s]$ over the $l$ with $w_l[o]\ne0$, which form a product of per-coordinate covering sets and are enumerated directly, divided by the scaling factor; the division is exact since the true value is an integer.
(e) \emph{Costs.} Preprocessing: $\RB^nO(N_{\mathrm{sub}}^2/D_A^{\gamma_A})=O(N^2K^{\delB}D_A^{-\gamma_A})=O(N^2D^{-(\gamma_A-x\delB)/(1+\alB x)})$. Encodings: $O(nkbNK^{1+\delB}D_A)$, negligible because $N_{\mathrm{sub}}\ge D_A^{1/\eps_A}$ and $1/\eps_A>5$. One query: at most $\cB^n\le K^{c+\eta_0}$ sub-queries of cost $O(D_A^{q_A})$ each, plus $O(n)$ for each coefficient, so $O(nD^{Q''})$ with $Q''=(x(c+\eta_0)+q_A)/(1+\alB x)<Q$ by (a) (using $\alB\ge\alpha-\eta_0$); $Q':=(Q+Q'')/2$ absorbs the factor $n=O(\log D)$.
(f) \emph{Range.} The sub-instances need $2\le D_A\le N_{\mathrm{sub}}^{\eps_A}$. Since $N_{\mathrm{sub}}\ge N/K\ge D^{1/\eps-x/(1+\alB x)}$ and $D_A\le2bD^{1/(1+\alB x)}$, it suffices that $1/(1+\alB x)<\eps_A(1/\eps-x/(1+\alB x))$, that is $\eps<(1+\alB x)/(x+1/\eps_A)$, which holds with the margin $\eta_0$ of (a); so the condition holds for $D\ge D_0$, and for $D<D_0$ the product is computed directly. All integers stay $N^{O(1)}$, so every arithmetic operation takes $O(1)$ word operations, and the sub-instances have entries of absolute value $N^{O(1)}=N_{\mathrm{sub}}^{O(1)}$, as Theorem~\ref{cp:P3} requires.
\end{proof}

With $Q=\frac12$ the queries of a wanted set of size $N^2/\sqrt D$ cost $O(N^2D^{Q'-1/2})$, so Proposition~\ref{cp:T8} gives $\Hyp(\eps,\gamma')$ with some $\gamma'>0$ for every $\eps<E_{1/2}(\alpha,c)=2c\es/(\es+2c-\alpha)$; with $Q=1$ it gives the per-entry hinted-OMv bounds of Section~\ref{cp:sec:composites}.

\begin{proposition}[column composite]\label{cp:T9}
If $\Tcol(\alpha,c')$ holds with $c'>\alpha>\es$, then for every $\eps<E_1(\alpha,c')=c'\es/(\es+c'-\alpha)$ there are $\gamma>0$, $Q'<1$ and a deterministic algorithm with preprocessing in time and space $O(N^2/D^\gamma)$ after which any whole column of $XY$ is computed in $O(ND^{Q'})$ time, for every column (worst case; time model; inputs as in Proposition~\ref{cp:T8}). By transposing the input, preprocessing the pair $(Y^{\mathsf T},X^{\mathsf T})$, the same holds for rows.
\end{proposition}

\begin{proof}
As for Proposition~\ref{cp:T8}, with the column cost in place of the entry cost. A column of $XY$ is a pair (top column block $J$, inner column $s$). For each of the $N_{\mathrm{sub}}$ inner rows $r$ it needs $P_l[r,s]$ for $l\in U_J=\prod_aU_B(J_a)$, at most $K^{1+c'+\eta}$ values, and then the $K$ outputs $(I,r)$, $I\in[K]$, are the image of these values under the column map of $B^{\otimes n}$, the Kronecker product of $n$ fixed $k\times\abs{U_B(J_a)}$ matrices. Applied one coordinate at a time it costs at most $nk\abs{U_J}$ operations, because every intermediate array has at most $\abs{U_J}$ entries; this uses $k\le\abs{U_B(\nu)}$, which holds because the $k$ outputs of a column $\nu$ are linearly independent bilinear forms, each a combination of the products of the terms in $U_B(\nu)$. So a column costs $O(nNK^{c'+\eta}D_A^{q_A})=O(ND^{Q'})$ with $Q'<1$ if and only if $x(c'+\eta-\alB)<1-q_A$, and the range is $\varphi(x)$ as in Proposition~\ref{cp:T8}; the supremum over $x<1/(c'-\alpha)$ is $E_1(\alpha,c')$. The union bound $\abs{U_J}$ holds for every column, so the bound is worst case.
\end{proof}

Lemma~\ref{cp:L10} shows that the standard level-1 family satisfies both hypotheses at every admissible point; Section~\ref{cp:sec:composites} draws the consequences.

\section{The level-1 laser family}\label{cp:sec:level1}

This section constructs, from Coppersmith's mixture of the tensors $F_6$ and $F_7$~\cite{Cop97}, the input family to which Section~\ref{cp:sec:range} applies \ThmDk, and the top algorithms of the per-entry and column composites. Two border decompositions of $F_q$ enter: the standard one $\Fstd$ of Section~\ref{cp:sec:laser}, and a merged one $\Fmrg$ with the same coverage at every output but a different distribution of the supports over the terms (Section~\ref{cp:sec:merged}). The laser restriction is the same for both; its count is Le Gall's Theorem~\ref{cp:P6:LG} (Section~\ref{cp:sec:laser1}). With the merged decomposition the family has typical support exponent $\cg=0$ (Section~\ref{cp:sec:moment}); with the standard one it gives the coverage and the column union that Propositions~\ref{cp:T8} and~\ref{cp:T9} need (Section~\ref{cp:sec:std}).

\subsection{The merged decomposition}\label{cp:sec:merged}

\begin{proposition}[merged decomposition]\label{cp:L1}
Write $S_x=\sum_{i=1}^qx_i$, and $S_y$, $S_z$ likewise, and define the $q+2$ products
\[
A_i=(x_0+\lambda x_i)(y_0+\lambda y_i)(z_0+\lambda z_i)\ (i\in[q]),\qquad\Oterm=x_0y_0z_0,
\]
\[
\Bp=(x_0+\lambda^2S_x-\lambda^3x_{q+1})(y_0+\lambda^2S_y-\lambda^3y_{q+1})(z_0+\lambda^2S_z-\lambda^3z_{q+1}).
\]
Over $\Z[\lambda]$, and hence over every commutative ring,
\[
\lambda\sum_{i=1}^qA_i-\Bp+(1-q\lambda)\Oterm=\lambda^3F_q+\lambda^4G_q(\lambda),
\]
where $G_q$ is a polynomial in $\lambda$ of degree at most $5$ whose coefficients are integer polynomials in the variables. So the terms $\lambda A_1,\dots,\lambda A_q$, $-\Bp$ and $(1-q\lambda)\Oterm$ form a border decomposition $\Fmrg$ of $F_q$ with $q+2$ terms and approximation degree~$3$.
\end{proposition}

\begin{proof}
\emph{The terms $A_i$.} $A_i=\Oterm+\lambda(x_iy_0z_0+x_0y_iz_0+x_0y_0z_i)+\lambda^2(x_iy_iz_0+x_iy_0z_i+x_0y_iz_i)+\lambda^3x_iy_iz_i$. Summing over $i$ and multiplying by $\lambda$,
\[
\lambda\sum_iA_i=q\lambda\Oterm+\lambda^2L_S+\lambda^3P+\lambda^4\sum_ix_iy_iz_i,
\]
with $L_S=S_xy_0z_0+x_0S_yz_0+x_0y_0S_z$ and $P=\sum_i(x_iy_iz_0+x_iy_0z_i+x_0y_iz_i)$.
\emph{The term $\Bp$.} Write its factors as $x_0+u_x$, $y_0+u_y$, $z_0+u_z$ with $u_x=\lambda^2S_x-\lambda^3x_{q+1}$ and so on; each $u$ is divisible by $\lambda^2$. Expanding, $\Bp=\Oterm+(u_xy_0z_0+x_0u_yz_0+x_0y_0u_z)+(u_xu_yz_0+u_xy_0u_z+x_0u_yu_z)+u_xu_yu_z$; the quadratic part is divisible by $\lambda^4$ and the cubic part by $\lambda^6$, and the linear part is $\lambda^2L_S-\lambda^3Q'$ with $Q'=x_{q+1}y_0z_0+x_0y_{q+1}z_0+x_0y_0z_{q+1}$. Hence $-\Bp=-\Oterm-\lambda^2L_S+\lambda^3Q'+\lambda^4(\cdots)$.
\emph{Adding} the term $(1-q\lambda)\Oterm$: the coefficient of $\Oterm$ is $q\lambda-1+1-q\lambda=0$, the terms $\pm\lambda^2L_S$ cancel, at order $\lambda^3$ there remains $P+Q'=F_q$, and everything else is divisible by $\lambda^4$. The highest power of $\lambda$ is $\lambda^9$, from the cubic part of $\Bp$ (each $u$ has degree $3$), so $G_q$ has degree at most $5$.
\end{proof}

$\Fmrg$ keeps the terms $A_i$ of $\Fstd$, moves the $\lambda^3$-parts of $C$ into the term $q+1$ ($\Bp$ is $B$ with $-\lambda^3x_{q+1}$ added in each factor), and leaves $\Oterm=x_0y_0z_0$ as the term $q+2$.

\begin{lemma}[supports of the merged decomposition]\label{cp:L2}
In $\Fmrg$ the $z$-forms of the terms $\lambda A_i$, $-\Bp$, $(1-q\lambda)\Oterm$ have the supports $\{z_0,z_i\}$, all $q+2$ variables $\{z_0,\dots,z_{q+1}\}$, and $\{z_0\}$; the same holds for the $x$- and $y$-forms. Consequently:
\begin{enumerate}[label=(\alph*)]
\item the coverage of $z_0$ is $q+2$, of $z_i$ ($1\le i\le q$) is $2$ (the terms $A_i$ and $\Bp$), and of $z_{q+1}$ is $1$ (the term $\Bp$); in $\Fstd$ the coverages are the same output by output ($z_i$ in $A_i$ and $B$, $z_{q+1}$ in $C$ only);
\item per coordinate the product degree is at most $9$: it is $4$ for $\lambda A_i$, $9$ for $-\Bp$ and $1$ for $(1-q\lambda)\Oterm$;
\item $z_0\mapsto\Oterm$, $z_i\mapsto A_i$, $z_{q+1}\mapsto\Bp$ is a system of distinct representatives of the outputs in their covering sets.
\end{enumerate}
\end{lemma}

\begin{proof}
The $z$-forms are $\lambda(z_0+\lambda z_i)$, $-(z_0+\lambda^2S_z-\lambda^3z_{q+1})$ and $(1-q\lambda)z_0$, and the coefficients $\lambda$, $\lambda^2$, $-\lambda^3$, $1$ are nonzero polynomials; this gives the supports and (a). For $\Fstd$ the $z$-forms are $\lambda(z_0+\lambda z_i)$, $-(z_0+\lambda^2S_z)$ and $(1-q\lambda)(z_0+\lambda^3z_{q+1})$. (b) The scalar coefficient and the three forms have $\lambda$-degrees $1+1+1+1$, $0+3+3+3$ and $1+0+0+0$. (c) Each representative covers its output, and the map is a bijection between the $q+2$ outputs and the $q+2$ terms.
\end{proof}

The supports of $\Bp$ and $\Oterm$ have sizes $q+2$ and $1$, against $q+1$ for $B$ and $2$ for $C$; the sums of the support sizes agree, as they must, since they equal the sums of the coverages. By the product rule, the coverage of every $z$-string under a tensor power is the same for both decompositions, and so is everything in the laser analysis that depends on the decomposition only through coverage. What changes is how the supports are distributed over the term strings, and this is what the profile sees (Section~\ref{cp:sec:moment}).

\subsection{Level-1 points and the laser count}\label{cp:sec:laser1}

\begin{definition}[level-1 points]\label{cp:L3}
Let $\Qs=\{6,7\}$. A \emph{level-1 point} is a rational $P=(\lambda_6,u_6,u_7)$ with $\lambda_6\in(0,1)$, $\lambda_7=1-\lambda_6$ and $0<u_q<w_q:=1/(q+2)$; we write points in the coordinates $(\lambda_6,\,u_6(6+2),\,u_7(7+2))$. Put $p_q=q/(2(q+2))$ and $v_q=(w_q-u_q)/2$. The \emph{joint type} of part $q$ is the distribution on the six block triples with
\[
p_{011}=p_{101}=p_q,\quad p_{002}=w_q,\quad p_{110}=u_q,\quad p_{200}=p_{020}=v_q;
\]
its $X$-marginal (the distribution of the $x$-block) is $(p_q+w_q+v_q,\,p_q+u_q,\,v_q)$, equal to its $Y$-marginal, and its $Z$-marginal is $(1,q,1)/(q+2)$. Let $\Ent^q_X$, $\Ent^q_Z$ and $\Ent^q_J$ be the Shannon entropies (natural logarithm) of the $X$-marginal, the $Z$-marginal and the joint type of part $q$, and
\[
\alpha(P)=\frac{\sum_q\lambda_qu_q\ln q}{\sum_q\lambda_qp_q\ln q},\qquad\slack(P)=\sum_q\lambda_q\bigl(\Ent^q_X-\Ent^q_Z\bigr).
\]
$P$ is \emph{admissible} if $\slack(P)>0$. Let $d(P)$ be the least positive integer such that for every multiple $M$ of $d(P)$ the numbers $M_q:=\lambda_qM$ and all per-part block counts ($M_q$ times the probabilities of the joint type and of its marginals) are integers. For such $M$ let $F^{(M)}=\bigotimes_{q\in\Qs}F_q^{\otimes M_q}$; its $M$ coordinates fall into two \emph{parts}. The \emph{kept output type} $\Ukept$ is the set of $z$-strings that have, on the coordinates of part $q$, exactly $M_q/(q+2)$ letters $z_0$, $qM_q/(q+2)$ letters in $\{z_1,\dots,z_q\}$ and $M_q/(q+2)$ letters $z_{q+1}$.
\end{definition}

The slack is a $\lambda$-weighted sum over the parts, and the laser count below uses only this sum: it is applied once to the whole mixture. A part may have $\Ent^q_X<\Ent^q_Z$; at the point $\Pzero$ of Lemma~\ref{cp:L11} the part $q=6$ has $\Ent^6_X\le\cpv{HXsixP0_hi}<\cpv{HZsixP0_lo}\le\Ent^6_Z$.

Let $T_X$ and $T_Z$ be the numbers of $x$-block strings and of $z$-block strings with the per-part marginal types of $P$ (equal for $x$ and $y$), and let $a_M=\prod_qq^{p_qM_q}$ and $b_M=\prod_qq^{u_qM_q}$.

\begin{lemma}[level-1 laser count]\label{cp:L4}
Let $P$ be a level-1 point with $\slack(P)\ge0$ and $M$ a multiple of $d(P)$. There is a zero-out of $F^{(M)}$ (setting variables to zero only) whose image is the direct sum of $s_M$ copies of $\mat{a_M}{b_M}{a_M}$ and which keeps only the variables of these copies, such that:
\begin{enumerate}[label=(\roman*)]
\item every copy is the block tensor $T_{IJK}$ of one block-triple string $(I,J,K)$ of the joint type of $P$; its outputs are all $a_M^2$ $z$-variables of the $z$-block $K$; and the kept outputs are $s_Ma_M^2$ distinct $z$-strings of the kept type $\Ukept$;
\item for every $\eps>0$ there is $C_\eps$ such that for all large $M$
\[
\ln s_M\ge M\sum_q\lambda_q\Ent^q_Z-2(1+\eps)\abs\Qs\ln(M+1)-\eps M\sum_q\lambda_q\bigl(\Ent^q_J-\Ent^q_Z\bigr)-C_\eps;
\]
in particular, for every $\eta>0$, $\ln s_M\ge M(\sum_q\lambda_q\Ent^q_Z-\eta)$ for all large $M$; and $s_M\le T_Z$;
\item every border decomposition of $F^{(M)}$ restricts, term by term, to a border decomposition of $s_M\cp\mat{a_M}{b_M}{a_M}$ in which every kept output is covered by a term of nonzero restricted $z$-form, and which has at least $s_Ma_M^2$ such terms.
\end{enumerate}
\end{lemma}

\begin{proof}
\emph{Blocks.} An $x$-variable $x_i$, $i\in\prod_\ell\{0,\dots,q+1\}$, lies in the $x$-block string $I$ with $I_\ell=0$, $1$ or $2$ according as $i_\ell=0$, $i_\ell\in[q]$ or $i_\ell=q+1$, and likewise for $y$ and $z$. Every monomial of $F_q$ has block sum~$2$, so $F^{(M)}=\sum_{(I,J,K)\in\Lambda}T_{IJK}$ in the notation of Theorem~\ref{cp:P6:LG} with $\tau=2$ and strings of length $M$, and $T_{IJK}=\bigotimes_\ell T_{I_\ell J_\ell K_\ell}$ with the constituents of Section~\ref{cp:sec:laser}.
\emph{Step 0 (first zero-out).} Zero every $x$-variable whose block string does not have, in every part $q$, exactly $M_q(p_q+w_q+v_q)$, $M_q(p_q+u_q)$ and $M_qv_q$ blocks $0$, $1$, $2$; the same for $y$; and zero every $z$-variable whose block string does not have the per-part $Z$-marginal $(1,q,1)M_q/(q+2)$. What remains is $\sum_{(I,J,K)\in U}T_{IJK}$, where $U$ is the set of triples of $\Lambda$ whose three block strings have these marginal types.
\emph{Step 1 ($U^*=U$).} At a coordinate the triple of blocks is one of $011$, $101$, $110$, $002$, $020$, $200$, and per part the six counts satisfy $n_{200}=\#\{I_\ell=2\}$, $n_{020}=\#\{J_\ell=2\}$, $n_{002}=\#\{K_\ell=2\}$, $n_{101}+n_{110}=\#\{I_\ell=1\}$, $n_{011}+n_{110}=\#\{J_\ell=1\}$ and $n_{011}+n_{101}=\#\{K_\ell=1\}$; this linear system determines all six counts. So every triple of $U$ has, in every part, the joint type of $P$, and in the notation of Theorem~\ref{cp:P6:LG} we may take $U^*=U$ and $N_i^*=N_i$. The regularity hypothesis of Theorem~\ref{cp:P6:LG} holds: the permutations of the coordinates inside the parts map $U$ onto $U$ and act transitively on the $x$-block strings of the right type, so $\abs{\{u\in U:f_1(u)=I\}}$ is the same number $N_1$ for all of them and $0$ for the others; likewise $N_2$ ($=N_1$, since the $X$- and $Y$-marginals agree) and $N_3$. With $T_1=T_2=T_X$ and $T_3=T_Z$, $\abs U=T_XN_1=T_ZN_3$. These counts are products over the parts of multinomial coefficients in numbers proportional to $M$, each count with at least one factor that has two positive arguments, so they are strictly increasing along the multiples of $d(P)$, as the setting of Theorem~\ref{cp:P6:LG} assumes.
\emph{Step 2 (Theorem~\ref{cp:P6:LG}).} It gives removal operations after which the surviving vertices form an edgeless graph with $s_M\ge c_\eps\abs U(N_1+N_2+N_3)^{-1-\eps}$ vertices, all in $U$. On the tensor: a $T_{IJK}$ with $(I,J,K)\in\Lambda$ survives if and only if its three blocks survive; it is then in $U$ by Step~0, hence a vertex of the final graph; no two surviving vertices share a block, so the surviving tensor is the direct sum of their $T_{IJK}$. Each has the joint type of $P$ and is therefore isomorphic to $\mat{\prod_qq^{n_{101}}}{\prod_qq^{n_{110}}}{\prod_qq^{n_{011}}}=\mat{a_M}{b_M}{a_M}$, since $p_{101}=p_{011}=p_q$ and $p_{110}=u_q$. The outputs of $T_{IJK}$ are the $z$-variables of block $K$, and there are $\prod_qq^{qM_q/(q+2)}=a_M^2$ of them ($2p_q=q/(q+2)$); their letter kinds are given by the blocks of $K$, so they lie in $\Ukept$, and distinct vertices have distinct $K$. This is (i), and it gives $s_M\le T_Z$. Finally zero every variable whose block lies in no surviving vertex. No monomial of the surviving tensor contains such a variable, so the image does not change, and every variable left is a variable of one copy; in particular every restricted $z$-form is a form in the $s_Ma_M^2$ kept outputs, as Definition~\ref{cp:T1} and Lemma~\ref{cp:T1:basic} require.
\emph{Step 3 (count).} Put $T_{\min}=\min(T_X,T_Z)$. Then $N_1+N_2+N_3=\abs U(2/T_X+1/T_Z)\le3\abs U/T_{\min}$, so $s_M\ge c_\eps3^{-1-\eps}T_{\min}(\abs U/T_{\min})^{-\eps}$. For a type class of length $n$ over three letters, $(n+1)^{-2}e^{n\Ent}\le\text{multinomial}\le e^{n\Ent}$ (the lower bound because the type class has the largest probability among the at most $(n+1)^2$ type classes under its own distribution). Hence $\ln T_X\ge M\sum_q\lambda_q\Ent^q_X-2\abs\Qs\ln(M+1)$, $\ln T_Z\ge M\sum_q\lambda_q\Ent^q_Z-2\abs\Qs\ln(M+1)$ and $\ln\abs U\le M\sum_q\lambda_q\Ent^q_J$. With $\slack(P)=\sum_q\lambda_q(\Ent^q_X-\Ent^q_Z)\ge0$ this gives $\ln T_{\min}\ge M\sum_q\lambda_q\Ent^q_Z-2\abs\Qs\ln(M+1)$ and $\ln(\abs U/T_{\min})\le M\sum_q\lambda_q(\Ent^q_J-\Ent^q_Z)+2\abs\Qs\ln(M+1)$, hence (ii) with $C_\eps=(1+\eps)\ln3-\ln c_\eps$. For the $\eta$-form choose $\eps$ with $\eps\sum_q\lambda_q(\Ent^q_J-\Ent^q_Z)<\eta/2$ and then $M$ large.
\emph{(iii).} A restriction maps the border identity $\sum_\tau u_\tau\otimes v_\tau\otimes w_\tau=\lambda^{h}F^{(M)}+\lambda^{h+1}(\cdots)$ to the same identity for the restricted forms, by linearity, and its $\lambda^h$-part is the direct sum. The slice of a kept output $o$ of a copy at order $\lambda^h$ is the bilinear form $(XY)_o\ne0$ (as $b_M\ge1$), so some restricted term has a $z$-form with nonzero coefficient at $o$. The count of terms is Lemma~\ref{cp:T1:basic}(a).
\end{proof}

\phantomsection\label{cp:rem:L4parts}%
Lemma~\ref{cp:L4} and its proof hold verbatim, with Definition~\ref{cp:L3} read for the part set, for any finite set $\Qs$ of integers $q\ge2$ with rational weights $\lambda_q>0$ summing to~$1$, in particular for $\Qs=\{q\}$ with $\lambda_q=1$, and for $0<u_q\le w_q$: the proof uses neither $u_q<w_q$ nor two parts.

\begin{remark}[on the citation]\label{cp:rem:LGcite}
Since $\slack(P)>0$ makes $T_X>T_Z$ for large $M$, the counts satisfy $N_1=N_2\ne N_3$, which the theorem of Coppersmith and Winograd~\cite{CW90}, the case $N_1=N_2=N_3$~\cite[Theorem~4.1]{LG12}, does not cover; Lemma~\ref{cp:L4} rests on~\cite[Theorem~4.2]{LG12}, and~\cite{CW90} is cited for the method. Coppersmith's level-1 analysis~\cite{Cop97} covers the same points; we cite it for the mixture and its value $\alpha>\cpv{alpha_cop97}$, not as an input of a proof. The constant $c_\eps$ of Theorem~\ref{cp:P6:LG} enters only the choice of $M$, never an exponent of a running time.
\end{remark}

\subsection{The input family}\label{cp:sec:family}

Fix an admissible point $P$. For every multiple $M$ of $d(P)$ with $b_M\ge2$ let $A_M$ be the restriction, by the zero-out of Lemma~\ref{cp:L4}, of the product decomposition of $F^{(M)}$ that carries $\Fmrg$ (with the $q$ of its part) on every coordinate, with the terms whose restricted $x$-, $y$- or $z$-form vanishes dropped. Its terms are term strings $\tau\in\prod_q\{A_1,\dots,A_q,\Bp,\Oterm\}^{M_q}$. The \emph{composition} of $\tau$ is the tuple $(n^q_A,n^q_{\Bp},n^q_{\Oterm})_{q\in\Qs}$ of the numbers of letters $A_\cdot$ (of any index), $\Bp$ and $\Oterm$ in part $q$, and the partition of the terms of $A_M$ is the partition by composition. Let $R_F=\prod_q(q+2)^{M_q}$ be the number of all term strings and $\ln a_1=\sum_q\lambda_qp_q\ln q$, so that $\ln a_M=M\ln a_1$.

\begin{lemma}[the input family]\label{cp:L5}
$(A_M)$ is an input family of decompositions of $s_M\cp\mat{a_M}{b_M}{a_M}$ that satisfies \Hcond1--\Hcond4:
\begin{itemize}
\item[\Hcond1] $\ln a_M=M\ln a_1$ with $\ln a_1>0$, $\alpha_M=\alpha(P)$ for every $M$, and $b_M\ge2$ for every such $M$, which holds for all large $M$;
\item[\Hcond2] $0\le\rho_M-2\le\delta_M:=\bigl(M\sum_q\lambda_q\Ent^q_Z-\ln s_M\bigr)/(M\ln a_1)\to0$;
\item[\Hcond3] the product degree is $e_M\le9M$, and the number of classes is $m_M\le(M+1)^{2\abs\Qs}=(M+1)^4$;
\item[\Hcond4] $C_M\le\prod_q\bigl((q+2)^{1/(q+2)}2^{q/(q+2)}\bigr)^{M_q}$, so $\cstar\le\cbar\le\co(P):=\sum_q\lambda_qc(q)/\ln a_1$ with $c(q)=(\ln(q+2)+q\ln2)/(q+2)$.
\end{itemize}
\end{lemma}

\begin{proof}
\Hcond1: $\ln a_M=\sum_qp_qM_q\ln q=M\ln a_1$ and $\ln b_M=M\sum_q\lambda_qu_q\ln q>0$, so $\alpha_M=\alpha(P)$, and $b_M\ge2$ once $M\sum_q\lambda_qu_q\ln q\ge\ln2$.
\Hcond2: the terms of $A_M$ are among all strings, so $r_M\le R_F$; and $r_M\ge s_Ma_M^2$ by Lemma~\ref{cp:L4}(iii). Since $\Ent\bigl((1,q,1)/(q+2)\bigr)+\frac q{q+2}\ln q=\ln(q+2)$ and $2\ln a_M=M\sum_q\lambda_q\frac q{q+2}\ln q$, we have $\ln R_F-2\ln a_M=M\sum_q\lambda_q\Ent^q_Z$ exactly. Hence $0\le\rho_M-2=\ln(r_M/(s_Ma_M^2))/\ln a_M\le\ln(R_F/(s_Ma_M^2))/\ln a_M=\delta_M$, and $\delta_M\le\bigl[2(1+\eps)\abs\Qs\ln(M+1)+\eps M\sum_q\lambda_q(\Ent^q_J-\Ent^q_Z)+C_\eps\bigr]/(M\ln a_1)$ by Lemma~\ref{cp:L4}(ii); letting $M\to\infty$ and then $\eps\to0$ gives $\delta_M\to0$.
\Hcond3: by Lemma~\ref{cp:L2}(b) every term string has product degree at most $9$ per coordinate, so $e_M\le9M$; the approximation degree is $3M$. The compositions of part $q$ are triples of nonnegative integers with sum $M_q$, at most $(M_q+1)^2$ of them, so $m_M\le\prod_q(M_q+1)^2\le(M+1)^4$. Both are $e^{O(\ln M)}$, against $\ln a_M=\Theta(M)$.
\Hcond4: by Lemma~\ref{cp:L2}(a) and the product rule, the coverage of an output $o\in\Ukept$ under the product decomposition is $\prod_\ell$ of $q+2$, $2$ or $1$ according as $o_\ell=z_0$, $o_\ell\in\{z_1,\dots,z_q\}$ or $o_\ell=z_{q+1}$, which is the bound stated, the same for all $o\in\Ukept$. The kept outputs lie in $\Ukept$ (Lemma~\ref{cp:L4}(i)), and restriction can only lower a coverage. Dividing by $\ln a_M=M\ln a_1$ gives the bound on $\cbar$ and $\cstar$.
\end{proof}

Only the upper bound $\cstar\le\co(P)$ is used below; restriction may kill covering terms, so equality is not claimed.

\subsection{The profile: symmetry and the moment bound}\label{cp:sec:moment}

The profile of $A_M$ is bounded through the supports of all term strings of the unrestricted product, by the following lemma for restrictions of products.

\begin{lemma}[restrictions of products]\label{cp:L6:lemma5}
Let $F$ be the product of border decompositions $F^{(1)},\dots,F^{(M)}$ over $\Q[\lambda]$ of base tensors, with term strings $\tau=(\beta_1,\dots,\beta_M)$, $R_F$ of them, and $z$-strings $o$, where $o\in\supp w_\tau$ if and only if $o_i\in\supp w^{(i)}_{\beta_i}$ for all $i$. Let $A$ be obtained from $F$ by a zero-out whose image is $s\cp\mat aba$, with its kept outputs in a set $Z^*$ of $z$-strings. Let $\Pi$ be a partition of all term strings such that for every class $c$ the number $\kappa^*_c$ of strings in $c$ whose support contains $o$ is the same for all $o\in Z^*$, and give each term of $A$ its class. Put $S^*(\tau)=\abs{\supp w_\tau\cap Z^*}$, let $\bar S_c$ be the mean of $S^*$ over $c$ and $\bar S(\tau)=\bar S_{c(\tau)}$. Then for every $t\in[0,1]$
\[
s^{t-1}Z_A(t)\le s^{-1}a^{-2t}R_F\,\mathrm{mean}_\tau\bar S(\tau)^t,\qquad\text{so}\qquad\mu(t)\le\frac{\ln(R_F/s)}{\ln a}-2+\frac{\ln\bigl(\mathrm{mean}_\tau\bar S(\tau)^t\bigr)}{\ln a},
\]
with $\mu(t)=\ln(s^{t-1}Z_A(t))/\ln a-2+2t$. If $S^*$ is constant on every class, $\bar S=S^*$.
\end{lemma}

\begin{proof}
A zero-out can only remove terms and remove outputs from supports, so the class $c$ of $A$ has at most $\abs{\Pi_c}$ terms and $\kappa_c\le\kappa^*_c$. Counting the pairs (string in $c$, element of $Z^*$ in its support) in two ways gives $\kappa^*_c\abs{Z^*}=\abs{\Pi_c}\bar S_c$. Hence $\abs{P_c}^{1-t}\kappa_c^t\le\abs{\Pi_c}^{1-t}(\kappa^*_c)^t=\abs{\Pi_c}\bar S_c^t\abs{Z^*}^{-t}$. The kept outputs are $sa^2$ distinct elements of $Z^*$, so $sa^2\le\abs{Z^*}$, and summing over $c$, $s^{t-1}Z_A(t)\le s^{t-1}(sa^2)^{-t}\sum_c\abs{\Pi_c}\bar S_c^t=s^{-1}a^{-2t}R_F\,\mathrm{mean}_\tau\bar S(\tau)^t$. Taking logarithms and dividing by $\ln a$ gives the bound on $\mu$.
\end{proof}

The constancy of $S^*$ on classes matters: by concavity of $y\mapsto y^t$, $\abs{\Pi_c}\bar S_c^t\ge\sum_{\tau\in c}S^*(\tau)^t$, which is the wrong direction for replacing $\bar S$ by $S^*$ without it. For the level-1 family both hypotheses hold by a symmetry group.

\begin{lemma}[the symmetry group]\label{cp:L6:lemmaC}
Let $\Gamma$ be the group generated by the permutations of the coordinates inside each part and, independently for every coordinate $\ell$, the permutations $\sigma_\ell$ of $\{1,\dots,q\}$ ($q$ of the part of $\ell$), acting on term letters by $A_i\mapsto A_{\sigma_\ell(i)}$, $\Bp\mapsto\Bp$, $\Oterm\mapsto\Oterm$ and on $z$-letters by $z_0\mapsto z_0$, $z_i\mapsto z_{\sigma_\ell(i)}$, $z_{q+1}\mapsto z_{q+1}$. Then
\begin{enumerate}[label=(\alph*)]
\item $\Gamma$ preserves the relation ``$o\in\supp w_\tau$'' for the product of the decompositions $\Fmrg$;
\item $\Gamma$ maps $\Ukept$ onto itself and is transitive on it;
\item the orbits of $\Gamma$ on term strings are exactly the composition classes.
\end{enumerate}
Hence for the composition partition the number $\kappa^*_c$ of class-$c$ strings covering $o$ is the same for all $o\in\Ukept$, and $S^*(\tau)=\abs{\supp w_\tau\cap\Ukept}$ is constant on every class.
\end{lemma}

\begin{proof}
(a) Per coordinate the supports of $A_i$, $\Bp$, $\Oterm$ are $\{z_0,z_i\}$, all letters, and $\{z_0\}$ (Lemma~\ref{cp:L2}), and $\sigma_\ell$ maps $\{z_0,z_i\}$ to $\{z_0,z_{\sigma_\ell(i)}\}$; coordinate permutations inside a part permute the factors of the product. (b) $\Ukept$ is defined by the numbers of letters of each kind per part, which $\Gamma$ preserves. Given $o,o'\in\Ukept$, permute the coordinates inside each part to move the letters $z_0$, the letters in $\{z_1,\dots,z_q\}$ and the letters $z_{q+1}$ of $o$ to the positions of those of $o'$, then permute the letters $1,\dots,q$ coordinate by coordinate. (c) The composition is invariant, and two strings with the same composition are mapped onto each other by a coordinate permutation inside the parts followed by per-coordinate letter permutations ($A_i\mapsto A_j$ at each $A$-coordinate).
For the consequence: an element $g\in\Gamma$ maps $\Cov(o)\cap\Pi_c$ bijectively onto $\Cov(go)\cap\Pi_c$ by (a) and (c), and $\Gamma$ is transitive on $\Ukept$ by (b); and $g$ maps $\supp w_\tau\cap\Ukept$ bijectively onto $\supp w_{g\tau}\cap\Ukept$, since $\Ukept$ is $\Gamma$-invariant, while every class is a single orbit.
\end{proof}

A group with one permutation of $\{1,\dots,q\}$ per part, applied to all coordinates of the part at once, is not transitive on $\Ukept$ (it preserves the multiset of letters at the positions of the letters in $\{z_1,\dots,z_q\}$); the per-coordinate group is the one that works.

\begin{corollary}[the profile of the family]\label{cp:L6}
For every $t\in[0,1]$ and every multiple $M$ of $d(P)$,
\[
\mu_M(t)\le\delta_M+\frac{\ln\bigl(\mathrm{mean}_\tau\Sall(\tau)^t\bigr)}{\ln a_M},\qquad\Sall(\tau)=\abs{\supp w_\tau\cap\Ukept},
\]
the mean over all $R_F$ term strings $\tau$ (alive or not, including those with letters $\Oterm$), with $0^t=0$ for $t>0$ and $0^0=1$.
\end{corollary}

\begin{proof}
Lemma~\ref{cp:L6:lemma5} with $Z^*=\Ukept$ applies by Lemma~\ref{cp:L6:lemmaC}, with $\bar S=S^*=\Sall$, and $\ln(R_F/s_M)/\ln a_M-2=\ln(R_F/(s_Ma_M^2))/\ln a_M=\delta_M$ (proof of Lemma~\ref{cp:L5}, \Hcond2).
\end{proof}

\begin{theorem}[moment bound]\label{cp:L7}
For every $t\in(0,1]$ and every multiple $M$ of $d(P)$,
\[
\frac1M\ln\mathbb E_\tau\,\Sall(\tau)^t\le\sum_q\lambda_q\,t^2\ln^2\frac{4eq}t,
\]
where $\tau$ is uniform over all $R_F$ term strings.
\end{theorem}

\begin{proof}
Fix $t$ and give the $z$-letters of a coordinate of part $q$ the weights $w(z_0)=1$, $w(z_i)=\beta:=1/t$ for $i\in[q]$, and $w(z_{q+1})=\zeta:=(t+q)/t^2$ (depending on $q$). Every $o\in\Ukept$ has the same weight $W=\prod_\ell w(o_\ell)=\prod_q(\beta^q\zeta)^{M_q/(q+2)}$, so, since $\supp w_\tau=\prod_\ell\supp w_{\tau_\ell}$,
\[
\Sall(\tau)\le\sum_{o\in\supp w_\tau}\frac{\prod_\ell w(o_\ell)}W=\frac{\prod_\ell f(\tau_\ell)}W,\qquad f(A_i)=1+\beta,\ f(\Bp)=1+q\beta+\zeta,\ f(\Oterm)=1.
\]
For uniform $\tau$ the letters are independent and uniform, so $\mathbb E\,\Sall^t\le W^{-t}\prod_\ell\mathbb Ef(\tau_\ell)^t$. Per coordinate, $Y=\ln f(\tau_\ell)$ lies in $[0,R]$ with $R=\ln(1+q\beta+\zeta)$, and Hoeffding's lemma~\cite{Hoe63} gives $\ln\mathbb Ee^{tY}\le t\,\mathbb EY+t^2R^2/8$. With the chosen weights $1+\beta=\beta(1+t)$ and $1+q\beta+\zeta=\zeta(1+t)$, so $t\,\mathbb EY\le t(q\ln\beta+\ln\zeta)/(q+2)+t\ln(1+t)\le t(q\ln\beta+\ln\zeta)/(q+2)+t^2$. Summed over the coordinates, the first part is $t\ln W$ and cancels against $W^{-t}$. And $R=\ln\bigl((t+q)(1+t)/t^2\bigr)\le\ln(4q/t^2)\le2\ln(4eq/t)$, so the rest per coordinate is at most $t^2\bigl(1+\ln^2(4eq/t)/2\bigr)\le t^2\ln^2(4eq/t)$, because $\ln(4eq/t)\ge\ln(4e)>\sqrt2$. Summing over the $M_q=\lambda_qM$ coordinates of each part and dividing by $M$ gives the claim.
\end{proof}

\begin{lemma}[$\cg=0$]\label{cp:L8}
Put $\Lampr(t)=\sum_q\lambda_qt^2\ln^2(4eq/t)/\ln a_1$. For the family $(A_M)$ of an admissible point:
\begin{enumerate}[label=(\roman*)]
\item $0\le\mu_M(t)\le\delta_M+\Lampr(t)$ for every $t\in(0,1]$ and every $M$;
\item $\Lam(0)=\lim_M\mu_M(0)=0$;
\item $0\le\Lam(t)\le\Lampr(t)$ for every $t\in(0,1]$;
\item $\lim_{t\to0+}\Lam(t)/t=0$; in particular $\cg=0$.
\end{enumerate}
\end{lemma}

\begin{proof}
(i) Corollary~\ref{cp:L6} and Theorem~\ref{cp:L7}, with $\ln a_M=M\ln a_1$: $\ln(\mathrm{mean}_\tau\Sall^t)/\ln a_M=[(1/M)\ln\mathbb E\Sall^t]/\ln a_1\le\Lampr(t)$; the lower bound is Lemma~\ref{cp:T1:basic}(b). (ii) $Z_M(0)=r_M$, so $\mu_M(0)=\rho_M-2\in[0,\delta_M]$, and $\delta_M\to0$. (iii) Take the limsup over $M$ in (i) at fixed $t$; the $M$-dependent term $\delta_M$ does not depend on $t$ and tends to $0$. (iv) $\Lam(t)/t\le t\sum_q\lambda_q\ln^2(4eq/t)/\ln a_1\to0$, and $\Lam\ge0$.
\end{proof}

\ThmDk{} uses $\cg$ only through the existence of arbitrarily small $t$ with $\Lam(t)\le ct$ for every $c>0$, which (iv) gives. An infimum over $t$ of $\limsup_M\mu_M(t)/t$ would not suffice in general, since it can be attained away from $0$; the explicit $o(t)$ bound is what closes this step.

\subsection{Rational points}\label{cp:sec:rational}

\begin{lemma}[rational admissible points]\label{cp:L9}
(a) For $q\ge4$ and $0<u_q<w_q$, $\partial\slack/\partial u_q<0$. (b) Every real point $P=(\lambda_6,u_6,u_7)$ with $\slack(P)\ge0$, $\lambda_6\in(0,1)$ and $0<u_q<w_q$ is a limit of rational admissible points. Hence the supremum of $\alpha(P)$ over the rational admissible points equals its supremum over all points with $\slack(P)\ge0$, $\lambda_6\in(0,1)$ and $0<u_q<w_q$.
\end{lemma}

\begin{proof}
(a) $\Ent^q_Z$ does not depend on $u_q$. For the $X$-marginal $(x_0,x_1,x_2)=(p+w+v,\,p+u,\,v)$ of part $q$ (indices $q$ dropped), $dx/du=(-\frac12,1,-\frac12)$, and the constant parts of $-\sum_i(\ln x_i+1)\,dx_i$ cancel, so $d\Ent^q_X/du=\frac12\ln(x_0x_2)-\ln x_1=\ln\bigl(\sqrt{x_0x_2}/x_1\bigr)$. Now $x_0x_2\le v<w/2=1/(2(q+2))$ and $x_1^2\ge p^2=q^2/(4(q+2)^2)$, so $x_0x_2/x_1^2\le2(q+2)/q^2<1$ for $q\ge4$; and $\partial\slack/\partial u_q=\lambda_q\,d\Ent^q_X/du$.
(b) Decreasing every $u_q$ by a small amount keeps $0<u_q<w_q$ and makes $\slack>0$ by (a); $\slack$ and $\alpha$ are continuous, so nearby rational points keep $\slack>0$ and $\alpha$ within any margin.
\end{proof}

The value of this supremum is computed, not certified, in Remark~\ref{cp:A4}.

\subsection{The standard family: coverage and column union}\label{cp:sec:std}

For an admissible point $P$ put
\[
\ccol(P)=\frac{\sum_q\lambda_q\bigl(p_q\ln(2(q+1))+\ln(q+2)/(q+2)\bigr)}{\ln a_1}.
\]

\begin{lemma}[standard level-1 tops]\label{cp:L10}
For every admissible rational point $P$ the hypotheses $\TPE(\alpha(P),\co(P))$ of Proposition~\ref{cp:T8} and $\Tcol(\alpha(P),\ccol(P)-1)$ of Proposition~\ref{cp:T9} hold. In the restriction of the standard decomposition by the zero-out of Lemma~\ref{cp:L4}, every kept output has coverage at most $\prod_q\bigl((q+2)^{1/(q+2)}2^{q/(q+2)}\bigr)^{M_q}$, and every column of every copy has a column union (the terms covering some output of the column) of at most $\prod_q\bigl((q+1)^{p_q}2^{p_q}(q+2)^{1/(q+2)}\bigr)^{M_q}$ terms; both bounds are the same for all kept outputs and all columns.
\end{lemma}

\begin{proof}
\emph{The family.} Let $A^{\mathrm{std}}_M$ be the restriction, by the zero-out of Lemma~\ref{cp:L4} (which depends on $F^{(M)}$ only), of the product decomposition carrying $\Fstd$ on every coordinate, with the trivial partition ($m=1$). Its product degree is at most $10M$ (Section~\ref{cp:sec:laser}: $4$, $6$, $10$ per coordinate), so \Hcond3{} holds; \Hcond1{} and \Hcond2{} hold as in Lemma~\ref{cp:L5}, whose proof uses only $R_F$, $s_M$, $a_M$, $b_M$ and Lemma~\ref{cp:L4}(iii).
\emph{Coverage.} In $\Fstd$, $z_0$ lies in all $q+2$ terms, $z_i$ in $A_i$ and $B$, and $z_{q+1}$ in $C$ only (Lemma~\ref{cp:L2}(a)); so the coverage of every $o\in\Ukept$ under the product decomposition is the bound stated, and restriction only lowers it. Its logarithm divided by $\ln a_M$ is $\co(P)$.
\emph{Column union.} A copy is $T_{IJK}$ with the joint type of $P$; its outputs are the $z$-strings of block $K$. The row of an output is given by its letters at the coordinates of type $101$, its column by its letters at the coordinates of type $011$, and at the other coordinates the $z$-letter is fixed: $z_0$ at the types $110$, $200$, $020$ and $z_{q+1}$ at $002$. A term string covers some output of a fixed column if and only if, coordinatewise, its letter covers some $z_i$, $i\in[q]$, at the $101$-coordinates ($q+1$ choices: $A_1,\dots,A_q,B$), covers the fixed letter at the $011$-coordinates ($2$ choices), covers $z_0$ at the $110$-, $200$- and $020$-coordinates (all $q+2$ terms), and covers $z_{q+1}$ at the $002$-coordinates ($C$ only). Since the conditions are coordinatewise, the column union is a product set of size $\prod_q(q+1)^{n_{101}}2^{n_{011}}(q+2)^{n_0}$ with $n_{101}=n_{011}=p_qM_q$ and $n_0=n_{110}+n_{200}+n_{020}=M_q/(q+2)$ per part; zero-outs only remove terms, so this bounds the restricted family for every column of every copy. Its logarithm divided by $\ln a_M$ is $\ccol(P)$.
\emph{From the family to one top algorithm.} Apply Lemma~\ref{cp:T5:lemma4} to $A^{\mathrm{std}}_M$ with $m=1$. At $t=0$ it gives $\delB\to0$ as $M$ and then $j$ grow, by \Hcond2{} and \Hcond3. At $t=1$, $\ZB(1)=\cB$ and $Z_A(1)=C_A$, so $\mu(1)=c$ and $\ln\cB/\ln k\le c+\ln(2(e+1))/\ln a+(1/j)\ln((e+1)r)/\ln a$, which is at most $\co(P)+\eta$ for suitable $(M,j)$. For column unions: in Step~1 of Lemma~\ref{cp:T5:lemma4} each term $(p,l)$ has its support inside that of $l$, so a column union grows by at most the factor $e+1$; in Step~2 a term $(\gamma,g)$ that covers some output of the column $(\sigma,\nu_1,\nu')$ of $s\cp T_0^{\otimes(j+1)}$ has a slot $\sigma'$ such that $m(\gamma,\sigma')$ covers an output of the column $(\sigma,\nu_1)$ of $A'$ and $g$ an output of the column $(\sigma',\nu')$ of $G_j$, and $(\gamma,g)\mapsto(m(\gamma,\sigma'),g)$, with the least such $\sigma'$, is injective, so the largest column union of $G_{j+1}$ is at most the product of those of $A'$ and of $G_j$; Step~4 only lowers it. Hence the column unions of $B$ are at most $\bigl((e+1)U_A\bigr)^j$ with $U_A$ the bound above, and $\ln\abs{U_B}/\ln k\le\ccol(P)+\ln(e+1)/\ln a_M\le\ccol(P)+\eta$. With $\alB=\alpha(P)$ this is $\TPE(\alpha(P),\co(P))$ and $\Tcol(\alpha(P),\ccol(P)-1)$, as $k^{1+(\ccol-1)+\eta}=k^{\ccol+\eta}$.
\end{proof}

Since every kept output has the same coverage bound and every column the same union bound, the maximal coverage equals the mean up to the restriction, and Propositions~\ref{cp:T8} and~\ref{cp:T9} need no relabeling of the outputs; this is why all composites of Section~\ref{cp:sec:composites} are deterministic.

\subsection{The point \texorpdfstring{$P_0$}{P0}}\label{cp:sec:P0}

\begin{lemma}[the constants at $\Pzero$]\label{cp:L11}
Let $\Pzero=(\cpv{P0_lam},\,\cpv{P0_u6},\,\cpv{P0_u7})$ in the coordinates $(\lambda_6,u_6(q+2),u_7(q+2))$. Then
\[
\cpv{slackP0_lo}<\slack(\Pzero)<\cpv{slackP0_hi},\qquad\cpv{alphaP0_lo}<\alpha(\Pzero)<\cpv{alphaP0_hi},
\]
and $\co(\Pzero)\le\cpv{cbarP0_hi}$;
so $\Pzero$ is admissible, $\alpha(\Pzero)>\cpv{thmD_range_P0}$, $1/(\co(\Pzero)+2)\ge\cpv{tcapP0_lo}$ and $2/(\co(\Pzero)+2)>\cpv{tpadP0_lo}$.
\end{lemma}

These values are certified: they are evaluated in exact rational arithmetic, with every logarithm of a rational enclosed by the series of $\operatorname{artanh}$ with a rigorous bound on its tail, by the script of Section~\ref{cp:sec:certificates} (the bound on $2/(\co(\Pzero)+2)$ follows from $\co(\Pzero)\le\cpv{cbarP0_hi}$), and that bound is certified a second time by an independent script in decimal interval arithmetic. By Lemma~\ref{cp:L5}, $\cstar\le\cbar\le\co(\Pzero)$ for the family of $\Pzero$, so on $(0,\cpv{tpadP0_lo}]$ the padding term of Corollary~\ref{cp:T5} is negative and $\Lamp=\Lam$ there.

\section{Composites with the data structure of~\texorpdfstring{\cite{v2}}{[Maj26a]}}\label{cp:sec:composites}

This section draws the consequences of the per-entry and column composites, Propositions~\ref{cp:T8} and~\ref{cp:T9}, for the standard level-1 tops of Lemma~\ref{cp:L10}. Their bottom is the data structure of~\cite[Theorem~\Vref{m:thm:general}]{v2} (Theorem~\ref{cp:P3}), so their queries are paid per entry or per column, and only the coverage and the column union of the top enter. Every running-time statement is in the time model, and every algorithm for which a running time is claimed is deterministic: the data structure of Theorem~\ref{cp:P3} is deterministic, the top is a fixed finite object, and at level~1 every kept output has the same coverage bound and every column the same union bound (Lemma~\ref{cp:L10}), so no output is relabelled and nothing is chosen at random. The exception is Lemma~\ref{cp:D:SD}, a statement on space and query time, whose hash maps are shown to exist by counting (Remark~\ref{cp:D:SDfalse}). We say ``deterministic'' at every threshold, because randomized refutations of the hinted OMv conjectures reach further, up to $\tau<\cpv{hinted_omv_mc}$ in~\cite[Corollary~\Vref{t:cor:omv}]{v2}. Each threshold of a level-1 top is certified at a named rational point, and every printed threshold is rounded down. Section~\ref{cp:sec:profile} explains why only the coverage of the top enters the per-entry thresholds: single entries lie on the plateau of the leaf-degree profile, where only its value at $t=1$ counts (Remark~\ref{cp:F9}).

\subsection{Two certified points}\label{cp:sec:D:points}

\begin{lemma}[the points $\Ptt$ and $\Pthree$]\label{cp:D:points}
In the coordinates $(\lambda_6,\,u_6(6+2),\,u_7(7+2))$ of Definition~\ref{cp:L3} let
\[
\Ptt=\bigl(\cpv{P23_lam},\,\cpv{P23_u6},\,\cpv{P23_u7}\bigr),\qquad\Pthree=\bigl(\cpv{P32_lam},\,\cpv{P32_u6},\,\cpv{P32_u7}\bigr).
\]
Both points are admissible: $\slack(\Ptt)>\cpv{P23_slack_lo}$ and $\slack(\Pthree)>\cpv{P32_slack_lo}$. Moreover
\begin{gather*}
\cpv{alphaP23_lo}<\alpha(\Ptt)<\cpv{alphaP23_hi},\qquad\cpv{coP23_lo}<\co(\Ptt)<\cpv{coP23_hi},\\
\cpv{alphaP32_lo}<\alpha(\Pthree)<\cpv{alphaP32_hi},\qquad\cpv{cprimeP32_lo}<\ccol(\Pthree)-1<\cpv{cprimeP32_hi}.
\end{gather*}
With $E_Q(\alpha,c)=c\es/(Q\es+c-Q\alpha)$ the range function of Proposition~\ref{cp:T8}, and $E_1(\alpha,c')=c'\es/(\es+c'-\alpha)$ that of Proposition~\ref{cp:T9}:
\begin{enumerate}[label=(\alph*)]
\item $E_Q(\alpha(\Ptt),\co(\Ptt))$ exceeds the value printed in row $Q$ of Table~\ref{cp:tab:D4}, for each of its seven rows; in particular $E_1(\alpha(\Ptt),\co(\Ptt))>\cpv{per_entry_P23}$ and $E_{1/2}(\alpha(\Ptt),\co(\Ptt))>\cpv{online_setdisj_R77}$;
\item $E_1(\alpha(\Pthree),\ccol(\Pthree)-1)>\cpv{hinted_omv_R32}$.
\end{enumerate}
\end{lemma}

These values are certified: they are evaluated in exact rational arithmetic, with every logarithm of a rational enclosed by the series of $\operatorname{artanh}$ with a rigorous bound on its tail, by the scripts of Section~\ref{cp:sec:certificates}. The margins are small (for $\cpv{online_setdisj_R77}$ about $4.5\cdot10^{-8}$), which is why the points are printed with all their digits; a point given to three or five digits can fail to be admissible.

\subsection{Single entries}\label{cp:sec:D:entries}

\begin{theorem}[per-entry composite]\label{cp:D1}
For every $\eps<\cpv{per_entry_P23}$ there are $\gamma>0$, $Q'<1$ and a deterministic algorithm that, for all $N\ge1$ and $2\le D\le N^\eps$ and all $X\in\Z^{N\times D}$, $Y\in\Z^{D\times N}$ with entries of absolute value $N^{O(1)}$, preprocesses $X$ and $Y$ in time and space $O(N^2/D^\gamma)$, after which any single entry of $XY$ is computed in $O(D^{Q'})$ time (time model). The threshold is certified at the point $\Ptt$.
\end{theorem}

\begin{proof}
$\Ptt$ is an admissible rational point (Lemma~\ref{cp:D:points}), so $\TPE(\alpha(\Ptt),\co(\Ptt))$ holds by Lemma~\ref{cp:L10}. By the intervals of Lemma~\ref{cp:D:points}, $\alpha(\Ptt)>\es$ and $\co(\Ptt)>\alpha(\Ptt)$, which are the hypotheses of Proposition~\ref{cp:T8} with $Q=1$. That proposition gives the structure for every $\eps<E_1(\alpha(\Ptt),\co(\Ptt))$, and this exceeds $\cpv{per_entry_P23}$ by Lemma~\ref{cp:D:points}(a).
\end{proof}

\begin{corollary}\label{cp:D2}
Conjecture~5.12 of~\cite{BNS19} fails, deterministically, for every $\tau_1<\cpv{per_entry_P23}$ and every $\tau_2\in(0,1)$: Phase~2 runs in $O(n^{2-\gamma\tau_1})$ time, Phase~3 only stores $J$, and Phase~4 runs in $O(n^{\tau_2+Q'\tau_1})$ time, with $\gamma$ and $Q'$ of Theorem~\ref{cp:D1}. Also Conjectures~5.2 and~5.7 fail for every $\tau<\cpv{per_entry_P23}$; Theorem~\ref{cp:D5} below gives more for these two.
\end{corollary}

\begin{proof}
Given $\tau_1<\cpv{per_entry_P23}$, apply Theorem~\ref{cp:D1} with $\eps=\tau_1$; it provides the single-entry interface of Lemma~\ref{cp:P5}(b) at $\eps$, and Lemma~\ref{cp:P5}(b) gives all three statements. Phase~4 reads $t_2$ entries of one row of $XY$ at the arbitrary columns $J$, so these are single-entry queries and nothing is shared between them.
\end{proof}

The same composite gives a single-entry structure for every query exponent, with a range that grows with the exponent.

\begin{proposition}[range against query exponent]\label{cp:D4}
Let $P$ be an admissible rational level-1 point with $\alpha(P)>\es$, and let $Q\in(0,1]$ with $\co(P)>Q\alpha(P)$. For every
\[
\eps<E_Q(\alpha(P),\co(P))=\frac{\co(P)\,\es}{Q\es+\co(P)-Q\alpha(P)}
\]
there are $\gamma>0$, $Q'<Q$ and a deterministic algorithm as in Theorem~\ref{cp:D1} that computes any single entry of $XY$ in $O(D^{Q'})$ time after a preprocessing in time and space $O(N^2/D^\gamma)$ (time model). At $P=\Ptt$ every $\eps$ below the value in Table~\ref{cp:tab:D4} qualifies.
\end{proposition}

\begin{table}[h]
\centering
\begin{tabular}{c|ccccccc}
$Q$ & $\cpv{D4_Qa}$ & $\cpv{D4_Qb}$ & $\cpv{D4_Qc}$ & $\cpv{D4_Qd}$ & $\cpv{D4_Qe}$ & $\cpv{D4_Qf}$ & $\cpv{D4_Qg}$\\\hline
every $\eps<$ & $\cpv{D4_a}$ & $\cpv{D4_b}$ & $\cpv{D4_c}$ & $\cpv{D4_d}$ & $\cpv{D4_e}$ & $\cpv{D4_f}$ & $\cpv{D4_g}$
\end{tabular}
\caption{Single entries in $O(D^{Q'})$ with $Q'<Q$ after a preprocessing in time and space $O(N^2/D^\gamma)$, $\gamma>0$: certified thresholds of Proposition~\ref{cp:D4} at the point $\Ptt$ (deterministic, time model).}\label{cp:tab:D4}
\end{table}

\begin{proof}
Lemma~\ref{cp:L10} gives $\TPE(\alpha(P),\co(P))$, and Proposition~\ref{cp:T8} applies with the given $Q$. At $\Ptt$ the hypotheses hold by the intervals of Lemma~\ref{cp:D:points}, and the table rows are Lemma~\ref{cp:D:points}(a).
\end{proof}

\begin{remark}[a superseded route to $\Hyp$]\label{cp:D11}
By the remark after Proposition~\ref{cp:T8}, the row $Q=\frac12$ of Table~\ref{cp:tab:D4} gives $\Hyp(\eps,\gamma)$ with some $\gamma>0$ for every $\eps<\cpv{online_setdisj_R77}$ (time model, certified at $\Ptt$). Corollary~\ref{cp:A2} supersedes this: it gives every $\eps\le\cpv{thmD_range_P0}$, by sharing the queries instead of paying them per entry. The number $\cpv{online_setdisj_R77}$ reappears in Theorem~\ref{cp:D3} because $\Hyp$, with its $N^2/\sqrt D$ wanted entries in time $N^2/D^\gamma$, asks of a per-entry structure exactly a query exponent below $\frac12$.
\end{remark}

\begin{remark}[the per-entry composite without the laser method]\label{cp:D10}
Proposition~\ref{cp:T8} needs of the top only its shape, its number of terms and its maximal coverage, and the maximal coverage is at most the number of terms. Suppose $\omega(1,a,1)=2$ for every $a<\alpha_0$. Then for every $a'<\alpha_0$, every $\eta>0$ and all large $k$ there is an exact algorithm over $\Q$ for $\mat k{\lfloor k^{a'}\rfloor}k$ with at most $k^{2+\eta}$ terms (by the definition of $\omega$ in Section~\ref{cp:sec:laser}), and its maximal coverage is at most its number of terms. So $\TPE(a',2)$ holds, and Proposition~\ref{cp:T8} with $Q=\frac12$ and $c=2$ gives $\Hyp(\eps,\gamma)$ with some $\gamma>0$ for every $\eps<4\es/(\es+4-a')$, hence, letting $a'\to\alpha_0$, for every
\[
\eps<\frac{4\es}{\es+4-\alpha_0}.
\]
With the published bound $\alpha\ge\AlphaVXXZ$~\cite{VXXZ24} this is every $\eps<\cpv{R11_range}$; if $\alpha\ge\cpv{alpha_preprints}$, as two preprints claim~\cite{AvdBZ26,Uzn26}, it would be every $\eps<\cpv{D10_pre}$. Both are below Corollary~\ref{cp:A2}, which uses only the level-1 analysis of Section~\ref{cp:sec:level1}. The bound $c=2$ is the worst-case coverage, which the level-1 tops improve to $\co(P)$.
\end{remark}

\subsection{Single entries below \texorpdfstring{$\sqrt D$}{sqrt D}: set disjointness}\label{cp:sec:D:setdisj}

In \emph{Online Set Disjointness} two families of $m$ sets over a universe of size $D$ are preprocessed in time and space polynomially in $D$ below $m^2$, after which a query $(I,J)$ is answered with $\abs{S_I\cap T_J}$ in time polynomially below $\sqrt D$, one query at a time~\cite{AV26}.

\begin{theorem}[Online Set Disjointness]\label{cp:D3}
For every $\eps<\cpv{online_setdisj_R77}$ there are $\gamma>0$, $Q<\frac12$ and a deterministic algorithm with the following property. Given two families $S_1,\dots,S_m$ and $T_1,\dots,T_m$ of subsets of $[D]$ with $2\le D\le m^\eps$, it preprocesses them in time and space $O(m^2/D^\gamma)$; afterwards, given any pair $(I,J)$, it outputs $\abs{S_I\cap T_J}$ in $O(D^Q)$ time. The query bound holds for every pair (worst case), and each query is answered on its own, without knowledge of later queries (online). Time model; certified at the point $\Ptt$.
\end{theorem}

\begin{proof}
Let $X\in\{0,1\}^{m\times D}$ have the incidence vectors of the $S_I$ as rows and $Y\in\{0,1\}^{D\times m}$ those of the $T_J$ as columns, built in $O(mD)$ time, which is below $m^2/D^\gamma$ once $\gamma\le1$, since $D\le m^\eps$ and $\eps<\frac12$. Then $(XY)[I,J]=\abs{S_I\cap T_J}$. Proposition~\ref{cp:D4} at $\Ptt$ with $Q=\frac12$ (row $\cpv{D4_Qe}$ of Table~\ref{cp:tab:D4}) gives the preprocessing and a query exponent $Q<\frac12$. The query algorithm of Proposition~\ref{cp:T8} reads only the stored sub-structures and the pair $(I,J)$, and its bound holds for every entry because every kept output of the level-1 top has the same coverage bound (Lemma~\ref{cp:L10}); so the bound is worst case and the algorithm is online.
\end{proof}

For comparison: the data structure of~\cite{v2} alone (Theorem~\ref{cp:P3} with $q_A<\frac12$) gives this for every $\eps<\es=\cpv{range_v1}$, the exact box count of~\cite[Theorem~\Vref{m:thm:fixedq}]{v2} for every $\eps<\cpv{hyp_half_v2}$, and~\cite{AV26} for $D\le m^{\cpv{osd_avw26}}$.

\begin{remark}[which conjectures need a query exponent below $\frac12$]\label{cp:D9}
Among the consequences of this section, only Online Set Disjointness, the set-disjointness construction below and the per-entry route to $\Hyp$ of Remark~\ref{cp:D11} need single entries with a query exponent below $\frac12$. None of the hinted OMv conjectures of Section~\ref{cp:sec:omv} does: Conjectures~5.2 and~5.7 ask for one column, $n$ entries within $n^{1+\tau}$ time, and Phase~4 of Conjecture~5.12 asks for one entry of $UN_{I,J}V$, which the reduction of Lemma~\ref{cp:P5} answers through $t_2$ entries of one row of a thin product within $n^{\tau_1+\tau_2}=t_2\cdot n^{\tau_1}$ time. In both cases a query exponent below~$1$ per entry suffices.
\end{remark}

The bounded-universe lower bound of Goldstein, Lewenstein and Porat~\cite[Theorem~4]{GLP19} reads: any solution to SetDisjointness with sets $S_1,\dots,S_m\subseteq[u]$ for any value of $u\in[N^\delta,N]$, where $N=\sum_i\abs{S_i}$ and $\delta>0$, must either use $\Omega(m^{2-o(1)})$ space or have $\tilde\Omega(u^{1/2-o(1)})$ query time, unless the Strong SetDisjointness Conjecture is false. Its conclusion fails for small universes:

\begin{corollary}\label{cp:D13}
For every $0<c<\cpv{D13_c}$ there are $\gamma>0$ and $Q<\frac12$ such that for every universe size $u=\lceil N^c\rceil$ the structure of Theorem~\ref{cp:D3}, applied to one family, answers SetDisjointness on every family $S_1,\dots,S_m\subseteq[u]$ with $N=\sum_i\abs{S_i}$ in deterministic preprocessing time and space $O(m^{2-c\gamma})$ and worst-case query time $O(u^Q)$ on a word RAM. Such $u$ lie in the range $[N^\delta,N]$ of~\cite[Theorem~4]{GLP19} for every $\delta\le c$, so for these universes both alternatives of its conclusion fail, by the fixed powers $m^{c\gamma}$ and $u^{1/2-Q}$.
\end{corollary}

\begin{proof}
Remove the empty sets; queries involving them are answered ``disjoint'' from an array of $m$ bits. Then $m\le N$, and $m\ge N/u$ since every set has at most $u$ elements. Since $c<\cpv{online_setdisj_R77}/(1+\cpv{online_setdisj_R77})$, that is $c/(1-c)<\cpv{online_setdisj_R77}$, we can choose $\eps$ with $c/(1-c)<\eps<\cpv{online_setdisj_R77}$. Then $c(1+\eps)<\eps$, so $u\le2N^c\le(N/u)^\eps\le m^\eps$ for all large $N$. Apply Theorem~\ref{cp:D3} at this $\eps$, with both families equal to $(S_i)$ and $D=u$; it gives $\gamma$ and $Q<\frac12$. Building the incidence matrix costs $O(N+mu)\le O(m^{1+\eps})$, within the bound below once $c\gamma<1-\eps$. The space is $O(m^2u^{-\gamma})\le O(m^2N^{-c\gamma})\le O(m^{2-c\gamma})$, because $u\ge N^c$ and $m\le N$, and the query time is $O(u^Q)$.
\end{proof}

The corollary is stated per universe size: it does not contradict Theorem~4 read as a statement about one solution that must work at every $u$ in the range at once, since Theorem~\ref{cp:D3} needs $u\le m^\eps$, which fails for the families of about $N/u$ sets of size $u$ as soon as $u>N^{\eps/(1+\eps)}$, and $\eps/(1+\eps)<\cpv{online_setdisj_R77}/(1+\cpv{online_setdisj_R77})$ for every $\eps<\cpv{online_setdisj_R77}$. Theorem~4 itself remains valid as an implication; its hypothesis, the Strong SetDisjointness Conjecture, is false (Remark~\ref{cp:D:SDfalse}).

The construction behind that remark turns any single-entry structure with a query exponent below $\frac12$ into a SetDisjointness structure with general universe. It is the construction of~\cite[Theorem~4]{GLP19}, re-run with a saving in powers of the universe instead of the number of sets, and with padding; it is also in~\cite[Proposition~\Vref{m:prop:setdisjone}]{v2}.

\begin{lemma}[from single entries to SetDisjointness]\label{cp:D:SD}
Let $0<\eps\le1$, $0<\gamma\le1$ and $0\le q<\frac12$, and suppose a deterministic algorithm stores, for all $m$ and $2\le D\le m^\eps$ and every $X\in\{0,1\}^{m\times D}$, a structure of $O(m^2/D^\gamma)$ words from which any entry of $XX^{\mathsf T}$ is computed in $O(D^q)$ time. Put
\[
g=\min(2\gamma,1-2q),\qquad\lambda=\frac2{4/\eps+2-2\gamma-g},\qquad\Esd(\eps,\gamma,q)=g\lambda=\frac{2g\eps}{4+(2-2\gamma-g)\eps}.
\]
Then there is a data structure for SetDisjointness (a family of sets over an arbitrary universe with total size $N$; a query asks whether two of the sets are disjoint) with worst-case query time $T=O(N^{\lambda(1-g)}\log^{O(1)}N)$ and space $S$ with $ST^2=O(N^{2-\Esd}\log^{O(1)}N)$, on a word RAM with $O(\log N)$-bit words. Its queries and its space are as stated for every instance; its preprocessing time is not bounded here.
\end{lemma}

\begin{proof}
Rename the elements into $[N]$ and keep each set as a sorted array, $O(N)$ words, so that a membership test costs $O(\log N)$. Store only the at most $N$ nonempty sets, in a sorted index; a query naming an absent index answers ``disjoint''. Put $\ell=N^\lambda$ and $s=\ell^{1-g}$; a set is \emph{small} if it has at most $s$ elements, \emph{large} if it has more than $\ell$, and \emph{medium} otherwise. There are $d\le N/\ell$ large and $e\le N/s$ medium sets. A query $(i,j)$ with $\abs{S_i}\le\abs{S_j}$ is answered as follows. If $S_i$ is small, look up its elements in $S_j$, in $O(s\log N)$ time. If $S_j$ is large, read a stored $d\times(d+e)$ bit matrix of the answers, which takes at most $2N^2/(\ell s)$ bits. If both are medium, use hashing into the universe $[D']$, $D'=8\ell^2$: for two medium sets and a random map $h$ into $[D']$, the union bound over at most $\ell^2$ pairs of elements shows that $h(S_i)\cap h(S_j)=\emptyset$ fails with probability at most $\frac18$ when $S_i\cap S_j=\emptyset$; so there are $K_h=O(\log N)$ maps $h_1,\dots,h_{K_h}$ such that every disjoint pair of medium sets stays disjoint under at least one of them, and these maps are fixed once for the instance. For each $k$ let $X_k$ be the $e'\times D'$ incidence matrix of the sets $h_k(S)$, $S$ medium, padded with empty rows to $e'=\max(e,\lceil D'^{1/\eps}\rceil)$ rows, so that $D'\le e'^\eps$, and store the given structure for $X_k$. The answer is ``disjoint'' if and only if $(X_kX_k^{\mathsf T})[i,j]=0$ for some $k$: a common element of $S_i$ and $S_j$ is a common element of their images under every $h_k$, and a disjoint pair is disjoint under some $h_k$. The queries never hash; they pass indices.

\emph{Query.} $T=O(s\log N+K_hD'^q)$, and $D'^q=8^q\ell^{2q}\le8s$ because $2q\le1-g$. So $T=O(\ell^{1-g}\log N)$.

\emph{Space.} $S=O(N)+2N^2/(\ell s)+K_h\,O(e'^2/D'^\gamma)$ words, and $e'^2\le N^2/s^2+4D'^{2/\eps}$. Hence, up to factors $\log^{O(1)}N$, $ST^2$ is at most the sum of the four terms
\[
N\ell^{2-2g},\qquad N^2\ell^{-g},\qquad N^2D'^{-\gamma}\le N^2\ell^{-2\gamma}\le N^2\ell^{-g},\qquad D'^{2/\eps-\gamma}\ell^{2-2g}=\Theta\bigl(\ell^{4/\eps-2\gamma+2-2g}\bigr),
\]
using $g\le2\gamma$. The first is at most $N^2\ell^{-g}$ if and only if $\lambda(2-g)\le1$, and the last if and only if $\lambda(4/\eps+2-2\gamma-g)\le2$. The choice of $\lambda$ makes the last condition an equality, and then the first holds, because $2(2-g)\le4/\eps+2-2\gamma-g$ is equivalent to $2-g+2\gamma\le4/\eps$, which is true for $\eps\le1$ and $\gamma\le1$. So $ST^2=O(N^2\ell^{-g}\log^{O(1)}N)=O(N^{2-g\lambda}\log^{O(1)}N)$.
\end{proof}

\begin{remark}[the Strong SetDisjointness Conjecture is false]\label{cp:D:SDfalse}
The Strong SetDisjointness Conjecture~\cite[Conjecture~3]{GKLP17}, \cite[Conjecture~2]{GLP19} asserts that every data structure for SetDisjointness with query time $T$ uses space $S=\tilde\Omega(N^2/T^2)$. The conjecture is false, and this is a one-line consequence of Lemma~\ref{cp:D:SD}: applied to the data structure of Theorem~\ref{cp:P3} with $Y=X^{\mathsf T}$, any $\eps_A<\es$, any $0<q_A<\frac12$ and its $\gamma_A$ (lowered to at most $1$), the lemma gives a data structure with $ST^2=O(N^{2-\Esd}\log^{O(1)}N)$ on every instance, where $\Esd=\Esd(\eps_A,\gamma_A,q_A)>0$ because $\min(2\gamma_A,1-2q_A)>0$. This consequence is also in~\cite[Theorem~\Vref{m:thm:setdisj}]{v2}, with explicit exponents. The model is the word RAM with $O(\log N)$-bit words, and the structure uses subtraction, as the identity of~\cite{v2} does; this is necessary, since the conjecture holds in the semigroup model~\cite{DMRU95}, as noted in~\cite{GLP19}. The maps $h_k$ of the proof of Lemma~\ref{cp:D:SD} are shown to exist by counting, so the preprocessing is not claimed to be deterministic; the conjecture concerns only space and query time, which are as stated.
\end{remark}

The composites of this section raise the range of a single-entry structure, but not the exponent $\Esd$:

\begin{proposition}[no composite raises $\Esd$]\label{cp:D8}
Let $0<\eps_A\le1$, $0<\gamma_A\le1$ and $0\le q_A<\frac12$, and suppose a deterministic algorithm preprocesses, for all $N$ and $2\le D\le N^{\eps_A}$, any $X\in\Z^{N\times D}$ and $Y\in\Z^{D\times N}$ with entries of absolute value $N^{O(1)}$ in time and space $O(N^2/D^{\gamma_A})$, after which any entry of $XY$ is computed in $O(D^{q_A})$ time, as Theorem~\ref{cp:P3} does; with $Y=X^{\mathsf T}$ it satisfies the hypothesis of Lemma~\ref{cp:D:SD}. Let $B$ be an exact top algorithm for $\mat kbk$ with $\alpha=\alB\ge0$, rank excess $\delta=\delB\ge0$ and coverage exponent $c=\ln\cB/\ln k\ge\alpha/2$. For $x\ge0$ the proof of Proposition~\ref{cp:T8}, with top $B$, split $x$ and this algorithm as the bottom, gives a single-entry structure whose exponents tend to
\[
\eps_x=\frac{1+\alpha x}{x+1/\eps_A},\qquad\gamma_x=\frac{\gamma_A-x\delta}{1+\alpha x},\qquad Q_x=\frac{xc+q_A}{1+\alpha x}.
\]
(range, preprocessing saving and query exponent) as its parameter $\eta_0$ tends to $0$. Put $\Esd:=0$ when $g\le0$. Then $\Esd(\eps_x,\gamma_x,Q_x)\le\Esd(\eps_A,\gamma_A,q_A)$, with equality only at $x=0$.
\end{proposition}

\begin{proof}
Write $u=1/\eps$, $g_A=\min(2\gamma_A,1-2q_A)$ and $\sigma=1/(1+\alpha x)\in(0,1]$. Then $\gamma_x=\sigma(\gamma_A-x\delta)\le\sigma\gamma_A$ and $1-2Q_x=\sigma(1+\alpha x-2xc-2q_A)\le\sigma(1-2q_A)$, because $2c\ge\alpha$; so $g_x:=\min(2\gamma_x,1-2Q_x)\le\sigma g_A$. If $g_x\le0$, then $\Esd=0$. Otherwise, with $u_x=\sigma(x+u_A)$, $\Esd=2g/(4u+2-2\gamma-g)$ increases in $g$ and decreases in the denominator, and $-2\gamma_x\ge-2\sigma\gamma_A$, $-g_x\ge-\sigma g_A$; hence
\begin{align*}
\Esd(\eps_x,\gamma_x,Q_x)&\le\frac{2\sigma g_A}{4\sigma(x+u_A)+2-2\sigma\gamma_A-\sigma g_A}=\frac{2g_A}{4x+4u_A+2/\sigma-2\gamma_A-g_A}\\
&\le\frac{2g_A}{4u_A+2-2\gamma_A-g_A}=\Esd(\eps_A,\gamma_A,q_A),
\end{align*}
using $x\ge0$ and $1/\sigma\ge1$; all denominators are positive since $u_A\ge1$, $\gamma_A\le1$ and $g_A\le1$. Equality forces $x=0$ when $g_A>0$. For $\eta_0>0$ the proof of Proposition~\ref{cp:T8} gives a smaller range, a smaller saving and a larger query exponent than these limits, and $\Esd$ does not decrease in $\eps$ and $\gamma$ and does not increase in $Q$, so the bound holds for those exponents too.
\end{proof}

The reason is visible in the proof: the composite multiplies $D$ by a power and divides $\gamma$ and $1-2Q$ by the same factor $1+\alpha x$, so it trades range for saving at an exchange rate of at most~$1$, while $\Esd$ behaves like the product $g\eps$. The range gain of Table~\ref{cp:tab:D4} lives at $\gamma\to0$, where $\Esd\to0$.

\subsection{Columns}\label{cp:sec:D:columns}

\begin{theorem}[column composite]\label{cp:D5}
For every $\eps<\cpv{hinted_omv_R32}$ there are $\gamma>0$, $Q'<1$ and a deterministic algorithm that preprocesses $X\in\Z^{N\times D}$ and $Y\in\Z^{D\times N}$, $2\le D\le N^\eps$, entries of absolute value $N^{O(1)}$, in time and space $O(N^2/D^\gamma)$, after which any whole column of $XY$, and any whole row, is computed in $O(ND^{Q'})$ time, for every column and every row (worst case; time model). The threshold is certified at the point $\Pthree$. Consequently Conjectures~5.2 and~5.7 of~\cite{BNS19} fail for every $\tau<\cpv{hinted_omv_R32}$, deterministically: Phase~2 runs in $O(n^{2-\gamma\tau})$ and Phase~3 in $O(n^{1+Q'\tau})$ time.
\end{theorem}

\begin{proof}
$\Pthree$ is an admissible rational point (Lemma~\ref{cp:D:points}), so $\Tcol(\alpha(\Pthree),\ccol(\Pthree)-1)$ holds by Lemma~\ref{cp:L10}, and the intervals of Lemma~\ref{cp:D:points} give $\ccol(\Pthree)-1>\alpha(\Pthree)>\es$. Proposition~\ref{cp:T9}, applied to $(X,Y)$ for columns and to $(Y^{\mathsf T},X^{\mathsf T})$ for rows, gives the structure for every $\eps<E_1(\alpha(\Pthree),\ccol(\Pthree)-1)$, which exceeds $\cpv{hinted_omv_R32}$ by Lemma~\ref{cp:D:points}(b). For $\tau<\cpv{hinted_omv_R32}$ take $\eps=\tau$; Lemma~\ref{cp:P5}(a) gives the refutations, with one Phase-3 column of the $n\times n$ product.
\end{proof}

\begin{corollary}\label{cp:D6}
Conjecture~5.12 of~\cite{BNS19} fails, deterministically,
\begin{enumerate}[label=(\alph*),ref=\thetheorem(\alph*)]
\item for every $\tau_1\in(0,1)$ and every $\tau_2<\cpv{hinted_omv_R32}\,\tau_1$;
\item\label{cp:D7} for every $\tau_2\in(0,1)$ and every $\tau_1<\cpv{hinted_omv_R32}\,\tau_2$.
\end{enumerate}
Part (b) adds to the region of Corollary~\ref{cp:D2} only points with $\tau_2>\cpv{D7_tau2}$.
\end{corollary}

\begin{proof}
(a) Given $\tau_2<\cpv{hinted_omv_R32}\tau_1$, take $\eps=\tau_2/\tau_1$ in Theorem~\ref{cp:D5}; it provides the column interface of Lemma~\ref{cp:P5}(a) at $\eps$, and Lemma~\ref{cp:P5}(a) gives the refutation, with Phase~3 in $O(n^{1+\tau_1-\gamma\tau_2})$ and Phase~4 in $O(n^{\tau_1+Q'\tau_2})$.
(b) Given $\tau_1<\cpv{hinted_omv_R32}\tau_2$, take $\eps=\tau_1/\tau_2$ in Theorem~\ref{cp:D5}. Phase~2 only stores $I$. In Phase~3 the matrices $U$ ($n\times t_1$) and $N_{I,J}$ ($t_1\times t_2$) are known; cut the $n$ rows of $U$ into $n/t_2$ blocks of $t_2$ rows and preprocess, for every block $U'$, the product $U'N_{I,J}$ with $t_2$ in the role of $N$ and $D=t_1\le t_2^\eps$, for rows. This costs $(n/t_2)\,O(t_2^2/t_1^\gamma)=O(n^{1+\tau_2-\gamma\tau_1})$, below $n^{\omega(\tau_2,\tau_1,1)}\ge n^{1+\tau_2}$ by the polynomial factor $n^{\gamma\tau_1}$. In Phase~4, $(UN_{I,J}V)_{i,j}=1$ if and only if $(UN_{I,J})_{i,b}>0$ and $V_{b,j}=1$ for some $b\le t_2$; row $i$ of $UN_{I,J}$ is one whole row of one block product, computed in $O(t_2t_1^{Q'})=O(n^{\tau_2+Q'\tau_1})$, below $n^{\tau_1+\tau_2}$, and the test against column $j$ of $V$ costs $O(t_2)$ more.
For the last sentence: a point of region (b) outside Corollary~\ref{cp:D2}'s region has $\cpv{per_entry_P23}\le\tau_1<\cpv{hinted_omv_R32}\tau_2$, so $\tau_2>\cpv{per_entry_P23}/\cpv{hinted_omv_R32}>\cpv{D7_tau2}$.
\end{proof}

\paragraph{Consequences for work that cites the conjectures.}\phantomsection\label{cp:D14}
The refutations of Theorem~\ref{cp:D5} remove the support of some conditional lower bounds. None of the lower bounds discussed below is refuted: we give no faster distance oracle and no faster attention algorithm. What is lost is the cited hypothesis at the stated parameters.

The lower bounds of~\cite{BNS19} that rest on Conjectures~5.2 and~5.7, namely their Theorems~5.3 and~5.8 and Corollaries~5.4, 5.5, 5.9 and~5.10, lose their support for every $\tau<\cpv{hinted_omv_R32}$; this is the end of their update--query trade-offs with the fastest queries.

For the distance oracles of Haeupler, Long and Saranurak~\cite{HLS24}: they cite~\cite[Corollary~5.10]{BNS19}, which rests on Conjecture~5.7, for the bound that no decremental or incremental distance oracle with approximation $5/3-\eps$ has worst-case update time $n^{2-\Omega(1)}$ and query time $n^{o(1)}$. An oracle with update time $n^{2-\delta}$ and query time $n^{o(1)}$ would solve Mv-hinted Mv with Phase~2 in $n^{2-(\delta-\tau)}$ and Phase~3 in $n^{1+o(1)}$ time, refuting Conjecture~5.7 at every $\tau<\delta$; so the bound at $\delta$ is supported by the conjecture at any $\tau\in(0,\delta)$. For $\delta\le\cpv{hinted_omv_R32}$ every such $\tau$ is refuted, and new evidence of hardness is needed; for larger $\delta$ the remaining cases $\tau\in[\cpv{hinted_omv_R32},\delta)$ of the conjecture still rule out update time $O(n^{\cpv{D14_HLS}})$, since $2-\cpv{hinted_omv_R32}=\cpv{D14_HLS_exact}$ (rounded down as in~\cite{v2}). Their setting has query time $n^{o(1)}$, so it needs only the form of Conjecture~5.7 with Phase~3 in $n^{1+o(1)}$; we do not refute that form.

For dynamic attention, van den Brand, Song and Zhou~\cite{BSZ24} state the hinted matrix--vector conjecture HMV (their Conjecture~3.1, restated as Conjecture~5.3): over the Boolean semiring, Phase~1 receives $M,V\in\{0,1\}^{n\times n}$, Phase~2 an $n\times n$ matrix $P$ with at most $n^\tau$ nonzero entries, Phase~3 an index $i$, and the answer is $MPV_{*,i}$; for every constant $0<\tau\le1$ no algorithm has polynomial Phase~1, Phase~2 in $O(n^{\omega(1,1,\tau)-\eps})$ and Phase~3 in $O(n^{1+\tau-\eps})$ time, for any constant $\eps>0$.

\begin{proposition}\label{cp:D:BSZ}
Conjecture~3.1 of~\cite{BSZ24} fails, deterministically, for every $\tau\in(0,\cpv{hinted_omv_R32})$.
\end{proposition}

\begin{proof}
Write $P=\sum_{k\le t}e_{r_k}e_{c_k}^{\mathsf T}$ with $t\le n^\tau$, padded with zero terms to $t=\lceil n^\tau\rceil$. Then $MPV=XY$ with $X=M_{*,(r_1,\dots,r_t)}$ ($n\times t$) and $Y=V_{(c_1,\dots,c_t),*}$ ($t\times n$), both known after Phase~2, and Phase~3 asks for one column of $XY$, the Boolean answer being the integer count tested for $>0$. Theorem~\ref{cp:D5} with $N=n$, $D=t$ and some $\eps\in(\tau,\cpv{hinted_omv_R32})$, so that $t\le n^\eps$ for all large $n$, gives Phase~2 in $O(n^{2-\gamma\tau})\le O(n^{\omega(1,1,\tau)-\gamma\tau})$ and Phase~3 in $O(n^{1+Q'\tau})=O(n^{1+\tau-(1-Q')\tau})$.
\end{proof}

So the lower bounds of~\cite{BSZ24} for dynamic attention (their Lemmas~5.5 and~5.9, and Lemma~1.4, the informal form of Lemma~5.9), each proved from Conjecture~3.1 at the same $\tau$, need a new hypothesis for $\tau<\cpv{hinted_omv_R32}$, and stand as before for $\tau\ge\cpv{hinted_omv_R32}$.

\section{Shared queries: \texorpdfstring{\ThmD}{Theorem D} and its range}\label{cp:sec:range}

This section draws the consequences of \ThmD{} for the range of $\eps$. It has two tiers.
From the level-1 analysis of Section~\ref{cp:sec:level1} alone, $\Hypk\kappa(\eps,\gamma)$ holds with some $\gamma>0$ for every $\kappa>0$ and every $\eps\le\cpv{thmD_range_P0}$ (Corollary~\ref{cp:A2}); the only cited input is Le Gall's Theorem~\ref{cp:P6:LG}. With a published fast rectangular product as the bottom, $\Hyp(\eps,\gamma)$ holds for every $\eps\le\cpv{A11_vxxz_half}$, using from~\cite{VXXZ24} only the bound $\alpha>\AlphaVXXZ$ (Corollary~\ref{cp:A11}).
Section~\ref{cp:sec:range-level1} proves the first tier: every $\eps$ up to the dual exponent $\alpha(P)$ \emph{of the top's family} is reached, at every $\kappa$, and a little more at finite~$x$. Section~\ref{cp:sec:range-converse} shows that no algebraic method that works for every $\kappa$ passes the true dual exponent $\alpha$, a statement in the multiplication model. Section~\ref{cp:sec:range-fast} replaces the inner products of \AlgD{} by a fast rectangular product on the dense terms; at $\kappa=1$ this reaches every $\eps<\alpha$, whatever the value of $\alpha$ is, and with the published bounds of~\cite{VXXZ24} even every $\eps\le\cpv{A11_tab1_one}$, above the published bound on $\alpha$ (Theorem~\ref{cp:A10}, Corollary~\ref{cp:A11}). The two-tier picture does not say that the range of the composite method is $\alpha$: the true $\alpha$ is claimed to be at least $\cpv{alpha_preprints}$~\cite{AvdBZ26,Uzn26}, which no top of this paper reaches; at $\kappa=1$ only the bottom of Section~\ref{cp:sec:range-fast} carries the range to~$\alpha$.

\subsection{The range at level 1}\label{cp:sec:range-level1}

\begin{theorem}[\ThmDk{} at level 1]\label{cp:A1}
Let $P$ be an admissible rational level-1 point (Definition~\ref{cp:L3}). For every $\kappa>0$ and every $0<\eps\le\alpha(P)$ there is $\gamma>0$ such that $\Hypk\kappa(\eps,\gamma)$ holds (time model; the algorithm is deterministic). For $0<\eps<\alpha(P)$ every $\gamma$ with $0<\gamma<\gbound{\mathrm D}\kappa(\eps)$ is admissible, where $\gbound{\mathrm D}\kappa(\eps)$ is the bound of Theorem~\ref{cp:T7}(b) with $\alpha=\alpha(P)$ and $\Lamp$ formed from the profile $\Lam$ of the input family of $P$ (Lemma~\ref{cp:L5}); here $\Lam\le\Lampr$ by Lemma~\ref{cp:L8}, and $\Lamp=\Lam$ for $t\le2/(\cstar+2)$.
\end{theorem}

\begin{proof}
By Lemma~\ref{cp:L5}, the family $(A_M)$ of $P$ satisfies \Hcond1--\Hcond4{} with $\alpha=\alpha(P)$. Here $\alpha(P)>0$ because every $u_q>0$, and $\alpha(P)<1$ because $u_q<w_q=1/(q+2)<p_q=q/(2(q+2))$ for $q\in\{6,7\}$. By Lemma~\ref{cp:L8}(iv) the family has $\cg=0$. Theorem~\ref{cp:T7}(a) gives the first statement and Theorem~\ref{cp:T7}(b) the second. The algorithm is \AlgD{} with the top algorithm of Corollary~\ref{cp:T5}.
\end{proof}

The proof uses no result of~\cite{v2}: its bottom is an inner product, and its only cited input is Le Gall's count, Theorem~\ref{cp:P6:LG}, through Lemma~\ref{cp:L4}. The constant $C$ of $\Hypk\kappa$ is astronomically large, because the top algorithm is a large fixed power; the saving $\gamma$ is not (Section~\ref{cp:sec:gamma}).

\begin{corollary}[the point $\Pzero$]\label{cp:A2}
For every $\kappa>0$ and every $0<\eps\le\cpv{thmD_range_P0}$ there is $\gamma>0$ such that $\Hypk\kappa(\eps,\gamma)$ holds (time model).
\end{corollary}

\begin{proof}
$\Pzero$ is admissible and $\alpha(\Pzero)>\cpv{alphaP0_lo}>\cpv{thmD_range_P0}$ (Lemma~\ref{cp:L11}, certified); apply Theorem~\ref{cp:A1}.
\end{proof}

Section~\ref{cp:sec:gamma} makes $\gamma$ explicit at $\Pzero$ (Theorem~\ref{cp:G3}).

\begin{remark}[the range beyond $\alpha(\Pzero)$]\label{cp:A3}
At finite $x$ the range of \ThmDk{} is $\Erange{\mathrm D}(x)=(1+\alpha x)/(1+x)>\alpha$, so it exceeds $\alpha(\Pzero)$ whenever \Dcond{} can be met at a finite $x$ (Theorem~\ref{cp:T7}(c)). At $\Pzero$ and $\kappa=\frac12$ this gives two certified values, which differ in the profile bound they use. With the bound $\Lam\le\Lampr$ of Lemma~\ref{cp:L8}, \Dcond{} holds at $t=\cpv{excess_crude_t}$ and $x=\cpv{excess_crude_x}$, so $\Hyp(\eps,\gamma)$ holds with some $\gamma>0$ for every $\eps<\cpv{excess_crude_P0}$. With the sharper bound $\Lam\le\Lamcert$ of Lemma~\ref{cp:G1}, it holds for every $\eps\le\cpv{excess_sharp_P0}$, at $x=\cpv{excess_sharp_x}$ with every $\gamma<\cpv{excess_sharp_gamma}$ (Table~\ref{cp:tab:G10}). Both are evaluated in exact rational arithmetic with rigorous logarithm enclosures (Section~\ref{cp:sec:certificates}).
\end{remark}

\begin{remark}[the supremum over level-1 points]\label{cp:A4}
Let $\alpha_C$ be the supremum of $\alpha(P)$ over the admissible rational level-1 points of Coppersmith's mixture of $F_6$ and $F_7$. By Lemma~\ref{cp:L9} it equals the supremum over all points with $\slack(P)\ge0$, $\lambda_6\in(0,1)$ and $0<u_q<w_q$. For every $\eps<\alpha_C$ some admissible rational $P$ has $\alpha(P)>\eps$, so Theorem~\ref{cp:A1} gives $\Hypk\kappa(\eps,\gamma)$ with some $\gamma>0$ for every $\kappa>0$ and every $\eps<\alpha_C$. The value of $\alpha_C$ is known only numerically: a floating-point optimisation gives $\alpha_C\approx\cpv{alphaC_computed}$ (computed, not certified). The certified value is $\cpv{thmD_range_P0}$ of Corollary~\ref{cp:A2}. Coppersmith~\cite{Cop97} reports $\alpha>\cpv{alpha_cop97}$ for this mixture; we cite it for the mixture, not as an input of a proof.
\end{remark}

\begin{remark}[per-entry composites have a different cost model]\label{cp:A5}
The per-entry composite of Proposition~\ref{cp:T8} pays one sub-query for every wanted entry and every top term that covers its block, $\abs W K^{\co}$ sub-queries in all, and its range $E_{1/2}(\alpha,c)=2c\es/(\es+2c-\alpha)$ depends on the top only through $\alpha$ and the coverage exponent~$c$. \ThmD{} pays for every pair (term, position) once, for all wanted entries at that position (Lemma~\ref{cp:T2}), and its bottom is an inner product, not the data structure of Theorem~\ref{cp:P3}, so $\es$ does not enter its range. The gain is seen at small $t$, through the way the supports are distributed over the terms, which coverage does not see: the standard and the merged decompositions of $F_q$ have the same coverage at every output (Lemma~\ref{cp:L2}(a)), and hence the same per-entry range, but the standard one has $\cg=\cpv{C12_cg_q7}$ at $q=7$ (computed; Proposition~\ref{cp:C12}) and the merged one $\cg=0$ (Lemma~\ref{cp:L8}). So bounds on per-entry composites in terms of $\alpha$ and coverage do not apply to \ThmD. One such bound is known. For the per-entry composite, $E_{1/2}(\alpha,c)=\es/(1-r/2)$ with $r=(\alpha-\es)/c$; if the top comes from a restriction of powers of $F_q$ with rank exponent~$2$ and its coverage is the product bound, then the logarithm of its maximal coverage is at least $(\ln(q+2)+q\ln2)/(q+2)$ per coordinate of $F_q$, and a floating-point computation with the quantum functionals of Christandl, Vrana and Zuiddam~\cite{CVZ18}, which are monotone under restriction and multiplicative, bounds $r$ by $\cpv{R27_r}$ for every $q\le12$, at any level of the laser analysis and for any mixture of values of~$q$ (we do not reproduce it). By that computation every such per-entry composite has range below $\cpv{R27_ceiling}$ (computed, not certified). This ceiling rests on the per-entry query count $\abs WK^{\co}$; \ThmD{} shares the queries, so the ceiling does not apply to it, and indeed Corollary~\ref{cp:A2} passes it.
\end{remark}

\subsection{No every-\texorpdfstring{$\kappa$}{kappa} algebraic method passes \texorpdfstring{$\alpha$}{alpha}}\label{cp:sec:range-converse}

The following converse is stated in the multiplication model of Section~\ref{cp:sec:models}; a time bound of an algorithm of this algebraic form implies the same bound on its non-scalar multiplications, not conversely. It is also proved in~\cite[Remark~\Vref{m:rem:converse}]{v2}; we restate it with its short proof, since Corollary~\ref{cp:A7} needs its exact hypotheses.

\begin{proposition}[the every-$\kappa$ converse; multiplication model]\label{cp:A6}
Let $\eps>0$ and $\kappa_0>0$. Suppose that for every $\kappa\in(0,\kappa_0]$ there are $\gamma=\gamma(\kappa)\ge0$ and $C$ such that for all $N$, all $D\le N^\eps$ and every $W\subseteq[N]^2$ with $\abs W\le N^2/D^\kappa$ there is a division-free straight-line program over $\Q$ in the entries of $X$ and $Y$ that computes the polynomials $(XY)[i,j]$, $(i,j)\in W$, as polynomials (equivalently, correctly for all integer matrices), with at most $CN^2D^{-\gamma}$ non-scalar multiplications. Then $\omega(1,\eps,1)=2$, and hence $\eps\le\alpha$.
\end{proposition}

\begin{proof}
Fix $\kappa\in(0,\kappa_0]$ with $\eps\kappa\le2$, let $N$ be large, and put $D=\lfloor N^\eps\rfloor\ge2$ and $s=\lfloor N^2/D^\kappa\rfloor$. Since $N^2/D^\kappa\ge1$, $s\ge N^2/(2D^\kappa)$, so $[N]^2$ is the union of at most $\lceil N^2/s\rceil\le2D^\kappa+1\le3D^\kappa$ sets of size at most~$s$. The programs for these sets, run one after the other, form one straight-line program that computes all entries of $XY$ with at most $3CN^2D^{\kappa-\gamma}\le3CN^{2+\eps\kappa}$ non-scalar multiplications. By Strassen's theorem~\cite{Str73} (Section~\ref{cp:sec:laser}), $R(\mat ND N)\le6CN^{2+\eps\kappa}$. Since $\lceil N^\eps\rceil\le D+1$ and one more inner index costs at most $N^2$ further rank-one terms, $R(\mat N{\lceil N^\eps\rceil}N)=O(N^{2+\eps\kappa})$, so $\omega(1,\eps,1)\le2+\eps\kappa$. This holds for every $\kappa\in(0,\kappa_0]$, so $\omega(1,\eps,1)\le2$; and $\omega(1,\eps,1)\ge2$, since the $N^2$ outputs are linearly independent. So $\omega(1,\eps,1)=2$, and $\eps\le\alpha$ by the definition of $\alpha$.
\end{proof}

Only small $\kappa$ and $\gamma\ge0$ are used: an algebraic algorithm with no saving at all, but working at every small $\kappa$, already forces $\eps\le\alpha$.

\begin{corollary}[sandwich]\label{cp:A7}
Let $\eps^{\mathrm{mult}}_{\mathrm{alg}}$ be the supremum of the $\eps$ for which the hypothesis of Proposition~\ref{cp:A6} holds for every $\kappa>0$, and $\eps^{\mathrm{time}}_{\mathrm{alg}}$ the supremum of the $\eps$ for which, for every $\kappa>0$, some $\gamma\ge0$ and some algorithm that is, for each $(N,D,W)$, a straight-line program as in Proposition~\ref{cp:A6} have running time $O(N^2D^{-\gamma})$ on the word RAM. Then for every admissible rational level-1 point~$P$
\[
\alpha(P)\le\eps^{\mathrm{time}}_{\mathrm{alg}}\le\eps^{\mathrm{mult}}_{\mathrm{alg}}\le\alpha .
\]
\end{corollary}

\begin{proof}
\AlgD{} is, for each $(N,D,W)$, a division-free straight-line program over $\Q$ that is correct as a polynomial identity (Remark~\ref{cp:T6:alg}), and by Theorem~\ref{cp:A1} it runs in time $O(N^2D^{-\gamma})$ with $\gamma>0$ for every $\kappa>0$ and every $\eps\le\alpha(P)$; this is the first inequality. The second holds because the number of non-scalar multiplications of such a program is at most its running time. The third is Proposition~\ref{cp:A6}.
\end{proof}

Since $\alpha>\AlphaVXXZ>\alpha(\Pzero)$ is known~\cite{VXXZ24}, the corollary says nothing new about $\alpha$. What it says is that $\alpha(P)$, the range of this construction, is a lower bound on $\eps^{\mathrm{time}}_{\mathrm{alg}}$, the best range of algebraic methods that work for every $\kappa$, and that no such method passes the true~$\alpha$. The lower end is a statement in the time model and the upper end one in the multiplication model.

\begin{remark}[one $\kappa$ caps nothing at $\kappa=1$]\label{cp:A8}
For one fixed $\kappa$, the covering argument in the proof of Proposition~\ref{cp:A6} gives only $\omega(1,\eps,1)\le2+\eps(\kappa-\gamma)$. So an algebraic algorithm with a saving $\gamma>0$ at a single $\kappa$ can exist only where $\omega(1,\eps,1)<2+\eps\kappa$. At $\kappa=1$ this excludes no $\eps\in(0,1]$: cutting the two factors of $\mat N{\lceil N^\eps\rceil}N$ into square blocks of side $\lceil N^\eps\rceil$ gives $\omega(1,\eps,1)\le2+\eps(\omega-2)<2+\eps$, as $\omega<3$. The cap at $\alpha$ appears only in the limit $\kappa\to0$, that is, for methods that work at every~$\kappa$. Theorem~\ref{cp:A10} below indeed passes $\alpha(\Pzero)$ at $\kappa=1$ and, with the published bounds of~\cite{VXXZ24}, also the published bound on~$\alpha$.
\end{remark}

\subsection{Fast rectangular bottoms}\label{cp:sec:range-fast}

\AlgD{} computes every needed sub-entry $p_l=(X_lY_l)[r,s]$ as an inner product of length $D_A$. A term $l$ that is needed at many positions is cheaper to handle by multiplying its whole sub-instance $X_lY_l$ with a fast rectangular algorithm, even one that is polynomially slower than $N_{\mathrm{sub}}^2$. This subsection makes that choice per term (Lemma~\ref{cp:A9}), with a fast rectangular product as the bottom (Lemma~\ref{cp:A10:bottom}), and converts the resulting costs, which contain powers of $N_{\mathrm{sub}}$, into savings in powers of $D$ by tiling the output (Lemma~\ref{cp:A10:tiling}). The result is Theorem~\ref{cp:A10}; its instances with the published bounds of~\cite{VXXZ24} are Corollary~\ref{cp:A11}.

\begin{lemma}[per-term choice]\label{cp:A9}
Let $B$ be an exact top algorithm for $\mat kbk$ with a partition (Definition~\ref{cp:T1}), and consider the composite of \AlgD{} with top $B^{\otimes n}$: sub-instances $P_l=X_lY_l$ of shape $N_{\mathrm{sub}}\times D_A\times N_{\mathrm{sub}}$ for $l\in[\RB]^n$, the needed sets $T_{rs}$ of Section~\ref{cp:sec:shared}, and $W_l:=\{(r,s):l\in T_{rs}\}$. Suppose that each term $l$ may use one of two options:
\begin{itemize}
\item[(F)] compute the whole product $P_l$ at cost at most $N_{\mathrm{sub}}^2\varphi$, after which each of its needed entries costs at most $D_A^{q_F}$;
\item[(S)] compute each of its needed entries directly, as an inner product, at cost at most $D_A$ each.
\end{itemize}
If every term takes the option with the smaller of the two bounds $N_{\mathrm{sub}}^2\varphi$ and $\abs{W_l}D_A$, then for every $W$ and every $t\in[0,1]$ the total cost of the sub-entries is at most
\[
N_{\mathrm{sub}}^2\,\varphi^{1-t}D_A^{\,t}\,\bar w^{\,t}\,\ZB(t)^n+S\,D_A^{q_F},\qquad\bar w=\abs W/N_{\mathrm{sub}}^2,\quad S=\sum_{(r,s)}\abs{T_{rs}} .
\]
The choice needs $W$ (the algorithm is offline in $W$), and the numbers $\abs{W_l}$ are computed in $O(S)$ time once the needed sets are built, with one direct-address table of size $\RB^n$.
\end{lemma}

\begin{proof}
Group the terms by class tuples $C=(c_1,\dots,c_n)$, $\abs C=\prod_a\abs{P_{c_a}}$. As in the proof of Lemma~\ref{cp:T2}, a wanted block $o$ has at most $\prod_a\kappa_{c_a}$ covering terms in $C$, so $\abs{T_{rs}\cap C}\le\abs{W_{rs}}\prod_a\kappa_{c_a}$. Counting the pairs (term of $C$, position at which it is needed) in two ways,
\[
\sum_{l\in C}\abs{W_l}=\sum_{(r,s)}\abs{T_{rs}\cap C}\le\abs W\prod_a\kappa_{c_a}=N_{\mathrm{sub}}^2\,\bar w\prod_a\kappa_{c_a}.
\]
Term $l$ costs at most $\min(N_{\mathrm{sub}}^2\varphi,\abs{W_l}D_A)+\abs{W_l}D_A^{q_F}$. A sum of minima is at most the minimum of the sums, so, using $\min(A,B)\le A^{1-t}B^t$, the first parts over $l\in C$ sum to at most
\[
\min\Bigl(\abs CN_{\mathrm{sub}}^2\varphi,\ D_AN_{\mathrm{sub}}^2\bar w\prod_a\kappa_{c_a}\Bigr)\le N_{\mathrm{sub}}^2(\abs C\varphi)^{1-t}\Bigl(D_A\bar w\prod_a\kappa_{c_a}\Bigr)^t .
\] Summed over all tuples, $\sum_C\abs C^{1-t}\prod_a\kappa_{c_a}^t=\ZB(t)^n$. The second parts sum to $\sum_l\abs{W_l}D_A^{q_F}=SD_A^{q_F}$. Nothing here uses $\varphi\le1$, so a bottom slower than $N_{\mathrm{sub}}^2$ is allowed. For the last sentence, increment a counter of $l$ for every element of every $T_{rs}$, in a table of the size of the tables of Lemma~\ref{cp:T4}.
\end{proof}

\begin{lemma}[fast rectangular bottoms]\label{cp:A10:bottom}
Let $0<e\le1$ and $g\ge0$ with $\omega(1,e,1)\le2+g$, and let $\delta>0$. There are a constant $C$ and an algorithm that, for all $N\ge1$ and $1\le D'\le N^e$, computes the product of $X\in\Z^{N\times D'}$ and $Y\in\Z^{D'\times N}$ with entries of absolute value $N^{O(1)}$ in at most $CN^{2+g+\delta}$ word-RAM steps; the transposed map $\Omega\mapsto\Omega Y^{\mathsf T}$, for $\Omega\in\Z^{N\times N}$ with entries of absolute value $N^{O(1)}$, is computed in the same time. For each $(N,D')$ the algorithm is a bilinear straight-line program over $\Q$, correct as a polynomial identity.
\end{lemma}

\begin{proof}
By the definition of $\omega(1,e,1)$ (Section~\ref{cp:sec:laser}; the rectangular exponents over $\Q$ and over $\mathbb C$ agree, by the argument of~\cite{Sch81}) there is $k\ge2$ with $R(\mat k{\lceil k^e\rceil}k)\le k^{2+g+\delta/2}$; fix an exact decomposition $B_0$ over $\Q$ with $r\le k^{2+g+\delta/2}$ terms, scaled to integer coefficients. Given $N$ and $D'$, let $n=\lceil\log_kN\rceil$, so that $k^n<kN$ and $\lceil k^e\rceil^n\ge k^{en}\ge N^e\ge D'$, and pad the product with zeros to $\mat{k^n}{\lceil k^e\rceil^n}{k^n}$. Evaluate $B_0^{\otimes n}$ recursively, one coordinate at a time. Level $j$ of the recursion handles $r^j$ products of size $k^{n-j}\times\lceil k^e\rceil^{n-j}\times k^{n-j}$ with $O(k\lceil k^e\rceil)$ operations per block entry; since $r\ge k^2\ge k\lceil k^e\rceil$, the total is $O(nr^n)\le O\bigl(n(kN)^{2+g+\delta/2}\bigr)\le CN^{2+g+\delta}$. The integers stay $N^{O(1)}$, and the final division by the scaling constant is exact. The transposed algorithm computes the $Y$-encodings as before, applies the transposed decoding to $\Omega$, multiplies by the $Y$-encodings and applies the transposed $X$-encoding; each is a Kronecker product of $n$ fixed matrices, applied one coordinate at a time with the transposed coefficient matrices, so it has the same cost. The choice of $k$ and $B_0$ depends on $(e,g,\delta)$ only. No continuity of $\omega(1,\cdot,1)$ is used.
\end{proof}

\begin{lemma}[tiling]\label{cp:A10:tiling}
Let $\eps\in(0,1)$ and $\kappa\in(0,1]$. Suppose there are an algorithm, constants $C$ and $\gamma_0>0$, an integer $J$ and exponents $t_j\in[0,1]$ ($j=0,\dots,J$) with the following property: for all $n$ and $D$ with $2\le D\le n^\eps$ and $n\le2D^{1/\eps}$ there are numbers $C_0=C_0(n,D)\ge0$ and $A_j=A_j(n,D)\ge0$ such that, for every $W'\subseteq[n]^2$ and all $X'\in\{0,1\}^{n\times D}$, $Y'\in\{0,1\}^{D\times n}$, the algorithm computes $(X'Y')|_{W'}$ in at most
\[
f_{n,D}(\abs{W'}):=C_0+\sum_jA_j\abs{W'}^{t_j}
\]
steps, and $f_{n,D}(n^2D^{-\kappa})\le Cn^2D^{-\gamma_0}$. Then for all $N_1,N_2\ge1$, $2\le D\le\min(N_1,N_2)^\eps$ and every $W\subseteq[N_1]\times[N_2]$ with $\abs W\le N_1N_2/D^\kappa$, the entries $(XY)|_W$ of $X\in\{0,1\}^{N_1\times D}$, $Y\in\{0,1\}^{D\times N_2}$ are computed in $O(N_1N_2D^{-\gamma})$ steps, with $\gamma=\min(\gamma_0,\kappa,2/\eps,1/\eps-1)$. In particular $\Hypk\kappa(\eps,\gamma)$ holds. The same holds for the transposed map $\Omega\mapsto\Omega Y^{\mathsf T}$, $\Omega$ supported on $W$, if the per-tile algorithm computes $\Omega'Y'^{\mathsf T}$ within a bound $f_{n,D}(\abs{W'})$ of the same form.
\end{lemma}

\begin{proof}
$W$ is given as a list of positions, as in Definition~\ref{cp:P1}.
\emph{Tile side.} Put $n=\lceil D^{1/\eps}\rceil$. Then $D\le n^\eps$, $n\le D^{1/\eps}+1\le2D^{1/\eps}$, and $n\le\min(N_1,N_2)$ because $D^{1/\eps}\le\min(N_1,N_2)$ and the $N_i$ are integers. For rational $\eps=p/q$, $n$ is the least integer with $n^p\ge D^q$ and is found by binary search with $O(\log N)$ arithmetic operations on numbers $N^{O(1)}$.
\emph{Padding.} Put $N_i'=n\lceil N_i/n\rceil\le2N_i$ and pad $X$ with zero rows and $Y$ with zero columns; $W$ is unchanged. There are $T=(N_1'/n)(N_2'/n)\le4N_1N_2n^{-2}\le4N_1N_2D^{-2/\eps}$ tiles.
\emph{Distribution of $W$.} Give each $(i,j)\in W$ the key $(\lfloor i/n\rfloor,\lfloor j/n\rfloor)$ and the local coordinates $(i\bmod n,\,j\bmod n)$. A counting sort on the $T$ keys produces the list $W_{\mathrm{tile}}$ of every tile in $O(T+\abs W)$ time and space; $T\le4N_1N_2D^{-\gamma}$ since $\gamma\le2/\eps$, and $\abs W\le N_1N_2D^{-\gamma}$ since $\gamma\le\kappa$. Returning the answers in the input order of $W$ is the inverse permutation, $O(\abs W)$.
\emph{Per tile.} Tile $(I,J)$ is the instance formed by rows $In,\dots,In+n-1$ of $X$, columns $Jn,\dots,Jn+n-1$ of $Y$, and $W_{\mathrm{tile}}$; its inner side is $D\le n^\eps$ and $n\le2D^{1/\eps}$, so it costs at most $f_{n,D}(\abs{W_{\mathrm{tile}}})$. The bound holds for every $W_{\mathrm{tile}}$, also for a tile that is entirely wanted.
\emph{Summing.} $f_{n,D}$ is concave and non-decreasing on $[0,\infty)$, and all tiles have the same $n$ and $D$, so by Jensen's inequality $\sum_{\mathrm{tiles}}f_{n,D}(\abs{W_{\mathrm{tile}}})\le Tf_{n,D}(\abs W/T)$. Here $\abs W/T\le N_1N_2D^{-\kappa}n^2/(N_1'N_2')\le n^2D^{-\kappa}$, so the sum is at most $TCn^2D^{-\gamma_0}=CN_1'N_2'D^{-\gamma_0}\le4CN_1N_2D^{-\gamma_0}$.
\emph{Total.} With the reading of $X$ and $Y$, $O((N_1+N_2)D)\le O(N_1N_2D^{1-1/\eps})$, all steps cost $O(N_1N_2D^{-\gamma})$, for every $D\ge2$. For the transposed map, the tile $(I,J)$ contributes $\Omega_{IJ}Y_J'^{\mathsf T}$ to the rows $In,\dots,In+n-1$ of the $N_1\times D$ output; adding these contributions costs $O(TnD)\le O(N_1N_2D^{1-1/\eps})$, and the output itself has $N_1D\le N_1N_2D^{-(1/\eps-1)}$ entries.
\end{proof}

Integer entries of absolute value $N^{O(1)}$ are handled in the same way, as in Definition~\ref{cp:P1}.

Write $a:=\alpha(\Pzero)$ and let $\Lamp$ be the function of Definition~\ref{cp:T1} for the input family of $\Pzero$ (Lemma~\ref{cp:L5}). For $e\in(a,1)$, $g\ge0$ and $\kappa\in(0,1]$ put $\varphi_e(x)=(1+ax)/(x+1/e)$ and, for $t\in(0,1)$ with $s:=t\kappa$ and $as>\Lamp(t)$,
\[
\xmin(t)=\frac{(t-s)+(1-t)g/e}{as-\Lamp(t)},\qquad F(\kappa;e,g):=\sup\bigl\{\varphi_e(\xmin(t)):t\in(0,1),\ at\kappa>\Lamp(t)\bigr\}.
\]
The set of these $t$ is not empty, since $\Lam(t)/t\to0$ (Lemma~\ref{cp:L8}) and $\Lamp=\Lam$ for small~$t$. Since $e>a$, $\varphi_e$ is strictly decreasing on $[0,\infty)$ from $\varphi_e(0)=e$ towards $a$, so $a<F(\kappa;e,g)\le e$.

\begin{theorem}[fast rectangular bottom]\label{cp:A10}
Let $e\in(\alpha(\Pzero),1)$ and $g\ge0$ with $\omega(1,e,1)\le2+g$. Then for every $\kappa\in(0,1]$ and every $0<\eps<F(\kappa;e,g)$ there is $\gamma>0$ such that $\Hypk\kappa(\eps,\gamma)$ holds (time model; the algorithm is deterministic). At $\kappa=1$,
\[
F(1;e,g)=\frac{\bar Ce+ag}{\bar C+g},\qquad\bar C:=\sup_{0<t<1}\frac{at-\Lamp(t)}{1-t}\in(0,\infty],
\]
read as $e$ if $\bar C=\infty$; so $F(1;e,0)=e$. Consequently $\Hypk1(\eps,\gamma)$ holds with some $\gamma>0$ for every $\eps<\alpha$, whatever the value of $\alpha$ is.
\end{theorem}

The algorithm is \AlgD{} with top algorithm $B$ from the level-1 family of $\Pzero$, except that every term chooses between its inner products and the fast bottom of Lemma~\ref{cp:A10:bottom} as in Lemma~\ref{cp:A9}; it runs on the tiles of Lemma~\ref{cp:A10:tiling}. Its saving comes from the sparse terms, through the factor $D_A^t\bar w^t\ZB(t)^n$; the dense terms pay $\varphi^{1-t}$, which is an excess, not a saving, and this excess is what the margin of $\xmin(t)$ pays for.

\begin{proof}
\emph{Reduction and choice of $x$.} For $\eps\le a$ Theorem~\ref{cp:A1} applies, so let $a<\eps<F(\kappa;e,g)$; since $\Hypk\kappa(\eps',\gamma)$ implies $\Hypk\kappa(\eps,\gamma)$ for $\eps\le\eps'$, we may take $\eps$ rational, as the tile side of Lemma~\ref{cp:A10:tiling} needs. Choose $t\in(0,1)$ with $as>\Lamp(t)$, $s=t\kappa$, and $\eps<\varphi_e(\xmin(t))$. Put
\[
x:=\frac{1-\eps/e}{\eps-a}>0,\qquad\text{so that}\qquad\varphi_e(x)=\eps .
\]
Since $\varphi_e$ is strictly decreasing, $x>\xmin(t)$, that is
\[
m_{\mathrm{sub}}:=x\bigl(as-\Lamp(t)\bigr)-(t-s)-(1-t)\frac ge>0 .
\]
Choose $\delta>0$ with $(1-t)\delta/e<m_{\mathrm{sub}}/3$, and $\eta>0$ with $x\eta<m_{\mathrm{sub}}/3$ and $\eta<(1/e-1)/(2x)$. Let $B$ be the top algorithm of Corollary~\ref{cp:T5} for the family of $\Pzero$, this $t$ and this $\eta$: $\alB=a$ (by Lemma~\ref{cp:L5}, $\alpha_M=\alpha(\Pzero)$ for every $M$), $\delB\le\eta$ and $\lamB(t)\le\Lamp(t)+\eta$. Let the bottom be the algorithm of Lemma~\ref{cp:A10:bottom} for $(e,g,\delta)$. The order of the choices is $\eps,\kappa$; then $t$, $x$; then $\delta$, $\eta$; then $B$ and the bottom; the algorithm then runs for all sizes. (Choosing $x$ with $\varphi_e(x)=\eps$ exactly, and not merely $\varphi_e(x)>\eps$, matters: it makes the bottom's range $D_A\le N_{\mathrm{sub}}^e$ tight, and only then is the excess of the bottom $(g+\delta)/e$ in powers of $D_A$.)

\emph{One tile.} Consider an instance $n\times D\times n$ with $2\le D\le n^\eps$ and $n\le2D^{1/\eps}$, so $D^{1/\eps}\le n\le2D^{1/\eps}$, and any $W'\subseteq[n]^2$. Run Steps~1--3 of \AlgD{} with top $B$ and the above $x$: $K$ and $D_A$ lie within constant factors of $D^{x/(1+ax)}$ and $D^{1/(1+ax)}$, and $N_{\mathrm{sub}}=n'/K$ with $n\le n'\le2n$. Since $1/\eps=(x+1/e)/(1+ax)$,
\[
D^{1/(e(1+ax))}=D^{1/\eps-x/(1+ax)}\le N_{\mathrm{sub}}\le4kD^{1/(e(1+ax))},
\]
so $D_A\le2bD^{1/(1+ax)}\le2bN_{\mathrm{sub}}^e$. In Step~4, a term $l$ with option (F) cuts the inner side of $P_l$ into at most $4b+1$ pieces of length at most $\lfloor N_{\mathrm{sub}}^e\rfloor$ and multiplies by the bottom, at cost $N_{\mathrm{sub}}^2\varphi$ with $\varphi=O(N_{\mathrm{sub}}^{g+\delta})=O(D^{(g+\delta)/(e(1+ax))})$, and then reads its needed entries ($q_F=0$); a term with option (S) computes its needed entries as inner products. The choice is that of Lemma~\ref{cp:A9}. The values $p_l=(X_lY_l)[r,s]$ are then decoded by Lemma~\ref{cp:T4}, and correctness is that of Theorem~\ref{cp:T6}, since every $p_l$ is computed exactly either way. The cost lines, in powers of $D$ relative to $n^2$ at $\abs{W'}\le n^2D^{-\kappa}$ (all constants depend on $B$, $x$, $\delta$ and the bottom only):
\begin{itemize}
\item \emph{Encodings} (Step~2), $O(n\log D\cdot K^{1+\delB}D_A)$: exponent $\bigl(1+x(1+\delB)\bigr)/(1+ax)-1/\eps=(1-1/e+x\delB)/(1+ax)<0$, because $\delB\le\eta<(1/e-1)/x$; here $e<1$ is used. The tables of Lemma~\ref{cp:T4} and of Lemma~\ref{cp:A9}, $O(\log D\cdot K^{2+\delB})$, are at most this, since $K\le n$.
\item \emph{Bucketing} (Step~3), $O(N_{\mathrm{sub}}^2+\abs{W'})$: savings $2x/(1+ax)$ and $\kappa$.
\item \emph{Sub-entries}, by Lemma~\ref{cp:A9} with $\bar w\le K^2D^{-\kappa}$ (up to a constant) and $\ZB(t)^n=K^{\lamB(t)+2-2t}$: the first term is $O(n^2D^{-e_{\mathrm{sub}}})$ with the saving
\[
\begin{split}
e_{\mathrm{sub}}&=\frac{x(as-\lamB(t))-(t-s)-(1-t)(g+\delta)/e}{1+ax}\\
&\ge\frac{m_{\mathrm{sub}}-x\eta-(1-t)\delta/e}{1+ax}>\frac{m_{\mathrm{sub}}}{3(1+ax)} .
\end{split}
\]
\item \emph{Needed entries and decoding}: $O(S\log D)$ with $S\le n'^2D^{-t\kappa}K^{\lamB(t)}$ (Lemma~\ref{cp:T2}(b)), saving $\bigl(s+x(as-\lamB(t))\bigr)/(1+ax)>0$.
\item \emph{Input and output}: $O(nD+\abs{W'})$, savings $1/\eps-1>0$ and $\kappa$.
\item \emph{Word size}: all numbers are $N^{O(1)}$, as in Theorem~\ref{cp:T6} and Lemma~\ref{cp:A10:bottom}.
\end{itemize}
As a function of $\abs{W'}$ each line is a coefficient, or a coefficient times $\abs{W'}^{t'}$ with $t'\in\{t,1\}$, and the bounds of Lemmas~\ref{cp:T2}(a) and~\ref{cp:A9} hold for every $W'$, not only for $\abs{W'}\le n^2D^{-\kappa}$. So the cost of a tile is at most $f_{n,D}(\abs{W'})$ of the form of Lemma~\ref{cp:A10:tiling}, whose coefficients depend on $n$, $D$ and the fixed parameters only, with $f_{n,D}(n^2D^{-\kappa})\le Cn^2D^{-\gamma_0}$ for a constant $C$ and some $\gamma_0>0$ (the factors $\log D$ are absorbed by a slightly smaller $\gamma_0$). For $D$ below a constant $D_0$ the entries are computed directly.

\emph{All sizes.} Lemma~\ref{cp:A10:tiling} gives $\Hypk\kappa(\eps,\gamma)$ with $\gamma=\min(\gamma_0,\kappa,2/\eps,1/\eps-1)>0$.

\emph{The case $\kappa=1$.} Then $s=t$, and for $t\in(0,1)$ with $at>\Lamp(t)$, $\xmin(t)=(1-t)g/(e(at-\Lamp(t)))=g/(eC_t)$ with $C_t=(at-\Lamp(t))/(1-t)>0$, and
\[
\varphi_e\Bigl(\frac g{eC_t}\Bigr)=\frac{e(1+ag/(eC_t))}{1+g/C_t}=\frac{C_te+ag}{C_t+g},
\]
which is non-decreasing in $C_t$ because $e>a$. It tends to $e$ as $C_t\to\infty$. Taking the supremum over $t$ gives the formula for $F(1;e,g)$, with $\bar C>0$ because the set of admissible $t$ is not empty. With $g=0$ the value is $e$ for every admissible $t$. Finally, let $\eps<\alpha$. If $\eps\le a$, Theorem~\ref{cp:A1} applies. Otherwise choose $e$ with $\eps<e<\min(\alpha,1)$; then $e>a$, and $\omega(1,e,1)=2$ because $e<\alpha$ and $\omega(1,\cdot,1)$ is non-decreasing; so $g=0$ is allowed and $\eps<e=F(1;e,0)$.
\end{proof}

\begin{corollary}[instances with the published bounds of~\cite{VXXZ24}]\label{cp:A11}
In the time model, by Theorem~\ref{cp:A10} with the top algorithm from the level-1 family of $\Pzero$, and with some $\gamma>0$ in each case:
\begin{enumerate}[label=(\alph*)]
\item since $\alpha>\AlphaVXXZ$~\cite{VXXZ24}: $\Hyp(\eps,\gamma)$ holds for every $\eps\le\cpv{A11_vxxz_half}$, and $\Hypk1(\eps,\gamma)$ for every $\eps<\AlphaVXXZ$;
\item since $\omega(1,\cpv{A11_tab1_e_half},1)\le\cpv{A11_tab1_om_half}$~\cite[Table~1]{VXXZ24}: $\Hyp(\eps,\gamma)$ holds for every $\eps\le\cpv{A11_tab1_half}$;
\item since $\omega(1,\cpv{A11_tab1_e_one},1)\le\cpv{A11_tab1_om_one}$~\cite[Table~1]{VXXZ24}: $\Hypk1(\eps,\gamma)$ holds for every $\eps\le\cpv{A11_tab1_one}$.
\end{enumerate}
\end{corollary}

\begin{proof}
The second half of (a) is the last statement of Theorem~\ref{cp:A10}. For the rest it suffices, by Theorem~\ref{cp:A10}, to show $v<F(\kappa;e,g)$ for the printed value $v$, with $(\kappa,e,g)=(\frac12,\AlphaVXXZ,0)$ in (a), $(\frac12,\cpv{A11_tab1_e_half},\cpv{A11_tab1_g_half})$ in (b) and $(1,\cpv{A11_tab1_e_one},\cpv{A11_tab1_g_one})$ in (c); in (a), $\omega(1,\AlphaVXXZ,1)=2$ because $\AlphaVXXZ<\alpha$. The certificate evaluates $\varphi_e(\xmin(t))$ in exact rational arithmetic at one point $t$ of the grid of Table~\ref{cp:tab:B2} ($t=\cpv{A11_t_half}$, $\cpv{A11_t_tab_half}$ and $\cpv{A11_t_one}$ respectively), with $a$ replaced by a certified lower bound $a'<a$ (the lower end of the enclosure of $\alpha(\Pzero)$ recorded with the certified grid, which lies below the enclosure of Lemma~\ref{cp:L11}, lowered by $\cpv{cert_tol}$) and $\Lamp(t)$ replaced by $M'_t=M_t+\cpv{cert_tol}$, where $M_t\ge\Lamp(t)$ is the certified value of Table~\ref{cp:tab:B2}. This gives a lower bound on $F$, because $\varphi_e(\xmin(t))$ increases when $a$ grows or $\Lamp(t)$ falls: $\xmin(t)$ then decreases, $\varphi_e$ is decreasing in $x$, and $\varphi_e$ increases with $a$. The values obtained exceed $v$ in all three cases. The script also checks every cost line of the proof of Theorem~\ref{cp:A10} directly at $\eps=v$ with the value of $x$ chosen there, and all margins are positive (Section~\ref{cp:sec:certificates}).
\end{proof}

Each printed value is the six-digit floor of the grid value: $v+\cpv{cert_step}$ is not certified at any grid point. In (a) the certificate does not use the strict inequality $\alpha>\AlphaVXXZ$: with $e=\AlphaVXXZ-\cpv{cert_tol_e}$ it gives the same printed value, so the reading $\alpha\ge\AlphaVXXZ$ of~\cite[Table~1]{VXXZ24} suffices. At $\kappa=\frac12$ the best grid point is the smallest one, so a finer grid near $t=0$ could raise the value slightly; $\cpv{A11_vxxz_half}$ is a certified lower bound, not the supremum of $F(\frac12;\AlphaVXXZ,0)$. At $\kappa=1$ the closed form of Theorem~\ref{cp:A10} gives the same value $\cpv{A11_tab1_one}$ in (c), with $\bar C\ge\cpv{A10_C}$ at $t=\cpv{A11_t_one}$ on the grid (certified: the script evaluates $(a't-M'_t)/(1-t)$ in exact rational arithmetic), and $\cpv{A11_grid_pos}$ of the $\cpv{A11_grid_n}$ grid points have $a't>M'_t$. Instance (a) rests on the same published input as the instance of Section~\ref{cp:C16} (from Proposition~\ref{cp:C14}), and its value is larger.

\begin{remark}[other inputs, and $\kappa$ slightly above $1$]\label{cp:A11:rem}
Theorem~\ref{cp:A10} transfers every future bound on any $\omega(1,e,1)$, and on $\alpha$, with no new proof. With Le Gall and Urrutia's earlier bound $\alpha>\cpv{alpha_lgu}$~\cite{LGU18} the grid value at $\kappa=\frac12$ is $\cpv{A11_lgu_half}$ (computed, not certified; superseded by Corollary~\ref{cp:A11}). Under the bound $\alpha\ge\cpv{alpha_preprints}$ claimed in two preprints of October 2026~\cite{AvdBZ26,Uzn26}, Theorem~\ref{cp:A10} gives $\Hypk1(\eps,\gamma)$ for every $\eps<\cpv{alpha_preprints}$, and the grid value at $\kappa=\frac12$ is $\cpv{A11_pre_half}$ (computed, not certified); both are conditional on that claim. For $\kappa$ slightly above~$1$ the same construction, with the bucketing done by a radix sort of $W$ and the encodings measured against $\abs WD$, gives positive savings on the strip $\kappa\in(1,\cpv{A10_strip})$ in computations on the certified grid; this strip is computed and argued, not proved, and no statement of this paper uses it. It is a different statement from Corollary~\ref{cp:B7}(a), which certifies the $\kappa$-batch statement with inner products as the bottom for every $\kappa\le\cpv{u8_plateau_end}$.
\end{remark}

\section{Explicit \texorpdfstring{$\gamma$}{gamma} and the product curve}\label{cp:sec:gamma}

\ThmDk{} gives, at every admissible level-1 point, some saving $\gamma>0$ for every $\eps\le\alpha(P)$; it does not say how large $\gamma$ is. This section makes the saving explicit at $\kappa=\frac12$, that is, for $\Hyp(\eps,\gamma)$. It has three parts. First, a bound on the profile of the level-1 family that is much sharper than the closed form $\Lampr$ of Lemma~\ref{cp:L8}, together with a certificate that evaluates it rigorously (Section~\ref{cp:sec:sharp}); with it \ThmDk{} gives an explicit $\gamma$ at $\Pzero$ and at level-1 points over other sets of parts (Section~\ref{cp:sec:explicit}). Second, \ThmDprod{}, the composite whose top algorithm is the tensor product of a level-1 top algorithm with the tile of~\cite{v2} (Section~\ref{cp:sec:prod}). Third, a table of certified points of the resulting curve $\gamma(\eps)$, and what the products do not change (Section~\ref{cp:sec:curve}). Every statement of this section is in the time model and deterministic. Every saving in it comes from an upper bound on a profile, so it is a lower bound on the saving of the construction, never a claim that the construction saves no more.

\subsection{The sharp level-1 profile}\label{cp:sec:sharp}

Fix an admissible rational level-1 point $P$ and its input family $(A_M)$ (Section~\ref{cp:sec:family}), and fix one part $q$ of $P$; write $\pi=1/(q+2)$, which is $w_q$ of Definition~\ref{cp:L3}. On the $M_q$ coordinates of part $q$ a term string has $n_A$ letters among $A_1,\dots,A_q$, $n_B$ letters $\Bp$ and $n_O$ letters $\Oterm$; its \emph{frequencies} are $x=(x_A,x_B,x_O)=(n_A,n_B,n_O)/M_q$. On part $q$ the kept output type $\Ukept$ has $n_0:=M_q\pi$ letters $z_0$, $n_1:=qM_q\pi$ letters in $\{z_1,\dots,z_q\}$ and $n_2:=M_q\pi$ letters $z_{q+1}$.

\begin{lemma}[exact support]\label{cp:G1:support}
Let $\tau$ be a term string, and let $(n_A,n_B,n_O)$ be the numbers of its letters $A_\cdot$, $\Bp$, $\Oterm$ on part $q$. The number of $z$-strings of part $q$ of the kept type whose letters lie in the supports of the letters of $\tau$ is
\[
S^{(q)}(\tau)=[n_O\le n_0]\,[n_B\ge n_2]\binom{n_B}{n_2}\sum_{j\ge0}\binom{n_B-n_2}{j}q^{\,n_B-n_2-j}\binom{n_A}{n_0-n_O-j},
\]
and $\Sall(\tau)=\prod_qS^{(q)}(\tau)$.
\end{lemma}

\begin{proof}
By Lemma~\ref{cp:L2} the supports of $A_i$, $\Bp$ and $\Oterm$ are $\{z_0,z_i\}$, all letters, and $\{z_0\}$. So a $z$-string $o$ of the kept type in the support has $o_\ell=z_0$ at every $\Oterm$-coordinate, $o_\ell\in\{z_0,z_i\}$ at an $A_i$-coordinate, and any letter at a $\Bp$-coordinate. Its $n_2$ letters $z_{q+1}$ can sit only at $\Bp$-coordinates: $\binom{n_B}{n_2}$ choices, none if $n_B<n_2$. Its $n_0$ letters $z_0$ include all $n_O$ letters at $\Oterm$-coordinates (none if $n_O>n_0$); the other $n_0-n_O$ sit at $j$ of the remaining $n_B-n_2$ $\Bp$-coordinates and at $n_0-n_O-j$ of the $A$-coordinates. Every remaining coordinate gets a letter in $\{z_1,\dots,z_q\}$: any of $q$ at a $\Bp$-coordinate, only $z_i$ at an $A_i$-coordinate. Multiplying the choices and summing over $j$ gives the formula. Since $\Ukept$ is a product over the parts and $\supp w_\tau=\prod_\ell\supp w_{\tau_\ell}$, $\Sall$ is the product of the per-part counts.
\end{proof}

Write $\varphi(s)=s\ln s$ with $\varphi(0)=0$, and let $\Hb$ be the binary entropy in nats. For real $d_B,d_O,y$ put
\[
x_B=\pi+d_B,\qquad x_O=\pi-d_O,\qquad x_A=q\pi-d_B+d_O,
\]
\[
\KL(x\,\|\,x^*)=x_A\ln\frac{x_A}{q\pi}+x_B\ln\frac{x_B}{\pi}+x_O\ln\frac{x_O}{\pi},\qquad x^*=(q\pi,\pi,\pi),
\]
\[
\sigma_q(d_B,d_O,y)=x_B\Hb\Bigl(\frac\pi{x_B}\Bigr)+d_B\Hb\Bigl(\frac y{d_B}\Bigr)+(d_B-y)\ln q+x_A\Hb\Bigl(\frac{d_O-y}{x_A}\Bigr)
\]
(with $w\Hb(s/w):=0$ for $w=0$), and on the polytope
\[
\Polq_q=\bigl\{(d_B,d_O,y):\ 0\le d_O\le\pi,\ 0\le y\le d_O,\ y\le d_B,\ d_B-y\le q\pi\bigr\}
\]
define, for $t\in(0,1]$,
\[
\Theta_q(t;d_B,d_O,y)=-\KL(x\,\|\,x^*)+t\,\sigma_q(d_B,d_O,y),\qquad\Theta^*_q(t)=\max_{\Polq_q}\Theta_q(t;\cdot),
\]
\[
\Lamcert(t)=\frac{\sum_q\lambda_q\Theta^*_q(t)}{\ln a_1}.
\]
With $w\Hb(s/w)=\varphi(w)-\varphi(s)-\varphi(w-s)$ for $0\le s\le w$,
\[
\sigma_q=\varphi(x_B)-\varphi(\pi)-\varphi(y)-\varphi(d_B-y)+(d_B-y)\ln q+\varphi(x_A)-\varphi(d_O-y)-\varphi(q\pi-d_B+y),
\]
the two terms $\varphi(d_B)$ cancelling. On $\Polq_q$ every argument of $\varphi$ is nonnegative ($x_A\ge q\pi-d_B+y\ge0$, $x_O\ge0$, $x_B\ge\pi$, $d_B-y\ge0$, $d_O-y\ge0$), and $\Polq_q\subseteq[0,1]^3$, since $d_B\le q\pi+y\le(q+1)\pi<1$.

\begin{lemma}[sharp profile]\label{cp:G1}
For the family of an admissible rational level-1 point $P$, every multiple $M$ of $d(P)$ and every $t\in(0,1]$,
\[
\frac1M\ln\mathbb E_\tau\,\Sall(\tau)^t\le\sum_q\lambda_q\Theta^*_q(t)+\frac{3\abs\Qs\ln(M+1)}M,
\]
with $\tau$ uniform over all term strings. Hence $\Lam(t)\le\Lamcert(t)$ for every $t\in(0,1]$.
\end{lemma}

\begin{proof}
Fix a part $q$ and write $M$ for $M_q$. For uniform $\tau$ the letters of part $q$ are independent and uniform over the $q+2$ letters, so the frequencies $x$ of $\tau$ follow the law of the type of $M$ independent letters of law $x^*$ over the three kinds $A_\cdot$, $\Bp$, $\Oterm$; by the type-class bound a given $x$ has probability at most $\exp(-M\KL(x\,\|\,x^*))$, and there are at most $(M+1)^2$ values of $x$. By Lemma~\ref{cp:G1:support}, $S^{(q)}(\tau)$ depends only on $x$, and it is $0$ unless $n_O\le n_0$ and $n_B\ge n_2$, that is, unless $d_O:=\pi-x_O\ge0$ and $d_B:=x_B-\pi\ge0$; also $d_O\le\pi$ since $n_O\ge0$. Otherwise $S^{(q)}(\tau)$ is a sum of at most $M+1$ terms indexed by $j$, and the term $j$ is nonzero only if $0\le j\le n_B-n_2$, $j\le n_0-n_O$ and $n_0-n_O-j\le n_A$; with $y:=j/M$ these conditions read $0\le y\le d_B$, $y\le d_O$ and $d_O-y\le x_A$, the last one being $d_B-y\le q\pi$. So every nonzero term has $(d_B,d_O,y)\in\Polq_q$. By $\binom nk\le\exp(n\Hb(k/n))$ its four factors are at most $\exp(Mx_B\Hb(\pi/x_B))$, $\exp(Md_B\Hb(y/d_B))$, $\exp(M(d_B-y)\ln q)$ and $\exp(Mx_A\Hb((d_O-y)/x_A))$, so the term is at most $\exp(M\sigma_q(d_B,d_O,y))$. Since $\KL$ depends on $(d_B,d_O)$ only and $t\le1$,
\[
\mathbb E\,S^{(q)}(\tau)^t\le(M+1)^2\,(M+1)^t\max_{\Polq_q}\exp\bigl(M[-\KL+t\sigma_q]\bigr)\le(M+1)^3\exp\bigl(M\Theta^*_q(t)\bigr).
\]
For the mixture, $\tau$ is uniform on a product over the parts and $\Sall$ is the product of the per-part counts (Lemma~\ref{cp:G1:support}), so $\mathbb E\,\Sall^t=\prod_q\mathbb E\,(S^{(q)})^t$; with $M_q=\lambda_qM\le M$ this gives the displayed bound. By Corollary~\ref{cp:L6} and $\ln a_M=M\ln a_1$,
\[
\mu_M(t)\le\delta_M+\Lamcert(t)+\frac{3\abs\Qs\ln(M+1)}{M\ln a_1}.
\]
Neither $\delta_M$ nor the last term depends on $t$, and both tend to $0$ (Lemma~\ref{cp:L5}, \Hcond2); taking the $\limsup$ over $M$ gives $\Lam(t)\le\Lamcert(t)$.
\end{proof}

Nothing in this proof restricts $t$; Section~\ref{cp:sec:batch} uses the lemma up to $t=1$. The closed form $\Lampr$ of Lemma~\ref{cp:L8} is a cruder upper bound on the same quantity; the maximum $\Theta^*_q(t)$ has to be evaluated, and the next lemma makes this rigorous.

\begin{lemma}[certificate]\label{cp:G2}
For $t>0$ the function $\Theta_q(t;\cdot)$ is concave and continuous on $\Polq_q$, and for every interior point $z^0$ of $\Polq_q$
\[
\Theta^*_q(t)\le\Theta_q(t;z^0)+\sum_{i=1}^3\abs{\partial_i\Theta_q(t;z^0)}.
\]
\end{lemma}

\begin{proof}
$-\KL$ is $-\varphi(x_A)-\varphi(x_B)-\varphi(x_O)$ plus an affine function of $(d_B,d_O)$, and $\varphi$ is convex while $x_A$, $x_B$, $x_O$ are affine; so $-\KL$ is concave. In the first form of $\sigma_q$ the terms $x_B\Hb(\pi/x_B)$, $d_B\Hb(y/d_B)$ and $x_A\Hb((d_O-y)/x_A)$ are perspectives $w\Hb(s/w)$ of the concave function $\Hb$, evaluated at affine $(s,w)$ with $0\le s\le w$; a perspective of a concave function is jointly concave, and $(d_B-y)\ln q$ is affine. So $\Theta_q(t;\cdot)$ is concave for $t>0$. It is continuous on $\Polq_q$, because the second forms of $\KL$ and $\sigma_q$ are sums of $\varphi$ at nonnegative affine arguments and affine terms. At an interior point $z^0$ all arguments of $\varphi$ are positive, so $\Theta_q$ is differentiable there, and concavity on the convex set $\Polq_q$ gives $\Theta_q(t;z)\le\Theta_q(t;z^0)+\nabla\Theta_q(t;z^0)\cdot(z-z^0)$ for every $z\in\Polq_q$. Since $\Polq_q\subseteq[0,1]^3$, $\abs{z_i-z^0_i}\le1$, which gives the claim.
\end{proof}

A certificate for $\Theta^*_q(t)$ is therefore one rational interior point $z^0$, which is found numerically and need only be close to the maximiser, and an enclosure of $\Theta_q(t;z^0)$ and of its gradient in interval arithmetic, starting from exact rationals; interiority of $z^0$ is checked exactly. The gradient is elementary. All certified values of the profile in this section and in Section~\ref{cp:sec:batch} are computed in this way; Section~\ref{cp:sec:certificates} names the scripts and the values that a second, independent implementation confirms.

\subsection{Explicit \texorpdfstring{$\gamma$}{gamma} at \texorpdfstring{$\Pzero$}{P0} and at other level-1 points}\label{cp:sec:explicit}

\begin{theorem}[explicit $\gamma$ at $\Pzero$]\label{cp:G3}
For every $0<\eps<\alpha(\Pzero)$ and every $\gamma$ with
\[
0<\gamma<\min\Bigl\{\cpv{G3_gamma_exact},\ \frac1\eps-\frac1{\alpha(\Pzero)}\Bigr\},
\]
$\Hyp(\eps,\gamma)$ holds (time model, deterministic); here $\cpv{alphaP0_lo}<\alpha(\Pzero)<\cpv{alphaP0_hi}$. In particular $\Hyp(\eps,\cpv{G3_gamma})$ holds for every $\eps\le\cpv{G3_eps}$, and for every $\eps\le\cpv{G3_eps_bind}$ the first term of the minimum is the smaller one.
\end{theorem}

\begin{proof}
Let $t_0=\cpv{G3_t}$. By Lemma~\ref{cp:L5}, $\cstar\le\co(\Pzero)$, and by Lemma~\ref{cp:L11}, $t_0<\cpv{tpadP0_lo}<2/(\co(\Pzero)+2)\le2/(\cstar+2)$; so $\Lamp(t_0)=\Lam(t_0)\le\Lamcert(t_0)$ by Corollary~\ref{cp:T5}(i) and Lemma~\ref{cp:G1}. Lemma~\ref{cp:G2}, applied at rational interior points for $q=6$ and $q=7$ and evaluated in interval arithmetic, certifies $\Lamcert(t_0)\le\cpv{u7_lamcert_t0}$ and
\[
\gamma_\infty:=\frac{t_0}2-\frac{\Lamcert(t_0)}{\alpha(\Pzero)}>\cpv{G3_gamma_exact}
\]
(upper end of $\Lamcert(t_0)$, lower end of $\alpha(\Pzero)$). By \ThmDk{}(b) with $\kappa=\frac12$, $\Hypk{1/2}(\eps,\gamma)$ holds for every $0<\eps<\alpha(\Pzero)$ and every $\gamma$ below $\min\{\sup_t(t/2-\Lamp(t)/\alpha),\,1/\eps-1/\alpha,\,2/\alpha\}$ with $\alpha=\alpha(\Pzero)$. The first term is at least $\gamma_\infty$, and the third exceeds $2>\frac12$; since every $\gamma$ of the theorem is below $\frac12$, this is $\Hyp(\eps,\gamma)$. For the particular statements: $1/\eps-1/\alpha(\Pzero)\ge1/\cpv{G3_eps}-1/\alpha(\Pzero)>\cpv{u7_G3_margin}>\cpv{G3_gamma}$ for $\eps\le\cpv{G3_eps}$; and $\eps_1:=1/(1/\alpha(\Pzero)+\gamma_\infty)$ is certified to be at least $\cpv{G3_eps_bind}$, so $1/\eps-1/\alpha(\Pzero)\ge\gamma_\infty>\cpv{G3_gamma_exact}$ for $\eps\le\cpv{G3_eps_bind}$.
\end{proof}

The value $\cpv{G3_gamma_exact}$ is the certified lower end of $\gamma_\infty$ rounded down; the floating-point optimum over $t$ of the same bound is within $10^{-8}$ of $\gamma_\infty$ at $t_0$, so a better tilt would gain nothing visible.

\begin{corollary}\label{cp:G4}
The supremum of $\eps\gamma$ over the pairs $(\eps,\gamma)$ of Theorem~\ref{cp:G3} is at least $\cpv{G4_epsgamma}$; values above it are reached at $\eps=\cpv{G3_eps_bind}$ with $\gamma$ close to $\cpv{G3_gamma_exact}$.
\end{corollary}

\begin{proof}
For $\eps\le\cpv{G3_eps_bind}$ every $\gamma<\cpv{G3_gamma_exact}$ is allowed, and the product $\cpv{G3_eps_bind}\cdot\cpv{G3_gamma_exact}$ exceeds $\cpv{G4_epsgamma}$.
\end{proof}

The tilt and the point enter the bound only through the part set and its weights. For a level-1 point $P$ write $\ln b_1(P)=\sum_q\lambda_qu_q\ln q$, so that $\alpha(P)=\ln b_1(P)/\ln a_1$.

\begin{lemma}[reduction]\label{cp:G5}
The quantity $\frac1M\ln\mathbb E_\tau\Sall(\tau)^t$ and the bound $\ln a_1\cdot\Lamcert(t)=\sum_q\lambda_q\Theta^*_q(t)$ depend on the part set $\Qs$ and the weights $(\lambda_q)$ only, not on the block parameters $(u_q)$. Consequently the first term of \ThmDk{}(b) at $\kappa=\frac12$, evaluated with $\Lamcert$ in place of $\Lamp$, is
\[
\sup_t\Bigl(\frac t2-\frac{\Lamcert(t)}{\alpha(P)}\Bigr)=\sup_t\Bigl(\frac t2-\frac{\sum_q\lambda_q\Theta^*_q(t)}{\ln b_1(P)}\Bigr),
\]
which at fixed $(\Qs,\lambda)$ increases with $\alpha(P)$.
\end{lemma}

\begin{proof}
The law of $\tau$ is uniform on all term strings, and $\Ukept$ is fixed by its per-part letter counts, which are those of the $Z$-marginal $(1,q,1)/(q+2)$ of Definition~\ref{cp:L3}; this marginal does not involve $u_q$, because $v_q=(w_q-u_q)/2$. The function $\Theta_q$ involves only $q$ and $t$, and $\ln a_1=\sum_q\lambda_qp_q\ln q$ only $(\Qs,\lambda)$. The identity follows from $\alpha(P)\ln a_1=\ln b_1(P)$.
\end{proof}

So within the level-1 family the only trade-off between range and explicit saving is the choice of the part set and its weights; the block parameters should be as large as admissibility allows. The points of Table~\ref{cp:tab:G6} below therefore take $u_q(q+2)=\cpv{u7_ubar}$ for every $q$, and four of them use other part sets than $\{6,7\}$.

\begin{remark}[other part sets]\label{cp:G6:parts}
Section~\ref{cp:sec:level1} fixes $\Qs=\{6,7\}$. Its Proposition~\ref{cp:L1} and Lemma~\ref{cp:L2} hold for every $q$; Definition~\ref{cp:L3}, Lemmas~\ref{cp:L4}--\ref{cp:L6:lemmaC}, Corollary~\ref{cp:L6}, Theorem~\ref{cp:L7}, Lemma~\ref{cp:L8} and Lemmas~\ref{cp:G1:support}--\ref{cp:G5} use the part set only through sums and products over its parts, and their proofs hold verbatim for every finite set $\Qs$ of integers $q\ge2$, rational weights $\lambda_q>0$ with $\sum_q\lambda_q=1$, and rational $0<u_q<w_q$. Lemma~\ref{cp:L9}, which needs $q\ge4$, is used only to pass from points with $\slack\ge0$ to rational admissible points, and it is not needed at a rational point whose slack is certified positive. So \ThmDk{} and the proof of Theorem~\ref{cp:G3} apply at every such point.
\end{remark}

\begin{table}[ht]
\centering\small
\begin{tabular}{l l l l}
\hline
part set $\Qs$ and weights & tilt $t$ & every $\gamma<$ & every $\eps\le$\\
\hline
$\{13,14\}$, $\lambda_{13}=\cpv{u7_G6_a_lam}$ & $\cpv{u7_G6_a_t}$ & $\cpv{G6_a_G}$ & $\cpv{G6_a_E}$\\
$\{11,12\}$, $\lambda_{11}=\cpv{u7_G6_b_lam}$ & $\cpv{u7_G6_b_t}$ & $\cpv{G6_b_G}$ & $\cpv{G6_b_E}$\\
$\{10\}$ & $\cpv{u7_G6_c_t}$ & $\cpv{G6_c_G}$ & $\cpv{G6_c_E}$\\
$\{8\}$ & $\cpv{u7_G6_d_t}$ & $\cpv{G6_d_G}$ & $\cpv{G6_d_E}$\\
$\{6,7\}$, $\lambda_6=\cpv{u7_G6_e_lam}$ & $\cpv{u7_G6_e_t}$ & $\cpv{G6_e_G}$ & $\cpv{G6_e_E}$\\
\hline
\end{tabular}
\caption{Certified level-1 points over other part sets (Theorem~\ref{cp:G6}); in every row $u_q(q+2)=\cpv{u7_ubar}$ for each $q$.}\label{cp:tab:G6}
\end{table}

\begin{theorem}[level-1 points over other part sets]\label{cp:G6}
At each of the five rational level-1 points of Table~\ref{cp:tab:G6}, $\Hyp(\eps,\gamma)$ holds (time model, deterministic) for every $\gamma$ below the third column and every $\eps$ up to the fourth. In particular $\Hyp(\eps,\cpv{G6_b_G})$ holds for every $\eps\le\cpv{G6_b_E}$, at the point with $\Qs=\{11,12\}$, $\lambda_{11}=\cpv{u7_G6_b_lam}$.
\end{theorem}

\begin{proof}
At each point the certificate checks $\slack(P)>0$ and $t\le1/(\co(P)+2)$ in interval arithmetic, so $P$ is admissible and $\Lamp(t)=\Lam(t)$ (Lemma~\ref{cp:L5} and Corollary~\ref{cp:T5}(i), read for the part set of the row by Remark~\ref{cp:G6:parts}). It certifies $\gamma_\infty=t/2-\Lamcert(t)/\alpha(P)$ to be larger than the third column and $1/(1/\alpha(P)+\gamma_\infty)$ to be at least the fourth. The proof of Theorem~\ref{cp:G3} then applies word for word.
\end{proof}

The best $\eps\gamma$ of the five rows, at least $\cpv{G6_epsgamma}$ for $\Qs=\{6,7\}$, $\lambda_6=\cpv{u7_G6_e_lam}$, is above that of $\Pzero$ (Corollary~\ref{cp:G4}); the gain of all level-1 points over $\Pzero$ is small, and the products of the next subsection gain more.

\subsection{Products with the tile: \texorpdfstring{\ThmDprod}{Theorem D\_x}}\label{cp:sec:prod}

\paragraph{The tile.} We use the squared construction of~\cite[Sections~\Vref{m:sec:tensor}--\Vref{m:sec:leaves}]{v2} with its notation, and call it the \emph{tile}. Fix a rational $c>1$; let $m$ run through the multiples of the denominator of $c$, and put $L=cm$. The tile $\Tile_{L,m}$ has $2L$ levels. At every level there are five terms $P_{11},P_{12},P_{21},P_{22},P_0$ with the output forms $\chi_{P_{ij}}=z_{ij}+z_0$ and $\chi_{P_0}=z_0$ \cite[Lemma~\Vref{m:lem:identity}]{v2}. Fixed re-indexing maps $\Gamma,\Delta$ and a list $\mathcal Q_g$ of $K_g\ge\mu_\infty^2\binom Lm$ good sets are found deterministically by~\cite[Lemma~\Vref{m:lem:hashes}]{v2}. For each run $s\in\{-m,\dots,m\}$ the input arrays $a_s$, $b_s$ are selections of entries of the input blocks \cite[Section~\Vref{m:sec:squared}]{v2}, and a \emph{leaf} is a string $\tau$ of $2L$ terms, with left and right forms $\Phi_\tau(a_s)$, $\Psi_\tau(b_s)$ and output form $\chi_\tau=\bigotimes_i\chi_{\tau_i}$ \cite[Section~\Vref{m:sec:leaves}]{v2}. Put $N_z=4^{L-m}$, $D_{\mathrm G}=4^m$ and $K_{\mathrm G}=K_gN_z$. The \emph{type class} $\UG$ is the set of output strings $w=w'w''$ with exactly $m$ letters $z_0$ in each half, so $\abs\UG=\beta_0^2$ with $\beta_0=\binom Lm4^{L-m}$, and the \emph{read set} $\RG\subseteq\UG$ consists of the strings with $Q(w'),Q(w'')\in\mathcal Q_g$, so $\abs\RG=K_{\mathrm G}^2\ge\mu_\infty^4\abs\UG$. By~\cite[Theorem~\Vref{m:thm:correct}]{v2} there is a bijection $o\mapsto w(o)$ from $[K_{\mathrm G}]^2$ onto $\RG$ such that the $2m+1$ runs together compute entry $o$ of the product of the $K_{\mathrm G}\times D_{\mathrm G}$ matrix $X_{\mathrm{band}}$, whose $k$-th row block is $X_{Q_k}$, with the $D_{\mathrm G}\times K_{\mathrm G}$ matrix $Y_{\mathrm{band}}$, whose $k'$-th column block is $Y_{Q_{k'}}$, as a single output at $w(o)$. We write $\Tile(c)$ for the sequence of tiles $\Tile_{cm,m}$.

\begin{lemma}[the tile is an exact bilinear algorithm on block entries]\label{cp:G8}
Let the entries of $X_{\mathrm{band}}$ and $Y_{\mathrm{band}}$ be matrices over $\Z$ of compatible sizes. For every run $s$ and leaf $\tau$ put $u_{s,\tau}(X_{\mathrm{band}})=\Phi_\tau(a_s)$, $v_{s,\tau}(Y_{\mathrm{band}})=\Psi_\tau(b_s)$ and $w_{s,\tau}[o]=\chi_\tau[w(o)]$ for $o\in[K_{\mathrm G}]^2$. Then
\[
(X_{\mathrm{band}}Y_{\mathrm{band}})[o]=\sum_{s=-m}^m\sum_\tau w_{s,\tau}[o]\,u_{s,\tau}(X_{\mathrm{band}})\,v_{s,\tau}(Y_{\mathrm{band}})\qquad\text{for every }o,
\]
with every product taken left factor first. So the triples $(u_{s,\tau},v_{s,\tau},w_{s,\tau})$ form an exact bilinear algorithm over $\Q$ for $\mat{K_{\mathrm G}}{D_{\mathrm G}}{K_{\mathrm G}}$ with at most $(2m+1)5^{2L}$ terms; terms whose output form vanishes on $\RG$ are dropped, and these include every leaf with more than $m$ letters $P_0$ in some half.
\end{lemma}

\begin{proof}
For integer entries this is~\cite[Theorem~\Vref{m:thm:correct}]{v2} combined with the expansion $\sum_s\sum_\tau\chi_\tau[w]\Phi_\tau(a_s)\Psi_\tau(b_s)$ of the runs at $w$ \cite[Section~\Vref{m:sec:leaves}]{v2}. Both sides are bilinear in the entries, with coefficients independent of them (the arrays are selections, the forms linear), so Lemma~\ref{cp:P4} gives the identity for block entries. A leaf with more than $m$ letters $P_0$ in a half covers no string of $\UG$, since $\chi_{P_0}=z_0$ and a string of $\UG$ has exactly $m$ letters $z_0$ per half.
\end{proof}

\paragraph{The product top algorithm.} Let $P$ be an admissible rational level-1 point, $\alpha=\alpha(P)$, and let $B$ be a top algorithm for $\mat kbk$ obtained from the family of $P$ by Corollary~\ref{cp:T5}, so that $b=k^\alpha$ (Lemma~\ref{cp:L5}, \Hcond1). The \emph{shares} are numbers $\fCW,\fG,\fdir\ge0$ with $\fCW+\fG+\fdir=1$. Given $D$, let $m$ be the largest multiple of the denominator of $c$ with $4^m\le D^{\fG}$, and put
\[
n=\Bigl\lfloor\frac{\fCW\ln D}{\ln b}\Bigr\rfloor,\qquad L=cm,\qquad D_A=\Bigl\lceil\frac D{b^n4^m}\Bigr\rceil,
\]
$K=k^nK_{\mathrm G}$ and $D'=b^n4^mD_A$, so that $b^n$, $4^m$ and $D_A$ lie within constant factors (depending on $b$ and $c$) of $D^{\fCW}$, $D^{\fG}$ and $D^{\fdir}$. The top algorithm is $C=B^{\otimes n}\otimes\Tile_{L,m}$, an exact bilinear algorithm over $\Q$ for $\mat K{b^n4^m}K$ on block matrices: a tensor product of exact bilinear algorithms acts on disjoint index coordinates, and block products are bilinear with the left blocks always on the left. Its terms (\emph{leaves}) are $l=(l_1,\dots,l_n,s,\tau)$ with $l_j\in[\RB]$; an output block $o=(o_B,o_G)$ is identified with the string $(o_1,\dots,o_n,w(o_G))$ of $n$ letters of $B$ and $2L$ letters $z$, and
\[
w_l[o]=\prod_{j\le n}w_{l_j}[o_j]\cdot\chi_\tau[w(o_G)],
\]
\[
\Cov(o)=\Cov_B(o_1)\times\dots\times\Cov_B(o_n)\times\{-m,\dots,m\}\times\Cov_{\mathrm G}(w(o_G)),
\]
where $\Cov_{\mathrm G}(w)$, the \emph{cube} of $w$ \cite[Section~\Vref{m:sec:leaves}]{v2}, consists of the leaves with an arbitrary term at the $2m$ levels where $w$ has $z_0$ and $P_{ij}$ at every level where $w$ has $z_{ij}$. The run index enters the input forms only.

For the tile and the bottom we need the following constants, with $\Hb$ the binary entropy and $\varphi(s)=s\ln s$:
\[
\gG(c)=c\Hb(1/c)+(c-1)\ln4,\qquad\AG(c)=2c\ln5-\gG(c),
\]
\[
\Ic_c(x)=\int_0^x\ln\frac{4(1-y)}{c-1+y}\,dy,\qquad\Jc_c(x)=\varphi(c-1+x)-\varphi(x)-\varphi(c-1),
\]
\[
\LamG(t,c)=\frac2{\ln4}\max_{0\le x\le1}\bigl[\Ic_c(x)+t\Jc_c(x)\bigr].
\]
Here $\AG(c)/\ln4$ is the encoding exponent of the tile per $\ln D_{\mathrm G}$ and $\LamG(t,c)$ its profile; $\gG(c)\le c\ln5\le\AG(c)$, because $\binom Lm4^{L-m}\le\sum_j\binom Lj4^{L-j}=5^L$ and $\binom Lm\ge e^{L\Hb(m/L)}/(L+1)$. In closed form $\Ic_c(x)=x\ln4-\varphi(1-x)-\varphi(c-1+x)+\varphi(c-1)$, so $\Ic_c+\Jc_c=x\ln4+\Hb(x)$. Both $\Ic_c$ and $\Jc_c$ are concave on $[0,1]$ ($\Ic_c'$ decreases, $\Jc_c''(x)=1/(c-1+x)-1/x<0$), so for $t\in[0,1]$ the maximum in $\LamG$ is certified by $\psi(x_0)+\abs{\psi'(x_0)}$ at a rational $x_0\in(0,1)$, $\psi=\Ic_c+t\Jc_c$.

\begin{theorem}[\ThmDprod]\label{cp:G7}
Let $P$ be an admissible rational level-1 point with family $(A_M)$, $\alpha=\alpha(P)$, profile $\Lam$ and coverage exponent $\cbar$, and let $c>1$ be rational. Let $\fCW,\fG,\fdir\ge0$ with $\fCW+\fG+\fdir=1$, and let $t\in(0,1]$, with $t\le1/(\cbar+2)$ if $\fCW>0$. Put
\[
\Qgain{\otimes}{}(t,f)=\frac t2-\fCW\frac{\Lam(t)}\alpha-\fG\LamG(t,c)-\fdir,\qquad\Eenc{\otimes}(\eps,f)=\frac1\eps-\frac{\fCW}\alpha-\fG\frac{\AG(c)}{\ln4}-\fdir .
\]
Then $\Hyp(\eps,\gamma)$ holds (time model, deterministic) for every $\eps\in(0,\frac12]$ and every $\gamma$ with $0<\gamma<\min\{\Qgain{\otimes}{}(t,f),\Eenc{\otimes}(\eps,f),\frac12\}$. Any upper bound on $\Lam(t)$, and any upper bound on $\LamG(t,c)$, may be used in place of it.
\end{theorem}

With $\fG=0$, $\fCW=\alpha x/(1+\alpha x)$ and $\fdir=1/(1+\alpha x)$ the first two terms are those of \ThmDk{} at $\kappa=\frac12$: $\Qgain{\otimes}{}=\Qgain{\mathrm D}{1/2}(x,t)$ and $\Eenc{\otimes}=1/\eps-(1+x)/(1+\alpha x)$. With $\fG=1$ the bound is $\min\{t/2-\LamG(t,c),\,1/\eps-\AG(c)/\ln4\}$, the tile alone in the cost model of this theorem; numerically it reproduces the exact leaf count of~\cite[Remark~\Vref{m:rem:exact}]{v2}, the values that the curve of Section~\ref{cp:sec:curve} is compared with (computed).

The restriction on $t$ comes from the choice of $B$ through Corollary~\ref{cp:T5}: for $t\le1/(\cbar+2)\le2/(\cstar+2)$ it gives $\lamB(t)\le\Lam(t)+\eta$. All certified points of Section~\ref{cp:sec:curve} satisfy it: their largest tilt is $\cpv{u7_G10_a_t}$, and the certificate checks $t\le1/(\co(P)+2)$ at each of them.

\paragraph{\AlgDprod.} This is \AlgD{} of Section~\ref{cp:sec:thmD} with $C$ in place of $B^{\otimes n}$.
\begin{enumerate}
\item Compute $n,m,L,D_A$; pad $N$ to $N'=K\lceil N/K\rceil$ and the inner dimension to $D'$, and view $X$ as a $K\times b^n4^m$ matrix of blocks of size $N_{\mathrm{sub}}\times D_A$, $N_{\mathrm{sub}}=N'/K$, and $Y$ likewise. Compute $\Gamma$, $\Delta$ and $\mathcal Q_g$ by~\cite[Lemma~\Vref{m:lem:hashes}]{v2}, and for each good set $Q$ the inverse of $S\mapsto\Delta[S]$ on the subsets of $[L]\setminus Q$.
\item Encode: for every $l_B\in[\RB]^n$ form $X_{l_B}=\sum u_{l_B}[I,J]X_{I,J}$, a $K_{\mathrm G}N_{\mathrm{sub}}\times4^mD_A$ block matrix, by the recursion of \AlgD{} (after scaling $B$ to integer coefficients); then for every run $s$ form $a_s$ from $X_{l_B}$ and compute $\Phi_\tau(a_s)$ for all leaves $\tau$ by the level-by-level recursion of~\cite[Section~\Vref{m:sec:squared}]{v2}. Do the same for $Y$.
\item Bucket $W$ by the inner position $(r,r')\in[N_{\mathrm{sub}}]^2$ by a two-pass radix sort with base $N_{\mathrm{sub}}$, and compute for every wanted entry its output string $(o_1,\dots,o_n,w(o_G))$ with the inverse tables of step~1.
\item For every position with $W_{rr'}\ne\emptyset$ and every run $s$: compute the stage sets over the $n+2L$ coordinates down to $T_{rr'}$ (Lemma~\ref{cp:G8:dec}); for every leaf $l\in T_{rr'}$ compute $p_l=(X_lY_l)[r,r']$ as an inner product of length $D_A$; decode by Lemma~\ref{cp:G8:dec}, add the $2m+1$ run results, and divide by the scaling factor of $B$ (the tile's coefficients are integers).
\end{enumerate}

\emph{The partition.} The terms of $B^{\otimes n}$ carry the product partition of Section~\ref{cp:sec:shared}. The live leaves $(s,\tau)$ of the tile are partitioned by $(s,e_1,e_2)$, where $e_h\le m$ is the number of letters $P_0$ of $\tau$ in half $h$, and $C$ carries the product of the two partitions. Since $\Cov(o)$ is a product set, the class $(c_1,c_2)$ of a tensor product $C_1\otimes C_2$ has $\abs{P_{c_1}}\abs{P_{c_2}}$ terms and coverage exactly $\kappa_{c_1}\kappa_{c_2}$ (the maximum over a product set of a product of nonnegative factors is the product of the maxima), so the profiles multiply: $Z_C(t)=\ZB(t)^nZ_{\mathrm G}(t)$.

\begin{lemma}[classes of the tile]\label{cp:G8:G}
Let $0\le e_1,e_2\le m$. The leaves with exactly $e_h$ letters $P_0$ in half $h$ number $\binom L{e_1}4^{L-e_1}\binom L{e_2}4^{L-e_2}$. Every $w\in\UG$ is covered by exactly
\[
\kappa^*_{e_1,e_2}=\binom m{e_1}4^{m-e_1}\binom m{e_2}4^{m-e_2}
\]
of them, every one of them covers exactly $\deg_{e_1,e_2}=\binom{L-e_1}{m-e_1}\binom{L-e_2}{m-e_2}$ strings of $\UG$, and $\binom L{e_1}4^{L-e_1}\binom L{e_2}4^{L-e_2}\deg_{e_1,e_2}=\abs\UG\,\kappa^*_{e_1,e_2}$.
\end{lemma}

\begin{proof}
A leaf $\tau$ covers $w$ if and only if $\chi_{\tau_i}$ contains $w_i$ at every level, that is, if $\tau$ has $P_0$ only where $w$ has $z_0$ and has $P_{ij}$ wherever $w$ has $z_{ij}$. For fixed $w\in\UG$, a covering leaf with $e_h$ letters $P_0$ in half $h$ chooses them among the $m$ levels of that half where $w$ has $z_0$, puts one of the four terms $P_{ij}$ at the other $m-e_h$ of these levels, and is forced elsewhere; this gives $\kappa^*$. For a fixed leaf, a covered $w\in\UG$ has $z_0$ at the $e_h$ levels with $P_0$ and at $m-e_h$ of the $L-e_h$ levels with some $P_{ij}$ in each half, and is forced elsewhere; this gives $\deg$. The identity counts the incidences (leaf of the class, string of $\UG$) in two ways.
\end{proof}

\begin{lemma}[profile of the tile]\label{cp:G8:Gp}
For every $t\in[0,1]$ and every $m\ge1$,
\[
\frac1{\abs\UG}\sum_\tau\deg(\tau)^t\le(m+1)^2(\mathrm e\,m)^2\exp\Bigl(2m\max_{0\le x\le1}\bigl[\Ic_c(x)+t\Jc_c(x)\bigr]\Bigr),
\]
the sum over the strings $\tau$ of $2L$ terms with at most $m$ letters $P_0$ in each half and $\deg(\tau)$ the number of strings of $\UG$ that $\tau$ covers. Consequently the tile's profile for the partition by $(s,e_1,e_2)$ satisfies
\[
K_{\mathrm G}^{-2+2t}Z_{\mathrm G}(t)\le(2m+1)\mu_\infty^{-4}(m+1)^2(\mathrm e\,m)^2D_{\mathrm G}^{\LamG(t,c)}=D_{\mathrm G}^{\LamG(t,c)+O(\ln m/m)}.
\]
\end{lemma}

\begin{proof}
With $d_h=m-e_h$ the sum factorises over the halves:
\[
\frac1{\abs\UG}\sum_\tau\deg(\tau)^t=\Bigl(\sum_{d=0}^m\frac{\beta_d}{\beta_0}\binom{L-m+d}d^t\Bigr)^2,
\]
where $\beta_d=\binom L{m-d}4^{L-m+d}$ is the number of half-strings with $m-d$ letters $P_0$ \cite[Section~\Vref{m:sec:leaves}]{v2}. First, with $x=d/m$ and $n'=L-m+d=m(c-1+x)$, $\binom{n'}d\le\exp(n'\Hb(d/n'))=\exp(m\Jc_c(x))$, the terms $m\ln m$ cancelling. Second, $\ln(\beta_d/\beta_0)=\sum_{e=1}^d\ln\frac{4(m-e+1)}{L-m+e}$, while $m\Ic_c(d/m)=\int_0^d\ln\frac{4(m-z)}{L-m+z}\,dz$. On $[e-1,e]$ the integrand is at least $\ln4+\ln(m-e)-\ln(L-m+e)$ when $e<m$, and $\int_{m-1}^m\ln(m-z)\,dz=-1$; so the $e$-th summand exceeds its integral by at most $\ln\frac{m-e+1}{m-e}$ for $e<m$ and by at most $1$ for $e=m$. These excesses telescope to at most $\ln m+1$, so $\beta_d/\beta_0\le\mathrm e\,m\exp(m\Ic_c(d/m))$. Each of the $m+1$ terms is therefore at most $\mathrm e\,m\exp(m[\Ic_c(x)+t\Jc_c(x)])$, and squaring gives the first claim.
For the profile: by Lemma~\ref{cp:G8:G}, a class $c=(s,e_1,e_2)$ has $\abs{P_c}\le\binom L{e_1}4^{L-e_1}\binom L{e_2}4^{L-e_2}$ live leaves and coverage $\kappa_c\le\kappa^*_{e_1,e_2}$ on $\RG\subseteq\UG$, and $\kappa^*=\abs{P}\deg/\abs\UG$ for the full count $\abs P$. Hence
\[
\abs{P_c}^{1-t}\kappa_c^t\le\sum_{\tau\in c}\deg(\tau)^t\,\abs\UG^{-t},
\]
and with $\abs\RG=K_{\mathrm G}^2\ge\mu_\infty^4\abs\UG$ and $-1+t\le0$,
\[
K_{\mathrm G}^{-2+2t}Z_{\mathrm G}(t)\le(2m+1)\abs\RG^{-1+t}\abs\UG^{-t}\sum_\tau\deg(\tau)^t\le(2m+1)\mu_\infty^{-4}\frac1{\abs\UG}\sum_\tau\deg(\tau)^t .
\]
Taking logarithms and dividing by $\ln D_{\mathrm G}=m\ln4$ gives the second claim.
\end{proof}

\begin{lemma}[decoding for products]\label{cp:G8:dec}
Fix a position $(r,r')$, a run $s$, and the values $p_l$ for the leaves $l\in T_{rr'}$ of run $s$. The values $y^{(s)}_o=\sum_lw_l[o]\,p_l$ over these leaves, $o\in W_{rr'}$, can be computed with at most $(n+2L+1)(c_{\max}+1)\abs{T_{rr'}}$ table operations and $(n+2L)c_{\max}\abs{T_{rr'}}$ arithmetic operations, $c_{\max}=\max(\cB,5)$, given direct-address tables on a universe of size $\RB^n5^{2L}$ allocated once, in $O((n+2L)\RB^n5^{2L})$ time. The wanted entry is $\sum_sy^{(s)}_o$.
\end{lemma}

\begin{proof}
The proof of Lemma~\ref{cp:T4} uses only that the output form of a term is a product of per-coordinate forms, that the covering set of an output is the product of per-coordinate covering sets, and that each coordinate has an injective map $\psi$ from its output letters to terms with $\psi(o)$ covering $o$. Here the coordinates are the $n$ coordinates of $B$, with $\psi_B$ from Lemma~\ref{cp:T3}, and the $2L$ levels of the tile, with $\psi(z_{ij})=P_{ij}$ and $\psi(z_0)=P_0$; at a tile level at most $5$ terms cover a letter. With the coordinates of $B$ first, the stage sets satisfy $\abs{M_i}\le\abs{T_{rr'}}$ by the same injection, and stage $i$ costs at most $c_{\max}\abs{M_i}$ insertions and multiply-adds. The run is fixed throughout and is not a coordinate; the stamp of a table cell is the pair (position, run). The sum over the runs is correct by Lemma~\ref{cp:G8}.
\end{proof}

\begin{proof}[Proof of Theorem~\ref{cp:G7}]
Fix $\eps$ and $\gamma$ as in the theorem, and let $3\xi>0$ be the smallest of the margins $\Qgain{\otimes}{}-\gamma$, $\Eenc{\otimes}-\gamma$ and $\frac12-\gamma$. Choose $\eta>0$ with $\fCW\eta/\alpha<\xi$. If $\fCW>0$, Corollary~\ref{cp:T5} gives a fixed exact top algorithm $B$ for $\mat kbk$ with $b=k^\alpha$, $\delB\le\eta$ and $\lamB(t)\le\Lamp(t)+\eta=\Lam(t)+\eta$, since $t\le1/(\cbar+2)\le2/(\cstar+2)$ (part (i)); if $\fCW=0$, take $n=0$ and no $B$; if $\fG=0$, take $m=L=0$ and no tile; then $C=B^{\otimes n}$, $K_{\mathrm G}=D_{\mathrm G}=1$, $Z_{\mathrm G}=1$, and every bound below holds with the tile's factors equal to~$1$. Run \AlgDprod{} for $D\ge D_0$, with a constant $D_0$ chosen below.

\emph{Correctness.} $C$ computes $\mat K{b^n4^m}K$ exactly on block matrices (Lemma~\ref{cp:G8} and the product rule), so block $o$ of $XY$ is $\sum_lw_l[o]X_lY_l$; Lemma~\ref{cp:G8:dec} computes exactly these sums at the wanted entries, and the inner products give $p_l=(X_lY_l)[r,r']$.

\emph{Needed leaves.} Lemma~\ref{cp:T2}(a), applied to $C$ with $n=1$ (its proof uses only the partition), gives $\sum\abs{T_{rr'}}\le N_{\mathrm{sub}}^2\bar w^tZ_C(t)$ with $\bar w=\abs W/N_{\mathrm{sub}}^2\le K^2D^{-1/2}$. With $Z_C=\ZB^nZ_{\mathrm G}$ (the runs are part of the classes of the tile) and Lemma~\ref{cp:G8:Gp},
\[
\sum\abs{T_{rr'}}\le N'^2D^{-t/2}\bigl[k^{n(-2+2t)}\ZB(t)^n\bigr]\bigl[K_{\mathrm G}^{-2+2t}Z_{\mathrm G}(t)\bigr]\le N'^2D^{-t/2}k^{n\lamB(t)}D_{\mathrm G}^{\LamG(t,c)+O(\ln m/m)} .
\]
Since $\ln k^n=n\ln b/\alpha=\fCW\ln D/\alpha+O(1)$ and $\ln D_{\mathrm G}=\fG\ln D+O(1)$, and each needed leaf costs an inner product of length $D_A=O(D^{\fdir})$ plus $O(n+L)=O(\log D)$ for decoding, this part costs $N^2D^{-t/2+\fCW\lamB(t)/\alpha+\fG\LamG(t,c)+\fdir+o(1)}\le N^2D^{-\Qgain{\otimes}{}+\fCW\eta/\alpha+o(1)}\le N^2D^{-\gamma-2\xi+o(1)}$.

\emph{Encodings.} Put $R=\RB^n5^{2L}$. The stage of $B$ costs $O(nkb\RB^nK_{\mathrm G}4^mN_{\mathrm{sub}}D_A)$, and $K_{\mathrm G}4^m=K_g4^L\le8^L\le25^L$. The stage of the tile costs, for every leaf of $B$ and every run, $O(L)$ per nonzero entry for forming the arrays (at most $12^L$ block entries) and $O(25^L)$ block operations for the recursion \cite[Section~\Vref{m:sec:squared}]{v2}. In total the encodings cost $O((n+L)(2m+1)RN_{\mathrm{sub}}D_A)$, and $N_{\mathrm{sub}}\le2N/K$. Now
\[
\ln\frac{RD_A}{KN}=n\ln\frac{\RB}k+\bigl[2L\ln5-\ln K_{\mathrm G}\bigr]+\ln D_A-\ln N .
\]
Here $n\ln(\RB/k)=n(1+\delB)\ln k=\fCW\frac{1+\delB}\alpha\ln D+O(1)$; $2L\ln5-\ln K_{\mathrm G}\le m\AG(c)+\ln(L+1)+2\ln(1/\mu_\infty)$, by $K_{\mathrm G}\ge\mu_\infty^2\binom Lm4^{L-m}$, $\ln\binom Lm\ge L\Hb(m/L)-\ln(L+1)$ and $2cm\ln5-cm\Hb(1/c)-(c-1)m\ln4=m\AG(c)$, which is $\fG\frac{\AG(c)}{\ln4}\ln D+O(\ln\ln D)$; and $\ln D_A=\fdir\ln D+O(1)$. With $N\ge D^{1/\eps}$ the encodings cost $N^2D^{-\Eenc{\otimes}(\eps,f)+\fCW\delB/\alpha+o(1)}\le N^2D^{-\gamma-2\xi+o(1)}$.

\emph{One-time costs, input and output.} The tables of Lemma~\ref{cp:G8:dec} cost $O((n+2L)R)$, at most the encodings. \cite[Lemma~\Vref{m:lem:hashes}]{v2} costs $\binom Lm(2^m+2^{L-m})L^{O(1)}\le8^LL^{O(1)}$ and needs $L\le\log_2N$; this holds since $L\le c\fG\ln D/(2\ln2)+O(1)$, $\log_2N\ge\ln D/(\eps\ln2)$, and $c\fG/2<c\fG\ln5/\ln4\le\fG\AG(c)/\ln4<1/\eps$ by $\AG(c)\ge c\ln5$ and $\Eenc{\otimes}>0$. The inverse tables have $K_g2^{L-m}\le R$ entries. The padding needs $K\le N$: $K=k^nK_{\mathrm G}\le D^{\fCW/\alpha+\fG\gG(c)/\ln4+o(1)}$ and $\gG(c)\le\AG(c)$, so $K\le D^{1/\eps-\Eenc{\otimes}+o(1)}\le N$ for $D\ge D_0$. Reading the input costs $O(ND)$; the radix sort costs $O(\abs W+N_{\mathrm{sub}})$ and visits only the positions with $W_{rr'}\ne\emptyset$; the output strings cost $O((n+L)\abs W)$. All of this is $O(N^2D^{-1/2}\log D)$, since $\eps\le\frac12$.

\emph{Word size.} The encodings of $B$ multiply the entries by at most $(kbH)^n=N^{O(1)}$, $H$ the largest coefficient of the scaled $B$; the per-level forms of the tile have coefficients in $\{0,\pm1\}$ and at most four nonzero entries, so a leaf form is a signed sum of at most $D^{O(1)}$ entries. Every number is $N^{O(1)}$ and fits in $O(1)$ words.

\emph{Conclusion.} The terms $o(1)$ are $O(\ln\ln D/\ln D)$, from the rounding of $n$, $m$ and $D_A$, the factors $2m+1$, $(m+1)^2(\mathrm e\,m)^2$ and $\mu_\infty^{-4}$, and the logarithmic factors; they depend on $B$, $c$ and $f$ only. So there is $D_0$ with total cost at most $CN^2D^{-\gamma}$ for all $D\ge D_0$; for $D<D_0$ the direct product costs $O(N^2D)=O(D_0^{1+\gamma}N^2D^{-\gamma})$. Replacing $\Lam(t)$ or $\LamG(t,c)$ by an upper bound only lowers $\Qgain{\otimes}{}$. The quantifiers are chosen in the order $(P,c,f,t)$ and $(\eps,\gamma)$, then $\eta$, then $B$, then $D_0$; for each instance $(n,m,L,D_A)$ are computed from $D$, and the tile from $(L,m)$, deterministically.
\end{proof}

\subsection{Two factors, the certified curve, and what it does not change}\label{cp:sec:curve}

The bound of \ThmDprod{} is a mixture: each factor $i\in\{\mathrm{CW},\mathrm G,\mathrm{dir}\}$ contributes $f_i(t/2-\Lambda_i(t))$ to $\Qgain{\otimes}{}$ and $f_ia_i$ to $1/\eps-\Eenc{\otimes}$, with
\[
(\Lambda_{\mathrm{CW}},a_{\mathrm{CW}})=\Bigl(\frac{\Lam(t)}\alpha,\frac1\alpha\Bigr),\qquad(\Lambda_{\mathrm G},a_{\mathrm G})=\Bigl(\LamG(t,c),\frac{\AG(c)}{\ln4}\Bigr),\qquad(\Lambda_{\mathrm{dir}},a_{\mathrm{dir}})=(1,1).
\]
Put $L_i=\sup_t(t/2-\Lambda_i(t))$, the supremum over the tilts allowed for that factor ($t\le1/(\cbar+2)$ for the level-1 top, $t\in(0,1]$ otherwise).

\begin{proposition}[two factors suffice, and the best pure factor bounds $\eps\gamma$]\label{cp:G9}
\begin{enumerate}[label=(\alph*)]
\item For fixed $P$, $c$, $t$ and $\eps$, the supremum over the share vectors $f$ of $\min\{\Qgain{\otimes}{}(t,f),\Eenc{\otimes}(\eps,f)\}$, when it is positive, is attained at a vector with at most two nonzero shares.
\item $a_i+L_i>0$ for every factor: $a_{\mathrm{dir}}+L_{\mathrm{dir}}=\frac12$, $a_{\mathrm{CW}}+L_{\mathrm{CW}}\ge1/\alpha$, and $a_{\mathrm G}+L_{\mathrm G}>\cpv{u7_aGL}$ for every $c>1$.
\item For every share vector $f$ and every pair $(\eps,\gamma)$ given by \ThmDprod{} with this $f$,
\[
\eps\gamma\le\max_{i:\,f_i>0}\frac{L_i}{a_i+L_i}.
\]
\end{enumerate}
\end{proposition}

\begin{proof}
(a) At fixed $t$ and $\eps$ both terms are affine in $f$. Maximising $\gamma$ subject to $\gamma\le\Qgain{\otimes}{}(t,f)$, $\gamma\le\Eenc{\otimes}(\eps,f)$, $\sum_if_i=1$ and $f\ge0$ is a linear program with two inequality rows and one equality row, so it has an optimal basic solution with at most three nonzero variables among $(f,\gamma)$; if the optimum is positive, $\gamma$ is one of them.
(b) $L_{\mathrm{dir}}=\sup_{t\le1}(t/2-1)=-\frac12$. For the level-1 top, $\Lam(t)/t\to0$ as $t\to0$ (Lemma~\ref{cp:L8}(iv)), so $L_{\mathrm{CW}}\ge0$ and $a_{\mathrm{CW}}+L_{\mathrm{CW}}\ge1/\alpha$. For the tile, $\Ic_c+\Jc_c=x\ln4+\Hb(x)$ has maximum $\ln5$ on $[0,1]$, so $\LamG(1,c)=2\ln5/\ln4$ and $L_{\mathrm G}\ge\frac12-2\ln5/\ln4>-\cpv{u7_LG_lo}$. And $\AG$ is convex on $(1,\infty)$ with $\AG'(c)=\ln\frac{25(c-1)}{4c}$, which vanishes at $c=\cpv{u7_cmin}$, so $\AG(c)/\ln4\ge\AG(\cpv{u7_cmin})/\ln4>\cpv{u7_aG_min}$. Hence $a_{\mathrm G}+L_{\mathrm G}>\cpv{u7_aG_min}-\cpv{u7_LG_lo}>\cpv{u7_aGL}$.
(c) Let $S_L=\sum_if_iL_i$ and $S_a=\sum_if_ia_i>0$. Since $\sum_if_i=1$, $\gamma<\Qgain{\otimes}{}=\sum_if_i(t/2-\Lambda_i(t))\le S_L$, and $\gamma<\Eenc{\otimes}$ gives $\eps<1/(S_a+\gamma)$. As $y\mapsto y/(S_a+y)$ increases on $[0,\infty)$, $\eps\gamma<\gamma/(S_a+\gamma)\le S_L/(S_a+S_L)=\sum_if_iL_i\big/\sum_if_i(a_i+L_i)$, and by (b) this mediant is at most the largest of the ratios $L_i/(a_i+L_i)$ with $f_i>0$.
\end{proof}

\begin{table}[ht]
\centering\small
\begin{tabular}{r l l l l}
\hline
& every $\eps$ & every $\gamma<$ & construction & parameters\\
\hline
1 & $<\cpv{G10_a_E}$ & $\cpv{G10_a_G}$ & $\Pprod\otimes\Tile(c)$ & $c=\cpv{u7_G10_a_c}$, $t=\cpv{u7_G10_a_t}$, $\fG=\cpv{u7_G10_a_f}$\\
2 & $<\cpv{G10_b_E}$ & $\cpv{G10_b_G}$ & $\Pprod\otimes\Tile(c)$ & $c=\cpv{u7_G10_b_c}$, $t=\cpv{u7_G10_b_t}$, $\fG=\cpv{u7_G10_b_f}$\\
3 & $<\cpv{G10_c_E}$ & $\cpv{G10_c_G}$ & $\Pprod\otimes\Tile(c)$ & $c=\cpv{u7_G10_c_c}$, $t=\cpv{u7_G10_c_t}$, $\fG=\cpv{u7_G10_c_f}$\\
4 & $<\cpv{G10_d_E}$ & $\cpv{G10_d_G}$ & $\Pprod\otimes\Tile(c)$ & $c=\cpv{u7_G10_d_c}$, $t=\cpv{u7_G10_d_t}$, $\fG=\cpv{u7_G10_d_f}$\\
5 & $<\cpv{G10_e_E}$ & $\cpv{G10_e_G}$ & $\Pprod\otimes\Tile(c)$ & $c=\cpv{u7_G10_e_c}$, $t=\cpv{u7_G10_e_t}$, $\fG=\cpv{u7_G10_e_f}$\\
6 & $<\cpv{G10_f_E}$ & $\cpv{G10_f_G}$ & $\Pprod\otimes\Tile(c)$ & $c=\cpv{u7_G10_f_c}$, $t=\cpv{u7_G10_f_t}$, $\fG=\cpv{u7_G10_f_f}$\\
7 & $<\cpv{G10_g_E}$ & $\cpv{G10_g_G}$ & $\Pprod\otimes\Tile(c)$ & $c=\cpv{u7_G10_g_c}$, $t=\cpv{u7_G10_g_t}$, $\fG=\cpv{u7_G10_g_f}$\\
8 & $<\cpv{G10_h_E}$ & $\cpv{G10_h_G}$ & $\Pprod\otimes\Tile(c)$ & $c=\cpv{u7_G10_h_c}$, $t=\cpv{u7_G10_h_t}$, $\fG=\cpv{u7_G10_h_f}$\\
9 & $<\cpv{G10_i_E}$ & $\cpv{G10_i_G}$ & $\Pprod\otimes\Tile(c)$ & $c=\cpv{u7_G10_i_c}$, $t=\cpv{u7_G10_i_t}$, $\fG=\cpv{u7_G10_i_f}$\\
10 & $<\cpv{G10_j_E}$ & $\cpv{G10_j_G}$ & $\Pprod\otimes\Tile(c)$ & $c=\cpv{u7_G10_j_c}$, $t=\cpv{u7_G10_j_t}$, $\fG=\cpv{u7_G10_j_f}$\\
11 & $<\cpv{G10_k_E}$ & $\cpv{G10_k_G}$ & $\Pprod\otimes\Tile(c)$ & $c=\cpv{u7_G10_k_c}$, $t=\cpv{u7_G10_k_t}$, $\fG=\cpv{u7_G10_k_f}$\\
12 & $<\cpv{G10_l_E}$ & $\cpv{G10_l_G}$ & $\Pprod\otimes\Tile(c)$ & $c=\cpv{u7_G10_l_c}$, $t=\cpv{u7_G10_l_t}$, $\fG=\cpv{u7_G10_l_f}$\\
13 & $\le\cpv{G10_m_E}$ & $\cpv{G10_m_G}$ & $\Pzero$, Theorem~\ref{cp:G3} & $t=\cpv{G3_t}$, $x\to\infty$\\
14 & $\le\cpv{G10_n_E}$ & $\cpv{G10_n_G}$ & $\Pzero$, \ThmDk & $t=\cpv{G3_t}$, $x=\cpv{u7_G10_n_x}$\\
15 & $\le\cpv{G10_o_E}$ & $\cpv{G10_o_G}$ & $\Pzero$, \ThmDk & $t=\cpv{G3_t}$, $x=\cpv{u7_G10_o_x}$\\
16 & $\le\cpv{G10_p_E}$ & $\cpv{G10_p_G}$ & $\Pzero$, \ThmDk & $t=\cpv{G3_t}$, $x=\cpv{u7_G10_p_x}$\\
17 & $\le\cpv{G10_q_E}$ & $\cpv{G10_q_G}$ & $\Pzero$, \ThmDk & $t=\cpv{G3_t}$, $x=\cpv{excess_sharp_x}$\\
\hline
\end{tabular}
\caption{Certified points of the curve $\gamma(\eps)$ (Theorem~\ref{cp:G10}). In rows 1--12, $\fCW=1-\fG$ and $\fdir=0$.}\label{cp:tab:G10}
\end{table}

\begin{theorem}[the certified curve]\label{cp:G10}
For each row of Table~\ref{cp:tab:G10}, $\Hyp(\eps,\gamma)$ holds (time model, deterministic) for every $\eps$ in the second column and every $\gamma$ below the third. Rows 1--12 are \ThmDprod{} at the level-1 point $\Pprod=(\cpv{u7_Pprod_lam},\,\cpv{u7_ubar},\,\cpv{u7_ubar})$ with the profile bound $\Lamcert$; rows 13--17 are \ThmDk{} at $\kappa=\frac12$ at $\Pzero$ with the same profile bound, row~13 being Theorem~\ref{cp:G3} and rows 14--17 using the finite $x$ shown.
\end{theorem}

\begin{proof}
At $\Pprod$ the certificate gives $\slack(\Pprod)>\cpv{u7_Pprod_slack}$ and $\alpha(\Pprod)>\cpv{u7_Pprod_alpha}$, and checks $t\le1/(\co(\Pprod)+2)$ for the $t$ of each row, so Theorem~\ref{cp:G7} applies with $\cbar\le\co(\Pprod)$ (Lemma~\ref{cp:L5}). In each of rows 1--12 it bounds $\Lam(t)\le\Lamcert(t)$ by Lemmas~\ref{cp:G1} and~\ref{cp:G2} and $\LamG(t,c)$ by the concavity certificate of Section~\ref{cp:sec:prod}, both in interval arithmetic, and so certifies a lower bound $Q_{\mathrm{lo}}$ on $\Qgain{\otimes}{}(t,f)$ with $Q_{\mathrm{lo}}>G$, the $\gamma$-bound of the row. The $\eps$-bound $E$ of the row is at most the lower end of the enclosure of $1/(a+Q_{\mathrm{lo}})$, with $a=\fCW/\alpha+\fG\AG(c)/\ln4$ taken at its upper end (the lower end of $\alpha$), so $\Eenc{\otimes}(\eps,f)=1/\eps-a>Q_{\mathrm{lo}}>G$ for every $\eps<E$. Theorem~\ref{cp:G7} with $\fdir=0$ gives the row. In rows 14--17 the certificate checks, at the $x$ of the row and $t=\cpv{G3_t}$, the condition \Dcond{} with $\Lamcert$, the range condition $\eps<(1+\alpha x)/(1+x)$ at the largest $\eps$ of the row, and that the encoding term $1/\eps-(1+x)/(1+\alpha x)$ at that $\eps$ and the gain $\Qgain{\mathrm D}{1/2}(x,t)$ both exceed $G$; the third term $2x/(1+\alpha x)$ exceeds $\frac12$. Both conditions persist for smaller $\eps$, and \ThmDk{} gives the row. All values are rounded in the safe direction: $G$ down to four significant digits and the $\eps$-bounds down to four decimals.
\end{proof}

Every row holds for all smaller $\eps$, so the table is a step function below the curve. Rows 15--17 lie beyond $\alpha(\Pzero)$, the range extension of \ThmDk(c), and row~14 improves on row~13 just below $\alpha(\Pzero)$: with the sharp profile the certified range at $\Pzero$ reaches $\cpv{G10_q_E}$, while the crude profile $\Lampr$ alone reaches $\cpv{excess_crude_P0}$ (Remark~\ref{cp:A3}).

\begin{remark}[the numerical curve]\label{cp:G11}
Optimising the bound of \ThmDprod{} numerically over $(c,t,f)$ at the level-1 point of the products gives a curve whose certified points are rows 1--12 of Table~\ref{cp:tab:G10}; the products exceed the tile alone from $\eps$ about $\cpv{G11_cross}$ on. Optimising \ThmDk{}(b) numerically over level-1 points with other part sets gives the level-1 envelope, which rises from about $\cpv{G11_env_lo}$ near the computed supremum $\cpv{alphaC_computed}$ of $\alpha(P)$ to about $\cpv{G11_env_hi}$ at $\eps=\cpv{G10_b_E}$; the numerical end of the level-1 range is $\cpv{G11_end}$. These are numerical optima, computed in floating point; they are not theorems, and only the certified rows of Tables~\ref{cp:tab:G6} and~\ref{cp:tab:G10} are.
\end{remark}

\begin{remark}[what the products do not change]\label{cp:G12}
The largest $\eps\gamma$ in Table~\ref{cp:tab:G10} is $\cpv{G12_epsgamma}$, at $\eps<\cpv{G10_a_E}$. By Proposition~\ref{cp:G9}(c), every pair given by \ThmDprod{} has $\eps\gamma$ at most the better of the values of its two pure factors, evaluated with the same profile bounds. With the bounds of this section these values are numerical optima: $\cpv{epsgamma_v1_opt}$ for the tile alone, the value of the exact leaf count of~\cite[Remark~\Vref{m:rem:exact}]{v2} (a numerical optimum there as well), and about $\cpv{G12_cw_epsgamma}$ for the level-1 top alone with $\Lamcert$ (computed; Table~\ref{cp:tab:G6} certifies $\cpv{G6_epsgamma}$); \cite[Corollary~\Vref{m:cor:prunedpoints}]{v2} reaches $\cpv{epsgamma_v2}$ by a different count. So, with these profile bounds, the products do not raise $\eps\gamma$, the quantity through which~\cite{v2} derives the exponents of their applications. What they raise is $\gamma$ at fixed $\eps$ between about $\cpv{G11_cross}$ and $\alpha(\Pprod)$, the dual exponent of the point of rows 1--12 (computed, Remark~\ref{cp:G11}): for example, at every $\eps\le\cpv{G6_b_E}$, from $\cpv{G6_b_G}$ (Theorem~\ref{cp:G6}) to $\cpv{G10_e_G}$ (Table~\ref{cp:tab:G10}, row~5).
\end{remark}

\phantomsection\label{cp:G13}%
With the closed-form bound $\Lampr$ of Lemma~\ref{cp:L8} in place of $\Lamcert$, so without Lemmas~\ref{cp:G1:support}--\ref{cp:G2}, \ThmDprod{} at the level-1 point $\Ppr=(\cpv{u7_Ppr_lam},\,\cpv{u7_Ppr_u6},\,\cpv{u7_Ppr_u7})$ with $\fdir=0$ still gives, certified, $\Hyp(\eps,\gamma)$ for every $\gamma<\cpv{G13_crude_a_G}$ and every $\eps<\cpv{G13_crude_a_E}$ ($c=\cpv{u7_cr_a_c}$, $t=\cpv{u7_cr_a_t}$, $\fG=\cpv{u7_cr_a_f}$) and for every $\gamma<\cpv{G13_crude_b_G}$ and every $\eps<\cpv{G13_crude_b_E}$ ($c=\cpv{u7_cr_b_c}$, $t=\cpv{u7_cr_b_t}$, $\fG=\cpv{u7_cr_b_f}$).

\section{Batches}\label{cp:sec:batch}

Definition~\ref{cp:P2} separates the two costs of a composite: phase~1 sees only $X$ and $Y$, and phase~2 answers a batch of at most $\bk$ positions, with the targets $CN^2D^{-g}$ for $\kappa\le1$ and $C\bk D^{1-g}$ for $\kappa\ge1$; $\fbat(\kappa)$ is the supremum of the $\eps$ for which some algorithm does this. This section proves the composite of Section~\ref{cp:sec:tops} in this two-cost form for every tilt $t\in(0,1]$ and with two bottoms, the direct inner product and the data structure of~\cite{v2} (Theorem~\ref{cp:B1}, \ThmB); certifies a grid of profile bounds at $\Pzero$ (Lemma~\ref{cp:B2}); gives a certified lower bound $\fcert(\kappa)$ on $\fbat(\kappa)$ for every $\kappa>0$, with the two monotonicity lemmas that carry certified values over intervals of $\kappa$ (Sections~\ref{cp:sec:batchcurve} and~\ref{cp:sec:mono}); says what the curve does not show (Section~\ref{cp:sec:batchnot}); and shows that batches of polynomially many full columns are easy without any composite (Section~\ref{cp:sec:fastcols}). All statements are in the time model and deterministic.

\subsection{\texorpdfstring{\ThmB}{Theorem B}}\label{cp:sec:thmB}

\begin{theorem}[\ThmB]\label{cp:B1}
Let $P$ be an admissible rational level-1 point with family $(A_M)$ and $\alpha=\alpha(P)$. Let $\kappa>0$, $t\in(0,1]$ and $L_t\ge\Lamp(t)$, and put $s=\min(0,1-\kappa)+t\kappa$. Let the bottom be either \emph{direct}, with $q_A=1$ and $\eps_A=1$, or the \emph{data structure}, with any $q_A>0$ and $\eps_A<\es$. If $x>0$ and
\[
xL_t+q_A<(1+\alpha x)\,s,\tag{B1}
\]
then the $\kappa$-batch statement holds at every $\eps<(1+\alpha x)/(x+1/\eps_A)$ (time model, deterministic). Consequently it holds at every $\eps$ below each of the numbers (a) and (b), where they are defined, and, by Theorem~\ref{cp:T7}(a) rather than by (B1), at every $\eps$ up to the number (c):
\begin{enumerate}[label=(\alph*)]
\item $\rdir(\kappa,t)=\alpha+\dfrac{(1-\alpha)(\alpha s-L_t)}{1-s+\alpha s-L_t}$, if $\alpha s>L_t$ and $s<1$;
\item $\rDS(\kappa,t)=\dfrac{\es L_t}{L_t-(\alpha-\es)s}$, if $L_t\ge\alpha s>0$;
\item $\alpha(P)$, if $\kappa\le1$; there every $\eps\le\alpha(P)$ is reached.
\end{enumerate}
The data-structure bottom uses Theorem~\ref{cp:P3}; the direct bottom uses no result of~\cite{v2}.
\end{theorem}

\begin{proof}
\emph{The top.} Fix $\eps$, $\kappa$, $t$, $x$ and the bottom with (B1) and $\eps<(1+\alpha x)/(x+1/\eps_A)$. For the data-structure bottom take $\gamma_A>0$ from Theorem~\ref{cp:P3} for $(\eps_A,q_A)$. Since both conditions are strict, there is $\eta>0$ such that (B1) holds with $L_t+\eta$ in place of $L_t$, $1/\eps>(1+x(1+\eta))/(1+\alpha x)$ (note $1/\eps>(x+1/\eps_A)/(1+\alpha x)\ge(1+x)/(1+\alpha x)$ as $\eps_A\le1$), and, for the data-structure bottom, $x\eta<\gamma_A$. Corollary~\ref{cp:T5} gives a top algorithm $B$ for $\mat kbk$ with $\alB=\alpha$ (Lemma~\ref{cp:L5}, \Hcond1), $\delB\le\eta$ and $\lamB(t)\le\Lamp(t)+\eta\le L_t+\eta$; this holds for every $t\in(0,1]$.

\emph{The algorithm.} Phase~1 runs steps~1 and~2 of \AlgD{} with this $B$ (the parameters $n$, $K$, $D_A$, the padding, the encodings $X_l$, $Y_l$ for all $l\in[\RB]^n$); for the data-structure bottom it also preprocesses each product $X_lY_l$, an instance with $N_{\mathrm{sub}}$ rows, inner dimension $D_A$ and $N_{\mathrm{sub}}$ columns, by Theorem~\ref{cp:P3}; and it allocates the tables of Lemma~\ref{cp:T4} and a zeroed array of $N_{\mathrm{sub}}^2$ bucket heads. Nothing in phase~1 depends on the batch. Phase~2, given a batch $\Bt$, buckets $\Bt$ by inner position through the bucket array, appending a position to a list when its bucket becomes nonempty; for every listed position $(r,r')$ it computes the needed sets down to $T_{rr'}$ (Lemma~\ref{cp:T4}), the values $p_l=(X_lY_l)[r,r']$ for $l\in T_{rr'}$ (an inner product of length $D_A$, or a single-entry query of Theorem~\ref{cp:P3}), and decodes; it outputs the wanted entries divided by the scaling factor, and finally resets the listed bucket heads.

\emph{Phase~1 cost.} As in the proof of Theorem~\ref{cp:T6}, the encodings cost
\[
O\bigl(nN^2D^{-1/\eps+(1+x(1+\delB))/(1+\alpha x)}\bigr),
\]
with a negative exponent by the choice of $\eta$; the tables cost $O(nK^{2+\delB})$ and the bucket array $O(N_{\mathrm{sub}}^2)=O(N^2D^{-2x/(1+\alpha x)})$, both polynomially below $N^2$ (exponents $e_3$, $e_4$ there); reading the input costs $O(ND)$. For the data-structure bottom the $\RB^n$ preprocessings cost $\RB^n\,O(N_{\mathrm{sub}}^2D_A^{-\gamma_A})=O(N'^2K^{\delB}D_A^{-\gamma_A})=O(N^2D^{(x\delB-\gamma_A)/(1+\alpha x)})$, with a negative exponent since $x\delB\le x\eta<\gamma_A$; the sub-instances satisfy $2\le D_A\le N_{\mathrm{sub}}^{\eps_A}$ for $D\ge D_0$ because $\eps<(1+\alpha x)/(x+1/\eps_A)$, as in part~(f) of the proof of Proposition~\ref{cp:T8}, and their entries are $N^{O(1)}=N_{\mathrm{sub}}^{O(1)}$. So phase~1 takes $O(N^2D^{-\gamma_{\mathrm{pre}}})$ steps for some $\gamma_{\mathrm{pre}}>0$.

\emph{Phase~2 cost.} By Lemma~\ref{cp:T2}(a) with the wanted set $\Bt$, $\bar w=\abs\Bt/N_{\mathrm{sub}}^2$ and $N_{\mathrm{sub}}K=N'\le2N$,
\[
\sum_{(r,r')}\abs{T_{rr'}}\le N_{\mathrm{sub}}^{2-2t}\abs\Bt^t\ZB(t)^n=N'^{2-2t}\abs\Bt^tK^{\lamB(t)}\le4\,\bk\Bigl(\frac{N^2}{\bk}\Bigr)^{1-t}K^{\lamB(t)}\le8\,\bk D^{\kappa(1-t)}K^{\lamB(t)},
\]
since $\abs\Bt\le\bk$ and $N^2/\bk\le2D^\kappa$ (if $N^2\ge D^\kappa$ then $\bk\ge N^2/(2D^\kappa)$, and otherwise $\bk=1$ and $N^2<D^\kappa$). Each element costs $O(D_A^{q_A})$ for its value ($q_A=1$ for the direct bottom) and $O(n\cB)$ for decoding. With $K\le D^{x/(1+\alpha x)}$, $D_A\le2bD^{1/(1+\alpha x)}$ and $\lamB(t)\ge0$ (Lemma~\ref{cp:T1:basic}), this part of phase~2 costs $O\bigl(n\,\bk D^{\kappa(1-t)+(x\lamB(t)+q_A)/(1+\alpha x)}\bigr)$, and by (B1) with $L_t+\eta$, $(x\lamB(t)+q_A)/(1+\alpha x)<s$. For $\kappa\ge1$, $\kappa(1-t)+s=1$, so the cost is polynomially below $\bk D$. For $\kappa\le1$, $s=t\kappa$ and $\bk\le N^2D^{-\kappa}$, so the cost is polynomially below $N^2$. The bucketing, the list, the reset, the output and the computation of the block and position of each wanted entry cost $O(n\abs\Bt)$, which is below the target once $g<1$ (for $\kappa\ge1$) and $g<\kappa$ (for $\kappa\le1$). So there are $g>0$ and $C$ with phase~2 within the target for $D\ge D_0$. For $D<D_0$, phase~1 does nothing and phase~2 computes every wanted entry directly in $O(D_0)$ steps each.

\emph{Correctness and word size} are as in the proof of Theorem~\ref{cp:T6}; the data structure returns the exact integers $(X_lY_l)[r,r']$.

\emph{The closed forms.} (B1) reads $x(L_t-\alpha s)<s-q_A$. (a) With the direct bottom it holds for every $x>x_{\min}=(1-s)/(\alpha s-L_t)$, and as $x$ decreases to $x_{\min}$ the range $(1+\alpha x)/(1+x)$ tends to $(\alpha-L_t)/(1-s+\alpha s-L_t)=\rdir(\kappa,t)$. (b) With the data-structure bottom and $L_t>\alpha s>0$ it holds for every $x<(s-q_A)/(L_t-\alpha s)$; as $q_A\to0$, $\eps_A\to\es$ and $x$ increases to $x_{\max}=s/(L_t-\alpha s)$, the range $(1+\alpha x)/(x+1/\eps_A)$ tends to $\es L_t/(L_t-(\alpha-\es)s)=\rDS(\kappa,t)$; if $L_t=\alpha s$, (B1) reads $q_A<s$ for every $x>0$, and the range tends to $\alpha=\rDS(\kappa,t)$ as $x\to\infty$. In both cases every $\eps$ below the limit is below the range at some admissible $(x,q_A,\eps_A)$. (c) For $\kappa\le1$ and $\eps\le\alpha(P)$, \ThmDk{}(a) gives $\Hypk\kappa(\eps,\gamma)$ by \AlgD; its steps~1--2 and the allocation of its tables do not depend on $W$ and form phase~1, its steps~3--4 form phase~2, and the cost bounds of the proof of Theorem~\ref{cp:T6} put both phases within $CN^2D^{-\gamma}$, the target for $\kappa\le1$.
\end{proof}

At $t=1$ and $\kappa\ge1$, $s=1$ for every $\kappa$, and (B1) with the data-structure bottom is $x(L_1-\alpha)<1-q_A$, the condition of Proposition~\ref{cp:T8} with $Q=1$ and $c=L_1$: the batch theorem then answers the batch entry by entry, and its range $\rDS(\kappa,1)=c\es/(\es+c-\alpha)$ is the per-entry range $E_1(\alpha,c)$ of that proposition.

\subsection{Profile bounds on a grid of tilts}\label{cp:sec:Mgrid}

\begin{lemma}[the certified grid]\label{cp:B2}
Let $\Tcert$ consist of $\cpv{B2_grid}$ tilts: $\cpv{u8_grid_lo}$ values in $[\cpv{u8_t_lo},\cpv{u8_t_cap}]$, the $\cpv{u8_grid_mid}$ values $\cpv{u8_t_a},\cpv{u8_t_b},\dots,\cpv{u8_t_z}$ other than $\cpv{u8_t_skip}$, and $t=1$. For every $t\in\Tcert$ a rational number $M_t$ is certified with
\[
M_t\ge\Lamp(t)
\]
for the family of $\Pzero$. For $t<1$, $M_t=\max\bigl(L_t,\,t\cdot\cpv{cbarP0_hi}-2+2t\bigr)$, where $L_t\ge\Lamcert(t)$ is certified by Lemma~\ref{cp:G2}; $M_1=\cpv{cbarP0_hi}$. Selected values are in Table~\ref{cp:tab:B2}.
\end{lemma}

\begin{proof}
$\Lamp(t)=\max(\Lam(t),t\cstar-2+2t)$. By Lemma~\ref{cp:G1}, $\Lam(t)\le\Lamcert(t)\le L_t$ for every $t\in(0,1]$, and $\cstar\le\co(\Pzero)\le\cpv{cbarP0_hi}$ by Lemmas~\ref{cp:L5} and~\ref{cp:L11}. At $t=1$, $\Lamp(1)=\cstar$ by Corollary~\ref{cp:T5}(ii).
\end{proof}

Every $L_t$ is certified by the method of Section~\ref{cp:sec:sharp}, in interval arithmetic; at 13 of the tilts a second implementation with an independent trust base reproduces all printed digits. On the whole grid $L_t$ is larger than the padding term $t\cdot\cpv{cbarP0_hi}-2+2t$, which is negative for $t<\cpv{tpadP0_lo}$; so $M_t=L_t$ at every $t<1$. Section~\ref{cp:sec:certificates} names the scripts and their outputs.

\begin{table}[ht]
\centering\small
\begin{tabular}{l | l l l l l}
\hline
$t$ & $\cpv{G3_t}$ & $\cpv{u8_t1}$ & $\cpv{u8_t2}$ & $\cpv{u8_t_cap}$ & $\cpv{u8_t4}$\\
$M_t\le$ & $\cpv{u7_lamcert_t0}$ & $\cpv{u8_M1}$ & $\cpv{u8_M2}$ & $\cpv{u8_M3}$ & $\cpv{u8_M4}$\\
\hline
$t$ & $\cpv{u8_t5}$ & $\cpv{u8_t6}$ & $\cpv{u8_t_z}$ & $1$ & \\
$M_t\le$ & $\cpv{u8_M5}$ & $\cpv{u8_M6}$ & $\cpv{u8_M7}$ & $\cpv{cbarP0_hi}$ & \\
\hline
\end{tabular}
\caption{Selected certified profile bounds at $\Pzero$ (Lemma~\ref{cp:B2}), rounded up.}\label{cp:tab:B2}
\end{table}

\subsection{The batch curve at \texorpdfstring{$\Pzero$}{P0}}\label{cp:sec:batchcurve}

With $\alpha=\alpha(\Pzero)$ and $M_t$ in place of $L_t$ in Theorem~\ref{cp:B1}, put
\[
\fcert(\kappa)=\max\Bigl(\alpha(\Pzero)\,[\kappa\le1],\ \max_{t\in\Tcert}\max\bigl(\rdir(\kappa,t),\rDS(\kappa,t)\bigr)\Bigr),
\]
where $\rdir$ and $\rDS$ count as $-\infty$ where their conditions fail.

\begin{theorem}[the batch curve at $\Pzero$]\label{cp:B3}
For every $\kappa>0$ and every $\eps<\fcert(\kappa)$ the $\kappa$-batch statement holds (time model, deterministic). The certified values of Table~\ref{cp:tab:B3} are lower bounds on $\fcert(\kappa)$: at each listed $\kappa$ the $\kappa$-batch statement holds for every $\eps$ below the value.
\end{theorem}

\begin{proof}
If $\eps<\fcert(\kappa)$, then $\kappa\le1$ and $\eps<\alpha(\Pzero)$, or there is $t\in\Tcert$ with $\eps<\rdir(\kappa,t)$ or $\eps<\rDS(\kappa,t)$; Theorem~\ref{cp:B1}(c), (a) or (b) with $L_t=M_t$ (Lemma~\ref{cp:B2}) applies. For the table, the certificate takes for each $\kappa$ a tilt $t\in\Tcert$, a bottom and a rational $x$ next to $x_{\min}$ or $x_{\max}$, checks (B1) strictly in interval arithmetic with the lower end of $\alpha(\Pzero)$, the upper end of $M_t$ and, for the data-structure bottom, $q_A\to0$ and the lower end of $\es$, and prints the lower end of the range at that $x$ rounded down to six digits. This is valid because (B1) and the range at fixed $x$ improve as $\alpha$ increases and $L_t$ decreases; and the printed values are below $\fcert(\kappa)$, since $\rdir$ and $\rDS$ increase in $\alpha$ and decrease in $L_t$.
\end{proof}

\begin{table}[ht]
\centering\small
\begin{tabular}{l l l | l l l}
\hline
$\kappa$ & every $\eps<$ & bottom & $\kappa$ & every $\eps<$ & bottom\\
\hline
$\cpv{u8_k_a}$ & $\cpv{u8_f_a}$ & direct & $\cpv{B3_kc}$ & $\cpv{B3_fc}$ & data structure\\
$\cpv{B3_ka}$ & $\cpv{B3_fa}$ & direct & $\cpv{B3_kd}$ & $\cpv{B3_fd}$ & data structure\\
$\cpv{u8_k_b}$ & $\cpv{u8_f_b}$ & direct & $\cpv{B3_ke}$ & $\cpv{B3_fe}$ & data structure\\
$\cpv{B3_kb}$ & $\cpv{B3_fb}$ & direct & $\cpv{B3_kf}$ & $\cpv{B3_ff}$ & data structure\\
$\cpv{u8_k_c}$ & $\cpv{u8_f_c}$ & direct & $\cpv{B3_kg}$ & $\cpv{B3_fg}$ & data structure\\
$\cpv{u8_k_d}$ & $\cpv{u8_f_d}$ & data structure & $\cpv{B3_kh}$ & $\cpv{B3_fh}$ & data structure\\
$\cpv{u8_k_e}$ & $\cpv{u8_f_e}$ & data structure & $\cpv{B3_ki}$ & $\cpv{B3_fi}$ & data structure\\
 & & & $\cpv{B3_kj}$ & $\cpv{B3_fj}$ & data structure\\
 & & & $\ge\cpv{B3_kk}$ & $\cpv{B3_floor}$ & data structure, $t=1$\\
\hline
\end{tabular}
\caption{Certified values of the batch curve at $\Pzero$ (Theorem~\ref{cp:B3}); ``data structure'' is the bottom of Theorem~\ref{cp:P3} with $q_A\to0$. The parameters $(t,x)$ of each value are in the certificate output.}\label{cp:tab:B3}
\end{table}

The curve has its largest certified value, $\cpv{B3_fb}$, at $\kappa=1$, that is at $\abs\Bt=N^2/D$, stays above $\alpha(\Pzero)$ for $\kappa$ slightly above~$1$ with the direct bottom, falls steeply once the data-structure bottom takes over, and from $\kappa=\cpv{B3_kk}$ on equals its $t=1$ value, at least $\cpv{B3_floor}$, the per-entry value at $\Pzero$ (Corollary~\ref{cp:B7}). $\fcert$ itself is non-increasing on $[1,\infty)$: $s=1-\kappa(1-t)$ decreases in $\kappa$, $\rdir$ and $\rDS$ do not decrease in $s$, $\rdir$ decreases to $\alpha$ as $\alpha s$ decreases to $L_t$, and below that point $\rDS$ continues from $\alpha$ downwards.

\begin{corollary}[explicit savings at $\kappa=\frac12$ and $\kappa=1$]\label{cp:B8}
At $\Pzero$:
\begin{enumerate}[label=(\alph*)]
\item $\Hypk1(\eps,\gamma)$ holds for every $0<\eps<\alpha(\Pzero)$ and every $0<\gamma<\min\{\cpv{u8_gD_one},\,1/\eps-1/\alpha(\Pzero)\}$; the corresponding statement for $\kappa=\frac12$ is Theorem~\ref{cp:G3}, with $\cpv{B8_gamma_half}$ in place of $\cpv{u8_gD_one}$ (rounded).
\item The range of $\Hypk\kappa$ exceeds $\alpha(\Pzero)$ by at least $\cpv{u8_exc_half}$ at $\kappa=\frac12$ and by at least $\cpv{u8_exc_one}$ at $\kappa=1$.
\end{enumerate}
\end{corollary}

\begin{proof}
(a) \ThmDk{}(b) with $\kappa=1$ and $t=\cpv{u8_t1}\in\Tcert$; since $t<\cpv{tpadP0_lo}$, $\Lamp(t)=\Lam(t)\le M_t$, and the certificate gives $t-M_t/\alpha(\Pzero)>\cpv{u8_gD_one}$ (lower end of $\alpha$). The third term $2/\alpha(\Pzero)$ exceeds $1\ge\gamma$. (b) Theorem~\ref{cp:B3} at $\kappa=\frac12$ and $\kappa=1$ (both with the direct bottom; for $\kappa\le1$ the $\kappa$-batch statement implies $\Hypk\kappa(\eps,\min(\gamma_{\mathrm{pre}},g,\kappa))$) and $\alpha(\Pzero)<\cpv{alphaP0_hi}$: $\cpv{B3_fa}-\cpv{alphaP0_hi}>\cpv{u8_exc_half}$ and $\cpv{B3_fb}-\cpv{alphaP0_hi}>\cpv{u8_exc_one}$.
\end{proof}

\subsection{Monotonicity in \texorpdfstring{$\kappa$}{kappa}}\label{cp:sec:mono}

\begin{lemma}[splitting, $\kappa\ge1$]\label{cp:B4}
Let $1\le\kappa_1\le\kappa_2$ and $\eps\le2$. If the $\kappa_2$-batch statement holds at $\eps$ with constants $\gamma_{\mathrm{pre}}$, $g$, $C$, where without loss of generality $g\le1$, then the $\kappa_1$-batch statement holds at $\eps$ with the same $\gamma_{\mathrm{pre}}$ and $g$ and the constant $8C+O(1)$.
\end{lemma}

\begin{proof}
Write $b_1=b_{\kappa_1}$ and $b_2=b_{\kappa_2}$, so $1\le b_2\le b_1$. Phase~1 is unchanged. The new phase~2 reads a batch $\Bt$ with $\abs\Bt\le b_1$ and cuts it into consecutive chunks of at most $b_2$ positions, in $O(\abs\Bt)$ steps. For each chunk it runs the old phase~2 with a \emph{write journal}: before every write to memory it appends the address and the old content to a journal kept in fresh memory (above every address the old algorithm can use, at the cost of one more bit per word), and after the chunk's output it undoes the journal in reverse order. Every chunk therefore starts from the memory left by phase~1 (cells never written are unchanged, written cells are restored), so the old worst-case bound $Cb_2D^{1-g}$ applies to each chunk. Journaling turns every step into at most three steps and the undo costs one step per journal entry, so a chunk costs at most $4Cb_2D^{1-g}+O(b_2)$. There are at most $\lceil b_1/b_2\rceil\le2b_1/b_2$ chunks, so the total is at most $8Cb_1D^{1-g}+O(b_1)\le(8C+O(1))\,b_1D^{1-g}$, using $g\le1$. At $\kappa_1=1$ the target $N^2D^{-g}$ is also met, since $b_1=\lfloor N^2/D\rfloor\le N^2/D$ when $D\le N^\eps\le N^2$.
\end{proof}

For \ThmB's own algorithm no journal is needed: phase~2 writes only bucket heads, which it resets through its list, and table cells, whose stamps it can take as the pair (batch, position).

\begin{lemma}[inclusion, $\kappa\le1$]\label{cp:B5}
If $\kappa_1\le\kappa_2\le1$ and the $\kappa_1$-batch statement holds at $\eps$, then the $\kappa_2$-batch statement holds at $\eps$.
\end{lemma}

\begin{proof}
A batch of at most $b_{\kappa_2}$ positions has at most $b_{\kappa_1}$ positions, and both targets are $CN^2D^{-g}$.
\end{proof}

\begin{corollary}\label{cp:B6}
\begin{enumerate}[label=(\alph*)]
\item $\fbat$ is non-decreasing on $(0,1]$, and $\min\{\fbat(\kappa_2),2\}\le\fbat(\kappa_1)$ for $1\le\kappa_1\le\kappa_2$; so $\min\{\fbat(\kappa),2\}\le\fbat(1)$ for every $\kappa>0$, and $\fbat(1)\ge\cpv{B3_fb}$.
\item Every certified lower bound $v<2$ on $\fcert(\kappa_i)$, in particular every value of Table~\ref{cp:tab:B3}, is a lower bound for $\fbat$ on $[1,\kappa_i]$ if $\kappa_i\ge1$, and on $[\kappa_i,1]$ if $\kappa_i\le1$.
\end{enumerate}
\end{corollary}

\begin{proof}
(a) Lemma~\ref{cp:B5} for $\kappa\le1$, and Lemma~\ref{cp:B4} for every $\eps\le2$ for $\kappa\ge1$; the value at $\kappa=1$ is Table~\ref{cp:tab:B3}. (b) The same lemmas, applied at every $\eps<v$, since $v\le\fcert(\kappa_i)\le\fbat(\kappa_i)$ by Theorem~\ref{cp:B3} and $v<2$.
\end{proof}

Splitting gives nothing for $\kappa<1$: cutting a batch of $N^2/D^{\kappa_1}$ positions into $D^{\kappa_2-\kappa_1}$ batches of $N^2/D^{\kappa_2}$ costs $D^{\kappa_2-\kappa_1}N^2D^{-g}$, because the target $N^2$ is not linear in the batch size; and by Lemma~\ref{cp:B5} $\fbat$ does not decrease there. A reading of the curve as non-increasing on all of $(0,\infty)$ would be wrong: the certified values rise from $\cpv{u8_f_a}$ at $\kappa=\cpv{u8_k_a}$ to $\cpv{B3_fb}$ at $\kappa=1$.

\begin{corollary}[plateau and floor at $\Pzero$]\label{cp:B7}
\begin{enumerate}[label=(\alph*)]
\item For every $0<\kappa\le\cpv{u8_plateau_end}$ the $\kappa$-batch statement holds for every $\eps<\alpha(\Pzero)$.
\item For every $\kappa\ge1$ the $\kappa$-batch statement holds for every $\eps<\cpv{B3_floor}$, the per-entry value at $\Pzero$.
\end{enumerate}
\end{corollary}

\begin{proof}
(a) For $\kappa\le1$ this is Theorem~\ref{cp:B1}(c). For $1<\kappa\le\cpv{u8_plateau_end}$ it is Corollary~\ref{cp:B6}(b) applied to the certified value $\fcert(\cpv{u8_plateau_end})\ge\cpv{u8_plateau_f}$ (direct bottom at $t=\cpv{u8_t1}$; Section~\ref{cp:sec:certificates}), which exceeds $\alpha(\Pzero)<\cpv{alphaP0_hi}$. The lower bound certified at $\kappa=\cpv{u8_k_d}$, $\cpv{u8_f_d}$, is below $\alpha(\Pzero)$, so the certificates of Table~\ref{cp:tab:B3} do not extend the plateau further; this says nothing about $\fbat$ there. (b) At $t=1$ and $\kappa\ge1$, $s=1$, so $\rDS(\kappa,1)=\es M_1/(M_1-\alpha+\es)$ does not depend on $\kappa$; its certified value is the entry $\cpv{B3_floor}$ of Table~\ref{cp:tab:B3}, which therefore holds for every $\kappa\ge1$.
\end{proof}

The value $\cpv{B3_floor}$ belongs to the point $\Pzero$; the per-entry threshold $\cpv{per_entry_P23}$ of Theorem~\ref{cp:D1} is larger because it is certified at a different point, $\Ptt$, chosen for the per-entry composite (Section~\ref{cp:sec:composites}). The two numbers do not contradict each other. Since Theorem~\ref{cp:D1} answers single entries after preprocessing, it gives the $\kappa$-batch statement for every $\kappa\ge1$ at every $\eps<\cpv{per_entry_P23}$; so $\fbat(\kappa)\ge\cpv{per_entry_P23}$ for every $\kappa\ge1$, above the values of Table~\ref{cp:tab:B3} from $\kappa=\cpv{B3_ki}$ on. Table~\ref{cp:tab:B3} and Corollary~\ref{cp:B7} describe the single point $\Pzero$.

\subsection{What the curve does not say}\label{cp:sec:batchnot}

\phantomsection\label{cp:B10}%
$\fcert$ is a lower bound on $\fbat$ obtained from one construction at one level-1 point. For $\kappa>1$ no upper bound on $\fbat(\kappa)$ is known, so the steep fall of $\fcert$ just above $\kappa=1$ is a property of the construction, not shown to be one of the problem. For $\kappa\le1$ the ceiling that is known is weak: at a single $\kappa$, an algebraic algorithm (multiplication model) gives only the inequality of Remark~\ref{cp:A8}, which bounds $\eps$ by the dual exponent only in the limit $\kappa\to0$. And $\fbat(\kappa)$ concerns every batch of the size: batches of a special shape can be answered beyond the curve. A whole column is a batch of $N$ positions; after preprocessing, Theorem~\ref{cp:D5} answers every whole column for every $\eps<\cpv{hinted_omv_R32}$ (at the point $\Pthree$), above the value $\cpv{per_entry_P23}$ that Theorem~\ref{cp:D1} guarantees for arbitrary batches of that size (and $\cpv{B3_floor}$ of the curve at $\Pzero$).

\subsection{Full columns by fast multiplication}\label{cp:sec:fastcols}

\begin{proposition}[full columns without preprocessing]\label{cp:B9}
Let $\omega<3$ be an upper bound on the exponent of square matrix multiplication; any published one will do, for instance $\omega\le\cpv{omega_vxxz}$~\cite{VXXZ24}. For every $\eps>0$, $\beta>0$ and $g<(3-\omega)\min(1,\beta/\eps)$ there is a constant $C$ such that, for all $N\ge1$ and $2\le D\le N^\eps$, all $X\in\Z^{N\times D}$ and $Y\in\Z^{D\times N}$ with entries of absolute value $N^{O(1)}$, and every set of $k\ge N^\beta$ column indices, the $k$ full columns of $XY$ are computed with no preprocessing in at most $C\abs\Bt D^{1-g}$ steps, $\abs\Bt=Nk$ (time model, deterministic).
\end{proposition}

\begin{proof}
We use the standard fact that for every $\delta>0$ there is $C_\delta$ such that two $m\times m$ integer matrices with entries $N^{O(1)}$, $m\le N$, are multiplied in at most $C_\delta m^{\omega+\delta}$ steps (a fixed exact bilinear algorithm of rank at most $m_0^{\omega+\delta}$, which exists since rank and border rank give the same exponent~\cite{Bin80}, applied recursively; all numbers stay $N^{O(1)}$). Let $b=\ln k/\ln D$; since $k\ge N^\beta\ge D^{\beta/\eps}$, $b\ge\beta/\eps$. The wanted product is $X$ times the $D\times k$ matrix of the wanted columns of $Y$; cutting the rows of $X$ into $N/k$ blocks of $k$ rows (padding by zero rows), it is $N/k$ products $\mat kDk$. If $D\le k$, cut each into $(k/D)^2$ products $\mat DDD$ (padding $k$ to a multiple of $D$); the cost per output entry is $C_\delta D^{\omega+\delta}/D^2=C_\delta D^{1-(3-\omega-\delta)}$. If $D>k$, cut the inner dimension into $D/k$ products $\mat kkk$ (padding $D$ to a multiple of $k$) and add them; the cost per entry is $(D/k)(C_\delta k^{\omega-2+\delta}+1)\le(C_\delta+1)D^{1-b(3-\omega-\delta)}$. The paddings cost a factor at most $4$, reading the input costs $O(ND)\le O(\abs\Bt D^{1-g})$ since $k\ge D^{\beta/\eps}\ge D^g$, and the constant $C_\delta$ is independent of $D$. So the cost per entry is $O(D^{1-(3-\omega-\delta)\min(1,\beta/\eps)})$; choosing $\delta$ with $(3-\omega-\delta)\min(1,\beta/\eps)>g$ gives the claim.
\end{proof}

So for batches of $N^\beta$ full columns, $\beta>0$ fixed, every $\eps$ is in the range without any preprocessing; the column threshold $\cpv{hinted_omv_R32}$ of Theorem~\ref{cp:D5} (at $\Pthree$) concerns single columns, or batches of $N^{o(1)}$ of them.

\begin{remark}[a column-union composite with a fast-multiplication bottom; the inequality is proved, the conclusion argued]\label{cp:B11}
One may try to combine the two mechanisms: run the column composite of Proposition~\ref{cp:T9} on a batch of $k=D^b$ columns at once, and compute the sub-columns that each top term needs by fast rectangular multiplication. Let $\gMM(b')$ be the per-entry saving of fast multiplication for $\mat{k'}{D_A}{k'}$ with $k'=D_A^{b'}$, and assume only that $\gMM$ is concave and non-decreasing on $[0,\infty)$ with $\gMM(0)=0$ and $\gMM\le1$. (For the fast-multiplication saving $\gMM(b)=1+2b-b\,\omega(1,1/b,1)=1+2b-\omega(b,1,b)$ concavity holds, since $\omega$ is convex by its subadditivity and homogeneity~\cite{LR83}, \cite[Facts~2.4 and~2.5]{ABFK24}; it is non-decreasing because $\omega(b+h,1,b+h)\le\omega(b,1,b)+2h$, and $\gMM(0)=0$ because $\omega(0,1,0)=1$.) With the parameters of Proposition~\ref{cp:T9} at a level-1 point with rank exponent $\alpha$ and column-union exponent $\ccol$, the union bound over the $k$ columns and Jensen's inequality over the top terms bound the cost of this composite per output entry by $D^{Q_{\mathrm{comb}}+o(1)}$ with
\[
Q_{\mathrm{comb}}=\frac{x(\ccol-1)+1-\gMM(b_A)}{1+\alpha x},\qquad b_A=b(1+\alpha x)-(2-\ccol)x,
\]
where $\gMM(b_A):=0$ if $b_A\le0$, while fast multiplication alone costs $D^{1-\gMM(b)}$ per entry. \emph{If $\alpha<\ccol-1$ and $\ccol\le2$, then $Q_{\mathrm{comb}}>1-\gMM(b)$ for every $x>0$ and $b>0$.} Indeed $(1+\alpha x)(Q_{\mathrm{comb}}-1+\gMM(b))=x(\ccol-1-\alpha)+\gMM(b)(1+\alpha x)-\gMM(b_A)$. If $b_A\le0$ this is at least $x(\ccol-1-\alpha)>0$. Otherwise, with a supergradient $\gMM'(b)\ge0$ at $b$, concavity gives $\gMM(b_A)\le\gMM(b)+\gMM'(b)(b_A-b)$ with $b_A-b=x(\alpha b-(2-\ccol))$, so the expression is at least $x\bigl[(\ccol-1-\alpha)+\alpha(\gMM(b)-b\gMM'(b))+(2-\ccol)\gMM'(b)\bigr]$, and $\gMM(b)-b\gMM'(b)\ge\gMM(0)=0$ by concavity. At the column point $\Pthree$ of Theorem~\ref{cp:D5}, $\ccol(\Pthree)\approx\cpv{u8_ccol_P32}$ and $\alpha(\Pthree)\approx\cpv{u8_alpha_P32}$ (computed), so the hypothesis holds there. So the cost \emph{bound} of the composite is never below that of plain fast multiplication. That the composite itself never beats fast multiplication would need the bound to be attained on some batch; for columns at distinct positions this is argued, not proved.
\end{remark}

\section{The transfer theorem and its converse}\label{cp:sec:transfer}

\ThmDk{} (Theorem~\ref{cp:T7}) uses an input family only through \Hcond1--\Hcond4{} and its profile. For the level-1 family, Section~\ref{cp:sec:level1} obtains $\cg=0$ from the merged decomposition, the symmetry group of the laser restriction and the moment bound. This section asks which decompositions and which restrictions of tensor powers give $\cg=0$, or at least $\cg<\alpha/2$, which already gives $\Hyp(\eps,\gamma)$ for every $\eps\le\min(\alpha,\frac12)$ (Corollary~\ref{cp:C2:half}).

The transfer theorem (Theorem~\ref{cp:C2}) answers the first question for every base tensor: if the tensor has a border decomposition with as many terms as $z$-variables whose supports are triangular with respect to an integer weight, then every family of restrictions of its powers by variable maps, with few preimages per output, rank exponent~$2$, linear copy size and a window on the weights of the preimages of the kept outputs, has $\cg=0$. Its proof needs no symmetry group and no type classes. Conversely (Theorem~\ref{cp:C6}), for decompositions with as many terms as $z$-variables, every partition of every such family has structural typical support exponent at least $\ktyp/\Acs$, where $\ktyp$ depends on the base decomposition only, and $\ktyp>0$ exactly when the decomposition is not triangular (Corollary~\ref{cp:C7}, Proposition~\ref{cp:C8}). With composition classes the exponent is pinned between $\ktyp/\Acs$ and a maximum over the limit types of the preimages of the kept outputs (Theorem~\ref{cp:C9}). Triangularity is therefore what gives $\cg=0$ (for the structural profile, and for the profile itself when every kept output has one preimage); it is not what \ThmD{} needs for the range $\alpha$: Example~\ref{cp:C13} is a non-triangular decomposition with $0<\cg<\alpha/2$, so \ThmD{} reaches its family's $\alpha$ all the same.

Two further statements concern top algorithms whose supports are not known: Proposition~\ref{cp:C14} mixes any family of exact algorithms with a family with $\cg=0$, and Proposition~\ref{cp:C17} shows that in the multiplication model the question is one about the rank of matrix products restricted to sparse output sets. The section ends with the instance of~\cite{VXXZ24}, which is conditional on a hypothesis on the parameters of their program, the one behind their published bound (Theorem~\ref{cp:C20}, Remark~\ref{cp:C20:F2}).

Every range statement of this section is in the time model of Definition~\ref{cp:P1} and goes through \ThmD{} or \ThmDk, except Proposition~\ref{cp:C17}, which is in the multiplication model; the statements about profiles are combinatorial.

\subsection{Variable-map families}\label{cp:sec:tr:setting}

\begin{definition}[variable-map families]\label{cp:C1:def}
Let $T$ be a tensor over $\Q$ whose set of $z$-variables is $Z$, $\abs Z=R\ge2$, and let $\mathcal D$ be a border decomposition of $T$ with terms $l\in[R']$, $z$-forms $w_l$, approximation degree $h$ and product degree $e_0$. For a term string $\tau\in[R']^M$ and a $z$-string $z\in Z^M$ put $\supp\tau=\prod_{i}\supp w_{\tau_i}$; so $z\in\supp\tau$ if and only if $z_i\in\supp w_{\tau_i}$ for every $i$, and $\supp\tau$ is the support of the $z$-form $w_\tau=w_{\tau_1}\otimes\dots\otimes w_{\tau_M}$ of the term $\tau$ of $\mathcal D^{\otimes M}$.

A \emph{variable-map family} for $(T,\mathcal D)$ consists of, for every $M$ in an infinite set of positive integers, an integer $Q_M\ge1$ and a restriction $(\varphi_X,\varphi_Y,\varphi_Z)$ of the direct sum of $Q_M$ copies of $T^{\otimes M}$ onto $s_M\cp\mat{a_M}{b_M}{a_M}$ that sends every variable to $0$ or to one variable with coefficient~$1$. Its \emph{terms} at level $M$ are the pairs $(j,\tau)$, $j\in[Q_M]$, $\tau\in[R']^M$: the term $\tau$ of $\mathcal D^{\otimes M}$ on copy $j$, pushed through the restriction. Its \emph{kept outputs} are the $s_Ma_M^2$ $z$-variables of $s_M\cp\mat{a_M}{b_M}{a_M}$. The \emph{preimage set} of a kept output $o$ is $\Pre(o)=\{(j,z):\varphi_Z\text{ maps the variable $z$ of copy $j$ to }o\}$, and $p_M=\max_o\abs{\Pre(o)}$. A term $(j,\tau)$ \emph{covers $o$ structurally} if $(j,z)\in\Pre(o)$ for some $z\in\supp\tau$. The \emph{structural profile} of a partition of a set of terms is defined as the profile of Definition~\ref{cp:T1} ($Z_M$, $\mu_M$, $\Lam$, $\cg$), with structural coverage in place of coverage; we write $\Lam^{\mathrm{str}}$ and $\cg^{\mathrm{str}}$.
\end{definition}

\begin{lemma}[pushing a decomposition through a variable-map restriction]\label{cp:C1}
In the setting of Definition~\ref{cp:C1:def}, let $A_M$ consist of the pushed terms of nonzero $z$-form at level $M$.
\begin{enumerate}[label=(\alph*)]
\item $A_M$ is a border decomposition of $s_M\cp\mat{a_M}{b_M}{a_M}$ with approximation degree $hM$ and product degree at most $e_0M$.
\item The support of the pushed $z$-form of $(j,\tau)$ consists of kept outputs that $(j,\tau)$ covers structurally.
\item The sets $\Pre(o)$ are nonempty and pairwise disjoint, every kept output lies in the support of some term of $A_M$, and $A_M$ has at least $s_Ma_M^2$ terms.
\end{enumerate}
Consequently, for every partition of all $Q_M(R')^M$ terms, the induced partition of $A_M$ has $\mu_M\le\mu^{\mathrm{str}}_M$; so every upper bound on $\Lam^{\mathrm{str}}$ or $\cg^{\mathrm{str}}$ holds for the profile of $(A_M)$.
\end{lemma}

\begin{proof}
(a) The $M$-th tensor power of $\mathcal D$ is a border decomposition of $T^{\otimes M}$ with the terms $\tau$, approximation degree $hM$ and product degree at most $e_0M$, since a coefficient of the product $u_\tau v_\tau w_\tau$ is a product of $M$ coefficients of products of $\mathcal D$; the direct sum over the copies is a border decomposition of the direct sum. The maps are $\Q$-linear; extended $\Q[\lambda]$-linearly they commute with taking the coefficient of a power of $\lambda$, so they map the border identity of the direct sum, term by term, to a border identity whose $\lambda^{hM}$-part is the image tensor $s_M\cp\mat{a_M}{b_M}{a_M}$. Terms of zero $z$-form contribute nothing and are dropped. A coefficient of a pushed product is a sum of coefficients of the source product, since every variable goes to $0$ or to one variable with coefficient~$1$; so its degree is at most $e_0M$.
(b) The coefficient of a kept output $o$ in the pushed $z$-form of $(j,\tau)$ is $\sum w_\tau[z]$ over the $z$ with $(j,z)\in\Pre(o)$, where $w_\tau[z]=\prod_iw_{\tau_i}[z_i]$; a nonzero summand needs $z\in\supp\tau$.
(c) Every variable has at most one image, so the preimage sets are disjoint. The image tensor contains $o$ with coefficient $1$ and is the image of the direct sum of the copies, so some variable is mapped to $o$. The $\lambda^{hM}$-coefficient of $\sum\varphi_Z(w_{(j,\tau)})[o]\,\varphi_X(u_{(j,\tau)})\otimes\varphi_Y(v_{(j,\tau)})$ is the bilinear form $(XY)_o\ne0$ of the copy of $o$; so some term has $o$ in its support. The count is Lemma~\ref{cp:T1:basic}(a).
For the consequence: by (b) the coverage of $o$ by the class-$c$ terms of $A_M$ is at most its structural coverage by all class-$c$ terms, the classes of $A_M$ are subsets of those of the partition, and $Z(t)=\sum_c\abs{P_c}^{1-t}\kappa_c^t$ is non-decreasing in every $\abs{P_c}$ and $\kappa_c$.
\end{proof}

\begin{definition}[hypotheses]\label{cp:C2:hyp}
For a variable-map family for $(T,\mathcal D)$ consider:
\begin{itemize}
\item[\Tcond1] ($z$-tight) $R'=R$, and a bijection $\pi\colon[R]\to Z$ with $\pi(l)\in\supp w_l$ for every $l$ is fixed; for a term string, $\pi(\tau)=(\pi(\tau_1),\dots,\pi(\tau_M))$.
\item[\Tcond2] (triangular) there is an integer weight $\wt\colon Z\to\{0,\dots,\wtop\}$, $\wtop\ge1$, with $\supp w_l\subseteq\{\pi(l)\}\cup\{z:\wt(z)<\wt(\pi(l))\}$ for every $l$; for a $z$-string, $\wt(z)=\sum_i\wt(z_i)$.
\item[\Tcond3] (few preimages) $\ln p_M=o(M)$.
\item[\Tcond4] (rank exponent $2$, linear copy size) $\zeta_M:=\ln\bigl(R^MQ_M/(s_Ma_M^2)\bigr)/M\to0$, $0<\Acs:=\liminf_M\ln a_M/M<\infty$, and $\alpha_M:=\ln b_M/\ln a_M\to\alpha$.
\item[\Tcond5] (weight window) there are numbers $\varpi_M$ and $\eta_M\to0$ such that $\varpi_M\le\wt(z)\le\varpi_M+\eta_MM$ for every kept output $o$ and every $(j,z)\in\Pre(o)$.
\end{itemize}
\end{definition}

By Lemma~\ref{cp:C1}(c), $s_Ma_M^2\le\abs{\bigcup_o\Pre(o)}\le Q_MR^M$, so $\zeta_M\ge0$. The number of copies $Q_M$ enters only through $\zeta_M$; it need not be $e^{o(M)}$. In \Tcond1, a bijection $\pi$ exists whenever $T$ is $z$-concise (its $z$-slices are linearly independent) and $R'=R$: the $R\times R$ matrix $W(\lambda)=(w_l[z])$ is invertible over $\Q(\lambda)$, since the $z$-flattening $\sum_l(u_l\otimes v_l)w_l^{\mathsf T}=U(\lambda)W(\lambda)$ has rank $R$ over $\Q(\lambda)$ by the argument of Lemma~\ref{cp:T1:basic}(a), and a nonzero term of the Leibniz expansion of $\det W(\lambda)$ is such a bijection.

\subsection{The transfer theorem}\label{cp:sec:tr:T}

\begin{theorem}[Theorem T, transfer for triangular decompositions]\label{cp:C2}
Let a variable-map family for $(T,\mathcal D)$ satisfy \Tcond1--\Tcond5. Partition the terms at level $M$ into the \emph{excess classes} $P_d=\{(j,\tau):(\wt(\pi(\tau))-\varpi_M)_+=d\}$, and put
\[
g(y)=y\ln\frac{e(R-1)}y\quad(g(0)=0),\qquad\cH=\frac2{\wtop^2},\qquad\PhiT(t)=\frac1\Acs\max_{y\in[0,1]}\bigl[t\,g(y)-(1-t)\cH y^2\bigr].
\]
\begin{enumerate}[label=(\alph*)]
\item For every $t\in[0,1]$, $\Lam^{\mathrm{str}}(t)\le\PhiT(t)$; hence also $\Lam(t)\le\PhiT(t)$ for the family $(A_M)$ of Lemma~\ref{cp:C1} with the induced partition.
\item $\PhiT(t)=O(t^{3/2}\ln(1/t))$ as $t\to0+$, so $\cg=0$ for $(A_M)$.
\item If $\alpha>0$, then $(A_M)$ with the induced partition, after finitely many $M$ are discarded, is an input family with \Hcond1--\Hcond4{} and limit $\alpha$.
\item If $0<\alpha<1$, then for every $\kappa>0$ and every $0<\eps\le\alpha$ there is $\gamma>0$ such that $\Hypk\kappa(\eps,\gamma)$ holds (time model); in particular $\Hyp(\eps,\gamma)$ holds with some $\gamma>0$ for every $0<\eps\le\min(\alpha,\frac12)$.
\end{enumerate}
\end{theorem}

\begin{proof}
Put $\wbar=\frac1R\sum_{z\in Z}\wt(z)$. Since all weights are nonnegative and every kept output has a preimage, we may replace $\varpi_M$ by $\max(\varpi_M,0)$ and assume $0\le\varpi_M\le\wtop M$; then the excess takes at most $\wtop M+1$ values in $[0,\wtop M]$.

\emph{(i) Deficit count.} Let $z\in\supp\tau$. At every coordinate either $z_i=\pi(\tau_i)$ or, by \Tcond2, $\wt(z_i)\le\wt(\pi(\tau_i))-1$, since weights are integers. Hence $\wt(z)\le\wt(\pi(\tau))$, and $z$ differs from $\pi(\tau)$ in at most $\wt(\pi(\tau))-\wt(z)$ coordinates. So for a fixed $z$ and $d\ge0$, every $\tau$ with $z\in\supp\tau$ and $\wt(\pi(\tau))\le\wt(z)+d$ agrees with $\pi^{-1}(z)$ outside at most $\min(d,M)$ coordinates, and the number of such $\tau$ is at most
\[
\sum_{i\le\min(d,M)}\binom Mi(R-1)^i\le(d+1)\,e^{M\tilde g(d/M)},\qquad\tilde g(y)=g(\min(y,1)),
\]
because $\binom Mi(R-1)^i\le(eM(R-1)/i)^i=e^{Mg(i/M)}$ and $g$ is increasing on $[0,R-1]\supseteq[0,1]$ ($g'(y)=\ln((R-1)/y)$, $R\ge2$). The minimum is needed: for $d>e(R-1)M$, $g(d/M)<0$.

\emph{(ii) Concentration.} For $\tau$ uniform in $[R]^M$, the letters $\pi(\tau_i)$ are independent and uniform on $Z$, so $\wt(\pi(\tau))$ is a sum of $M$ independent variables with values in $[0,\wtop]$ and mean $\wbar$. Hoeffding's inequality~\cite{Hoe63} gives, for every $y\ge0$, at most $R^Me^{-\cH y^2M}$ strings $\tau$ with $\wt(\pi(\tau))\ge M\wbar+yM$, and likewise at most $R^Me^{-\cH y^2M}$ strings $z\in Z^M$ with $\wt(z)\le M\wbar-yM$.

\emph{(iii) Where the window lies.} Put $u_M=\wbar-\varpi_M/M$ and $y_M=u_M-\eta_M$. Every preimage of a kept output has weight at most $\varpi_M+\eta_MM=M(\wbar-y_M)$. If $y_M>0$, Lemma~\ref{cp:C1}(c) and (ii) give $s_Ma_M^2\le\abs{\bigcup_o\Pre(o)}\le Q_MR^Me^{-\cH y_M^2M}$, that is $\cH y_M^2\le\zeta_M$. Hence $u_M\le\eta'_M:=\eta_M+(\zeta_M/\cH)^{1/2}$, and $\eta'_M\to0$.

\emph{(iv) Classes.} A term of class $d>0$ has $\wt(\pi(\tau))=\varpi_M+d=M\wbar+M(d/M-u_M)$, so by (ii) $\abs{P_d}\le Q_MR^Me^{-\cH((d/M-u_M)_+)^2M}$; for $d=0$ we use $\abs{P_0}\le Q_MR^M$. A class-$d$ term $(j,\tau)$ that covers a kept output $o$ structurally has $(j,z)\in\Pre(o)$ with $z\in\supp\tau$; then $\wt(z)\ge\varpi_M$ by \Tcond5{} and $\wt(\pi(\tau))\le\varpi_M+d$, so by (i) there are at most $(d+1)e^{M\tilde g(d/M)}$ such $\tau$ for each of the at most $p_M$ preimages. So the structural class coverage is $\kappa_d\le p_M(d+1)e^{M\tilde g(d/M)}$. (The width of the window is not used here.)

\emph{(v) Profile.} With at most $\wtop M+1$ classes and $\mu_M(t)=[\ln Z(t)-(1-t)\ln(s_Ma_M^2)]/\ln a_M$, (iv) gives
\[
\mu^{\mathrm{str}}_M(t)\ln a_M\le(1+t)\ln(\wtop M+1)+(1-t)\zeta_MM+t\ln p_M+M\max_dB_M(d/M),
\]
with $B_M(0)=0$ and $B_M(y)=t\tilde g(y)-(1-t)\cH((y-u_M)_+)^2$ for $y>0$. The function $\tilde g$ is concave and non-decreasing on $[0,\infty)$ with $\tilde g(0)=0$ ($g$ is concave on $[0,1]$ and $g'(1)=\ln(R-1)\ge0$), hence subadditive. Put $x=(u_M)_+\le\eta'_M$. For $0<y\le x$, $B_M(y)\le t\tilde g(x)$. For $y>x$, $(y-u_M)_+\ge y-x>0$ and $\tilde g(y)\le\tilde g(y-x)+\tilde g(x)$, so $B_M(y)\le t\tilde g(x)+t\tilde g(y-x)-(1-t)\cH(y-x)^2$. For $v>1$, $t\tilde g(v)-(1-t)\cH v^2$ is below its value at $v=1$; so in all cases $B_M(y)\le t\tilde g(\eta'_M)+\Acs\PhiT(t)$, and the right side is nonnegative. Let $\eta>0$; for all large $M$, $\ln a_M\ge(\Acs-\eta)M$, and dividing,
\[
\mu^{\mathrm{str}}_M(t)\le\frac{(1+t)\ln(\wtop M+1)+\zeta_MM+\ln p_M}{(\Acs-\eta)M}+\frac{t\tilde g(\eta'_M)+\Acs\PhiT(t)}{\Acs-\eta}.
\]
The first term tends to $0$ by \Tcond3{} and \Tcond4, and $\tilde g(\eta'_M)\to0$. So $\Lam^{\mathrm{str}}(t)\le\Acs\PhiT(t)/(\Acs-\eta)$ for every $\eta>0$, which is (a); the statement for $(A_M)$ is Lemma~\ref{cp:C1}.

\emph{(vi) $\PhiT(t)=O(t^{3/2}\ln(1/t))$.} Let $t\le\frac12$ and $\delta_t=(4t\,g(1)/\cH)^{1/2}$, with $g(1)=1+\ln(R-1)$, and let $t$ be so small that $\delta_t\le1$. For $y\in(\delta_t,1]$ the bracket is at most $t\,g(1)-\frac12\cH\delta_t^2=-t\,g(1)<0$; for $y\in[0,\delta_t]$ it is at most $t\,g(\delta_t)$. So $\PhiT(t)\le t\,g(\delta_t)/\Acs=O(t^{3/2}\ln(1/t))$. By (a), $\Lam(t)/t\to0$, and $\Lam\ge0$ by Lemma~\ref{cp:T1:basic}(b), whose proof uses only that $A_M$ is a border decomposition of $s_M\cp\mat{a_M}{b_M}{a_M}$; so $\cg=0$. This is (b).

\emph{(vii) \Hcond1--\Hcond4.} Discard the finitely many $M$ with $\ln a_M<\frac12\Acs M$. \Hcond1: $a_M\to\infty$, $\alpha_M\to\alpha$, and $\ln b_M=\alpha_M\ln a_M\to\infty$ when $\alpha>0$, so $b_M\ge2$ after discarding finitely many more $M$. \Hcond2: the number $r_M$ of terms of $A_M$ satisfies $s_Ma_M^2\le r_M\le Q_MR^M$ (Lemma~\ref{cp:C1}(c)), so $0\le\rho_M-2\le\zeta_MM/\ln a_M\to0$. \Hcond3: the product degree is at most $e_0M$ (Lemma~\ref{cp:C1}(a)) and there are at most $\wtop M+1$ classes, both $e^{O(\ln M)}$ against $\ln a_M\ge\frac12\Acs M$. \Hcond4: a kept output $o$ is covered only by terms $(j,\tau)$ with $(j,z)\in\Pre(o)$ and $z\in\supp\tau$ for some $z$, and for each of the at most $p_M$ preimages there are at most $R^M$ strings $\tau$; so $C_M\le p_MR^M$ and $\ln C_M/\ln a_M\le2(\ln p_M/M+\ln R)/\Acs$, which is bounded by \Tcond3. This is (c).

(d) By (b) and (c), Theorem~\ref{cp:T7}(a) applies.
\end{proof}

\begin{corollary}[$\cg<\alpha/2$ suffices]\label{cp:C2:half}
Let $(A_M)$ be an input family with \Hcond1--\Hcond4, limit $\alpha\in(0,1)$ and $\cg<\alpha/2$. Then $\Hyp(\eps,\gamma)$ holds with some $\gamma>0$ for every $0<\eps\le\min(\alpha,\frac12)$ (time model).
\end{corollary}

\begin{proof}
By the definition of $\cg$ as a lower limit there are arbitrarily small $t>0$ with $\Lam(t)<t\alpha/2$; take one with $t\le2/(\cstar+2)$, so that $\Lamp(t)=\Lam(t)$ (Corollary~\ref{cp:T5}(i)). At $\kappa=\frac12$ the numerator of $\Qgain{\mathrm D}{1/2}(x,t)$ is $x(t\alpha/2-\Lam(t))+t/2-1$, which is positive for all large $x$, so \Dcond{} holds. Since $\alpha<1$, $\Erange{\mathrm D}(x)=(1+\alpha x)/(1+x)>\alpha$, so every $0<\eps\le\alpha$ satisfies $\eps<\Erange{\mathrm D}(x)$, and Theorem~\ref{cp:T7} gives $\Hypk{1/2}(\eps,\gamma)$ for every $0<\gamma<\gbound{\mathrm D}{1/2}(x,t,\eps)$, which is positive. For $\eps\le\frac12$ this is $\Hyp(\eps,\gamma)$.
\end{proof}

\begin{remark}[\Tcond5{} for zero-outs]\label{cp:C3}
Suppose that $Z$ is partitioned into blocks and $\wt$ is constant on every block, that the restriction is a zero-out followed by a bijection from the surviving variables of all copies onto the variables of $s_M\cp\mat{a_M}{b_M}{a_M}$ (so $p_M=1$), and that at every level the block strings of the preimages of all kept outputs have one common type (the number of coordinates in each block). Then all these preimages have the same weight $\sum_\beta n_\beta\wt_\beta$, where $n_\beta$ is the number of coordinates in block $\beta$ and $\wt_\beta$ the weight on it, and \Tcond5{} holds with $\eta_M=0$. This is the case of the merged decomposition $\Fmrg$ of $F_q$ with the weight $\wt(z_0)=0$, $\wt(z_i)=1$ ($1\le i\le q$), $\wt(z_{q+1})=2$, which is the block index of Section~\ref{cp:sec:laser}, under the zero-out of Lemma~\ref{cp:L4} with one value of $q$, whose kept outputs have the kept output type (Section~\ref{cp:sec:std:Fq}).
\end{remark}

\begin{remark}[a block-constant weight is needed; argued, not proved]\label{cp:C3:nec}
If $\wt$ is not constant on the blocks of the restriction, \Tcond5{} generally fails, and the following argument suggests that the conclusion does too: a single copy then contains kept outputs of many weights; a wanted set filled with outputs of slightly lower weight than the typical one, which form a fraction $e^{-M\iota}$ of the outputs for a small $\iota>0$ and so far more than $K^{-\alpha/2}$ of them ($K$ the side of the top algorithm of \AlgD{} built from the family), is covered by $e^{Mg(\eta)}$ typical terms per output. So the transfer appears to need triangularity with respect to a weight that is constant on blocks. We do not use this remark.
\end{remark}

\subsection{Triangular decompositions}\label{cp:sec:tr:tri}

For a decomposition with $R'=R$ and a bijection $\pi$ as in \Tcond1, let $\Dpi$ be the directed graph on the terms with an edge $l\to l'$ whenever $l'\ne l$ and $\pi(l')\in\supp w_l$, and call $\pi$ a \emph{perfect matching} of the support graph (the bipartite graph between terms and $z$-variables with the edges $(l,z)$, $z\in\supp w_l$). The condition $l'\ne l$ matters: since $\pi(l)\in\supp w_l$, without it every vertex would carry a loop and every graph would have a cycle.

\begin{proposition}[triangular $\Leftrightarrow$ acyclic]\label{cp:C4}
Let $\mathcal D$ have $R'=R$ terms and let $\pi$ be a perfect matching of its support graph. The following are equivalent:
\begin{enumerate}[label=(\alph*)]
\item $\mathcal D$ is triangular with respect to $\pi$, that is, \Tcond2{} holds for some integer weight;
\item $\pi$ is the only perfect matching of the support graph;
\item $\Dpi$ has no directed cycle;
\item for one, equivalently for every, $m\ge1$, the term string $\tau$ in which every term occurs exactly $m$ times (in any order) has exactly one $z$-string of the uniform type (every $z$-variable exactly $m$ times) in $\supp\tau$.
\end{enumerate}
In particular, whether \Tcond2{} holds does not depend on the choice of $\pi$.
\end{proposition}

\begin{proof}
(a)$\Rightarrow$(b): if $\sigma$ is a perfect matching, then $\sum_l\wt(\sigma(l))=\sum_z\wt(z)=\sum_l\wt(\pi(l))$, while $\wt(\sigma(l))\le\wt(\pi(l))$ for every $l$ with equality only if $\sigma(l)=\pi(l)$; so $\sigma=\pi$.
(b)$\Rightarrow$(c): a simple cycle $l_1\to\dots\to l_k\to l_1$, $k\ge2$, gives a second perfect matching $\sigma$ with $\sigma(l_j)=\pi(l_{j+1})$ (indices mod $k$) and $\sigma=\pi$ elsewhere.
(c)$\Rightarrow$(a): put $\wt(\pi(l))$ equal to the length of a longest directed path starting at $l$; an edge $l\to l'$ gives $\wt(\pi(l'))\le\wt(\pi(l))-1$, so every $z=\pi(l')\in\supp w_l$ with $l'\ne l$ has $\wt(z)<\wt(\pi(l))$. (If all these lengths are $0$, add $1$ to every weight to have $\wtop\ge1$.)
(a)$\Rightarrow$(d), for every $m$ and indeed for every term string $\tau$: $\pi(\tau)\in\supp\tau$, and a $z$-string in $\supp\tau$ of the same type as $\pi(\tau)$ has the same weight, so by step (i) of the proof of Theorem~\ref{cp:C2} it differs from $\pi(\tau)$ in no coordinate. For the uniform string $\tau$, $\pi(\tau)$ has the uniform type.
not (c)$\Rightarrow$ not (d), for every $m$: for a simple cycle $l_1\to\dots\to l_k\to l_1$, $k\ge2$, choose for each $j$ one coordinate $i_j$ with $\tau_{i_j}=l_j$, and change $\pi(\tau)$ at $i_j$ to $\pi(l_{j+1})$. The new string lies in $\supp\tau$ (edges), differs from $\pi(\tau)$, and has the same type, since each $\pi(l_j)$ loses one occurrence and gains one.
Since (d) for one $m$ implies (c), and (c) implies (d) for every $m$, the two readings of (d) agree; and (b) does not refer to a choice of $\pi$.
\end{proof}

\begin{example}[the two decompositions of $F_q$]\label{cp:C4:ex}
Both $\Fmrg$ and $\Fstd$ have $q+2$ terms, as many as $F_q$ has $z$-variables. For $\Fmrg$ take $\pi(A_i)=z_i$, $\pi(\Bp)=z_{q+1}$, $\pi(\Oterm)=z_0$ and the weight $(\wt(z_0),\wt(z_i),\wt(z_{q+1}))=(0,1,2)$: by Lemma~\ref{cp:L2} the supports are $\{z_0,z_i\}$, all of $Z$, and $\{z_0\}$, so \Tcond2{} holds, and $\wt$ is the block index of Section~\ref{cp:sec:laser}. For $\Fstd$ the supports are $\{z_0,z_i\}$ for $A_i$, $\{z_0,\dots,z_q\}$ for $B$ and $\{z_0,z_{q+1}\}$ for $C$; a perfect matching must send $C$ to $z_{q+1}$, and then either $B$ to $z_0$ and every $A_i$ to $z_i$, or one $A_j$ to $z_0$, $B$ to $z_j$ and the other $A_i$ to $z_i$. So there are $q+1$ perfect matchings, and $\Fstd$ is not triangular by Proposition~\ref{cp:C4}.
\end{example}

\begin{proposition}[upper sets of the weight]\label{cp:C5}
Let $T$ be $z$-concise and let $\mathcal D$ be a border decomposition of $T$ with $R'=R$ terms that is triangular with respect to $\pi$ and $\wt$. For a threshold $v$ let $Z_{\ge v}=\{z:\wt(z)\ge v\}$ and let $T|_{Z_{\ge v}}$ be $T$ with the $z$-variables outside $Z_{\ge v}$ set to $0$. Then $\bR(T|_{Z_{\ge v}})=\abs{Z_{\ge v}}$ for every $v$, and every $z$-variable of the largest weight has a $z$-slice of rank $1$. Consequently, a $z$-concise tensor all of whose $z$-slices have rank at least $2$ has no triangular decomposition with as many terms as $z$-variables.
\end{proposition}

\begin{proof}
A term $l$ has a $z$-form with nonzero restriction to $Z_{\ge v}$ if and only if $\wt(\pi(l))\ge v$: if $\wt(\pi(l))\ge v$ then $\pi(l)\in\supp w_l\cap Z_{\ge v}$, and if $\wt(\pi(l))<v$ every variable of $\supp w_l$ has weight at most $\wt(\pi(l))<v$. Since $\pi$ is a bijection, exactly $\abs{Z_{\ge v}}$ terms survive the zero-out, and they form a border decomposition of $T|_{Z_{\ge v}}$. Conversely, the border rank of $T|_{Z_{\ge v}}$ is at least the rank of its $z$-flattening (as in Lemma~\ref{cp:T1:basic}(a)), which is $\abs{Z_{\ge v}}$ because the $z$-slices of $T$ are linearly independent. For $z$ of the largest weight, $\pi^{-1}(z)$ is the only term with $z$ in its support, so the $z$-slice of $T$ is the $\lambda^h$-coefficient of $w_{\pi^{-1}(z)}[z]\,u\otimes v$, a limit of rank-one matrices, of rank at most $1$, and nonzero by conciseness.
\end{proof}

\subsection{Mixing top algorithms}\label{cp:sec:tr:mix}

The next proposition uses a family of exact algorithms of which nothing is known but their size, rank and largest coverage, and mixes it by tensor products with a top algorithm of a family with $\cg=0$. It rests on \ThmD{} for one top algorithm (Theorem~\ref{cp:T6}), not on \ThmDk, because the top algorithms it builds do not come from one input family.

\begin{lemma}[products of top algorithms]\label{cp:C14:prod}
Let $B$ and $B'$ be exact top algorithms over $\Q$ for $\mat kbk$ and $\mat{k'}{b'}{k'}$ with partitions. Then $B\otimes B'$ is an exact top algorithm over $\Q$ for $\mat{kk'}{bb'}{kk'}$, and with the pair classes $P_c\times P'_{c'}$ and $\theta=\ln k/\ln(kk')$,
\[
\alpha_{B\otimes B'}=\theta\alB+(1-\theta)\alpha_{B'},\qquad\lambda_{B\otimes B'}(t)=\theta\lamB(t)+(1-\theta)\lambda_{B'}(t)\quad(t\in[0,1]),
\]
in particular $\delta_{B\otimes B'}=\theta\delB+(1-\theta)\delta_{B'}$. For tensor powers, $B^{\otimes n}$ with the classes $P_{c_1}\times\dots\times P_{c_n}$ has the same $\alB$, $\lamB$ and $\delB$ as $B$, and maximal coverage $\cB^n$.
\end{lemma}

\begin{proof}
The tensor product of bilinear algorithms for $\mat kbk$ and $\mat{k'}{b'}{k'}$ is a bilinear algorithm for $\mat kbk\otimes\mat{k'}{b'}{k'}=\mat{kk'}{bb'}{kk'}$, with the $z$-forms $w_l\otimes w'_{l'}$, whose supports are $\supp w_l\times\supp w'_{l'}$. So $\Cov((o,o'))=\Cov(o)\times\Cov'(o')$, the class sizes multiply, and the class coverages $\max_{(o,o')}\abs{\Cov(o)\cap P_c}\,\abs{\Cov'(o')\cap P'_{c'}}$ are the products $\kappa_c\kappa'_{c'}$. Hence $Z_{B\otimes B'}(t)=\ZB(t)Z_{B'}(t)$, and $\ln Z=(\lamB(t)+2-2t)\ln k+(\lambda_{B'}(t)+2-2t)\ln k'$ gives the formula for $\lambda$; $\ln(bb')=\ln b+\ln b'$ gives the one for $\alpha$. Powers are the case $B'=B$, repeated.
\end{proof}

\begin{proposition}[mixing top algorithms]\label{cp:C14}
Let $(B_k)_{k\in\mathcal K}$, $\mathcal K$ an infinite set of integers, be exact top algorithms over $\Q$ for $\mat k{b_k}k$ with $\ln b_k/\ln k\to\alpha_{\mathcal B}$ and $R_{B_k}=k^{2+\delta_k}$, $\delta_k\to0$, and let $\cbar_{\mathcal B}=\limsup_k\ln C_k/\ln k$, where $C_k$ is the maximal coverage of $B_k$. Let $\mathcal C$ be an input family with \Hcond1--\Hcond4, limit $\alpha_{\mathcal C}$ and $\cg=0$, with $0<\alpha_{\mathcal C}<\alpha_{\mathcal B}\le1$. Then $\alpha_{\mathcal B}\le\cbar_{\mathcal B}\le2$, and $\Hypk{1/2}(\eps,\gamma)$ holds with some $\gamma>0$ (time model) for every
\[
0<\eps<\alpha_{\mathcal C}+\theta^*(\alpha_{\mathcal B}-\alpha_{\mathcal C}),\qquad\theta^*=\frac{\alpha_{\mathcal C}}{2\cbar_{\mathcal B}-\alpha_{\mathcal B}+\alpha_{\mathcal C}};
\]
for $\eps\le\frac12$ this is $\Hyp(\eps,\gamma)$.
\end{proposition}

\begin{proof}
\emph{Coverage bounds.} An output $o=(i,j)$ of $\mat kbk$ is the bilinear form $\sum_mx_{im}y_{mj}$, of rank $b$ as a bilinear form, and in an exact algorithm it is $\sum_{l\in\Cov(o)}w_l[o](u_l\cdot X)(v_l\cdot Y)$, a sum of $\abs{\Cov(o)}$ forms of rank at most $1$; so $C_k\ge b_k$ and $\cbar_{\mathcal B}\ge\alpha_{\mathcal B}$. Also $C_k\le R_{B_k}$, so $\cbar_{\mathcal B}\le2$. In particular the denominator of $\theta^*$ is at least $\alpha_{\mathcal B}+\alpha_{\mathcal C}>0$, and $\theta^*\le\alpha_{\mathcal C}/(\alpha_{\mathcal B}+\alpha_{\mathcal C})<1$.

\emph{Choices.} Write $a_{\mathcal B}=\alpha_{\mathcal B}$, $a_{\mathcal C}=\alpha_{\mathcal C}$, $\bar c=\cbar_{\mathcal B}$, and for $\theta\in[0,1]$ put $\bar\alpha(\theta)=\theta a_{\mathcal B}+(1-\theta)a_{\mathcal C}$ and $m(\theta)=\bar\alpha(\theta)/2-\theta\bar c$; then $\bar\alpha$ is increasing and $m(\theta)>0$ if and only if $\theta<\theta^*$. Fix $\eps<\bar\alpha(\theta^*)$, choose $\theta_0\in(0,\theta^*)$ with $\eps<\bar\alpha(\theta_0)$, put $m_0=m(\theta_0)>0$ and $g_0=\bar\alpha(\theta_0)-\eps>0$, and choose $\xi>0$ with $\theta_0+\xi<1$ such that $m(\theta)\ge m_0/2$ and $\bar\alpha(\theta)\ge\bar\alpha(\theta_0)-g_0/4$ whenever $\abs{\theta-\theta_0}\le\xi$.
Since $\cg(\mathcal C)=0$, there is $t\in(0,2/(\cstar_{\mathcal C}+2)]$ with $\Lam_{\mathcal C}(t)\le m_0t/8$; there $\Lamp_{\mathcal C}(t)=\Lam_{\mathcal C}(t)$ (Corollary~\ref{cp:T5}(i)). Put $\nu=t\min\{m_0/16,\ g_0/8,\ (1-\theta_0-\xi)(1-a_{\mathcal C})/4\}$.
Corollary~\ref{cp:T5}, applied to $\mathcal C$ at this $t$ with accuracy $\nu$, gives an exact top algorithm $B_{\mathcal C}$ for $\mat{k_{\mathcal C}}{b_{\mathcal C}}{k_{\mathcal C}}$, $b_{\mathcal C}\ge2$, with $\abs{\alpha_{B_{\mathcal C}}-a_{\mathcal C}}\le\nu$, $\delta_{B_{\mathcal C}}\le\nu$ and $\lambda_{B_{\mathcal C}}(t)\le m_0t/8+\nu$. Give each $B_k$ a single class; then $\ZB(t)=R_{B_k}^{1-t}C_k^t$ and $\lambda_{B_k}(t)=(1-t)\delta_k+tc_k$ with $c_k=\ln C_k/\ln k$. Choose $k\in\mathcal K$ with $\abs{\alpha_{B_k}-a_{\mathcal B}}\le\nu$, $c_k\le\bar c+\nu$ and $\delta_k\le\nu$. The numbers $n\ln k/(n\ln k+n'\ln k_{\mathcal C})$, $n,n'\ge1$, are dense in $(0,1)$; choose $n,n'$ with $\theta':=n\ln k/(n\ln k+n'\ln k_{\mathcal C})$ in $[\theta_0-\xi,\theta_0+\xi]$, and let $B=B_k^{\otimes n}\otimes B_{\mathcal C}^{\otimes n'}$ with pair classes.

\emph{The product.} By Lemma~\ref{cp:C14:prod}, $\alB=\theta'\alpha_{B_k}+(1-\theta')\alpha_{B_{\mathcal C}}$, $\delB=\theta'\delta_k+(1-\theta')\delta_{B_{\mathcal C}}\le\nu$, and
\begin{align*}
\frac{t\alB}2-\lamB(t)&\ge\theta'\Bigl[\frac{t(a_{\mathcal B}-\nu)}2-\nu-t(\bar c+\nu)\Bigr]+(1-\theta')\Bigl[\frac{t(a_{\mathcal C}-\nu)}2-\frac{m_0t}8-\nu\Bigr]\\
&\ge t\,m(\theta')-\tfrac52\nu-\frac{m_0t}8\ \ge\ \frac{m_0t}2-\frac{5m_0t}{32}-\frac{m_0t}8\ >\ 0.
\end{align*}
Also $\alB\ge\bar\alpha(\theta')-\nu\ge\eps+5g_0/8$, so $\alB/(1+\delB)\ge\alB-\delB\ge\eps+g_0/2>\eps$; and $\alB\le\theta'+(1-\theta')a_{\mathcal C}+\nu\le1-(1-\theta_0-\xi)(1-a_{\mathcal C})+\nu<1$, so the middle side $b_k^nb_{\mathcal C}^{n'}$ of $B$ is at least $b_{\mathcal C}\ge2$ and at most its outer side $k^nk_{\mathcal C}^{n'}$.

\emph{\ThmD.} At $\kappa=\frac12$ the numerator of $\Qgain B{1/2}(x,t)$ is $x(t\alB/2-\lamB(t))+t/2-1$, positive for all large $x$. Since $\alB<1$, $\Erange B(x)-\alB/(1+\delB)=(1+\delB-\alB)/\bigl((1+x(1+\delB))(1+\delB)\bigr)>0$, so $\eps<\Erange B(x)$ for every $x$. Theorem~\ref{cp:T6} with $B$, $\kappa=\frac12$, a large $x$ and this $t$ gives $\Hypk{1/2}(\eps,\gamma)$ for every $0<\gamma<\gbound B{1/2}(x,t,\eps)$, which is positive.
\end{proof}

\begin{proposition}[optimality of $\theta^*$]\label{cp:C15}
Let $B$ be an exact top algorithm for $\mat kbk$ with a single class and coverage exponent $c=\ln\cB/\ln k$, and $B'$ an exact top algorithm for $\mat{k'}{b'}{k'}$ with any partition and $\alpha_{B'}>0$. If $\theta=\ln k/\ln(kk')\ge\alpha_{B'}/(2c-\alB+\alpha_{B'})$, then $\Qgain{B\otimes B'}{1/2}(x,t)<0$ for every $x>0$ and $t\in(0,1]$: \ThmD{} at $\kappa=\frac12$ gives nothing for $B\otimes B'$ with pair classes.
\end{proposition}

\begin{proof}
As in the proof of Proposition~\ref{cp:C14}, $c\ge\alB$, so the denominator is positive. With one class, $\lamB(t)=(1-t)\delB+tc\ge tc$, and $\lambda_{B'}\ge0$ (Lemma~\ref{cp:T1:basic}). By Lemma~\ref{cp:C14:prod}, $t\alpha_{B\otimes B'}/2-\lambda_{B\otimes B'}(t)\le\frac t2\bigl[\theta\alB+(1-\theta)\alpha_{B'}-2\theta c\bigr]=\frac t2\bigl[\alpha_{B'}-\theta(2c-\alB+\alpha_{B'})\bigr]\le0$, so the numerator of $\Qgain{B\otimes B'}{1/2}(x,t)$ is at most $t/2-1<0$.
\end{proof}

With $B=B_k$, $c=c_k$ and $B'$ a top algorithm of $\mathcal C$, the threshold $\alpha_{B'}/(2c-\alB+\alpha_{B'})$ tends to $\theta^*$ along the $k$ with $c_k\to\cbar_{\mathcal B}$ as $\alpha_{B'}\to\alpha_{\mathcal C}$; so $\theta^*$ is, in the limit, the largest share of a single-class $B_k$ for which \ThmD{} at $\kappa=\frac12$ applies to such products.

\phantomsection\label{cp:C16}%
With the level-1 family of $\Pzero$ as $\mathcal C$ ($\cg=0$ by Lemma~\ref{cp:L8}, $\alpha(\Pzero)>\cpv{alphaP0_lo}$ by Lemma~\ref{cp:L11}), with exact algorithms $B_k$ of rank $k^{2+o(1)}$ for $\mat k{\lceil k^{\alpha_{\mathcal B}}\rceil}k$ for every $\alpha_{\mathcal B}<\AlphaVXXZ$, which exist because $\omega(1,\alpha_{\mathcal B},1)=2$ by the bound $\alpha\ge\AlphaVXXZ$ of~\cite{VXXZ24} on the dual exponent $\alpha$ of matrix multiplication, and with $\cbar_{\mathcal B}\le2$, Proposition~\ref{cp:C14} gives $\Hyp(\eps,\gamma)$ with some $\gamma>0$ for every $\eps<\cpv{C16_range}$ (the range is increasing in $\alpha_{\mathcal C}$ and $\alpha_{\mathcal B}$ and decreasing in $\cbar_{\mathcal B}$, so these bounds may be inserted, and the value is checked in exact rational arithmetic); this uses only the printed bound of~\cite{VXXZ24}, and Section~\ref{cp:sec:range} obtains more, up to $\cpv{A11_vxxz_half}$, from the same printed bound with a fast rectangular product as the bottom (Corollary~\ref{cp:A11}).

\subsection{The multiplication model}\label{cp:sec:tr:mult}

For a tensor $T$ and a set $W$ of its $z$-variables, $T|_W$ denotes $T$ with the $z$-variables outside $W$ set to $0$; for $\mat NDN$ and a set $W$ of positions these are the bilinear forms $(XY)[a,b]$, $(a,b)\in W$.

\begin{proposition}[the multiplication model]\label{cp:C17}
In the multiplication model of Section~\ref{cp:sec:models} the following hold.
\begin{enumerate}[label=(\alph*)]
\item If, for given $N$, $D$ and $W$, a straight-line program of the multiplication model computes the entries of $XY$ at $W$ with $L$ non-scalar multiplications, then $R(\mat NDN|_W)\le2L$.
\item Conversely, a bilinear algorithm with $r$ terms for $\mat NDN|_W$ computes these entries with $r$ non-scalar multiplications.
\item If an exact top algorithm $B$ for $\mat kbk$ satisfies $\tau:=t\alB/2-\lamB(t)>0$ for some $t\in[0,1]$, then $R(\mat kbk|_S)\le k^{2-\tau}$ for every set $S$ of outputs with $\abs S\le k^{2-\alB/2}$.
\end{enumerate}
Consequently, in the multiplication model, $\Hyp(\eps,\gamma)$ (for all $N$, all $2\le D\le N^\eps$ and all $W$ with $\abs W\le N^2/\sqrt D$, at most $CN^2D^{-\gamma}$ non-scalar multiplications) holds if and only if $R(\mat NDN|_W)\le C'N^2D^{-\gamma}$ for all such $N$, $D$, $W$ and some constant $C'$; so the supremum of the $\eps$ for which some $\gamma>0$ works is the same for both statements.
\end{proposition}

\begin{proof}
(a) is Strassen's bound (Section~\ref{cp:sec:laser}) for the tensor $\mat NDN|_W$, whose bilinear forms are the wanted entries. (b) Each term of a bilinear algorithm is one product of a linear form in the entries of $X$ with one in the entries of $Y$; linear operations are free in the multiplication model. (c) Lemma~\ref{cp:T2}(a) with $n=1$, $N_{\mathrm{sub}}=1$ and the wanted set $S$ gives $\abs{\bigcup_{o\in S}\Cov(o)}\le\abs S^t\ZB(t)=\abs S^tk^{\lamB(t)+2-2t}\le k^{(2-\alB/2)t+\lamB(t)+2-2t}=k^{2-\tau}$, and the terms in $\bigcup_{o\in S}\Cov(o)$, with their $z$-forms restricted to $S$, form an exact decomposition of $\mat kbk|_S$. The consequence follows from (a) with $C'=2C$ and from (b) with $C=C'$.
\end{proof}

The equivalence holds in the multiplication model only. In the time model nothing about rank follows from $\Hyp$, since a word-RAM algorithm need not be algebraic; the direction from $\Hyp$ to the rank bound holds for algorithms of the algebraic form of Section~\ref{cp:sec:models}, whose running time bounds their number of non-scalar multiplications, such as \AlgD{} (Remark~\ref{cp:T6:alg}). A rank bound does not by itself give $\Hyp$ in the time model: a bilinear algorithm with about $N^2$ terms for $\mat NDN|_W$ may have linear forms with up to $N^3D\le N^{3+\eps}$ coefficients in all, and forming them is the cost that the composite methods of this paper avoid by structure.

\begin{remark}[discussion; argued, not proved]\label{cp:C17:disc}
By Remark~\ref{cp:T6:alg} and Theorem~\ref{cp:A1}, the range of Proposition~\ref{cp:C17} in the multiplication model is at least $\alpha(P)$ for every admissible level-1 point $P$, and by Proposition~\ref{cp:C17}(c) every top algorithm with $\lamB(t)<t\alB/2$ for some $t$ gives a rank saving on sparse output sets. We do not know whether the range in the multiplication model equals $\alpha$. A proof that it is smaller would need, for some $\eps<\alpha$ and every $\gamma>0$, sets $W$ with $R(\mat N{N^\eps}N|_W)>N^2D^{-\gamma}$: a lower bound on the rank of a tensor whose largest dimension is $N^2/\sqrt D$ that exceeds this dimension by the factor $D^{1/2-\gamma}$. No lower bound on the rank or border rank of an explicit family of tensors that is superlinear in the largest dimension is known, and counting does not give one here, since there are at most $2^{N^2}$ sets $W$ while, over a finite field, the tensors of rank $N^2D^{-\gamma}$ in these dimensions are far more numerous.
\end{remark}

\subsection{The converse: plans and the typical support rate}\label{cp:sec:tr:C}

\begin{definition}[plans]\label{cp:C6:def}
Let $\mathcal D$ satisfy $R'=R$, and let $\Sr=\{(l,z):z\in\supp w_l\}$ be its support relation. For distributions $c$ on the terms and $p$ on $Z$, a \emph{plan} with marginals $(c,p)$ is a distribution $J$ on $\Sr$ with term marginal $c$ and $z$-marginal $p$; $\Pl(c,p)$ is the set of these plans. Put
\[
\Hzl(c,p)=\max_{J\in\Pl(c,p)}\Ent_J(Z\mid L),\qquad\Hlz(c,p)=\max_{J\in\Pl(c,p)}\Ent_J(L\mid Z),
\]
with the value $-\infty$ if $\Pl(c,p)$ is empty, where $(L,Z)$ is distributed according to $J$ and $\Ent$ is the entropy with the natural logarithm. With $\unif$ the uniform distribution, the \emph{typical support rate} of $\mathcal D$ is $\ktyp=\Hzl(\unif,\unif)$. For $J\in\Pl(\unif,\unif)$, $\Ent_J(Z\mid L)=\Ent(J)-\ln R=\Ent_J(L\mid Z)$, so also $\ktyp=\Hlz(\unif,\unif)$. Under \Tcond1{} the diagonal plan $J(l,\pi(l))=1/R$ lies in $\Pl(\unif,\unif)$, so $0\le\ktyp\le\ln R$.
\end{definition}

\begin{theorem}[Theorem C, the converse]\label{cp:C6}
Let a variable-map family for $(T,\mathcal D)$ satisfy \Tcond1, \Tcond3{} and \Tcond4, with $\Acs=\liminf_M\ln a_M/M$. Fix at every level $M$ a partition of a set of terms that contains all terms covering some kept output structurally. Then for every plan $J\in\Pl(\unif,\unif)$ and every $t\in(0,1]$,
\[
\Lam^{\mathrm{str}}(t)\ge t\,\Ent_J(Z\mid L)/\Acs;
\]
hence $\Lam^{\mathrm{str}}(t)\ge t\ktyp/\Acs$ and $\cg^{\mathrm{str}}\ge\ktyp/\Acs$. If $p_M=1$ for all $M$, the same bounds hold for the profile of the family $(A_M)$ of Lemma~\ref{cp:C1} with every partition.
\end{theorem}

\begin{proof}
Fix $J$, $t$ and $\delta\in(0,1]$, and write $H=\Ent_J(Z\mid L)$. For $(l,z)$ with $J(l,z)>0$ put $\phi(l,z)=-\ln(RJ(l,z))$; since the term marginal of $J$ is uniform, $H=\sum_{l,z}J(l,z)\phi(l,z)$. Let $\Delta_J=1+\max\abs{H-\phi(l,z)}$ over these $(l,z)$.

\emph{Step 1 (a tail bound).} Let $J^M(\tau,z)=\prod_iJ(\tau_i,z_i)$ on $[R]^M\times Z^M$. Both marginals of $J^M$ are uniform: $\sum_\tau J^M(\tau,z)=R^{-M}$ for every $z$ and $\sum_zJ^M(\tau,z)=R^{-M}$ for every $\tau$. Where $J^M(\tau,z)>0$ put $X(\tau,z)=\sum_i\phi(\tau_i,z_i)$, so that $J^M(\tau,z)=R^{-M}e^{-X(\tau,z)}$. For $0\le\vartheta\le1/\Delta_J$, $\abs{\vartheta(H-\phi)}\le1$, and $\abs{e^y-1-y}\le y^2$ for $\abs y\le1$ gives $\sum J\,e^{\vartheta(H-\phi)}\le1+\vartheta^2\Delta_J^2\le e^{\vartheta^2\Delta_J^2}$, because $\sum J(H-\phi)=0$. With $\vartheta=\delta/(2\Delta_J^2)$ and Markov's inequality,
\[
J^M\bigl(X\le M(H-\delta)\bigr)\le e^{-\vartheta\delta M}\Bigl(\sum J\,e^{\vartheta(H-\phi)}\Bigr)^M\le e^{-\delta^2M/(4\Delta_J^2)}.
\]
\emph{Step 2 (heavy terms).} Let $K'=\bigcup_o\Pre(o)$ be the set of preimages of kept outputs; $\abs{K'}\ge s_Ma_M^2$ by Lemma~\ref{cp:C1}(c). Put $\theta_M=e^{-\zeta_MM}/4=s_Ma_M^2/(4Q_MR^M)$. For a term $x=(j,\tau)$ let $g(x)$ be the sum of $J^M(\tau,z)$ over the $z$ with $(j,z)\in K'$, $J^M(\tau,z)>0$ and $X(\tau,z)>M(H-\delta)$. Then $g(x)\le R^{-M}$, and by the $z$-marginal and Step 1, applied once per copy,
\[
\sum_xg(x)\ge\abs{K'}R^{-M}-Q_Me^{-\delta^2M/(4\Delta_J^2)}.
\]
Since $\zeta_M\to0$, for all large $M$ the last exponential is at most $\theta_M$, and then $\sum_xg(x)\ge\frac34s_Ma_M^2R^{-M}$. Let $H_M$ be the set of terms with $g(x)\ge\theta_MR^{-M}$. There are $Q_MR^M$ terms, so $\sum_xg(x)\le\abs{H_M}R^{-M}+Q_MR^M\theta_MR^{-M}=\abs{H_M}R^{-M}+\frac14s_Ma_M^2R^{-M}$; hence $\abs{H_M}\ge s_Ma_M^2/2$.

\emph{Step 3 (heavy terms cover many outputs).} Let $x=(j,\tau)\in H_M$. Every $z$ counted in $g(x)$ has $J(\tau_i,z_i)>0$ for all $i$, so $z_i\in\supp w_{\tau_i}$ and $z\in\supp\tau$; and $J^M(\tau,z)<R^{-M}e^{-M(H-\delta)}$. So more than $g(x)R^Me^{M(H-\delta)}\ge\theta_Me^{M(H-\delta)}$ preimages $(j,z)\in K'$ have $z\in\supp\tau$. Each is a preimage of one kept output, which $x$ then covers structurally, and each kept output has at most $p_M$ preimages; so $x$ covers structurally at least $S_M:=\theta_Me^{M(H-\delta)}/p_M$ kept outputs.

\emph{Step 4 (the profile).} Let $(P_c)$ be the partition at level $M$; its terms include $H_M$. The incidences between $P_c\cap H_M$ and the $s_Ma_M^2$ kept outputs number at least $\abs{P_c\cap H_M}S_M$, so the structural class coverage satisfies $\kappa_c\ge\abs{P_c\cap H_M}S_M/(s_Ma_M^2)$. With $\abs{P_c}\ge\abs{P_c\cap H_M}$ and $t\in(0,1]$,
\[
Z^{\mathrm{str}}_M(t)\ge\sum_c\abs{P_c\cap H_M}\Bigl(\frac{S_M}{s_Ma_M^2}\Bigr)^t=\abs{H_M}\Bigl(\frac{S_M}{s_Ma_M^2}\Bigr)^t\ge\frac{s_Ma_M^2}2\Bigl(\frac{S_M}{s_Ma_M^2}\Bigr)^t,
\]
and with $\mu_M(t)=[\ln Z(t)-(1-t)\ln(s_Ma_M^2)]/\ln a_M$,
\[
\mu^{\mathrm{str}}_M(t)\ln a_M\ge t\ln S_M-\ln2=t\bigl[M(H-\delta-\zeta_M)-\ln4-\ln p_M\bigr]-\ln2 .
\]
\emph{Step 5 (the limit).} Let $\eta\in(0,\Acs)$. For all large $M$, $\ln a_M\ge(\Acs-\eta)M$, and for infinitely many $M$, $\ln a_M\le(\Acs+\eta)M$. By \Tcond3{} and \Tcond4, $(\ln4+\ln p_M)/M\to0$ and $\zeta_M\to0$, so the terms $t(\ln4+\ln p_M)+\ln2$, divided by $\ln a_M$, tend to $0$. If $H>0$, take $\delta<H$: along the infinitely many $M$ of the second kind, $H-\delta-\zeta_M>0$ and $\mu^{\mathrm{str}}_M(t)\ge t(H-\delta-\zeta_M)/(\Acs+\eta)-o(1)$, so $\Lam^{\mathrm{str}}(t)\ge t(H-\delta)/(\Acs+\eta)$. If $H=0$, for all large $M$, $\mu^{\mathrm{str}}_M(t)\ge-t(\delta+\zeta_M)/(\Acs-\eta)-o(1)$, so $\Lam^{\mathrm{str}}(t)\ge-t\delta/(\Acs-\eta)$. Letting $\delta\to0$ and $\eta\to0$ gives $\Lam^{\mathrm{str}}(t)\ge tH/\Acs$. The maximum over $J$ is $\ktyp$, and $\cg^{\mathrm{str}}=\liminf_{t\to0+}\Lam^{\mathrm{str}}(t)/t\ge\ktyp/\Acs$.

\emph{Actual supports.} If $p_M=1$, the coefficient of a kept output $o$ in the pushed $z$-form of $(j,\tau)$ is $w_\tau[z]$ if the single preimage of $o$ is $(j,z)$, and $0$ otherwise (proof of Lemma~\ref{cp:C1}(b)). So the terms of $A_M$ are exactly the terms that cover some kept output structurally, and their coverage is their structural coverage; the profile of a partition of $A_M$ is its structural profile, and the theorem applies to it.
\end{proof}

The structural profile bounds the profile of Definition~\ref{cp:T1} from above (Lemma~\ref{cp:C1}), so Theorem~\ref{cp:C6} limits what the structural profile can certify; it does not bound the actual profile from below when $p_M>1$, since the contributions of several preimages of an output may cancel. It concerns profiles only and does not say that $\Hyp$ fails above any range.

\begin{corollary}[Corollary C1: $\cg=0$ forces triangularity]\label{cp:C7}
Let $\mathcal D$ have $R'=R$ terms and a perfect matching $\pi$. For every simple directed cycle $l_1\to\dots\to l_k\to l_1$ of $\Dpi$ with $k\ge2$, $\ktyp\ge k\ln2/R$. So if $\mathcal D$ is not triangular, every variable-map family for $(T,\mathcal D)$ with \Tcond1, \Tcond3, \Tcond4{} has $\cg^{\mathrm{str}}\ge k\ln2/(R\Acs)>0$ for every partition, and the same holds for the profile of $(A_M)$ when $p_M=1$. Together with Theorem~\ref{cp:C2}: for a family with \Tcond1, \Tcond3, \Tcond4{} and $p_M=1$, if some partition has $\cg=0$ then $\mathcal D$ is triangular; and if $\mathcal D$ is triangular and \Tcond5{} holds, the excess classes have $\cg=0$.
\end{corollary}

\begin{proof}
The plan $J$ with mass $1/(2R)$ on each of $(l_j,\pi(l_j))$ and $(l_j,\pi(l_{j+1}))$ (indices mod $k$) and mass $1/R$ on $(l,\pi(l))$ for every $l$ off the cycle lies on $\Sr$ (edges and $\pi(l)\in\supp w_l$), and its marginals are uniform: the row of each $l_j$ has mass $1/R$, and the column of $\pi(l_j)$ receives $1/(2R)$ from $l_j$ and from $l_{j-1}$. Each row on the cycle has conditional entropy $\ln2$, the others $0$, so $\Ent_J(Z\mid L)=k\ln2/R$. If $\mathcal D$ is not triangular, $\Dpi$ has such a cycle by Proposition~\ref{cp:C4}; Theorem~\ref{cp:C6} gives the bounds.
\end{proof}

A \emph{height} on $\Dpi$ is a function $\hgt$ on the terms that is constant on every strongly connected component and satisfies $\hgt(l')<\hgt(l)$ for every edge $l\to l'$ between different components; one exists, since the components form an acyclic graph (take the length of a longest path in that graph). Put $\wt(\pi(l))=\hgt(l)$.

\begin{lemma}[plans respect heights]\label{cp:C8:H}
Let $J\in\Pl(c,p)$ with $\sum_zp(z)\wt(z)\ge\sum_lc(l)\hgt(l)$. Then equality holds, and $J$ vanishes on every $(l,z)\in\Sr$ with $z=\pi(l')$ for $l'$ in another component than $l$. In particular this holds for every $J\in\Pl(p\circ\pi,p)$, where $(p\circ\pi)(l)=p(\pi(l))$.
\end{lemma}

\begin{proof}
$\sum_zp(z)\wt(z)-\sum_lc(l)\hgt(l)=\sum_{(l,z)\in\Sr}J(l,z)(\wt(z)-\hgt(l))$. Every $(l,z)\in\Sr$ has $z=\pi(l')$ with $l'=l$ or $l\to l'$, so $\wt(z)=\hgt(l')\le\hgt(l)$, with strict inequality when $l'$ lies in another component. So the sum is at most $0$, and it is $0$ only if $J$ vanishes on the pairs between components. For $J\in\Pl(p\circ\pi,p)$ the two sides are equal.
\end{proof}

\begin{proposition}[Proposition K: components]\label{cp:C8}
Let $\mathcal D$ have $R'=R$ terms and a perfect matching $\pi$. For a strongly connected component $\mathsf C$ of $\Dpi$ let $d_{\mathsf C}(l)=\abs{\{l'\in\mathsf C:\pi(l')\in\supp w_l\}}$ (this counts $l'=l$), let $\nu(\mathsf C)$ be the largest number of vertices of $\mathsf C$ covered by vertex-disjoint simple cycles of length at least $2$ inside $\mathsf C$, call a doubly stochastic $\mathsf C\times\mathsf C$ matrix $\beta$ \emph{admissible} if $\beta_{ll'}=0$ unless $\pi(l')\in\supp w_l$, and let
\[
\kappa_{\mathsf C}=\frac1R\max\Bigl\{\sum_{l\in\mathsf C}\Ent(\beta_{l\cdot}):\beta\text{ admissible}\Bigr\},
\]
where $\beta_{l\cdot}$ is the row of $l$.
Then $\ktyp=\sum_{\mathsf C}\kappa_{\mathsf C}$ over the components with at least two vertices, and for each of them
\[
\nu(\mathsf C)\ln2/R\le\kappa_{\mathsf C}\le\frac1R\sum_{l\in\mathsf C}\ln d_{\mathsf C}(l)\le\abs{\mathsf C}\ln\abs{\mathsf C}/R.
\]
In particular $\ktyp=0$ if and only if $\Dpi$ is acyclic, that is, if and only if $\mathcal D$ is triangular.
\end{proposition}

\begin{proof}
By Lemma~\ref{cp:C8:H} with $c=p=\unif$, every $J\in\Pl(\unif,\unif)$ is supported on the pairs $(l,\pi(l'))$ with $l$, $l'$ in one component, so $J$ is a direct sum of blocks $J_{\mathsf C}$, and $RJ_{\mathsf C}$, indexed by $(l,l')\in\mathsf C^2$, is doubly stochastic and vanishes unless $\pi(l')\in\supp w_l$. Conversely every choice of such blocks is a plan. The conditional entropy $\Ent_J(Z\mid L)=\frac1R\sum_l\Ent(RJ(l,\cdot))$ adds over the blocks. A component with one vertex $l$ contributes only $(l,\pi(l))$, a row of entropy $0$. The upper bounds hold since a row with at most $d_{\mathsf C}(l)$ nonzero entries has entropy at most $\ln d_{\mathsf C}(l)\le\ln\abs{\mathsf C}$. For the lower bound apply the construction of the proof of Corollary~\ref{cp:C7} to a family of vertex-disjoint cycles covering $\nu(\mathsf C)$ vertices. A component with at least two vertices contains a cycle of length at least $2$, so it contributes at least $2\ln2/R>0$; with Proposition~\ref{cp:C4} this gives the last sentence.
\end{proof}

\begin{lemma}[additivity]\label{cp:C8:add}
Let $\mathcal D_1$ and $\mathcal D_2$ be decompositions with as many terms as $z$-variables ($R_1$ and $R_2$) and perfect matchings $\pi_1$, $\pi_2$. Then $\mathcal D_1\otimes\mathcal D_2$ has $R_1R_2$ terms and $z$-variables, the perfect matching $\pi_1\times\pi_2$, and $\ktyp(\mathcal D_1\otimes\mathcal D_2)=\ktyp(\mathcal D_1)+\ktyp(\mathcal D_2)$.
\end{lemma}

\begin{proof}
The support of $w_{l_1}\otimes w_{l_2}$ is $\supp w_{l_1}\times\supp w_{l_2}$, so the support relation is the product of the two. For a plan $J$ of the product with uniform marginals, the marginal of $(L_1,Z_1)$ is a plan of $\mathcal D_1$ with uniform marginals, likewise for $(L_2,Z_2)$, and $\Ent(Z_1Z_2\mid L_1L_2)\le\Ent(Z_1\mid L_1)+\Ent(Z_2\mid L_2)$, since conditioning does not increase entropy. Product plans attain the sum.
\end{proof}

\subsection{Composition classes}\label{cp:sec:tr:G}

The \emph{composition} $\cmpn(\tau)$ of a term string is the distribution of its letters, and the \emph{type} $\typ(z)$ of a $z$-string is the distribution of its letters; an $M$-type is a distribution with values in $\frac1M\Z$. For an $M$-type $c$ let $T_c$ be the number of term strings of composition $c$; then $T_c\le e^{M\Ent(c)}$.

\begin{lemma}[counts]\label{cp:C9:count}
For an $M$-type $c$ on the terms and a $z$-string $z$ of type $p$, the number $N(c,z)$ of term strings $\tau$ of composition $c$ with $z\in\supp\tau$ is at most $(M+1)^{R^2}e^{M\Hlz(c,p)}$.
\end{lemma}

\begin{proof}
Group these $\tau$ by the joint type of the pairs $(\tau_i,z_i)$. It is an $M$-type in $\Pl(c,p)$, since every pair lies in $\Sr$; there are at most $(M+1)^{R^2}$ of them. For a fixed joint type $J$ and the fixed $z$, the number of $\tau$ is $\prod_{z'}(Mp(z'))!/\prod_{l,z'}(MJ(l,z'))!\le e^{M\Ent_J(L\mid Z)}$, by the bound $\binom{n}{n_1,\dots,n_r}\le e^{n\Ent(n_1/n,\dots,n_r/n)}$ applied for each $z'$.
\end{proof}

\begin{lemma}[semicontinuity]\label{cp:C9:usc}
$\Hlz$ and $\Hzl$ are upper semicontinuous on the product of the simplices of distributions on the terms and on $Z$, with values in $\{-\infty\}\cup[0,\ln R]$.
\end{lemma}

\begin{proof}
The set of triples $(c,p,J)$ with $J\in\Pl(c,p)$ is closed, and $J$ ranges over a compact simplex; $\Ent_J(L\mid Z)=\Ent(J)-\Ent(p)$ is continuous on it. If $(c_n,p_n)\to(c,p)$ and $\Hlz(c_n,p_n)\ge h_0$, maximizers $J_n$ have a convergent subsequence whose limit lies in $\Pl(c,p)$ and has $\Ent_J(L\mid Z)\ge h_0$. If $\Pl(c,p)$ is empty, it is empty near $(c,p)$, by closedness. The same argument applies to $\Hzl$.
\end{proof}

\begin{theorem}[Theorem G: composition classes]\label{cp:C9}
Let a variable-map family for $(T,\mathcal D)$ satisfy \Tcond1, \Tcond3{} and \Tcond4, and partition the terms by composition: $P_c=\{(j,\tau):\cmpn(\tau)=c\}$. Let $\mathcal P$ be the set of limit points of the types of kept preimages, that is, of the $p$ for which there are levels $M_i\to\infty$ and preimages $(j_i,z_i)$ of kept outputs at level $M_i$ with $\typ(z_i)\to p$, and put $\PhiG=\max_{p\in\mathcal P}\Hlz(\unif,p)$.
\begin{enumerate}[label=(\alph*)]
\item If $\alpha>0$, then $(A_M)$ with the induced partition, after finitely many $M$ are discarded, is an input family with \Hcond1--\Hcond4.
\item $\cg^{\mathrm{str}}\le\PhiG/\Acs$; hence also $\cg\le\PhiG/\Acs$ for $(A_M)$.
\item $\unif\in\mathcal P$, so $\PhiG\ge\ktyp$, and $\ktyp/\Acs\le\cg^{\mathrm{str}}\le\PhiG/\Acs$.
\item Assume \Tcond{5$^*$}: $\PhiG=\ktyp$. Then $\cg^{\mathrm{str}}=\ktyp/\Acs$, and $\cg=\ktyp/\Acs$ for $(A_M)$ when $p_M=1$.
\item If \Tcond{5$^*$} holds, $0<\alpha<1$ and $2\ktyp<\Acs\alpha$, then $\Hyp(\eps,\gamma)$ holds with some $\gamma>0$ for every $0<\eps\le\min(\alpha,\frac12)$ (time model).
\end{enumerate}
\end{theorem}

\begin{proof}
(a) As in step (vii) of the proof of Theorem~\ref{cp:C2}, with at most $(M+1)^R$ classes in \Hcond3.

(b) Fix $M$. Then $\abs{P_c}=Q_MT_c\le Q_Me^{M\Ent(c)}$ by the bound on $T_c$ before Lemma~\ref{cp:C9:count}, and by that lemma the structural coverage of a kept output $o$ by $P_c$ is at most $\sum_{(j,z)\in\Pre(o)}N(c,z)\le p_M(M+1)^{R^2}e^{MF_M(c)}$, where $F_M(c)$ is the largest $\Hlz(c,p)$ over the types $p$ of kept preimages at level $M$; classes with $\kappa_c=0$ contribute nothing for $t>0$. With at most $(M+1)^R$ classes and $\ln(s_Ma_M^2)=M\ln R+\ln Q_M-\zeta_MM$,
\begin{align*}
\mu^{\mathrm{str}}_M(t)\ln a_M&\le R\ln(M+1)+tR^2\ln(M+1)+t\ln p_M+(1-t)\zeta_MM\\
&\quad+M\max_c\bigl[(1-t)(\Ent(c)-\ln R)+tF_M(c)\bigr].
\end{align*}
Fix $e>0$. By Lemma~\ref{cp:C9:usc} and the compactness of $\{\unif\}\times\mathcal P$ there is $\delta>0$ with $\Hlz(c,p)\le\PhiG+e$ whenever $\norm{c-\unif}_1<\delta$ and $p$ has distance less than $\delta$ from $\mathcal P$ (otherwise a sequence violating this would have a limit point $(\unif,p^*)$, $p^*\in\mathcal P$, with $\Hlz(\unif,p^*)\ge\PhiG+e$). By compactness of the simplex, for all large $M$ all types of kept preimages are within $\delta$ of $\mathcal P$. For $\norm{c-\unif}_1\ge\delta$, Pinsker's inequality gives $\Ent(c)-\ln R\le-\delta^2/2$, and with $\Hlz\le\ln R$ the bracket is at most $-(1-t)\delta^2/2+t\ln R<0$ for $t<t_0(\delta)$. For $\norm{c-\unif}_1<\delta$ the bracket is at most $t(\PhiG+e)$, since $\Ent(c)\le\ln R$. Since $\PhiG\ge0$ (see (c)), with $\ln a_M\ge(\Acs-\eta)M$ for all large $M$ this gives $\Lam^{\mathrm{str}}(t)\le t(\PhiG+e)/(\Acs-\eta)$ for $0<t<t_0$, so $\cg^{\mathrm{str}}\le(\PhiG+e)/(\Acs-\eta)$ for every $e,\eta>0$. The bound for $(A_M)$ is Lemma~\ref{cp:C1}.

(c) Some copy $j$ and some type $p$ carry at least $\abs{K'}/(Q_M(M+1)^R)\ge R^Me^{-\zeta_MM}(M+1)^{-R}$ of the preimages of kept outputs, $K'=\bigcup_o\Pre(o)$, by Lemma~\ref{cp:C1}(c). There are at most $e^{M\Ent(p)}\le R^M$ $z$-strings of type $p$, so $\Ent(p)\ge\ln R-\zeta_M-R\ln(M+1)/M$, and Pinsker's inequality gives $\norm{p-\unif}_1\to0$. So $\unif\in\mathcal P$ and $\PhiG\ge\Hlz(\unif,\unif)=\ktyp\ge0$. The lower bound $\cg^{\mathrm{str}}\ge\ktyp/\Acs$ is Theorem~\ref{cp:C6}. (d) follows from (b), (c) and, for $p_M=1$, the last part of Theorem~\ref{cp:C6}. (e) By (b) and \Tcond{5$^*$}, $\cg\le\ktyp/\Acs<\alpha/2$ for $(A_M)$, which is an input family by (a); apply Corollary~\ref{cp:C2:half}.
\end{proof}

The bound (e) does not need $p_M=1$: it uses only the upper bound (b). The equality case $2\ktyp=\Acs\alpha$ is not covered.

\begin{proposition}[sufficient conditions for \Tcond{5$^*$}]\label{cp:C10}
In the setting of Theorem~\ref{cp:C9}, \Tcond{5$^*$} holds in each of the following cases.
\begin{enumerate}[label=(\alph*)]
\item $\mathcal P=\{\unif\}$.
\item There are a finite group $\Gamma$ of pairs $(\sigma,\rho)$ of permutations of the terms and of $Z$ with $(\sigma(l),\rho(z))\in\Sr$ for all $(l,z)\in\Sr$, and a convex set $\mathcal Q$ of distributions on $Z$ with $\mathcal P\subseteq\mathcal Q$ and $\rho_*\mathcal Q=\mathcal Q$ for $(\sigma,\rho)\in\Gamma$, whose only $\Gamma$-fixed point is $\unif$.
\item $\mathcal D$ is triangular (\Tcond2) and \Tcond5{} holds.
\end{enumerate}
\end{proposition}

\begin{proof}
By Theorem~\ref{cp:C9}(c) it suffices to show $\Hlz(\unif,p)\le\ktyp$ for $p\in\mathcal P$. (a) $\Hlz(\unif,\unif)=\ktyp$. (b) $\Ent_J(L\mid Z)$ is concave in $J$, and the set of pairs $(J,p)$ with $J\in\Pl(\unif,p)$ is convex, so $p\mapsto\Hlz(\unif,p)$ is concave. A pair $(\sigma,\rho)\in\Gamma$ maps $\Pl(\unif,p)$ onto $\Pl(\unif,\rho_*p)$ and preserves the conditional entropy, so $\Hlz(\unif,\cdot)$ is $\Gamma$-invariant. For $p\in\mathcal Q$ the average $\bar p$ of its $\Gamma$-orbit lies in $\mathcal Q$ and is $\Gamma$-fixed, so $\bar p=\unif$, and concavity gives $\Hlz(\unif,\unif)\ge\Hlz(\unif,p)$. (c) By step (iii) of the proof of Theorem~\ref{cp:C2}, which uses only \Tcond4, \Tcond5{} and Lemma~\ref{cp:C1}(c), every kept preimage at level $M$ has $\wt(z)\ge\varpi_M\ge M(\wbar-\eta'_M)$ with $\eta'_M\to0$, so every $p\in\mathcal P$ has $\sum_zp(z)\wt(z)\ge\wbar=\sum_l\unif(l)\wt(\pi(l))$. A triangular decomposition has only one-vertex components, and $\hgt=\wt\circ\pi$ is a height (it decreases strictly along every edge); for the merged decomposition of $F_q$ this is the weight $(0,1,2)$ of Example~\ref{cp:C4:ex}, not a $0/1$ height. By Lemma~\ref{cp:C8:H}, every $J\in\Pl(\unif,p)$ vanishes off the pairs $(l,\pi(l))$, so $L$ is a function of $Z$ under $J$ and $\Ent_J(L\mid Z)=0=\ktyp$ (Proposition~\ref{cp:C8}).
\end{proof}

\subsection{The standard decomposition of \texorpdfstring{$F_q$}{F\_q}}\label{cp:sec:std:Fq}

\begin{proposition}[the typical support rate of $\Fstd$]\label{cp:C11}
For $q\ge2$ let $x_q=\bigl((q+1)-\sqrt{q^2-2q+5}\bigr)/(2(q-1))$, the root in $(0,1/q)$ of $(q-1)x^2-(q+1)x+1=0$. The standard decomposition $\Fstd$ of $F_q$ has
\[
\ktyp(\Fstd)=\frac{-2\ln x_q-(q-1)\ln(1-x_q)}{q+2},
\]
and $\ktyp(\Fstd)=\cpv{C11_q5}$ for $q=5$.
\end{proposition}

\begin{proof}
Take $\pi(A_i)=z_i$, $\pi(B)=z_0$, $\pi(C)=z_{q+1}$; then $\Dpi$ has the edges $A_i\to B$, $B\to A_i$ and $C\to B$, so its components are $\{B,A_1,\dots,A_q\}$ and $\{C\}$ (Example~\ref{cp:C4:ex}). In a plan with uniform marginals, normalised so that each row is a distribution, the row of $C$ is concentrated on $z_{q+1}$, since only $C$ covers $z_{q+1}$; the row of $B$ is a distribution $x=(x_0,x_1,\dots,x_q)$ on $(z_0,\dots,z_q)$; the row of $A_j$ is $(1-y_j,y_j)$ on $(z_j,z_0)$. The column of $z_j$ forces $y_j=x_j$, and then the column of $z_0$ has mass $x_0+\sum_jx_j=1$. So
\[
\ktyp(\Fstd)=\frac1{q+2}\max_x\Bigl[\Ent(x)+\sum_{j=1}^q\Hb(x_j)\Bigr],
\]
with $\Hb$ the binary entropy. The objective is concave and invariant under permutations of $x_1,\dots,x_q$, so a maximiser has $x_j=x$, $x_0=1-qx$. On $[0,1/q]$ the function $f(x)=-(1-qx)\ln(1-qx)-2qx\ln x-q(1-x)\ln(1-x)$ is strictly concave with $f'(x)=q\ln\bigl((1-qx)(1-x)/x^2\bigr)$, which tends to $+\infty$ at $0$ and to $-\infty$ at $1/q$; its zero satisfies $(1-qx)(1-x)=x^2$, that is $(q-1)x^2-(q+1)x+1=0$, and $1-qx=x^2/(1-x)$. Substituting $\ln(1-qx)=2\ln x-\ln(1-x)$ gives $f(x_q)=-2\ln x_q-(q-1)\ln(1-x_q)$. For $q=5$, $x_5=(3-\sqrt5)/4=1/(2\varphi^2)$ and $1-x_5=\varphi/2$ with $\varphi$ the golden ratio, so $f(x_5)=2\ln2+4\ln\varphi+4\ln2-4\ln\varphi=6\ln2$.
\end{proof}

\paragraph{The standard level-1 family of one $F_q$.}\label{cp:C12:fam}
Fix $q\ge2$ and a rational $u\in(0,w]$, $w=1/(q+2)$, and put $p=q/(2(q+2))$ and $v=(w-u)/2$, the joint type $(p_{011},p_{101},p_{110},p_{002},p_{020},p_{200})=(p,p,u,w,v,v)$ of Definition~\ref{cp:L3} with one value of $q$. Assume $\Ent_X>\Ent_Z$ for the entropies of its $X$-marginal $(p+w+v,\,p+u,\,v)$ and of its $Z$-marginal $(1,q,1)/(q+2)$. For $M$ in a set of multiples of a fixed integer, the zero-out of Lemma~\ref{cp:L4} with $\Qs=\{q\}$ (see the paragraph after Lemma~\ref{cp:L4}) maps $F_q^{\otimes M}$ onto $s_M\cp\mat{a_M}{b_M}{a_M}$ with $a_M=q^{pM}$, $b_M=q^{uM}$ and $\ln s_M\ge M(\Ent_Z-\eta)$ for every $\eta>0$ and all large $M$. It keeps the block tensors $T_{IJK}$ of this joint type, which are pairwise distinct in each of $X$, $Y$, $Z$, and each $T_{IJK}$ is the tensor product of $pM$ copies of $\mat11q$, $pM$ copies of $\mat q11$, $uM$ copies of $\mat1q1$ and $(w+2v)M$ copies of $\mat111$, hence $T_{IJK}\cong\mat{q^{pM}}{q^{uM}}{q^{pM}}$ by a bijective relabeling of its variables. The composite map sends every variable to $0$ or to itself and then relabels, so every kept output has exactly one preimage: $Q_M=1$ and $p_M=1$. The \emph{standard level-1 family} at $(q,u)$ is the variable-map family for $(F_q,\Fstd)$ given by these restrictions. It satisfies \Tcond1{} ($R=q+2$ terms and $z$-variables), \Tcond3{} ($p_M=1$) and \Tcond4: since $\Ent_Z=\ln(q+2)-\frac q{q+2}\ln q$ and $a_M^2=q^{qM/(q+2)}$, $\ln\bigl((q+2)^M/(s_Ma_M^2)\bigr)=M\Ent_Z-\ln s_M\le\eta M$; $\ln a_M/M=p\ln q$ for every $M$, so
\[
\Acs=A_q:=\frac{q\ln q}{2(q+2)},
\]
a limit, not only a lower limit; and $\alpha_M=u/p$ for every $M$. The kept outputs are the $z$-strings of the kept blocks $K$, of type $(\#z_0,\#\{z_1,\dots,z_q\},\#z_{q+1})=M(1,q,1)/(q+2)$, and their number per block, $q^{qM/(q+2)}=a_M^2$, is the number of outputs of $\mat{a_M}{b_M}{a_M}$. Since the restriction is a zero-out with $p_M=1$, the supports of the restricted terms and the structural supports agree on the kept outputs. The same restriction carries $\Fmrg$ in place of $\Fstd$; that is the merged level-1 family at $(q,u)$, to which Remark~\ref{cp:C3} applies with $\eta_M=0$.

\begin{proposition}[the typical support exponent of the standard level-1 family]\label{cp:C12}
For the standard level-1 family at $(q,u)$, the partition by composition of Theorem~\ref{cp:C9} (the distribution of all $q+2$ letters of the term string) has $\cg=\ktyp(\Fstd)/A_q$, for the profile of Definition~\ref{cp:T1}, which here equals the structural one; and every partition has $\cg\ge\ktyp(\Fstd)/A_q$. At $q=7$, $\cg=\ktyp(\Fstd)/A_7=\cpv{C12_cg_q7}$ (computed).
\end{proposition}

\begin{proof}
Every kept preimage is a kept output and has the type $M(1,q,1)/(q+2)$ in the coarse letters, so $\mathcal P$ lies in the closed convex set $\mathcal Q=\{p:p(z_0)=p(z_{q+1})=1/(q+2)\}$. The group $\Gamma=S_q$, acting simultaneously on $(A_1,\dots,A_q)$ and $(z_1,\dots,z_q)$ and fixing $B$, $C$, $z_0$, $z_{q+1}$, preserves the support relation of $\Fstd$ and $\mathcal Q$, and the only $\Gamma$-fixed point of $\mathcal Q$ is $\unif$. So \Tcond{5$^*$} holds by Proposition~\ref{cp:C10}(b), and Theorem~\ref{cp:C9}(d) and Theorem~\ref{cp:C6} with $p_M=1$ give both statements, with $\Acs=A_q$.
\end{proof}

\begin{example}[a non-triangular decomposition with $0<\cg<\alpha/2$]\label{cp:C13}
Let $T=F_7^{\otimes6}$ with the decomposition $\mathcal D=\Fstdq{7}\otimes{\Fmrgq{7}}^{\otimes5}$; it has $9^6$ terms and $z$-variables and the perfect matching $\pi_{\mathrm{std}}\times\pi_{\mathrm{mrg}}^{\times5}$ (Lemma~\ref{cp:C8:add}). Take $u=w=\frac19$ at $q=7$: then $\alpha=u/p=\cpv{C13_alpha}$, and $\Ent_X=\ln2>\Ent_Z=\ln9-\frac79\ln7$ because $2^9\cdot7^7>3^{18}$. For $M$ a multiple of $6d$, where $d$ is the integer of the family (so that $M/6$ and $5M/6$ are multiples of~$d$), take the standard level-1 family at $(7,\frac19)$ at level $M/6$, the merged level-1 family at $(7,\frac19)$ at level $5M/6$, and the tensor product of their restrictions; grouping the coordinates into $M/6$ blocks of one standard and five merged coordinates, a renaming of positions, makes it a variable-map family for $(T,\mathcal D)$ at level $M/6$, with $Q=p=1$. Give its terms the pair classes: the composition classes of Theorem~\ref{cp:C9} on the standard factor and the excess classes of Theorem~\ref{cp:C2} on the merged factor. Then the family is an input family with \Hcond1--\Hcond4{} and limit $\alpha=\cpv{C13_alpha}$, $\mathcal D$ is not triangular, and
\[
\cg=\frac{\ktyp(\Fstdq{7})}{6A_7}=\frac{\cg(\text{std})}6=\cpv{C13_cg}<\cpv{C13_half_alpha}=\frac\alpha2,
\]
where $\cg(\text{std})=\cpv{C12_cg_q7}$ is the value of Proposition~\ref{cp:C12} at $q=7$; the values are computed, and the inequality is equivalent to $x_7^2(1-x_7)^6>7^{-3}$ with $x_7=(4-\sqrt{10})/6$, which holds with a wide margin (checked in exact rational arithmetic). By Corollary~\ref{cp:C2:half}, $\Hyp(\eps,\gamma)$ holds with some $\gamma>0$ for every $0<\eps\le\cpv{C13_alpha}$ (time model).
\end{example}

\begin{proof}
\emph{The family.} Both factors satisfy \Tcond1, \Tcond3{} and \Tcond4{} with $\Acs=A_7$ per coordinate of $F_7$ and $\alpha_M=\frac27$ (the restriction does not depend on the decomposition). For the product, $s=s_1s_2$, $a=a_1a_2$ and $b=b_1b_2$, so the quantity of \Tcond4{} is the sum of those of the factors, $\ln a=MA_7=6(M/6)A_7$ and $\alpha_M=\frac27$; so \Tcond4{} holds at level $M'=M/6$ with $\Acs=6A_7$. \Hcond1--\Hcond4{} follow as in step (vii) of the proof of Theorem~\ref{cp:C2}: the number of pair classes is polynomial in $M$, and the product degree is at most $10M$.
\emph{Upper bound.} By Lemma~\ref{cp:C14:prod}, applied to the decompositions (class sizes and class coverages multiply, by the product form of the supports), $Z_M(t)=Z_{1}(t)Z_{2}(t)$ for the two factors, so with $\theta=\ln a_1/\ln a=\frac16$,
\[
\mu_M(t)=\theta\mu_{1,M/6}(t)+(1-\theta)\mu_{2,5M/6}(t),\qquad\Lam(t)\le\tfrac16\Lam_1(t)+\tfrac56\Lam_2(t),
\]
where $\Lam_1$ and $\Lam_2$ are the profiles of the two factor families (an upper limit along a subsequence is at most the upper limit along the sequence). By Theorem~\ref{cp:C2} and Remark~\ref{cp:C3}, $\Lam_2(t)/t\to0$; by Proposition~\ref{cp:C12}, $\liminf_{t\to0+}\Lam_1(t)/t=\cg(\text{std})$. Hence $\cg\le\cg(\text{std})/6$.
\emph{Lower bound.} By Lemma~\ref{cp:C8:add}, $\ktyp(\mathcal D)=\ktyp(\Fstdq{7})+5\ktyp(\Fmrgq{7})=\ktyp(\Fstdq{7})$, since $\Fmrgq{7}$ is triangular (Example~\ref{cp:C4:ex}, Proposition~\ref{cp:C8}). Theorem~\ref{cp:C6}, with $p=1$, gives $\cg\ge\ktyp(\Fstdq{7})/(6A_7)=\cg(\text{std})/6$. As $\ktyp(\mathcal D)>0$, $\mathcal D$ is not triangular (Proposition~\ref{cp:C8}).
\emph{The inequality.} $\ktyp(\Fstdq{7})/(6A_7)<\frac17$ with $6A_7=\frac73\ln7$ and Proposition~\ref{cp:C11} reads $-2\ln x_7-6\ln(1-x_7)<3\ln7$, that is $x_7^2(1-x_7)^6>7^{-3}$; with $3.1622<\sqrt{10}<3.1623$ the left side exceeds $0.0079$, while $7^{-3}<0.0030$.
\end{proof}

The example shows that triangularity, which is what gives $\cg=0$ (Corollary~\ref{cp:C7}), is not needed for \ThmD{} to reach the family's $\alpha$: a single $z$-tight non-triangular base decomposition can have $0<\cg<\alpha/2$, and no assumption on the hashing of a laser construction is used. The same product construction works for $a$ standard and $b$ merged coordinates per block, with $\cg=\frac a{a+b}\cg(\text{std})$.

\subsection{The instance of~\texorpdfstring{\cite{VXXZ24}}{VXXZ24}}\label{cp:sec:tr:vxxz}

Vassilevska Williams, Xu, Xu and Zhou~\cite{VXXZ24} bound $\omega(1,\kappa,1)$ by a degeneration of $2^{o(n)}$ copies of the $n$-th power of the fourth power $CW_5^{\otimes4}$ of the Coppersmith-Winograd tensor, that is of $CW_5^{\otimes4n}$, into copies of a matrix product $\mat m{m^\kappa}m$ (their Algorithm~1, with $q=5$ and $\ell^*=3$), and obtain the bound $\alpha\ge\AlphaVXXZ$ on the dual exponent of matrix multiplication. To avoid a clash with the density exponent $\kappa$ of $\Hypk\kappa$ we write $\beta$ for their $\kappa$, and $\bz=\AlphaVXXZ$. Their $CW_q$ is the tensor $F_q$ of Section~\ref{cp:sec:laser}, their level-1 partition of the $z$-variables is $\{z_0\}\sqcup\{z_1,\dots,z_q\}\sqcup\{z_{q+1}\}$, and a decomposition of $CW_q$ enters their analysis only through the bound $\tilde R(CW_q)\le q+2$ on the right side of their condition~(27)~\cite[Sections~3.6, 3.7 and~8]{VXXZ24}; their restriction acts on the tensor and involves no decomposition. Condition~(27) reads
\[
\lim_{\eps\to0}\lim_{n\to\infty}V^{1/n}\min\bigl(A,B^{1/\beta},C\bigr)^{\omega'/n}\ge(q+2)^{2^{\ell^*-1}},
\]
where $V$ copies of $\mat ABC$ are the output of Algorithm~1 for a given choice of its parameters, and $\eps>0$ is the slack of their construction (not the exponent of $\Hyp$); with $\omega'=2$ it gives $\omega(1,\beta,1)=2$.

\begin{lemma}[Lemma V: the structure of the degeneration of~\cite{VXXZ24}]\label{cp:C18}
For every feasible choice of the parameters of Algorithm~1 of~\cite{VXXZ24} at $q=5$, $\ell^*=3$ (one that satisfies the constraints of~\cite[Theorems~5.3, 6.1 and~6.3]{VXXZ24}), every $\eps>0$, every $n$ and every $\beta\in(0,1)$, the degeneration of Algorithm~1, followed by the zero-out of each output copy of $\mat ABC$ to $\mat m{\lfloor m^\beta\rfloor}m$ with $m=\min(A,B^{1/\beta},C)$, is a restriction of the direct sum of $Q(n)$ copies of $F_5^{\otimes N}$, $N=4n$, with $\ln Q(n)=o(n)$ as $n\to\infty$ at fixed parameters and $\eps$, that sends every variable to $0$ or to one variable with coefficient~$1$. Moreover, every step of its global stage (step~1 of Algorithm~1) sends a variable to a variable whose index string is a permutation of the index string of the source.
\end{lemma}

\begin{proof}
By inspection of the proofs of~\cite[Proposition~5.1, Theorem~5.3, Theorem~6.1, Proposition~6.2, Theorem~6.3, Theorem~7.2 and Corollary~4.2]{VXXZ24} and of Algorithm~1; Appendix~\ref{cp:app:vxxz} lists every step with its place in~\cite{VXXZ24} and the operation it performs. Each step is a zero-out, a renaming of positions, a permutation of chunks of positions, or an identification $x^{(i)}\mapsto x$ of variables of copies whose summands have pairwise disjoint supports (so that no coefficient other than $1$ arises), or a tensor product or direct sum of such maps; compositions of such maps are of the same kind. No step uses the parameter of the definition of degeneration~\cite[Section~3.3]{VXXZ24}. The random choices of the construction (hash weights, permutations) only show that a good choice exists; a member of the family is one fixed restriction. The number of input copies is $2^{o(n)}$ by step~1 of Algorithm~1, and step~2 consumes $2^{o(n)}$ copies per output copy at each of the $\ell^*-1$ constituent stages.
\end{proof}

The second statement is not a statement of~\cite{VXXZ24}, whose theorems say ``can be degenerated into'' with their general definition of degeneration; it is a reading of their proofs, given in full in Appendix~\ref{cp:app:vxxz}.

\begin{lemma}[window, preimages and copy size]\label{cp:C19}
Fix a feasible choice of the parameters of Algorithm~1 at $q=5$, $\ell^*=3$, and $\eps>0$, $n$ and $\beta$. For a $z$-string of length $N$ let its level-1 composition be the fractions of its letters in $\{z_0\}$, $\{z_1,\dots,z_5\}$ and $\{z_6\}$.
\begin{enumerate}[label=(\alph*)]
\item There is a vector $\Gamma$, determined by the parameters, such that every $z$-variable of the input copies whose image under the restriction of Lemma~\ref{cp:C18} is nonzero has level-1 composition within $\cpv{C19_window}\eps$ of $\Gamma$ in the maximum norm; $\cpv{C19_window}=\cpv{C19_window_c}=3^{2^{\ell^*-1}}3^{\ell^*}$.
\item Every variable of the output has at most $\abs{\tilde{\mathcal D}}^{3}r^{9}$ preimages, where $r=\lceil2^{C_1N/\log N}\rceil$ with the constant $C_1$ of~\cite[Corollary~4.2]{VXXZ24} and $\abs{\tilde{\mathcal D}}=2^{o(n)}$ bounds the numbers of complete split distributions of~\cite[Theorems~5.3 and~6.3]{VXXZ24}; this is $e^{o(N)}$ as $n\to\infty$ at fixed $\eps$.
\item For every restriction of the direct sum of $Q$ copies of $F_5^{\otimes N}$ onto $s\cp\mat aba$ (by any linear maps), $s\le Q(N+1)^{18}3^{5N/3}$. Consequently, for every variable-map family for $(F_5,\mathcal D)$ (any decomposition $\mathcal D$ with $R'=7$) with $\zeta_M\to0$ as in \Tcond4, $\liminf_M\ln a_M/M\ge\frac16\ln\frac{343}{243}>\cpv{C19_copy_lo}$.
\end{enumerate}
\end{lemma}

\begin{proof}
(a) The degeneration begins with step~1 of Algorithm~1, so a variable $z'$ with nonzero image has a nonzero image $z$ in some copy of the level-$\ell^*$ $(3^{\ell^*}\eps)$-interface tensor that step~1 produces. By the second statement of Lemma~\ref{cp:C18}, $z'$ and $z$ have the same letters, hence the same level-1 composition. By~\cite[Definitions~3.6 and~4.1]{VXXZ24}, $z$ lies in a level-1 block whose restriction $\hat K_t$ to the $n_t$ chunks of each term $t$ of the interface tensor satisfies $\norm{\mathrm{split}(\hat K_t)-\beta_{Z,t}}_\infty\le3^{\ell^*}\eps$, where $\beta_{Z,t}$ is the prescribed complete split distribution on $\{0,1,2\}^L$, $L=2^{\ell^*-1}$. The fraction of the $Ln_t$ positions of term $t$ with level-1 index $v$ is $\mathrm{comp}(\mathrm{split}(\hat K_t))(v)$, where $\mathrm{comp}(\xi)(v)=\frac1L\sum_\sigma\xi(\sigma)\#\{i:\sigma_i=v\}$, and $\abs{\mathrm{comp}(\xi)(v)-\mathrm{comp}(\xi')(v)}\le\sum_\sigma\abs{\xi(\sigma)-\xi'(\sigma)}\le3^L\norm{\xi-\xi'}_\infty$. Averaging over the terms with the weights $n_t/n$ given by the parameters gives the claim with $\Gamma=\sum_t(n_t/n)\,\mathrm{comp}(\beta_{Z,t})$ and the constant $3^L3^{\ell^*}$. Variables filled in by the hole fixing have their preimages among such variables, by permutations that keep letters, so they are covered.
(b) By Lemma~\ref{cp:C18} every variable has at most one image, so the number of preimages of a composed map is at most the product over its steps. A zero-out or a bijection contributes $1$. In~\cite[Theorems~5.3 and~6.3]{VXXZ24} an output copy is the sum of one copy per complete split distribution $\xi$, each from its own input copy, so a variable gets at most $\abs{\tilde{\mathcal D}}$ preimages there. In the hole fixing one may use exactly $r$ broken copies, and in the proof of~\cite[Theorem~7.2]{VXXZ24} each broken copy is consumed once, relabeled by a bijection and zeroed out, so a variable gets at most one preimage per broken copy. Each of the $\ell^*$ stages applies one hole fixing per region, and the three regions are tensor factors, whose counts multiply. The product is at most $\abs{\tilde{\mathcal D}}^{\ell^*}r^{3\ell^*}$.
(c) We use the slice rank $\mathrm{SR}$ of tensors~\cite{Tao16}: a diagonal tensor with $s$ nonzero entries has slice rank $s$, a restriction does not increase it, and $\mathrm{SR}(\bigoplus_{i\le Q}T)\le Q\,\mathrm{SR}(T)$. Group the monomials of $F_5^{\otimes N}$ into the parts of fixed joint type $P$ of their support triples, a distribution on the $18$ support triples of $F_5$ with denominators $N$; there are at most $(N+1)^{18}$ parts. In the part of type $P$ choose a coordinate $i$ with the smallest entropy $\Ent(P_i)$ of the $i$-th marginal of $P$, and write the part as a sum of slices along $i$, one per string of type $P_i$; that is at most $e^{N\Ent(P_i)}$ slices. Every support triple of $F_5$ has a coordinate equal to the index $0$, so $\sum_iP_i(0)\ge1$ and some marginal has $P_i(0)\ge\frac13$; the largest entropy of a distribution on $7$ letters with mass at least $\frac13$ on one of them is $\Hb(\frac13)+\frac23\ln6=\frac53\ln3$. So $\mathrm{SR}(F_5^{\otimes N})\le(N+1)^{18}3^{5N/3}$. Keeping one monomial $x_{11}y_{11}z_{11}$ in each copy of $\mat aba$ is a further zero-out, onto a diagonal tensor with $s$ entries, so $s\le Q(N+1)^{18}3^{5N/3}$. For a variable-map family, $\ln(s_Ma_M^2)=M\ln7+\ln Q_M-\zeta_MM$, so $2\ln a_M\ge M(\ln7-\frac53\ln3)-18\ln(M+1)-\zeta_MM$, and $\liminf\ln a_M/M\ge\frac12(\ln7-\frac53\ln3)=\frac16\ln\frac{343}{243}$. The last value exceeds $\cpv{C19_copy_lo}$ because $\frac{343}{243}>e^{0.3444}$, checked with a rational upper bound on the exponential series.
\end{proof}

\begin{theorem}[the instance of~\cite{VXXZ24}, conditional]\label{cp:C20}
For $0<\beta<\bz$ let \Vhyp\beta{} be the statement: for some choice of the parameters of Algorithm~1 of~\cite{VXXZ24} at $q=5$ and $\ell^*=3$ that satisfies the constraints of~\cite[Theorems~5.3, 6.1 and~6.3]{VXXZ24} (a feasible point of their program~(28)), condition~(27) of~\cite{VXXZ24} holds with $\omega'=2$ at $\kappa=\beta$. If \Vhyp\beta{} holds for every $\beta\in(0,\bz)$, then for every $0<\eps<\bz$ there is $\gamma>0$ such that $\Hyp(\eps,\gamma)$ holds (time model). More precisely, \Vhyp\beta{} alone gives $\Hyp(\eps,\gamma)$ with some $\gamma>0$ for every $0<\eps\le\beta$.
\end{theorem}

\begin{proof}
Fix $\beta$ and a choice of parameters as in \Vhyp\beta. Take the base $T=F_5$ with the merged decomposition $\mathcal D=\Fmrgq{5}$ (Proposition~\ref{cp:L1} at $q=5$; seven terms and seven $z$-variables). By Example~\ref{cp:C4:ex}, \Tcond1{} and \Tcond2{} hold with $\pi(\Oterm)=z_0$, $\pi(A_i)=z_i$, $\pi(\Bp)=z_6$ and the weight $(0,1,2)$, which is the level-1 index of~\cite{VXXZ24}.
\emph{The family.} In logarithmic form, (27) with $\omega'=2$ says that $L(\eps):=\lim_{n\to\infty}\frac1n\ln(Vm^2)$ satisfies $\lim_{\eps\to0}L(\eps)\ge4\ln7$, with $m=\min(A,B^{1/\beta},C)$. For $k\ge1$ choose $\eps_k\le1/k$ with $L(\eps_k)\ge4\ln7-1/k$, and then $n_k\ge k$, with $n_k>n_{k-1}$, so large that at $(\eps_k,n_k)$: $\frac1{n_k}\ln(Vm^2)\ge4\ln7-2/k$; the number $Q$ of input copies satisfies $\frac1{n_k}\ln Q\le1/k$; and the preimage bound of Lemma~\ref{cp:C19}(b) is at most $e^{N_k/k}$, $N_k=4n_k$. The $k$-th member of the family is the restriction of Lemma~\ref{cp:C18} at $(\eps_k,n_k)$, onto $s=V$ copies of $\mat m{\lfloor m^\beta\rfloor}m$, at level $M=N_k$.
\emph{\Tcond3.} By Lemmas~\ref{cp:C18} and~\ref{cp:C19}(b) it is a variable-map family, and $\ln p_M\le M/k$.
\emph{\Tcond4.} $\ln(7^MQ/(sm^2))\le4n_k\ln7+n_k/k-n_k(4\ln7-2/k)=3n_k/k$, so $\zeta_M\to0$. By Lemma~\ref{cp:C19}(c), $\Acs\ge\frac16\ln\frac{343}{243}>0$; in particular $m\to\infty$, and $\alpha_M=\ln\lfloor m^\beta\rfloor/\ln m\to\beta$. (The zero-out to $\mat m{\lfloor m^\beta\rfloor}m$ is possible since $m\le A$, $m\le C$ and $m^\beta\le B$.)
\emph{\Tcond5.} Every preimage of a kept output is a variable of the input copies with nonzero image, so by Lemma~\ref{cp:C19}(a) its level-1 composition is within $\cpv{C19_window}\eps_k$ of $\Gamma$, and its weight $\wt(z)=M(\mathrm{comp}_1(z)+2\,\mathrm{comp}_2(z))$ lies in an interval of length $6\cdot\cpv{C19_window}\,\eps_kM$; so \Tcond5{} holds with $\eta_M=6\cdot\cpv{C19_window}\,\eps_k\to0$.
Theorem~\ref{cp:C2}(d), with $\alpha=\beta\in(0,1)$, gives $\Hyp(\eps,\gamma)$ with some $\gamma>0$ for every $0<\eps\le\min(\beta,\frac12)=\beta$. If \Vhyp\beta{} holds for every $\beta<\bz$, every $\eps<\bz$ is at most some such $\beta$.
\end{proof}

\begin{remark}[the hypothesis and the published claim]\label{cp:C20:F2}
The bound $\alpha\ge\bz$ of~\cite{VXXZ24} on the dual exponent rests on a numerical solution of their program~(28), whose constraints include condition~(27); Table~1 of~\cite{VXXZ24} lists the upper bound $2$ on $\omega(1,\kappa,1)$ at $\kappa=\bz$. A point returned by a numerical solver satisfies the constraints up to the solver's tolerance, and the paper prints neither the point nor a certificate; we do not verify the claim here. The hypothesis of Theorem~\ref{cp:C20} is what that bound asks of the construction near $\bz$: in~(27), $\kappa$ enters only through $\min(A,B^{1/\kappa},C)$, which does not decrease when $\kappa$ decreases, and the constraints of~\cite[Theorems~5.3, 6.1 and~6.3]{VXXZ24} do not involve $\kappa$, so a feasible choice of parameters that satisfies (27) with $\omega'=2$ at some $\kappa$ satisfies it at every smaller $\kappa$. The hypothesis asks no more than this: the parameters may depend on $\beta$, and (27) is not needed at $\bz$ itself. Without the reading of the construction in Lemma~\ref{cp:C18}, the printed bound alone gives every $\eps<\cpv{C16_range}$ through Proposition~\ref{cp:C14}, and more through the fast rectangular bottom of Section~\ref{cp:sec:range}.
\end{remark}

\begin{remark}[the standard decomposition does not certify the range]\label{cp:C21}
Under \Vhyp\beta, the restrictions of the proof of Theorem~\ref{cp:C20} with $\Fstdq{5}$ in place of $\Fmrgq{5}$ form a variable-map family for $(F_5,\Fstdq{5})$ with \Tcond1, \Tcond3, \Tcond4{} and $\ln Q_M=o(M)$. By Lemma~\ref{cp:C1}(c), $s_Ma_M^2\le Q_M7^M$, so $\Acs\le\frac12\ln7$, and by Theorem~\ref{cp:C6} and Proposition~\ref{cp:C11} every partition has
\[
\cg^{\mathrm{str}}\ge\frac{\ktyp(\Fstdq{5})}{\Acs}\ge\frac{6\ln2/7}{\frac12\ln7}=\frac{12\ln2}{7\ln7}>\cpv{C21_cg}>\cpv{C21_half_alpha}=\frac\bz2.
\]
So no partition certifies $\cg<\beta/2$ through the structural profile. Since these families may have $p_M>1$ (the identifications in rows 13, 16 and 21 of Appendix~\ref{cp:app:vxxz} give a variable several preimages), this does not bound the profile of Definition~\ref{cp:T1} from below: coefficients of several preimages of an output could cancel.
\end{remark}

\section{Unequal outer sides and ABFK's Hypothesis 1.9}\label{cp:sec:unequal}

This section treats products $XY$ with $X\in\Z^{n_1\times D}$, $Y\in\Z^{D\times n_2}$ and unequal outer sides $n_1\le n_2$, with $D\le n_1^\eps$ and wanted sets $W\subseteq[n_1]\times[n_2]$ of size at most $n_1n_2/D^\kappa$; write $\sigma_P=\ln n_1/\ln n_2$ for the shape of an instance. With a top algorithm whose own outer sides are unequal, \ThmD{} reaches, against the smaller side, inner dimensions up to $n_1^{\cpv{E13_top}}$ for very unbalanced shapes, and in general up to a step function $E(\sigma_P)$ of the shape (Theorem~\ref{cp:E8}); the converse of Section~\ref{cp:sec:range} has an analogue for shapes (Proposition~\ref{cp:E1}). The application is to the hypothesis of Abboud, Bringmann, Fischer and Künnemann~\cite{ABFK24} on partially sparse triangle listing, whose instances are exactly such products, transposed (Section~\ref{cp:sec:unequal-abfk}).

\subsection{Rectangular exponents and the converse for shapes}\label{cp:sec:unequal-conv}

For $\sigma\in(0,1]$ let $\alstar(\sigma)$ be as in Section~\ref{cp:sec:laser}, and put $\eps^*(\sigma):=\alstar(\sigma)/\sigma$; this is not $\es$, the range of~\cite{v2}. Since $\beta\mapsto\omega(\sigma,\beta,1)$ is non-decreasing and $\omega(\sigma,\beta,1)\ge1+\sigma$ (the $n^\sigma\cdot n$ outputs are linearly independent), $\{\beta:\omega(\sigma,\beta,1)=1+\sigma\}$ is an interval starting at~$0$. The model of the next proposition is the multiplication model of Section~\ref{cp:sec:models}.

\begin{proposition}[the converse for shapes; multiplication model]\label{cp:E1}
Fix $\sigma\in(0,1]$, $\eps>0$ and $\kappa_0>0$. Suppose that for every $\kappa\in(0,\kappa_0]$ there are $\gamma=\gamma(\kappa)\ge0$ and $C$ such that for every $N$, with $n_1=\lceil N^\sigma\rceil$, every $1\le D\le n_1^\eps$ and every $W\subseteq[n_1]\times[N]$ with $\abs W\le n_1N/D^\kappa$, there is a division-free straight-line program over $\Q$ in the entries of $X\in\Q^{n_1\times D}$ and $Y\in\Q^{D\times N}$, with control flow depending on $(n_1,D,N,W)$ only, that computes $(XY)|_W$ as polynomials with at most $Cn_1ND^{-\gamma}$ non-scalar multiplications. Then $\omega(\sigma,\sigma\eps,1)=1+\sigma$; hence $\sigma\eps\le\alstar(\sigma)$, that is, $\eps\le\eps^*(\sigma)$.
\end{proposition}

\begin{proof}
Fix $\kappa\in(0,\kappa_0]$ with $\eps\kappa\le1$, put $N$ large, $n_1=\lceil N^\sigma\rceil\le2N^\sigma$ and $D=\lfloor n_1^\eps\rfloor$. Cut $[n_1]\times[N]$ into batches of size $s=\lfloor n_1N/D^\kappa\rfloor$; since $1\le D^\kappa\le n_1\le n_1N$, there are at most $3D^\kappa$ batches. Their programs, run one after the other, form one straight-line program for all entries of $XY$ with at most $3D^\kappa Cn_1ND^{-\gamma}\le6C2^{\eps\kappa}N^{1+\sigma+\sigma\eps\kappa}$ non-scalar multiplications. By Strassen's theorem~\cite{Str73}, $R(\mat{n_1}DN)=O(N^{1+\sigma+\sigma\eps\kappa})$. Since $D\ge n_1^\eps/2\ge N^{\sigma\eps}/2$, splitting the inner dimension into two parts bounds $R(\mat{\lceil N^\sigma\rceil}{\lceil N^{\sigma\eps}\rceil}N)$ by the same quantity times a constant. This holds for every $\kappa\in(0,\kappa_0]$, so $\omega(\sigma,\sigma\eps,1)\le1+\sigma$, and the reverse inequality holds because the $n_1N$ outputs are linearly independent. Hence $\sigma\eps\le\alstar(\sigma)$.
\end{proof}

As in Remark~\ref{cp:A8}, the every-$\kappa$ form is essential: one fixed $\kappa$ gives only $\omega(\sigma,\sigma\eps,1)\le1+\sigma+\sigma\eps\kappa$. Algorithm~D$_\sigma$ below is of the required form (Remark~\ref{cp:E4:alg}).

\subsection{\texorpdfstring{\ThmDsig}{Theorem D\_sigma}: unequal tops}\label{cp:sec:unequal-Ds}

\begin{definition}[input families for unequal tops]\label{cp:E3:def}
An \emph{input family for unequal tops} is a sequence $(A_M)$, $M$ in an infinite set of integers, where $A_M$ is a border decomposition over $\Q[\lambda]$ of $s_M\cp\mat{e_M}{h_M}{l_M}$ with $r_M$ terms of nonzero $z$-form, product degree at most $d_M$, and a partition into $m_M$ classes with profile $Z_M$ and maximal coverage $C_M$ (as in Definition~\ref{cp:T1}; each copy has the $e_Ml_M$ outputs of an $e_M\times l_M$ matrix). Put $L_M=\ln(e_Ml_M)$. Its hypotheses are:
\begin{itemize}
\item[\Hscond1] $e_M,l_M\to\infty$, $h_M\ge2$, $\ln e_M/\ln l_M\to\sigma$ and $\ln h_M/\ln l_M\to\beta$, with $0<\beta<\sigma\le1$;
\item[\Hscond2] $\rho_M:=\ln(r_M/s_M)/L_M\to1$;
\item[\Hscond3] $(\ln(d_M+1)+\ln m_M)/L_M\to0$;
\item[\Hscond4] $\cbar:=\sup_M\ln C_M/L_M<\infty$.
\end{itemize}
Its profile is $\Lam_\sigma(t)=\limsup_M\mu^\sigma_M(t)$ with $\mu^\sigma_M(t)=\ln\bigl(s_M^{t-1}Z_M(t)\bigr)/L_M-1+t$, and $\cg=\liminf_{t\to0+}\Lam_\sigma(t)/t$. For an exact top algorithm $B$ over $\Q$ for $\mat{k_1}b{k_2}$ with a partition we write $\sigma_B=\ln k_1/\ln k_2$, $\beta_B=\ln b/\ln k_2$, $\RB=(k_1k_2)^{1+\delB}$ (for $k_1=k_2$, half the $\delB$ of Definition~\ref{cp:T1}) and $\lambda^\sigma_B(t)=\ln\ZB(t)/\ln(k_1k_2)-1+t$.
\end{definition}

For $e_M=l_M$ the normalisation $L_M=2\ln a_M$ halves every exponent of Definition~\ref{cp:T1}. As in Lemma~\ref{cp:T1:basic}, $\mu^\sigma_M(t)\ge\rho_M-1\ge0$: the $z$-flattening of $s\cp\mat ehl$ has rank $sel$, so $r_M\ge s_Me_Ml_M$, and $\kappa_c\ge\abs{P_c}/(s_Me_Ml_M)$.

\begin{lemma}[top algorithms of unequal shape]\label{cp:E3:top}
Let $(A_M)$ satisfy \Hscond1--\Hscond4. For every $t\in(0,1/(2(\cbar+1))]$ and every $\eta>0$ there are $M$ and $j$ such that the construction of Lemma~\ref{cp:T5:lemma4}, applied to $A_M$ as a decomposition of $s_M\cp\mat{e_M}{h_M}{l_M}$, gives an exact top algorithm $B$ over $\Q$ for $\mat{k_1}b{k_2}=\mat{e_M^j}{h_M^j}{l_M^j}$ with a partition such that $\sigma_B=\ln e_M/\ln l_M$, $\beta_B=\ln h_M/\ln l_M$, $\abs{\sigma_B-\sigma}\le\eta$, $\abs{\beta_B-\beta}\le\eta$, $\delB\le\eta$ and $\lambda^\sigma_B(t)\le\Lam_\sigma(t)+\eta$.
\end{lemma}

\begin{proof}
The proof of Lemma~\ref{cp:T5:lemma4} never uses that the two outer sides are equal: Step~1 (interpolation) is shape-free; in Step~2, $s\cp T_0^{\otimes(j+1)}=(s\cp T_0)\otimes T_0^{\otimes j}$ with $T_0=\mat ehl$, the matrices $X^{(m)}$ and $Y^{(m)}$ have shapes $e^j\times h^j$ and $h^j\times l^j$, and the grouping and the injection behind $\kappa^{j+1}_{(c,d)}\le\kappa'_c\kappa^j_d$ do not use $e=l$; Steps~3 and~4 are shape-free. Dividing by $jL$ in place of $2j\ln a$, it gives
\[
\lambda^\sigma_B(t)\le\Bigl(1-\frac1j\Bigr)\Bigl[\max\Bigl(\mu^\sigma(t),\,tc-1+t+\frac{\ln m}L\Bigr)+\frac{\ln(2(d+1))}L\Bigr]+\frac1j\,\frac{\ln((d+1)r)}L,\qquad c=\frac{\ln C_A}L .
\]
For $t\le1/(2(\cbar+1))$ the padding term satisfies $tc_M-1+t\le t(\cbar+1)-1\le-\frac12$, so for large $M$ it is below $0\le\mu^\sigma_M(t)$. The rest is the proof of Corollary~\ref{cp:T5}, with \Hscond1--\Hscond4{} in place of \Hcond1--\Hcond4. The shape exponents are constant along the powers, so $\sigma_B=\ln e_M^j/\ln l_M^j=\ln e_M/\ln l_M$, and likewise $\beta_B$.
\end{proof}

The restriction $t\le1/(2(\cbar+1))$ is conservative (the padding term is already $\le0$ for $t\le1/(\cbar+1)$); only small $t$ is used below.

\begin{lemma}[the square family as a family for \ThmDsig]\label{cp:E8:square}
Let $(A_M)$ be an input family in the sense of Definition~\ref{cp:T1} with $\alpha\in(0,1)$, and read $A_M$ as a decomposition of $s_M\cp\mat{e_M}{h_M}{l_M}$ with $e_M=l_M=a_M$ and $h_M=b_M$, with $d_M$ the product degree (called $e_M$ in Definition~\ref{cp:T1}). Then $(A_M)$ satisfies \Hscond1--\Hscond4{} with $\sigma=1$ and $\beta=\alpha$, $L_M=2\ln a_M$, its profile is $\Lam_\sigma=\Lam/2$, and its constant in \Hscond4{} is $\cbar/2$; in particular its typical support exponent is $\cg/2$.
\end{lemma}

\begin{proof}
With $L_M=2\ln a_M$: $\ln e_M/\ln l_M=1$ and $\ln h_M/\ln l_M=\alpha_M\to\alpha$, so \Hscond1{} holds with $\sigma=1>\beta=\alpha>0$; $\ln(r_M/s_M)/L_M=\rho_M/2\to1$ by \Hcond2; $(\ln(d_M+1)+\ln m_M)/L_M\to0$ by \Hcond3; and $\sup_M\ln C_M/L_M=\cbar/2$. The profile is $\ln(s_M^{t-1}Z_M(t))/L_M-1+t=\bigl(\mu_M(t)+2-2t\bigr)/2-1+t=\mu_M(t)/2$, and the limsup over $M$ gives $\Lam/2$.
\end{proof}

For this family the $t$-range of Lemma~\ref{cp:E3:top} is $t\le1/(\cbar+2)$, and the quantities of Theorem~\ref{cp:E3} below at $\sigma=1$ are those of \ThmDk.

\paragraph{Algorithm D$_\sigma$.}
Given $n_1\le n_2$, $D$, $X\in\Z^{n_1\times D}$, $Y\in\Z^{D\times n_2}$ and $W\subseteq[n_1]\times[n_2]$, a top algorithm $B$ for $\mat{k_1}b{k_2}$ and $x>0$:
\begin{enumerate}
\item Put $n=\lfloor x\ln D/((1+\beta_Bx)\ln k_2)\rfloor$, $K_1=k_1^n$, $K_2=k_2^n$, $B'=b^n$ and $D_A=\lceil D/B'\rceil$; pad to $n_1'=K_1\lceil n_1/K_1\rceil$, $n_2'=K_2\lceil n_2/K_2\rceil$ and $D'=B'D_A$; put $N_1=n_1'/K_1$ and $N_2=n_2'/K_2$.
\item Cut $X$ into $K_1\times B'$ blocks of size $N_1\times D_A$ and $Y$ into $B'\times K_2$ blocks of size $D_A\times N_2$, and compute the encodings $X_l$, $Y_l$ of $B^{\otimes n}$ for all $l\in[\RB]^n$, one coordinate at a time, after scaling $B$ to integer coefficients.
\item Bucket $W$ by position $(r,s)\in[N_1]\times[N_2]$ and build the needed sets $M_n,\dots,M_0=T_{rs}$ (Lemma~\ref{cp:T4}).
\item For every position and every $l\in T_{rs}$ compute $p_l=\sum_{j<D_A}X_l[r,j]Y_l[j,s]$; decode by Lemma~\ref{cp:T4}; divide by the scaling constant.
\end{enumerate}
Lemmas~\ref{cp:T2}--\ref{cp:T4} hold for $B^{\otimes n}$ with outputs $[K_1]\times[K_2]$ and positions $[N_1]\times[N_2]$, with the same proofs: Lemma~\ref{cp:T2} is a bound over the set of positions and a product over the $n$ coordinates, with $\ZB(t)^n=(K_1K_2)^{1-t+\lambda^\sigma_B(t)}$ and $\bar w=\abs W/(N_1N_2)$; Lemmas~\ref{cp:T3} and~\ref{cp:T4} are statements about the terms of $B$ and one position.

\begin{theorem}[\ThmDsig]\label{cp:E3}
Let $(A_M)$ satisfy \Hscond1--\Hscond4{} with limits $\sigma$, $\beta$, profile $\Lam_\sigma$ and constant $\cbar$. Let $\kappa>0$, $x>0$ and $t\in(0,1/(2(\cbar+1))]$ satisfy
\[
\Qgain\sigma\kappa(x,t):=\frac{t\kappa(1+\beta x)-(1+\sigma)x\Lam_\sigma(t)-1}{1+\beta x}>0 .
\]
Let $0<\eps<\Erange\sigma(x):=(1+\beta x)/(1+\sigma x)$, let
\[
0<\gamma<\gbound\sigma\kappa(x,t,\eps):=\min\Bigl\{\Qgain\sigma\kappa(x,t),\ \frac1\eps-\frac{1+\sigma x}{1+\beta x},\ \frac{(1+\sigma)x}{1+\beta x}\Bigr\},
\]
and let $c_0\ge1$. Then there is $C$ such that for all $n_1\le n_2$ with $n_1\le c_0n_2^\sigma$, all $2\le D\le n_1^\eps$, all $X\in\Z^{n_1\times D}$ and $Y\in\Z^{D\times n_2}$ with entries of absolute value $n_2^{O(1)}$ and every $W\subseteq[n_1]\times[n_2]$ with $\abs W\le n_1n_2/D^\kappa$, Algorithm~D$_\sigma$ with a suitable top algorithm $B$ computes $(XY)|_W$ in at most $Cn_1n_2D^{-\gamma}$ word-RAM steps (time model; deterministic).
\end{theorem}

\begin{proof}
\emph{Step 1 (top).} Lemma~\ref{cp:E3:top} gives, for every $\eta>0$, a top algorithm $B$ with $(\beta_B,\sigma_B,\delB)$ within $\eta$ of $(\beta,\sigma,0)$ and $\lambda^\sigma_B(t)\le\Lam_\sigma(t)+\eta$; $B$ does not depend on $W$. The order of the choices is $\kappa,\eps,t,x,\gamma$, then $\eta$, then $B$.

\emph{Step 2 (sizes).} $K_2$ lies within a factor $k_2$ of $D^{x/(1+\beta_Bx)}$, $K_1=K_2^{\sigma_B}$, $B'=K_2^{\beta_B}\le D$ (since $\beta_Bx/(1+\beta_Bx)<1$), and $D_A\le2bD^{1/(1+\beta_Bx)}$. The padding needs $K_1\le n_1$ and $K_2\le n_2$; both follow, for $D\ge D_0$, from the encoding bounds of Step~3, which give $K_1D_A\le n_1D^{-\gamma'}$ and $K_2D_A\le c\,n_2D^{-\gamma'}$ for some $\gamma'>0$.

\emph{Step 3 (encodings).} The $X$-encoding has $n$ stages; stage $a$ holds $\RB^a(k_1b)^{n-a}$ blocks of size $N_1\times D_A$, each made with at most $k_1b$ operations per entry. Since the $x$-flattening of $\mat{k_1}b{k_2}$ has rank $k_1b\le\RB$, the total is $O(n\RB^nN_1D_A)=O(n\,n_1K_2(K_1K_2)^{\delB}D_A)$. Likewise the $Y$-encoding costs $O(n\,n_2K_1(K_1K_2)^{\delB}D_A)$, using $k_2b\le\RB$. The instance gives $n_1\ge D^{1/\eps}$ and $n_2\ge(n_1/c_0)^{1/\sigma}\ge c_0^{-1/\sigma}D^{1/(\sigma\eps)}$. So the two encodings cost $O(n_1n_2\log D\cdot D^{-g})$ with the gains
\[
g_Y=\frac1\eps-\frac{\sigma_Bx+1+(1+\sigma_B)x\delB}{1+\beta_Bx},\qquad g_X=\frac1{\sigma\eps}-\frac{x+1+(1+\sigma_B)x\delB}{1+\beta_Bx}
\]
(the latter with the constant $c_0^{1/\sigma}$). At the limit point $(\beta_B,\sigma_B,\delB)=(\beta,\sigma,0)$,
\[
g_X-g_Y=(1-\sigma)\Bigl(\frac1{\sigma\eps}-\frac x{1+\beta x}\Bigr)\ge0,
\]
because $\eps<(1+\beta x)/(1+\sigma x)<(1+\beta x)/(\sigma x)$ and $\sigma\le1$. So at the limit the $X$-encoding gain is at least the $Y$-encoding gain, which is at least the second term of $\gbound\sigma\kappa$. The gains are continuous in $(\beta_B,\sigma_B,\delB)$ near the limit (as $1+\beta x>0$) and non-increasing in $\delB$, so with the strict inequality $\gamma<\gbound\sigma\kappa$ every top within a small enough $\eta$ of the limit has both gains at least some $\gamma'>\gamma$. In words: the smaller side $n_1$ is paired with the smaller top dimension $K_1$, and the binding condition is $K_1D_A<n_1$; an instance with $\sigma_P<\sigma$ only makes $n_2$ larger.

\emph{Step 4 (bucketing).} $O(N_1N_2+\abs W)$ with $N_1N_2\le4n_1n_2/(K_1K_2)$: gain $(1+\sigma_B)x/(1+\beta_Bx)$, and $\kappa$.

\emph{Steps 5--6 (queries).} With $\bar w=\abs W/(N_1N_2)\le K_1K_2D^{-\kappa}$ and $\ZB(t)^n=(K_1K_2)^{1-t+\lambda^\sigma_B(t)}$, Lemma~\ref{cp:T2} bounds the number of needed sub-entries by $n_1'n_2'D^{-t\kappa}(K_1K_2)^{\lambda^\sigma_B(t)}$; each costs $D_A$. The gain is $t\kappa-[(1+\sigma_B)x\lambda^\sigma_B(t)+1]/(1+\beta_Bx)$, which tends to $\Qgain\sigma\kappa(x,t)$ or more as $\eta\to0$; here $\lambda^\sigma_B(t)\ge0$, as in Lemma~\ref{cp:T1:basic}.

\emph{Step 7 (decoding).} Lemma~\ref{cp:T4}: a factor $O(n\cB)=O(\log D)$.

\emph{Step 8 (tables).} $O(n(K_1K_2)^{1+\delB})$: relative to $n_1n_2$ this is the $Y$-encoding term times $K_2/(n_2D_A)\le1$.

\emph{Step 9 (input and output).} $(n_1+n_2)D\le2n_2D$, so the relative cost is at most $2D/n_1\le2D^{1-1/\eps}$, with gain $1/\eps-1\ge g_Y$ at the limit because $(1+\sigma x)/(1+\beta x)\ge1$. The output has $\abs W\le n_1n_2D^{-\kappa}$ numbers, gain $\kappa>\Qgain\sigma\kappa(x,t)$.

\emph{Steps 10--14.} $\gamma<\kappa$, since $\Qgain\sigma\kappa(x,t)<t\kappa$ by $\Lam_\sigma\ge0$. For $D<D_0$ the product is computed directly in $O(n_1n_2D_0^{1+\gamma}D^{-\gamma})$. There is no sub-instance range condition. After scaling, every number is $n_2^{O(1)}$, since the encodings multiply entries by at most $(k_1bH)^n=D^{O(1)}$, $H$ the largest scaled coefficient. The factors $\log D$ are absorbed by the margin $\gamma'>\gamma$.

\emph{The constant $c_0$.} The bound $n_1\le c_0n_2^\sigma$ enters only through the lower bound on $n_2$ in $g_X$ and in Step~2, and weakens it by the constant factor $c_0^{1/\sigma}$; so $C$ depends on $c_0$.

\emph{Correctness.} $B^{\otimes n}$ is an exact bilinear algorithm for $\mat{K_1}{B'}{K_2}$; it applies to block matrices with compatible blocks ($X$-blocks $N_1\times D_A$ on the left, $Y$-blocks $D_A\times N_2$ on the right), and Lemma~\ref{cp:T4} decodes exactly.
\end{proof}

\begin{corollary}[\ThmDsig, consequences]\label{cp:E4}
In the setting of Theorem~\ref{cp:E3}:
\begin{enumerate}[label=(\alph*)]
\item if $\cg=0$, then for every $\kappa>0$ and every $0<\eps\le\beta/\sigma$ some $\gamma>0$ has the property of Theorem~\ref{cp:E3} (for every $c_0$); and for every $t\le1/(2(\cbar+1))$ with $t\kappa\beta>(1+\sigma)\Lam_\sigma(t)$ and $t\kappa<1$, with $\xmin(t)=(1-t\kappa)/(t\kappa\beta-(1+\sigma)\Lam_\sigma(t))$, so does every $\eps<\beta/\sigma+(1-\beta/\sigma)/(1+\sigma\xmin(t))$;
\item for $0<\eps<\beta/\sigma$, every $\gamma$ below $\min\{\sup_t(t\kappa-(1+\sigma)\Lam_\sigma(t)/\beta),\ 1/\eps-\sigma/\beta,\ (1+\sigma)/\beta\}$, the supremum over $t\in(0,1/(2(\cbar+1))]$, is admissible.
\end{enumerate}
\end{corollary}

\begin{proof}
(a) Since $\cg=0$ is a liminf, there are arbitrarily small $t$ with $(1+\sigma)\Lam_\sigma(t)\le t\kappa\beta/2$; then $x>2/(t\kappa\beta)$ gives $\Qgain\sigma\kappa(x,t)>0$. Also $(1+\beta x)/(1+\sigma x)-\beta/\sigma=(\sigma-\beta)/(\sigma(1+\sigma x))>0$, so every $\eps\le\beta/\sigma$ lies in the range. The second half: $\Qgain\sigma\kappa(x,t)>0$ if and only if $x(t\kappa\beta-(1+\sigma)\Lam_\sigma(t))>1-t\kappa$, and $(1+\beta x)/(1+\sigma x)=\beta/\sigma+(1-\beta/\sigma)/(1+\sigma x)$ decreases in $x$. (b) Let $x\to\infty$: $\Qgain\sigma\kappa(x,t)\to t\kappa-(1+\sigma)\Lam_\sigma(t)/\beta$, $1/\eps-(1+\sigma x)/(1+\beta x)$ decreases to $1/\eps-\sigma/\beta$, and $(1+\sigma)x/(1+\beta x)$ increases to $(1+\sigma)/\beta$; the range condition holds for every $x$, since $(1+\beta x)/(1+\sigma x)>\beta/\sigma>\eps$.
\end{proof}

\begin{remark}[the algebraic form of Algorithm D$_\sigma$]\label{cp:E4:alg}
As for \AlgD{} (Remark~\ref{cp:T6:alg}), for each $(n_1,D,n_2,W)$ Algorithm~D$_\sigma$ is a division-free straight-line program over $\Q$ whose control flow depends on $(n_1,D,n_2,W)$ only, consisting of integer encodings, inner products of encoded entries, decodings with the fixed coefficients of $B$ and one multiplication by a rational constant; it is correct as a polynomial identity, because $B^{\otimes n}$ is an exact bilinear algorithm over~$\Q$.
\end{remark}

The transposed products are needed for the application in Section~\ref{cp:sec:unequal-abfk}.

\begin{proposition}[transposition]\label{cp:E5}
Let $\mathcal A$ be \AlgD{} of Theorem~\ref{cp:T6}, Algorithm~D$_\sigma$ of Theorem~\ref{cp:E3}, or the per-tile algorithm of the proof of Theorem~\ref{cp:A10}, each with its parameters fixed. There is an algorithm $\mathcal A^{\mathsf T}$ that, given $W$, $Y$ and $\Omega\in\Z^{n_1\times n_2}$ supported on $W$ with entries of absolute value $n_2^{O(1)}$, computes $\Omega Y^{\mathsf T}\in\Z^{n_1\times D}$ in at most the running time of $\mathcal A$ on $(W,Y)$ plus $O(\abs W+n_1D)$ steps, with a control flow that is the same function of the sizes and $W$. Consequently Theorems~\ref{cp:T6}, \ref{cp:T7}, \ref{cp:E3} and Corollary~\ref{cp:E4} hold for $\Omega Y^{\mathsf T}$ in place of $(XY)|_W$, with the same ranges and the same $\gamma$, since $\abs W\le n_1n_2D^{-\kappa}$, $n_1D\le n_1n_2D^{1-1/\eps}$ and $\gamma<\min(\kappa,1/\eps-1)$; and with Lemma~\ref{cp:A10:tiling}, whose per-tile bound keeps its form under the extra cost $O(\abs{W'}+nD)$ since $nD\le n^2D^{1-1/\eps}$, so does Theorem~\ref{cp:A10}, also for rectangular $W$.
\end{proposition}

\begin{proof}
For fixed sizes, $W$ and $Y$, \AlgD{} computes the linear map $L_Y=d^{-1}\,\mathrm{Dec}\circ P_Y\circ E_X$ in $X$, where $E_X:X\mapsto(X_l)_l$ is the $X$-encoding (Step~2), $P_Y$ multiplies by the fixed $Y$-encodings ($p_{l,r,s}=\sum_jX_l[r,j]Y_l[j,s]$), $\mathrm{Dec}$ is the decoding of Lemma~\ref{cp:T4} and $d$ the scaling constant. $\mathcal A^{\mathsf T}$ computes $L_Y^{\mathsf T}=d^{-1}E_X^{\mathsf T}P_Y^{\mathsf T}\mathrm{Dec}^{\mathsf T}$:
(1$'$) the same $n$, $K$, $D_A$ and padding (the transpose of zero-padding is truncation); (2$'$) the $Y$-encodings $Y_l$, exactly as in Step~2; (3$'$) the same bucketing and needed sets $M_n,\dots,M_0=T_{rs}$; (5$'$) the transposed decoding: per position, start with $u_n(o)=\Omega_{o,r,s}$ on $M_n=W_{rs}$ and for $i=n,\dots,1$ scatter $u_{i-1}(o_{<i},l,l_{>i})\mathrel{+}=w_l[o_i]\,u_i(o_{\le i},l_{>i})$ for each element of $M_i$ and each $l\in\Cov(o_i)$; the result $\pi_{l,r,s}=u_0(l)$ on $T_{rs}$ equals $\sum_{o\in W_{rs},\,l\in\Cov(o)}w_l[o]\Omega_{o,r,s}$, and the scatter traverses exactly the pairs that Lemma~\ref{cp:T4} traverses, at the same cost; (4$'$) $G_l[r,j]\mathrel{+}=\pi_{l,r,s}Y_l[j,s]$ for every position, every $l\in T_{rs}$ and $j<D_A$, at the cost of Step~4's inner products; (2$''$) the transposed $X$-encoding $G=E_X^{\mathsf T}((G_l)_l)$, by running Step~2's recursion backwards with every stage's coefficient matrix transposed, at the cost of the encoding; (6$'$) output $d^{-1}G$ on the unpadded $n_1\times D$ part; for $D<D_0$ the direct sparse product.
Correctness: $L_Y(X)=(XY)|_W$ is a polynomial identity (Remark~\ref{cp:T6:alg}), so for all $X$, $\langle L_Y^{\mathsf T}(\Omega),X\rangle=\langle\Omega,(XY)|_W\rangle=\sum_{(i,j)\in W}\Omega_{ij}\sum_kX_{ik}Y_{kj}=\langle\Omega Y^{\mathsf T},X\rangle$, hence $L_Y^{\mathsf T}(\Omega)=\Omega Y^{\mathsf T}$; $\Omega Y^{\mathsf T}$ is an integer matrix, so the final division is exact. Every number is bounded as in Theorem~\ref{cp:T6}, by $(\cB H)^n\max\abs\Omega$, $N_{\mathrm{sub}}D^{O(1)}\max\abs\pi\max\abs{Y_l}$ and $(\RB H)^n\max\abs{G_l}$, all $n_2^{O(1)}$.
Algorithm~D$_\sigma$ uses the same three linear maps between spaces of other sizes, so the same reverse pass applies. If the smaller side is $n_2$, $(XY)|_W=((Y^{\mathsf T}X^{\mathsf T})|_{W^{\mathsf T}})^{\mathsf T}$ is computed by Algorithm~D$_\sigma$ on $(Y^{\mathsf T},X^{\mathsf T})$, linear in its right factor, and the same reverse pass with the roles of the two factors exchanged computes $\Omega Y^{\mathsf T}$.
For the per-tile algorithm of Theorem~\ref{cp:A10}, the terms with option~(S) are transposed as above. A term with option~(F) forms the $N_{\mathrm{sub}}\times N_{\mathrm{sub}}$ matrix $\pi_l$, zero outside $W_l$, at cost $N_{\mathrm{sub}}^2$, below the cost of its bottom product, and computes $\pi_lY_l^{\mathsf T}$ with the transposed bottom of Lemma~\ref{cp:A10:bottom}, at the same cost as the forward bottom. The choice between the options depends on $W$ only, so it fixes one linear program.
\end{proof}

\subsection{The unequal level-1 family and the step function}\label{cp:sec:unequal-level1}

Fix one $q\ge2$ ($\Qs=\{q\}$, $\lambda_q=1$). A \emph{single-$q$ unequal point} is a rational $P=(p_e,p_l,u,v_x,v_y)$ with
\[
p_e+p_l=\frac q{q+2},\qquad u+v_x+v_y=w:=\frac1{q+2},\qquad\text{all five}>0,\qquad u<p_e\le p_l .
\]
Its joint type is $p_{101}=p_e$, $p_{011}=p_l$, $p_{002}=w$, $p_{110}=u$, $p_{200}=v_x$, $p_{020}=v_y$; its $X$-, $Y$- and $Z$-marginals are $X=(p_l+w+v_y,\,p_e+u,\,v_x)$, $Y=(p_e+w+v_x,\,p_l+u,\,v_y)$ and $Z=(w,\,q/(q+2),\,w)$, the last one the same as in Definition~\ref{cp:L3}. Put $\slack_X=\Ent(X)-\Ent(Z)$ and $\slack_Y=\Ent(Y)-\Ent(Z)$. Let $d(P)$ be the least positive integer such that all block counts $M$ times the joint and marginal probabilities are integers for every multiple $M$ of $d(P)$. As in Section~\ref{cp:sec:family}, $A_M$ is the restriction of the product decomposition carrying $\Fmrg$ on every coordinate by the zero-out of Proposition~\ref{cp:E6} below, with the terms of vanishing restricted forms dropped and the composition partition.

\begin{proposition}[the unequal level-1 family]\label{cp:E6}
Let $P$ be a single-$q$ unequal point with $\slack_X>0$ and $\slack_Y>0$, and let $M$ run through the multiples of $d(P)$. There is a zero-out of $F_q^{\otimes M}$ whose image is the direct sum of $s_M$ copies of $\mat{e_M}{h_M}{l_M}$ with
\[
\ln e_M=Mp_e\ln q,\qquad\ln h_M=Mu\ln q,\qquad\ln l_M=Mp_l\ln q,
\]
whose kept outputs are $s_Me_Ml_M$ distinct $z$-strings of the kept output type $\Ukept$ of Definition~\ref{cp:L3} (with $\Qs=\{q\}$), each covered by a term of nonzero restricted $z$-form, and with $\ln s_M\ge M(\Ent(Z)-\eta)$ for every $\eta>0$ and all large $M$. The restricted merged family $(A_M)$ satisfies \Hscond1--\Hscond4{} with
\[
\sigma=\frac{p_e}{p_l},\qquad\beta=\frac u{p_l},\qquad\Lam_\sigma(t)\le\frac{t^2\ln^2(4eq/t)}{\frac q{q+2}\ln q}\ \ (t\in(0,1]),\qquad\cbar\le\frac{(\ln(q+2)+q\ln2)/(q+2)}{\frac q{q+2}\ln q},
\]
and $\cg=0$. The only cited input is Theorem~\ref{cp:P6:LG}.
\end{proposition}

\begin{proof}
\emph{The count.} This is the proof of Lemma~\ref{cp:L4} with unequal $X$- and $Y$-marginals. After the first zero-out of the block strings of wrong marginal types, the set $U$ of remaining block triples consists of the triples of the joint type of $P$, because at level~1 the joint type is determined by the three marginals: $p_{200}$, $p_{020}$ and $p_{002}$ are the block-$2$ frequencies, and $p_{110}=(X_1+Y_1-Z_1)/2$, $p_{101}=(X_1+Z_1-Y_1)/2$, $p_{011}=(Y_1+Z_1-X_1)/2$. So $U^*=U$ in Theorem~\ref{cp:P6:LG}. The coordinate permutations map $U$ onto itself and act transitively on the block strings of each marginal type, so every such $x$-, $y$- or $z$-block string lies in the same number $N_1$, $N_2$ or $N_3$ of triples of $U$; the theorem allows $N_1$, $N_2$, $N_3$ to differ. With $T_X$, $T_Y$, $T_Z$ the numbers of block strings of the three marginal types, $N_1+N_2+N_3\le3\abs U/\min(T_X,T_Y,T_Z)$, and Theorem~\ref{cp:P6:LG} leaves, by removal operations only, $s_M\ge c_\eps3^{-1-\eps}\min(T_X,T_Y,T_Z)\,(\abs U/\min(T_X,T_Y,T_Z))^{-\eps}$ pairwise block-disjoint triples of $U$. Since $\ln T_X\ge M\Ent(X)-2\ln(M+1)$, $\ln T_Y\ge M\Ent(Y)-2\ln(M+1)$ and $\ln T_Z\ge M\Ent(Z)-2\ln(M+1)$, and both slacks are positive, $\ln\min(T_X,T_Y,T_Z)\ge M\Ent(Z)-2\ln(M+1)$; with $\ln\abs U\le M\Ent(\text{joint type})$ this gives the bound on $s_M$ as in Lemma~\ref{cp:L4}(ii). A surviving triple $(I,J,K)$ is $T_{IJK}$, the tensor product of $Mp_e$ copies of $\mat q11$ (type $101$), $Mp_l$ copies of $\mat11q$ (type $011$), $Mu$ copies of $\mat1q1$ (type $110$) and copies of $\mat111$, that is $\mat{e_M}{h_M}{l_M}$; its outputs are the $q^{M(p_e+p_l)}=e_Ml_M$ $z$-variables of block $K$, which lie in $\Ukept$, and distinct triples have distinct $K$. The covering statement is Lemma~\ref{cp:L4}(iii) with $s_Ma_M^2$ replaced by $s_Me_Ml_M$.
\emph{\Hscond1.} $\ln e_M/\ln l_M=p_e/p_l$ and $\ln h_M/\ln l_M=u/p_l$ for every $M$, and $0<u<p_e\le p_l$ gives $0<\beta<\sigma\le1$; $h_M\ge2$ for large $M$.
\emph{\Hscond2.} $\Ent(Z)+\frac q{q+2}\ln q=\ln(q+2)$ and $L_M=\ln(e_Ml_M)=M\frac q{q+2}\ln q$, so $\ln s_M+L_M\ge M\ln(q+2)-o(M)=\ln R_F-o(M)$ with $R_F=(q+2)^M$ the number of term strings. With $s_Me_Ml_M\le r_M\le R_F$ this gives $\rho_M\to1$.
\emph{\Hscond3.} The product degree is at most $9M$ (Lemma~\ref{cp:L2}(b)) and there are at most $(M+1)^2$ composition classes; both are $e^{O(\ln M)}$, against $L_M=\Theta(M)$.
\emph{\Hscond4.} The coverage of an output in $\Ukept$ is the product over the coordinates of $q+2$, $2$ or $1$ (Lemma~\ref{cp:L2}(a)); restriction only lowers it; and $L_M=M\frac q{q+2}\ln q$ does not depend on how $q/(q+2)$ is split between $p_e$ and $p_l$.
\emph{Profile.} Lemma~\ref{cp:L6:lemma5} holds for a zero-out with image $s\cp\mat ehl$, with $sa^2$ replaced by $sel$ (its proof uses the shape only through the number of kept outputs), and with $Z^*=\Ukept$ it applies by Lemma~\ref{cp:L6:lemmaC}: $\Ukept$ and the term strings are those of the square case, since the $Z$-marginal is unchanged. It gives $\mu^\sigma_M(t)\le\ln(R_F/(s_Me_Ml_M))/L_M+\ln(\mathrm{mean}_\tau\Sall(\tau)^t)/L_M$. The first term tends to $0$ by the computation in \Hscond2, and Theorem~\ref{cp:L7} with $\lambda_q=1$ bounds the second by $t^2\ln^2(4eq/t)/(\frac q{q+2}\ln q)$. So $\Lam_\sigma(t)\le t^2\ln^2(4eq/t)/(\frac q{q+2}\ln q)=o(t)$, and $\cg=0$.
\end{proof}

The restriction keeps other $x$- and $y$-blocks than in the square case. That does not matter, because Lemma~\ref{cp:L6:lemma5} uses only upper bounds: restriction only removes terms and supports.

\begin{table}[ht]\centering\small
\begin{tabular}{lllll}
$i$ & point $(q,\sigma,u(q+2),v_x(q+2))$ & $E_i=u/p_e$ & $\beta=u/p_l$ & $\slack_X,\ \slack_Y>$\\\hline
$1$ & $\cpv{E7pt_a}$ & $\cpv{E7E_a}=\cpv{E7Ed_a}$ & $\cpv{E7b_a}$ & $\cpv{E7sl_a}$\\
$2$ & $\cpv{E7pt_b}$ & $\cpv{E7E_b}>\cpv{E7Ed_b}$ & $\cpv{E7b_b}$ & $\cpv{E7sl_b}$\\
$3$ & $\cpv{E7pt_c}$ & $\cpv{E7E_c}=\cpv{E7Ed_c}$ & $\cpv{E7b_c}$ & $\cpv{E7sl_c}$\\
$4$ & $\cpv{E7pt_d}$ & $\cpv{E7E_d}=\cpv{E7Ed_d}$ & $\cpv{E7b_d}$ & $\cpv{E7sl_d}$\\
$5$ & $\cpv{E7pt_e}$ & $\cpv{E7E_e}>\cpv{E7Ed_e}$ & $\cpv{E7b_e}$ & $\cpv{E7sl_e}$\\
$6$ & $\cpv{E7pt_f}$ & $\cpv{E7E_f}>\cpv{E7Ed_f}$ & $\cpv{E7b_f}$ & $\cpv{E7sl_f}$\\
$7$ & $\cpv{E7pt_g}$ & $\cpv{E7E_g}>\cpv{E7Ed_g}$ & $\cpv{E7b_g}$ & $\cpv{E7sl_g}$
\end{tabular}
\caption{The seven certified single-$q$ unequal points, rows $i=1,\dots,7$: $w=1/(q+2)$, $p_l=q/((1+\sigma)(q+2))$, $p_e=\sigma p_l$, $v_y=w-u-v_x$; the slacks are certified lower bounds.}\label{cp:E7:table}
\end{table}

\begin{lemma}[seven certified points]\label{cp:E7}
For each row of Table~\ref{cp:E7:table}, the point with the listed $q$, $\sigma$, $u$ and $v_x$, with $w$, $p_l$, $p_e$ and $v_y$ as in the caption, is a single-$q$ unequal point (all five probabilities are positive and $u<p_e\le p_l$) with $\slack_X$ and $\slack_Y$ above the listed lower bounds, so Proposition~\ref{cp:E6} applies to it; and $p_e/p_l=\sigma$, $u/p_l=\beta$ and $u/p_e=E_i$ exactly.
\end{lemma}

The identities are exact rational arithmetic and the slacks are enclosed in exact rational arithmetic with rigorous logarithm enclosures; Section~\ref{cp:sec:certificates} names the script. The smallest slack exceeds $\cpv{E7_slack_min}$, so these points are far from the boundary; the square point $\Pzero$ has slack above $\cpv{slackP0_lo}$ (Lemma~\ref{cp:L11}).

For $s\in(0,1]$ let
\[
E(s):=\max\bigl(\{\alpha(\Pzero)\}\cup\{E_i:\sigma_i\ge s\}\bigr),
\]
with $(\sigma_i,E_i)$ the rows of Table~\ref{cp:E7:table}. $E$ is a non-increasing step function, constant on the intervals $(\sigma_{i+1},\sigma_i]$, with $E=\alpha(\Pzero)$ on $(\frac34,1]$ and largest value $E_7=\cpv{E7E_g}>\cpv{E13_top}$ on $(0,\frac1{10}]$.

\begin{theorem}[$\gamma$ uniform in the shape]\label{cp:E8}
For every $\kappa>0$, $\eps>0$ and $c_0\ge1$ there are $\gamma>0$ and $C$ with the following property. For all integers $2\le n_1\le n_2$ and every $s\in(0,1]$ with $n_1\le c_0n_2^s$ and $\eps\le E(s)$ (for instance $s=\ln n_1/\ln n_2$ with $c_0=1$), every $2\le D\le n_1^\eps$, all $X\in\Z^{n_1\times D}$ and $Y\in\Z^{D\times n_2}$ with entries of absolute value $n_2^{O(1)}$ and every $W$ with $\abs W\le n_1n_2/D^\kappa$, the entries $(XY)|_W$, and $\Omega Y^{\mathsf T}$ for every $\Omega$ supported on $W$ with entries $n_2^{O(1)}$, are computed in at most $Cn_1n_2D^{-\gamma}$ word-RAM steps (time model; deterministic). The same holds with the roles of $n_1$ and $n_2$ exchanged. The only cited input is Theorem~\ref{cp:P6:LG}.
\end{theorem}

\begin{proof}
There are finitely many families: the seven of Lemma~\ref{cp:E7} (Proposition~\ref{cp:E6}), with $\beta/\sigma=E_i$, and the square family of $\Pzero$ read as a family for \ThmDsig{} with $\sigma=1$ and $\beta=\alpha(\Pzero)$ (Lemmas~\ref{cp:E8:square}, \ref{cp:L5} and~\ref{cp:L8}). All have $\cg=0$. For fixed $(\kappa,\eps,c_0)$, Corollary~\ref{cp:E4}(a) gives for each family $F$ with $\eps\le E_F$ a pair $(\gamma_F,C_F)$ valid for all instances with $n_1\le c_0n_2^{\sigma_F}$, and Proposition~\ref{cp:E5} gives the same for $\Omega Y^{\mathsf T}$. Take $\gamma=\min\gamma_F$ and $C=\max C_F$ over these families. An instance with $n_1\le c_0n_2^s$ and $\eps\le E(s)$ is served by a family with $\sigma_F\ge s$ and $E_F\ge\eps$, since then $n_1\le c_0n_2^s\le c_0n_2^{\sigma_F}$. For $n_1>n_2$ use $(XY)|_W=((Y^{\mathsf T}X^{\mathsf T})|_{W^{\mathsf T}})^{\mathsf T}$ and the last part of the proof of Proposition~\ref{cp:E5}.
\end{proof}

The constant $\gamma$ is uniform in the shape; this is what the applications of Proposition~\ref{cp:E1} (in Propositions~\ref{cp:E2} and~\ref{cp:E16}(b)) and Corollary~\ref{cp:E13} use. The range $E=\alpha(\Pzero)$ on $(\frac34,1]$ comes from \ThmDsig{} at $\sigma=1$ directly; the bounds of Proposition~\ref{cp:E2} on rank exponents give no algorithm.

Measured against the smaller side, the inner dimension can thus go up to $n_1^{\cpv{E7Ed_b}}$ when $n_1=n_2^{1/2}$ and up to $n_1^{\cpv{E7Ed_g}}$ when $n_1\le n_2^{1/10}$, against $n_1^{\alpha(\Pzero)}$ for the square top. The unequal top does not reach ``$\alpha$ against the larger side'': it moves $\eps$ from $\alpha(\Pzero)$ towards the input-size ceiling~$1$.

\begin{proposition}[bounds on $\eps^*(\sigma)$]\label{cp:E2}
For every $\sigma\in(0,1]$,
\[
\max\bigl(\alpha,E(\sigma)\bigr)\le\eps^*(\sigma)\le\min\Bigl\{1,\ \frac{\alpha(1+\sigma)}{2\sigma}\Bigr\},
\]
$\eps^*$ is non-increasing on $(0,1]$, and $\omega=2$ implies $\eps^*\equiv1$.
\end{proposition}

\begin{proof}
\emph{$\eps^*\le1$.} The $y$-flattening of $\mat{n^\sigma}{n^\beta}n$ has rank $n^{\beta+1}$, so $\omega(\sigma,\beta,1)\ge1+\beta$, and equality with $1+\sigma$ forces $\beta\le\sigma$.
\emph{Monotonicity and $\eps^*\ge\alpha$.} Let $\sigma'<\sigma$, $\beta$ admissible at $\sigma$, and $n=m^{\sigma'/\sigma}$. Then $\mat{m^{\sigma'}}{m^{\sigma'\beta/\sigma}}m=\mat{n^\sigma}{n^\beta}n\otimes\mat11{m^{1-\sigma'/\sigma}}$, whose rank is at most $m^{(\sigma'/\sigma)(1+\sigma)+1-\sigma'/\sigma+o(1)}=m^{1+\sigma'+o(1)}$. So $\sigma'\beta/\sigma$ is admissible at $\sigma'$, i.e. $\eps^*(\sigma')\ge\eps^*(\sigma)$; at $\sigma=1$ this gives $\eps^*(\sigma')\ge\alstar(1)=\alpha$.
\emph{The symmetrisation bound.} $\mat ehl$ and $\mat lhe$ have the same rank, and $\mat ehl\otimes\mat lhe=\mat{el}{h^2}{el}$. So $\beta$ admissible at $\sigma$ gives $\omega(1,2\beta/(1+\sigma),1)=2$, hence $\alstar(\sigma)\le\alpha(1+\sigma)/2$; divide by $\sigma$.
\emph{$\omega=2$.} $\mat{n^\sigma}{n^\sigma}n=\mat{n^\sigma}{n^\sigma}{n^\sigma}\otimes\mat11{n^{1-\sigma}}$ has rank $n^{1+\sigma+o(1)}$, so $\beta=\sigma$ is admissible.
\emph{$\eps^*\ge E$.} For $\eps\le E(\sigma)$, Theorem~\ref{cp:E8} with $s=\sigma$ and $c_0=2$ gives, for every $\kappa>0$, an algorithm on the instances with $n_1=\lceil N^\sigma\rceil\le2N^\sigma$ which by Remark~\ref{cp:E4:alg} has the form required in Proposition~\ref{cp:E1}, and whose number of non-scalar multiplications is at most its running time. Proposition~\ref{cp:E1} gives $\eps\le\eps^*(\sigma)$.
\end{proof}

The two upper bounds are separate: $\alpha(1+\sigma)/(2\sigma)$ is the smaller one only for $\sigma>\alpha/(2-\alpha)$, and the chain ``$\eps^*(\sigma)\le\alpha(1+\sigma)/(2\sigma)\le1$'' would be false for small~$\sigma$. With $\alpha\ge\AlphaVXXZ$~\cite{VXXZ24} the lower bound $\eps^*\ge\alpha$ improves on $E$ for $\sigma\in(\frac34,1]$; the claimed $\alpha\ge\cpv{alpha_preprints}$~\cite{AvdBZ26,Uzn26} would improve it further. An unequal-shape value is never cheaper evidence about $\alpha$ than the square one: by the symmetrisation bound, a value $\eps$ at shape $\sigma$ certifies only $\alpha\ge2\sigma\eps/(1+\sigma)$, which for the rows of Table~\ref{cp:E7:table} is below $\alpha(\Pzero)$.

\begin{remark}[rates]\label{cp:E9}
As $x\to\infty$, Corollary~\ref{cp:E4}(b) allows every $\gamma$ below $\min\{\sup_t(t\kappa-(1+\sigma)\Lam_\sigma(t)/\beta),\ 1/\eps-1/E_i\}$ for the family of row $i$ (the third term $(1+\sigma)/\beta$ of Corollary~\ref{cp:E4}(b) exceeds $1$ and is omitted). With the profile bound of Proposition~\ref{cp:E6}, evaluated in floating point on a grid of $t$, the first term is about $\cpv{E9_half_half}$ at $\kappa=\frac12$ and $\cpv{E9_half_one}$ at $\kappa=1$ for the row $\sigma=\frac12$, and about $\cpv{E9_tenth_half}$ and $\cpv{E9_tenth_one}$ for the row $\sigma=\frac1{10}$ (computed, not certified). These rates apply to $\eps<E_i$, with $\gamma$ also below $1/\eps-1/E_i$; they are of the order of the square case.
\end{remark}

\subsection{The partially sparse triangle hypothesis of ABFK}\label{cp:sec:unequal-abfk}

Abboud, Bringmann, Fischer and Künnemann~\cite{ABFK24} define the problem PS-AE-Triangle$(x,y,z,m)$: given a tripartite graph $G=(X,Y,Z,E)$ with $\abs X\le x$, $\abs Y\le y$, $\abs Z\le z$ and $\abs{E\cap(Y\times Z)}\le m$, decide for each edge in $(X\times Z)\cap E$ whether it lies in a triangle~\cite[Definition~1.8]{ABFK24}. Only the $Y$--$Z$ edges are bounded by $m$. Their Hypothesis~1.9 reads: ``For all $a,b,c\ge0$, the PS-AE-Triangle$(m^a,m^b,m^c,m)$ problem cannot be solved in time $O(m^{\min\{1+a,\omega(a,b,c)\}-\epsilon})$, for any $\epsilon>0$''~\cite[Hypothesis~1.9]{ABFK24}. Their model is the word RAM, $\omega(a,b,c)$ is the exponent of arithmetic circuits over fields of characteristic~$0$, which over $\Q$ is the rank exponent of Section~\ref{cp:sec:laser}, and they use $\max\{a+b,a+c,b+c\}\le\omega(a,b,c)\le a+b+c$~\cite[Fact~2.2]{ABFK24}. The hypothesis does not say whether it concerns deterministic or randomized algorithms; all algorithms below are deterministic.

We say that Hypothesis~1.9 \emph{fails at} $(a,b,c)$ if its instance at this triple is false, that is, if PS-AE-Triangle$(m^a,m^b,m^c,m)$ is solved in time $O(m^{\min\{1+a,\omega(a,b,c)\}-\delta})$ for some $\delta>0$. Since the hypothesis is one statement quantified over all triples, one such triple refutes it, and it is already refuted at triples in the range of~\cite[Theorem~\Vref{m:thm:general}]{v2} (\cite[Corollary~\Vref{m:cor:abfk}]{v2}); what matters below is at which triples it fails. For a triple with $a>0$ and $b+c>1$ put
\[
\kappa:=\frac{b+c-1}a,\qquad\eps:=\frac a{\min(b,c)},\qquad s:=\frac{\min(b,c)}{\max(b,c)} .
\]

\begin{lemma}[the reduction]\label{cp:E10}
Let $a>0$, $b+c>1$ and $a+\max(b,c)<\min\{1+a,b+c\}$, and let $\kappa$ and $\eps$ be as above. Suppose that for some $\eps'>\eps$ and some $\gamma>\max(0,\kappa-1)$ there is a deterministic algorithm that, for all $n_1,n_2\ge1$ with $\min(n_1,n_2)\le2\max(n_1,n_2)^s$, where $s=\min(b,c)/\max(b,c)$, and all $2\le D\le\min(n_1,n_2)^{\eps'}$, every $Y\in\{0,1\}^{D\times n_2}$ and every $\Omega\in\{0,1\}^{n_1\times n_2}$ with support $W$, $\abs W\le n_1n_2/D^\kappa$, computes $\Omega Y^{\mathsf T}$ in $O(n_1n_2D^{-\gamma})$ word-RAM steps. Then Hypothesis~1.9 fails at $(a,b,c)$. The proof uses only $\omega(a,b,c)\ge b+c$; no value of $\alpha$ or $\omega$ enters.
\end{lemma}

\begin{proof}
Given an instance, add isolated vertices so that $\abs X=D:=\lfloor m^a\rfloor$, $\abs Y=n_2:=\lfloor m^b\rfloor$ and $\abs Z=n_1:=\lfloor m^c\rfloor$; this changes neither the triangles nor the number of $Y$--$Z$ edges. The padded sizes satisfy $\min(n_1,n_2)\le2\max(n_1,n_2)^s$: if $c\le b$, then $n_1\le m^c=(m^b)^s\le(2n_2)^s\le2n_2^s$, and likewise if $c>b$. Let $\Omega\in\{0,1\}^{n_1\times n_2}$ be the $Z$--$Y$ biadjacency matrix (at most $m$ ones) and $Z'\in\{0,1\}^{n_2\times D}$ the $Y$--$X$ one, and put $Y:=Z'^{\mathsf T}$. Then $(\Omega Y^{\mathsf T})[z,x]$ is the number of paths $z$--$y$--$x$, and a query edge $xz$ lies in a triangle if and only if this number is positive. The condition $a+\max(b,c)<\min\{1+a,b+c\}$ forces $a<\min(b,c)$, so $b,c>0$. For large $m$, $n_1n_2\ge m^{b+c}/4$ and $D^\kappa\le m^{a\kappa}=m^{b+c-1}$, so $\abs W\le m\le4n_1n_2/D^\kappa$; split $W$ into at most four parts of size at most $n_1n_2/D^\kappa$ and add the four products. Also $D\le m^a\le(2\min(n_1,n_2))^{\eps}\le\min(n_1,n_2)^{\eps'}$ for large $m$. So $\Omega Y^{\mathsf T}$ costs $O(n_1n_2D^{-\gamma})=O(m^{b+c-a\gamma})$. Reading the graph costs $O(m+m^{a+b}+m^{a+c})$ and answering the queries $O(m^{a+c})$, so the total is $O(m^{\max\{b+c-a\gamma,\,1,\,a+\max(b,c)\}})$. Here $b+c-a\gamma<b+c$ as $\gamma>0$, and $b+c-a\gamma<1+a$ as $\gamma>\kappa-1$; $1<\min\{1+a,b+c\}$; and $a+\max(b,c)<\min\{1+a,b+c\}$ by assumption. By~\cite[Fact~2.2]{ABFK24}, $\min\{1+a,\omega(a,b,c)\}\ge\min\{1+a,b+c\}$, so the hypothesis fails at $(a,b,c)$ with some $\delta>0$. Every step is deterministic, on a word RAM with $O(\log m)$-bit words.
\end{proof}

The same reduction proves~\cite[Corollary~\Vref{m:cor:abfk}]{v2}. For $\kappa\le1$, $\min\{1+a,b+c\}=b+c$, and the condition $a+\max(b,c)<b+c$ is $a<\min(b,c)$, that is $\eps<1$; the orientation (inner dimension $D=\abs X$, outer sides $\abs Z$ and $\abs Y$) is the only one that works, since in the other the queried part is unbounded.

Let $F_\square(\kappa)$ be the supremum of the $\eps$ for which Theorem~\ref{cp:T7}, applied to the input family of $\Pzero$ (with part (a) or (c)), gives $\Hypk\kappa(\eps,\gamma)$ with some $\gamma>0$. Then $F_\square(\kappa)\ge\alpha(\Pzero)$ for every $\kappa>0$, and at $\kappa=\cpv{B3_ka}$ and $\kappa=\cpv{B3_kb}$ the certified values of Theorem~\ref{cp:B3}, $\cpv{B3_fa}$ and $\cpv{B3_fb}$, are values of Theorem~\ref{cp:T7}(c) with the certified grid of Table~\ref{cp:tab:B2}: Table~\ref{cp:tab:B3} lists their bottom as direct, that is inner products, and for $\kappa\le1$ the direct-bottom range $\rdir(\kappa,t)$ of Theorem~\ref{cp:B1}(a) equals $\alpha+(1-\alpha)/(1+\xmin(t))$ of Theorem~\ref{cp:T7}(c) at the same tilt, with $L_t$ in place of $\Lamp(t)$, so $F_\square(\frac12)\ge\cpv{B3_fa}$ and $F_\square(1)\ge\cpv{B3_fb}$.

\begin{proposition}[the square region]\label{cp:E11}
Hypothesis~1.9 fails at every $(a,b,c)$ with $a>0$, $b+c>1$, $\kappa\in(0,1]$ and $a/\min(b,c)<F_\square(\kappa)$; in particular at every such triple with $a/\min(b,c)<\alpha(\Pzero)$, and at $\kappa=1$ at every such triple with $a/\min(b,c)<\cpv{B3_fb}$.
\end{proposition}

\begin{proof}
Let $\eps<\eps'<F_\square(\kappa)$, $\eps'$ rational, with $\eps=a/\min(b,c)<1$. Theorem~\ref{cp:T7} gives \AlgD{} with a top algorithm $B$ such that on every $n\times D\times n$ instance with $D\le n^{\eps'}$ and $n\le2D^{1/\eps'}$ the cost is $C_0(n,D)$ (input, encodings, tables, bucket array) plus $A_1(n,D)\abs{W'}$ (bucketing, output) plus $A_t(n,D)\abs{W'}^t$ (queries and decoding, by Lemma~\ref{cp:T2}(a), for every $W'$), with coefficients that depend on $n$, $D$ and the fixed parameters only, and it is $O(n^2D^{-\gamma_0})$ at $\abs{W'}=n^2D^{-\kappa}$. Proposition~\ref{cp:E5} transposes it at an extra cost $O(\abs{W'}+nD)$, which keeps this form, since $nD\le n^2D^{1-1/\eps'}$. Lemma~\ref{cp:A10:tiling}, in its rectangular and transposed form, gives $\Omega Y^{\mathsf T}$ in $O(n_1n_2D^{-\gamma})$ with $\gamma>0$ for all $n_1,n_2$ and $D\le\min(n_1,n_2)^{\eps'}$. Since $\kappa\le1$, Lemma~\ref{cp:E10} applies.
\end{proof}

For $\kappa$ slightly above $1$, the total time of \AlgD{} stays below $\abs WD$ on a strip $\kappa\in(1,\cpv{A10_strip})$, with values computed in exact arithmetic on the certified grid; this strip is computed and argued, not proved, and no statement below uses it. It is a different statement from Corollary~\ref{cp:B7}(a): the strip needs a saving $\gamma>\kappa-1$ of the whole computation, the plateau only $\gamma_{\mathrm{pre}}>0$ and a saving in phase~2.

\begin{proposition}[the region with a fast rectangular bottom]\label{cp:E12}
Let $e\in(\alpha(\Pzero),1)$ and $g\ge0$ with $\omega(1,e,1)\le2+g$, and let $F(\kappa;e,g)$ be as in Theorem~\ref{cp:A10}. Then Hypothesis~1.9 fails at every $(a,b,c)$ with $a>0$, $b+c>1$, $\kappa\in(0,1]$ and $a/\min(b,c)<F(\kappa;e,g)$. In particular it fails at every triple with $\kappa=1$ and $a/\min(b,c)<\alpha$, whatever the value of $\alpha$ is, and, with $\omega(1,\cpv{A11_tab1_e_one},1)\le\cpv{A11_tab1_om_one}$~\cite[Table~1]{VXXZ24}, at every triple with $\kappa=1$ and $a/\min(b,c)\le\cpv{A11_tab1_one}$.
\end{proposition}

\begin{proof}
As for Proposition~\ref{cp:E11}, with the per-tile algorithm of the proof of Theorem~\ref{cp:A10} in place of \AlgD: its cost in a tile has the form of Lemma~\ref{cp:A10:tiling}, with coefficients that depend on $n$, $D$ and the fixed parameters only, and Proposition~\ref{cp:E5} transposes it at an extra cost $O(\abs{W'}+nD)$, which keeps this form. For $\eps<\eps'<F(\kappa;e,g)$ (and $\eps'<1$, as $F\le e<1$) Lemma~\ref{cp:A10:tiling} gives $\Omega Y^{\mathsf T}$ in $O(n_1n_2D^{-\gamma})$ for all $D\le\min(n_1,n_2)^{\eps'}$, and Lemma~\ref{cp:E10} applies. The case $\kappa=1$, $a/\min(b,c)<\alpha$ follows as in the last part of the proof of Theorem~\ref{cp:A10}, and the last statement from Corollary~\ref{cp:A11}(c), whose proof shows $\cpv{A11_tab1_one}<F(1;\cpv{A11_tab1_e_one},\cpv{A11_tab1_g_one})$.
\end{proof}

\begin{corollary}[unequal sides]\label{cp:E13}
Let $a>0$, $b+c>1$, $\kappa\le1$ and $a/\min(b,c)<E(s)$, with $s=\min(b,c)/\max(b,c)$ and $E$ the step function of Theorem~\ref{cp:E8}. Then PS-AE-Triangle$(m^a,m^b,m^c,m)$ is solved deterministically in time $O(m^{b+c-\delta})$ for some $\delta>0$, and Hypothesis~1.9 fails at $(a,b,c)$. Equivalently, the condition is $a<\min(b,c)E(s)$, with $E$ up to $E_7=\cpv{E7E_g}$.
\end{corollary}

\begin{proof}
Pad as in Lemma~\ref{cp:E10}, with $n_1=\lfloor m^c\rfloor$ and $n_2=\lfloor m^b\rfloor$. If $c\le b$, then $n_1\le m^c=(m^b)^{c/b}\le(2n_2)^s\le2n_2^s$; if $c>b$, the same holds with the roles of $n_1$ and $n_2$ exchanged. Take $\eps'$ with $a/\min(b,c)<\eps'<E(s)$. Theorem~\ref{cp:E8} with $c_0=2$, in its transposed form, gives $\Omega Y^{\mathsf T}$ in $O(n_1n_2D^{-\gamma})$ for every $D\le\min(n_1,n_2)^{\eps'}$, with $\gamma>0=\max(0,\kappa-1)$. Since $E(s)\le E_7=\cpv{E7E_g}<1$, $a<\min(b,c)$, and Lemma~\ref{cp:E10} applies; its proof gives the time bound, because for $\kappa\le1$ the target is $m^{b+c}$.
\end{proof}

\begin{remark}[the ceiling of the method]\label{cp:E14}
The route of Lemma~\ref{cp:E10} refutes the hypothesis only where $a<\min(b,c)$ (and $b+c>1$): it reads an input of size $m^{a+\max(b,c)}$ and aims at the bound $m^{\min\{1+a,b+c\}}\le m^{b+c}$. This is the same ceiling as $\eps^*\le1$ in Proposition~\ref{cp:E2}, so no unequal top, of any level, moves the refuted wedge past $a=\min(b,c)$; it can only move $E(s)$ towards~$1$. This is a limit of the method, not a statement about the hypothesis: for $a\ge\min(b,c)$ with $\max(b,c)<1$, and likewise for $b+c\le1$, whether PS-AE-Triangle can be solved faster than the hypothesis allows is open.
\end{remark}

\begin{remark}[what Corollary~\ref{cp:E13} adds]\label{cp:E15}
There are four refuted regions, all deterministic: that of~\cite{v2}, $a/\min(b,c)<\es$ for $\kappa\le1$ (and a positive bound for every $\kappa>1$; \cite[Corollary~\Vref{m:cor:abfk}]{v2}); the square region of Proposition~\ref{cp:E11}; the region of Proposition~\ref{cp:E12}; and that of Corollary~\ref{cp:E13}. The first three do not depend on the shape $s$. For $s\in(\frac34,1]$, $E(s)=\alpha(\Pzero)\le F_\square(\kappa)$, so Corollary~\ref{cp:E13} adds nothing there, and it adds nothing for $\kappa>1$. What it adds is the set of triples with $\kappa\le1$, $s\le\frac34$ and $a/\min(b,c)<E(s)$ that lie in none of the other three regions. Table~\ref{cp:E15:table} shows samples. Its square column uses $\alpha(\Pzero)$ only: with the larger $F_\square(\kappa)$ near $\kappa=1$, and with the region of Proposition~\ref{cp:E12}, the other regions' upper ends are larger, so the table overstates what Corollary~\ref{cp:E13} adds at these samples. The extremes of the wedge do not move: on the step of $E$ with right end $s$, $\kappa\le1$ forces $\min(b,c)<1/(1+1/s-E(s))$ and $a<E(s)/(1+1/s-E(s))$, and both are largest at $s=1$; so every triple of the region of Corollary~\ref{cp:E13} has $a<\alpha(\Pzero)/(2-\alpha(\Pzero))<\cpv{E15_a}$ and $\min(b,c)<1/(2-\alpha(\Pzero))<\cpv{E15_bc}$, the same extremes as for the square region at $\alpha(\Pzero)$.
\end{remark}

\begin{table}[ht]\centering\small
\begin{tabular}{llllll}
$(b,c)$ & $s$ & $E(s)$ & Cor.~\ref{cp:E13} & square, at $\alpha(\Pzero)$ & \cite{v2}, at $\es$\\\hline
$\cpv{E15bc_1}$ & $\cpv{E15s_1}$ & $\cpv{E7Ed_f}$ & $(\cpv{E15lo_1},\cpv{E15r_1})$ & $(\cpv{E15lo_1},\cpv{E15q_1})$ & $(\cpv{E15lo_1},\cpv{E15v_1})$\\
$\cpv{E15bc_2}$ & $\cpv{E15s_2}$ & $\cpv{E7Ed_g}$ & $(\cpv{E15lo_2},\cpv{E15r_2})$ & $(\cpv{E15lo_2},\cpv{E15q_2})$ & empty\\
$\cpv{E15bc_3}$ & $\cpv{E15s_3}$ & $\cpv{E7Ed_d}$ & $(\cpv{E15lo_3},\cpv{E15r_3})$ & $(\cpv{E15lo_3},\cpv{E15q_3})$ & empty\\
$\cpv{E15bc_4}$ & $\cpv{E15s_4}$ & $\cpv{E7Ed_b}$ & $(\cpv{E15lo_4},\cpv{E15r_4})$ & $(\cpv{E15lo_4},\cpv{E15q_4})$ & $(\cpv{E15lo_4},\cpv{E15v_4})$\\
$\cpv{E15bc_5}$ & $\cpv{E15s_5}$ & $\cpv{E7Ed_a}$ & $(\cpv{E15lo_5},\cpv{E15r_5})$ & $(\cpv{E15lo_5},\cpv{E15q_5})$ & $(\cpv{E15lo_5},\cpv{E15v_5})$\\
$\cpv{E15bc_6}$ & $\cpv{E15s_6}$ & $\cpv{E7Ed_f}$ & $(\cpv{E15lo_6},\cpv{E15r_6})$ & empty & empty
\end{tabular}
\caption{Refuted intervals of $a$ at $\kappa\le1$ for fixed $(b,c)$: the lower end is $b+c-1$ (from $\kappa\le1$), the upper ends are $\min(b,c)$ times $E(s)$, $\alpha(\Pzero)$ and $\es$; the column $E(s)$ and the upper ends are truncated to six decimals. The square column uses $\alpha(\Pzero)$, not $F_\square(\kappa)$ and not the region of Proposition~\ref{cp:E12}.}\label{cp:E15:table}
\end{table}

ABFK's exponent $\sigma(r)$, for $r\in[0,2]$, is the unique solution $\sigma$ of $\omega(\sigma-1,2-\sigma,1+r-\sigma)=\sigma$~\cite[Definition~1.5, Lemma~4.4]{ABFK24}, and their Theorem~1.10 states that sparse Boolean matrix multiplication with input size $m_{\mathrm{in}}$ and output size $m_{\mathrm{in}}^r$ cannot be solved in time $O(m_{\mathrm{in}}^{\sigma(r)-\epsilon})$ unless Hypothesis~1.9 fails. Its proof uses the hypothesis only at the triple
\[
T(r):=\bigl(\sigma(r)-1,\ 2-\sigma(r),\ 1+r-\sigma(r)\bigr).
\]
Let $\mu$ be the solution of $\omega(\mu,1,1)=1+2\mu$; Le Gall and Urrutia~\cite{LGU18} show $\mu<\cpv{E16_mu}$, and~\cite{VXXZ24} shows $\mu<\cpv{mu_vxxz}$; either bound works below.

\begin{proposition}[ABFK's Theorem~1.10 triples]\label{cp:E16}
\begin{enumerate}[label=(\alph*)]
\item For every $r\in(0,\cpv{E16_r_one}]$, $T(r)$ lies in the region of Corollary~\ref{cp:E13}, whatever the value of $\sigma(r)$; so the instance of Hypothesis~1.9 at $T(r)$ that the proof of~\cite[Theorem~1.10]{ABFK24} uses is false.
\item For every $r\in(0,\cpv{E16_r_two}]$, $\sigma(r)=1+r/3$, and $T(r)=(r/3,1-r/3,2r/3)$ lies in the region of Corollary~\ref{cp:E13}; so the same holds for these $r$.
\item For every $r\in(\cpv{E16_r_out},2]$, $T(r)$ lies outside the region of Corollary~\ref{cp:E13}.
\end{enumerate}
So for $r\in(0,\cpv{E16_r_two}]$ Theorem~1.10's lower bound loses its support; Theorem~1.10 itself remains valid as an implication.
\end{proposition}

\begin{proof}
Write $\sigma=\sigma(r)$ and $T(r)=(a,b,c)$. Four facts hold for every $r\in(0,2]$.
(1) $\omega(a,b,c)\ge b+c=3+r-2\sigma$~\cite[Fact~2.2]{ABFK24} gives $\sigma\ge1+r/3$.
(2) Hence $\kappa=(b+c-1)/a=(2+r-2\sigma)/(\sigma-1)\le1$; also $a>0$, and $b+c>1$ if and only if $\sigma<1+r/2$.
(3) For $r\le1$, $c\le b$, so $\eps=a/c=(\sigma-1)/(1+r-\sigma)$, which is increasing in $\sigma$ and equals $\frac12$ at $\sigma=1+r/3$; the shape $s=c/b=(1+r-\sigma)/(2-\sigma)$ has $ds/d\sigma=(r-1)/(2-\sigma)^2<0$ for $r<1$, so $s\le2r/(3-r)$.
(4) $\sigma$ is convex~\cite[Lemma~4.5]{ABFK24}, $\sigma(0)=1$~\cite[Lemma~4.6]{ABFK24} and $\sigma(1)=1+\mu/(1+\mu)$~\cite[Lemma~4.7]{ABFK24}, so $\sigma(r)\le1+r\mu/(1+\mu)$ on $[0,1]$. At this bound $\eps=\mu$, and $\eps$ is increasing in $\sigma$, so $\eps\le\mu<\cpv{E16_mu}$; and since $\mu<1$, $\sigma<1+r/2$.
(a) For $r\le\frac15$, $s\le2r/(3-r)\le\frac17$, so $E(s)\ge E(\frac17)=\cpv{E7E_f}>\cpv{E7Ed_f}>\cpv{E16_mu}>\eps$. With $\kappa\le1$, $a>0$ and $b+c>1$ from (2) and (4), $T(r)$ is in the region.
(b) At $\sigma_0=1+r/3$ the triple is $(r/3,1-r/3,2r/3)$, with $\eps=\frac12$ and $s_0=2r/(3-r)\le\frac15$ for $r\le\frac3{11}$ (with equality at $r=\frac3{11}$). Since $E(s_0)\ge E(\frac15)=\cpv{E7E_e}>\frac12$, Theorem~\ref{cp:E8} with Remark~\ref{cp:E4:alg} and Proposition~\ref{cp:E1} (at shape $s_0$, with $c_0=2$) gives $\omega(s_0,s_0/2,1)=1+s_0$. By homogeneity, with the factor $1-r/3$, and symmetry, $\omega(r/3,1-r/3,2r/3)=(1-r/3)(1+s_0)=1+r/3=\sigma_0$. So $\sigma_0$ solves ABFK's equation, and by uniqueness~\cite[Lemma~4.4]{ABFK24}, $\sigma(r)=\sigma_0$. Then $T(r)$ has $\kappa=1$ and $\eps=\frac12<E(s_0)$, so it lies in the region.
(c) For $r<1$: by (1) and (4), $c=1+r-\sigma\ge r/(1+\mu)$ and $b=2-\sigma\le1-r/3$, so $s=c/b\ge3r/((1+\mu)(3-r))$, which exceeds $\frac15$ for $r>3(1+\mu)/(16+\mu)$; this threshold increases in $\mu$ and is below $\cpv{E16_r_out}$ for $\mu<\cpv{E16_mu}$. Then $E(s)\le E_4=\cpv{E7E_d}<\frac12\le\eps$ by (3). For $r\ge1$, suppose that $T(r)$ were inside. Then $b+c>1$, so $\sigma<1+r/2\le2$ by (2) and $b>0$; $s=b/c$ and $\eps=a/b$, and by (1), $2\eps s=2a/c\ge1$, so $\eps\ge\frac12$ as $s\le1$. Now $\eps<E(s)\le E_7=\cpv{E7E_g}$, hence $s>1/(2E_7)=9000/10703>\cpv{E16_s}>\frac34$, hence $E(s)=\alpha(\Pzero)<\frac12\le\eps$, a contradiction.
\end{proof}

On $(\cpv{E16_r_two},\cpv{E16_r_out}]$ the answer depends on the unknown $\sigma(r)$. The triples $T(r)$ have $a+b=1$, hence $\eps(\kappa+1)=c/\min(b,c)\ge1$; so they lie in a region of the form $\{a/\min(b,c)<F(\kappa)\}$, $\kappa\le1$, only where $F(\kappa)>\frac12$. This is why the sentence ``ABFK's lower bounds keep their stated support'' is true for the region of~\cite{v2} and for the square region at $\alpha(\Pzero)$, but must not be carried over to Corollary~\ref{cp:E13}, where $E(s)>\frac12$ for $s\le\frac15$.

\begin{remark}[the value of $\sigma(r)$ for small $r$]\label{cp:E17}
Theorem~\ref{cp:E8} and Proposition~\ref{cp:E1}, which use no bound on $\omega$ or $\alpha$, give $\sigma(r)=1+r/3$ for $r\in[0,\cpv{E16_r_two}]$ (proof of Proposition~\ref{cp:E16}(b); $\sigma(0)=1$ by~\cite[Lemma~4.6]{ABFK24}). This is the value that~\cite[Lemma~4.10]{ABFK24} gives under $\omega=2$. Under the claimed bound $\alpha\ge\cpv{alpha_preprints}$~\cite{AvdBZ26,Uzn26}, which gives $\mu=\frac12$, the same follows for $r\in[0,1]$ from the lower bound $\sigma(r)\ge1+r/3$ (fact (1) in the proof of Proposition~\ref{cp:E16}) and~\cite[Lemma~4.9]{ABFK24}, whose first term $1+r\mu/(1+\mu)=1+r/3$ is the largest of its three terms on $[0,1]$.
\end{remark}

\begin{remark}[ABFK's balanced triple at the boundary]\label{cp:E18}
Suppose $\alpha\ge\frac12$, as claimed in~\cite{AvdBZ26,Uzn26}. Then $\omega(\frac12,1,1)=2$, so $\mu=\frac12$, and the triple of~\cite[Theorem~1.3]{ABFK24} and $T(1)$ are both $(\frac13,\frac23,\frac23)$, with $\kappa=1$ and $\eps=\frac12$. By Proposition~\ref{cp:E12} with $g=0$, Hypothesis~1.9 fails at every triple with $\kappa=1$ and $\eps<\frac12$; this one lies exactly on the boundary. By the closed form of Theorem~\ref{cp:A10}, $F(1;e,g)>\frac12$ if and only if $g<\bar C(e-\frac12)/(\frac12-\alpha(\Pzero))$, and with the certified lower bounds of $\alpha(\Pzero)$ and of $\bar C$ (at $t=\cpv{A11_t_one}$ on the grid) the factor $\bar C/(\frac12-\alpha(\Pzero))$ exceeds $\cpv{E18_g}$. So a bound $\omega(1,e,1)\le2+g$ with $e\in(\frac12,1)$ and $g<\cpv{E18_g}(e-\frac12)$ would put $(\frac13,\frac23,\frac23)$ inside the region of Proposition~\ref{cp:E12}. Every AE-Triangle$(n^{1/2},n,n,n^{3/2})$ instance is a PS-AE-Triangle$(m^{1/3},m^{2/3},m^{2/3},m)$ instance with $m=n^{3/2}$, since all its edges number at most $m$; so the resulting algorithm solves AE-Triangle$(n^\mu,n,n,n^{1+\mu})$ in time $n^{1+2\mu-\Omega(1)}$, and by~\cite[Theorem~1.3]{ABFK24} sparse Boolean matrix multiplication would then run in time $O(m^{4/3-\epsilon})$ for some $\epsilon>0$. This is a conditional statement: we do not know whether a bound of that kind above $\frac12$ is published.
\end{remark}

\section{The degree profile}\label{cp:sec:profile}

The savings of the composites that share queries depend on the top through one function, the distribution of the supports of its terms, measured by a log-moment. This section states what that function determines, in three regimes. A batch of wanted entries saves a value of its Legendre function (Section~\ref{cp:sec:F:batch}). A data structure with a fixed query exponent follows a frontier built from the same function (Section~\ref{cp:sec:F:frontier}). A range is not a value of the Legendre function: it is set by a second parameter, the encoding cost, at a zero of it (Remark~\ref{cp:F8}). For the tile of~\cite{v2} and for tensor powers of one fixed decomposition the profile is an exact log-moment, and the batch statements below are identities (the frontier of Theorem~\ref{cp:F6} is a construction whose optimality is not claimed); for the level-1 family of Section~\ref{cp:sec:level1}, \ThmD{}'s profile is only bounded, and the statements give lower bounds on the savings (Section~\ref{cp:sec:F:wording}). Section~\ref{cp:sec:F:small} treats small budgets.

\subsection{Profiles and their Legendre function}\label{cp:sec:F:def}

\begin{definition}[leaf-degree profile, Legendre function]\label{cp:F1}
Let $\mathcal A$ be an exact decomposition over $\Q$ of a tensor whose $z$-variables are its outputs, and $R$ a set of \emph{read} outputs, $M=\abs R$. The \emph{degree} of a term $\tau$ is $\deg(\tau)=\abs{\supp w_\tau\cap R}$; terms of degree $0$ are dropped. For a sequence $(\mathcal A_n,R_n)$ of such pairs whose inner dimensions $D_n$ tend to infinity, the \emph{leaf-degree profile} in units of $\ln D$ is
\[
\Lleaf(t)=\lim_{n\to\infty}\frac{\ln\bigl(\frac1{M_n}\sum_\tau\deg(\tau)^t\bigr)}{\ln D_n}\qquad(t\in[0,1]),
\]
where the limit exists. For a function $\Lambda$ on $[0,1]$ its \emph{Legendre function} is $\Leg(\beta)=\sup_{0\le t\le1}(t\beta-\Lambda(t))$ for $\beta\ge0$.
\end{definition}

For the powers $B^{\otimes n}$ of one exact decomposition $B$ of $\mat kbk$, read on all outputs, no limit is needed: degrees multiply, so $\frac1M\sum_\tau\deg(\tau)^t=(S(t)/k^2)^n$ with $S(t)=\sum_\rho\deg(\rho)^t$, and with $D=b^n$
\[
\Lleaf(t)=\frac{\lambda(t)}{\alB},\qquad\lambda(t)=\log_kS(t)-2,
\]
at every $n$. The parameter $\alB$ only converts units: in units of $\ln K$, $K=k^n$, the Legendre function is $\alB\Leg(\beta/\alB)$.

\begin{lemma}\label{cp:F1:basic}
\begin{enumerate}[label=(\alph*)]
\item For every bounded $\Lambda$, $\Leg$ is convex and nondecreasing, and $\Leg(\beta')-\Leg(\beta)\le\beta'-\beta$ for $0\le\beta\le\beta'$.
\item If $\Lambda$ is nondecreasing, then $-\Lambda(0)\le\Leg(\beta)\le\beta-\Lambda(0)$.
\item $\Lleaf$ is convex and nondecreasing where it exists, and $\Lleaf(0)\ge0$ for decompositions of matrix products read on all outputs.
\end{enumerate}
\end{lemma}

\begin{proof}
(a) $\Leg$ is a supremum of affine functions of $\beta$ with slopes $t\in[0,1]$; so it is convex, nondecreasing, and $t\beta'-\Lambda(t)\le t\beta-\Lambda(t)+(\beta'-\beta)$ gives the increment bound. (b) $t=0$ gives the lower bound; $\Lambda(t)\ge\Lambda(0)$ and $t\le1$ give the upper. (c) $t\mapsto\ln\sum_\tau e^{t\ln\deg(\tau)}$ is convex, and nondecreasing since $\deg\ge1$; pointwise limits keep both properties. A decomposition of a matrix product has at least as many terms of nonzero $z$-form as outputs (Lemma~\ref{cp:T1:basic}(a)), so the value at $t=0$ is $\log_D(\#\text{terms}/M)\ge0$.
\end{proof}

\ThmD's profile is a different function, built from class maxima, and it bounds the leaf-degree profile from above.

\begin{remark}[class maxima against leaf degrees]\label{cp:F1:thmD}
Let $B$ be an exact top algorithm for $\mat kbk$ with a partition, read on all outputs, and $S_l=\abs{\supp w_l}$. Inside a class, $\sum_{l\in P_c}S_l=\sum_o\abs{\Cov(o)\cap P_c}\le k^2\kappa_c$, so Hölder's inequality gives $\sum_{l\in P_c}S_l^t\le\abs{P_c}^{1-t}(k^2\kappa_c)^t$, and summing over the classes, $\log_k\sum_lS_l^t-2\le\lamB(t)$. In the same way $\mu_M(t)\ge\log_{a_M}\bigl(\frac1{s_Ma_M^2}\sum_lS_l^t\bigr)$ for an input family. Equality holds when all terms of a class have the same support size and every output has the same class coverage. So every saving computed from \ThmD's profile is at most the value the leaf-degree profile would give.
\end{remark}

\subsection{Batches}\label{cp:sec:F:batch}

The first term of the saving of \ThmDk{} at density $D^{-\kappa}$ is the Legendre function at $\kappa$, as long as the best tilt lies where \ThmDk's profile $\Lamp$ equals the family's profile $\Lam$.

\begin{proposition}[the batch value of \ThmDk]\label{cp:F2}
Let $(A_M)$ satisfy \Hcond1--\Hcond4 with $\alpha\in(0,1)$, profile $\Lam$ and $\cstar$ as in Definition~\ref{cp:T1}, let $\tpad=2/(\cstar+2)$, and let $\hat\Lambda$ be a continuous convex function on $[0,1]$ with $\Lam\le\hat\Lambda$ on $(0,1]$. Put $\hat\Lambda_D=\hat\Lambda/\alpha$ and let $\hat\Leg$ be its Legendre function. Then for every $\kappa>0$ the first term of the minimum in Theorem~\ref{cp:T7}(b) satisfies
\[
\sup_{0<t\le1}\Bigl(t\kappa-\frac{\Lamp(t)}\alpha\Bigr)\ \ge\ \sup_{0\le t\le\tpad}\bigl(t\kappa-\hat\Lambda_D(t)\bigr),
\]
and the right side equals $\hat\Leg(\kappa)$ whenever $\kappa\le\hat\Lambda_D'(\tpad)$ (the right derivative). So for $0<\eps<\alpha$ and such $\kappa$, $\Hypk\kappa(\eps,\gamma)$ holds (time model) for every $0<\gamma<\min\{\hat\Leg(\kappa),\,1/\eps-1/\alpha,\,2/\alpha\}$.
\end{proposition}

\begin{proof}
On $(0,\tpad]$, $\Lamp=\Lam\le\hat\Lambda$ (Corollary~\ref{cp:T5}(i)), which gives the inequality; the value at $t=0$ is the limit $t\to0+$, by continuity. The function $\varphi(t)=t\kappa-\hat\Lambda_D(t)$ is concave. If its right derivative at $\tpad$, $\kappa-\hat\Lambda_D'(\tpad)$, is at most $0$, then $\varphi(t)\le\varphi(\tpad)$ for $t\ge\tpad$, so its supremum over $[0,1]$ is attained on $[0,\tpad]$. The last sentence is Theorem~\ref{cp:T7}(b).
\end{proof}

For the level-1 family at $\Pzero$ take $\hat\Lambda=\Lamcert$ of Lemma~\ref{cp:G1}, which is continuous and convex, as a maximum over a compact set of functions affine in $t$. There $\tpad>\cpv{tpadP0_lo}$ (Lemma~\ref{cp:L11}), so the identity of Proposition~\ref{cp:F2} holds for every $\kappa\le(\Lamcert/\alpha)'(\tpad)$. The derivative of a convex function does not decrease, and the right derivative of $\Lamcert/\alpha$ at $\cpv{tcapP0_lo}<\tpad$ is about $\cpv{F2_kappa}$ (computed by finite differences, not certified); so by this computation the identity covers every $\kappa$ at which Sections~\ref{cp:sec:gamma} and~\ref{cp:sec:batch} print a value of this saving, all of which are at most $1$. The printed level-1 savings are values of $\hat\Leg$ computed from the upper bound $\Lamcert$ (or from the certified bounds $M_t$ of Table~\ref{cp:tab:B2}) and not from $\Lam$ itself: they are lower bounds on what \ThmDk{} saves, and \ThmDk{} itself is an upper bound on the work. Whether the work of \AlgD{} on the level-1 family is governed exactly by the true profile is open; for tensor powers of a fixed base it is:

\begin{proposition}[exactness for tensor powers]\label{cp:F3}
Let $B$ be an exact decomposition over $\Q$ of $\mat kbk$ whose terms all have nonzero $z$-form, read on all outputs, with $\lambda$ as above, and let $\kappa>0$. For $n\ge1$ and $K=k^n$ let $W$ be a random set of outputs of $B^{\otimes n}$ containing each output independently with probability $p=K^{-\kappa}$, and let $N_W$ be the number of terms of $B^{\otimes n}$ whose support meets $W$. Then
\[
\mathbb E\,N_W=K^{\,2-\Leg_K(\kappa)+o(1)}\quad(n\to\infty),\qquad\Leg_K(\kappa)=\sup_{0\le t\le1}\bigl(t\kappa-\lambda(t)\bigr).
\]
\end{proposition}

\begin{proof}
$\mathbb E N_W=\sum_\tau\bigl(1-(1-p)^{\deg\tau}\bigr)$.
\emph{Upper bound.} $1-(1-p)^d\le\min(1,pd)\le(pd)^t$ for $t\in[0,1]$, and $\sum_\tau(p\deg\tau)^t=p^tS(t)^n=K^{2+\lambda(t)-t\kappa}$; take the infimum over $t$.
\emph{Lower bound.} $1-(1-p)^d\ge1-e^{-pd}\ge(1-1/e)\min(1,pd)$. Let $d_1<\dots<d_r$ be the distinct degrees of the terms of $B$ and $m_i$ their multiplicities. For $\nu\in\N^r$ with $\sum_i\nu_i=n$, exactly $\binom n\nu\prod_im_i^{\nu_i}$ terms of $B^{\otimes n}$ have $\nu_i$ coordinates of degree $d_i$, and each has degree $\prod_id_i^{\nu_i}$. Keeping the largest of these at most $(n+1)^r$ groups and writing $x=\nu/n$,
\[
\frac{\ln\mathbb EN_W}{n\ln k}\ge\max_{x\in\Delta_n}\phi(x)-O\Bigl(\frac{\log n}n\Bigr),\qquad\phi(x)=\frac{\Ent(x)+\sum_ix_i\ln m_i+\min\bigl(0,\sum_ix_i\ln d_i-\kappa\ln k\bigr)}{\ln k},
\]
where $\Delta_n$ is the set of points of the simplex $\Delta$ with denominator $n$; since $\phi$ is continuous on $\Delta$, $\max_{\Delta_n}\phi\to\max_\Delta\phi$. Now $\min(0,y)=\min_{t\in[0,1]}ty$, and $\Ent(x)+\sum_ix_i\ln m_i+t(\sum_ix_i\ln d_i-\kappa\ln k)$ is concave in $x$ and linear in $t$ on compact convex sets, so by Sion's minimax theorem the maximum over $x$ of the minimum over $t$ is the minimum over $t$ of the maximum over $x$. The inner maximum is $\ln\sum_im_id_i^t=\ln S(t)=(2+\lambda(t))\ln k$ (Gibbs' variational principle). Hence $\max_\Delta\phi=\min_{t\in[0,1]}(2+\lambda(t)-t\kappa)=2-\Leg_K(\kappa)$, which matches the upper bound.
\end{proof}

\AlgD{} with top algorithm $B^{\otimes n}$ computes at each inner position exactly the products of the terms in $T_{rs}$, so for a random wanted set its number of inner products per position is $N_W$. The upper bound in the proof is the Hölder step of Lemma~\ref{cp:T2}; so for tensor powers that step is tight in the exponent, once the class-maximum profile is replaced by the leaf-degree profile. The restriction $t\le1$ cannot be dropped: $t\kappa-\lambda(t)$ is concave, so whenever $\kappa>\lambda'(1)$ its supremum over all $t\ge0$ exceeds $\Leg_K(\kappa)$, and a formula with unrestricted $t$ would claim fewer needed terms than there are.

The batch statement of Definition~\ref{cp:P2} also has a saving for sparser batches, $\kappa>1$, where the budget per wanted entry stays at $D$, the cost of an inner product.

\begin{proposition}[sparse batches]\label{cp:F4}
Let $P$ be an admissible rational level-1 point, $\alpha=\alpha(P)$, and for $t\in(0,1]$ let $L_t\ge\Lamp(t)$ be the numbers of Theorem~\ref{cp:B1}. For $\kappa>0$ put $s=t\kappa-(\kappa-1)_+$ and
\[
\Ebat(\kappa)=\sup_{0<t\le1}\Bigl(t\kappa-\frac{L_t}\alpha\Bigr)-(\kappa-1)_+ .
\]
\begin{enumerate}[label=(\alph*)]
\item With the direct bottom, the proof of Theorem~\ref{cp:B1} at tilt $t$ and split $x$ with $s>(xL_t+1)/(1+\alpha x)$ gives, for every $\eps<(1+\alpha x)/(1+x)$, phase~2 within the target of Definition~\ref{cp:P2} with every saving $0<g<s-(xL_t+1)/(1+\alpha x)$ (when $\bk\ge N^2/(2D^\kappa)$, this is $O(N^2D^{\min(0,1-\kappa)-g})$ steps). As $x\to\infty$ the saving bound tends to $s-L_t/\alpha$ while the range bound stays above $\alpha$. So if $\Ebat(\kappa)>0$, then for every $\eps\le\alpha$ and every $0<g<\Ebat(\kappa)$ the $\kappa$-batch statement holds at $\eps$ with this $g$ (time model).
\item $\Ebat$ is nondecreasing on $(0,1]$ and nonincreasing on $[1,\infty)$, so it is largest at $\kappa=1$.
\item Let $P=\Pzero$ and $L_t=M_t$ (Lemma~\ref{cp:B2}). For $\kappa>1$, $\fcert(\kappa)>\alpha(\Pzero)$ if and only if $\Ebat(\kappa)>0$ with the supremum taken over $\Tcert$.
\end{enumerate}
\end{proposition}

\begin{proof}
(a) is a reading of the proof of Theorem~\ref{cp:B1} with $q_A=1$ and $\eps_A=1$: its phase-2 work is $N^2D^{-t\kappa+(x\lamB(t)+1)/(1+\alB x)+o(1)}$ against the target $N^2D^{\min(0,1-\kappa)}$, and $\min(0,1-\kappa)=-(\kappa-1)_+$; the range $(1+\alpha x)/(1+x)$ exceeds $\alpha$ for every $x>0$, since $\alpha<1$.
(b) $\Ebat=\tilde\Leg-(\kappa-1)_+$, where $\tilde\Leg(\kappa)=\sup_t(t\kappa-L_t/\alpha)$ is nondecreasing with $\tilde\Leg(\kappa')-\tilde\Leg(\kappa)\le\kappa'-\kappa$ by Lemma~\ref{cp:F1:basic}(a). On $(0,1]$, $\Ebat=\tilde\Leg$; on $[1,\infty)$, $\Ebat(\kappa')-\Ebat(\kappa)=\tilde\Leg(\kappa')-\tilde\Leg(\kappa)-(\kappa'-\kappa)\le0$.
(c) For $\kappa>1$, $\fcert(\kappa)$ is the maximum over the grid of $\rdir(\kappa,t)$ and $\rDS(\kappa,t)$. Where $\rdir$ is defined ($\alpha s>L_t$ and $s<1$), $\rdir(\kappa,t)-\alpha=(1-\alpha)(\alpha s-L_t)/(1-s+\alpha s-L_t)>0$. Where $\rDS$ is defined ($L_t\ge\alpha s>0$), $\rDS(\kappa,t)\le\alpha$, with equality only if $L_t=\alpha s$: $\rDS(\kappa,t)<\alpha$ is equivalent to $(\alpha-\es)(L_t-\alpha s)>0$. At $t=1$, $s=1$ and $\rdir$ is not defined, and $s-M_1/\alpha=1-\cpv{cbarP0_hi}/\alpha(\Pzero)<0$; for $t<1$, $s<1$. So $\fcert(\kappa)>\alpha$ if and only if some grid tilt has $s-L_t/\alpha>0$.
\end{proof}

So the plateau of Corollary~\ref{cp:B7}(a), the $\kappa>1$ at which the certified curve is at least $\alpha(\Pzero)$, ends at the zero of $\Ebat$ taken over the certified grid with $L_t=M_t$; by Corollary~\ref{cp:B7}(a) this zero is not below $\cpv{u8_plateau_end}$ (certified). That the true range $\fbat(\kappa)$ of the problem is largest at $\kappa=1$ is a different statement, Corollary~\ref{cp:B6}, proved by splitting and inclusion and not by the Legendre form.

\begin{remark}[the plateau $t=1$ and single entries]\label{cp:F9}
For a convex profile $\Lambda$ and $\kappa\ge\Lambda'(1)$ (left derivative) the slope $\kappa-\Lambda'(t)$ of $t\kappa-\Lambda(t)$ is nonnegative on $[0,1]$, so the best tilt is $t=1$ and $\Leg(\kappa)=\kappa-\Lambda(1)$; for $\kappa\ge1$ the sparse-batch saving $\Leg(\kappa)-(\kappa-1)$ is the constant $1-\Lambda(1)$. At $t=1$ the leaf-degree profile and the $\Sall$ moment are mean coverages. For the merged level-1 family the $\Sall$ moment equals the coverage exponent exactly: every output of the kept type is covered by $(q+2)^{M_q/(q+2)}2^{qM_q/(q+2)}$ terms on the part $q$ (the bound of Lemma~\ref{cp:L5} \Hcond4, the same for every such output), so the moment $\frac1{M}\ln\mathbb E_\tau\Sall(\tau)$ tends to $\sum_q\lambda_q(\ln(q+2)+q\ln2)/(q+2)$, which is $\co(P)$ in units of $\ln a_1$, the coverage exponent that Proposition~\ref{cp:T8} uses. A single wanted entry is a batch of density $K^{-2}=D^{-2/\alpha}$, and at $\Pzero$, $2/\alpha(\Pzero)>\cpv{F9_entry_lo}$, beyond the slope $\Lamcert'(1)/\alpha$, which is about $\cpv{F9_kappa}$ (computed). So single entries lie on the plateau, where the profile enters only through its value at $t=1$, and the thresholds of Theorem~\ref{cp:D1}, Theorem~\ref{cp:D3} and Table~\ref{cp:tab:D4} are functions of $\alpha(P)$ and $\co(P)$ alone. This uses that the largest coverage equals the mean, which holds at level~1 (Lemma~\ref{cp:L10}); where only a worst-case coverage is known, as in Remark~\ref{cp:D10}, the profile gives only the mean-coverage value.
\end{remark}

\subsection{Data structures at a fixed query exponent}\label{cp:sec:F:frontier}

A data structure that answers each query on its own cannot share queries, but it can store the products of the terms of high degree and compute the others at query time. For tensor powers of a fixed base this gives a frontier between preprocessing and query exponent that is again built from the Legendre function. Worst-case outputs are handled by the following averaging lemma.

\begin{lemma}[symmetrisation]\label{cp:F22}
Let a finite group $\Gamma$ act transitively on a finite set $Z$, and let $f:Z\to\R$ have mean $\bar f$. Then for every $n$ and every $o\in Z^n$ there is $g\in\Gamma$ with $\sum_{j=1}^nf(go_j)\le n\bar f$.
\end{lemma}

\begin{proof}
For each $z\in Z$ the map $g\mapsto gz$ is $\abs{\mathrm{Stab}(z)}$-to-one onto $Z$, by transitivity, so $\frac1{\abs\Gamma}\sum_gf(gz)=\bar f$. Summing over $j$, the average over $g$ of $\sum_jf(go_j)$ is $n\bar f$, and the minimum is at most the average.
\end{proof}

The model counts linear operations on given term products, as in Section~\ref{cp:sec:tops}, where the products are computed from the encodings and their cost is accounted for separately.

\begin{definition}[box structures]\label{cp:F6:model}
Let $\mathcal A$ be an exact decomposition with terms $\tau$, outputs $o$, coefficients $c(\tau,o)=w_\tau[o]$ and covering sets $\Cov(o)$, and let the products $\Pi(\tau)$ of its terms be given. A \emph{box} is a pair $(R,C)$ of a set $R$ of terms and a set $C$ of outputs with $c(\tau,o)\ne0$ and $c(\tau,o)=f(\tau)g(o)$ for all $(\tau,o)\in R\times C$; its value is $\sum_{\tau\in R}f(\tau)\Pi(\tau)$. A \emph{box structure} computes in a preprocessing the values of a set of boxes, each as a linear combination of \emph{children}, which are boxes or term products whose term sets partition $R$; the preprocessing cost is the number of term products used plus the total number of children. The query for an output $o$ returns $y_o=\sum_{\tau\in\Cov(o)}c(\tau,o)\Pi(\tau)$ as a sum of $g(o)$ times the values of some boxes containing $o$ (the \emph{top boxes} of $o$) and of terms $c(\tau,o)\Pi(\tau)$ (the \emph{singletons} of $o$), where the term sets of the top boxes and the singletons partition $\Cov(o)$; its cost is their number.
\end{definition}

\begin{theorem}[the frontier for tensor powers]\label{cp:F6}
Let $B$ be an exact decomposition over $\Q$ of $\mat kbk$ with $b\ge2$ whose $\RB$ terms all have nonzero $z$-form, with maximal coverage $\cB$ and $\lambda(s)=\log_kS(s)-2$, $S(s)=\sum_\rho\deg(\rho)^s$ (all outputs read). Let $n\ge1$, $K=k^n$, $u>0$ and $s\in(0,1]$. There are an exact decomposition $\tilde B_n$ of $\mat K{b^n}K$ with $k^2\RB^n$ terms and a box structure for it with preprocessing cost at most $k^2(\cB+1)\,n\,K^{2+\lambda(s)-su}$ and query cost at most $(1+n\cB)K^{u(1-s)+\lambda(s)+\log_k2}$ for every output. Both are explicit and deterministic. The costs are those of Definition~\ref{cp:F6:model}: linear operations on given term products; they are not word-RAM times, and the cost of computing the term products is not included.
Consequently, with $D=b^n$ and $\Leg$ the Legendre function of $\lambda/\alB$: for every $u'>0$ with $\Leg(u')>0$ there is a box structure with preprocessing cost $K^2D^{-\Leg(u')+o(1)}$ and query cost $D^{\,u'-\Leg(u')+\log_k2/\alB+o(1)}$ for every output, as $n\to\infty$.
\end{theorem}

\begin{proof}
\emph{One copy.} For a term string $\tau$ of $B^{\otimes n}$ let $\mathrm{ord}(\tau)=\sum_j\ln\deg(\tau_j)\ge0$, and put $\theta=u\ln K$. For $0\le j<n$, a suffix $\sigma$ of length $n-j$ (term letters at the coordinates $j+1,\dots,n$) with $\mathrm{ord}(\sigma)\ge\theta$ and a prefix $\zeta$ of outputs of length $j$, let the \emph{generalised box} $G(j,\sigma,\zeta)$ have the terms $(\rho,\sigma)$, $\rho\in\Cov(\zeta_1)\times\dots\times\Cov(\zeta_j)$, and the outputs $(\zeta,o'')$ with $c(\sigma,o'')\ne0$, and value $\sum_\rho c(\rho,\zeta)\Pi(\rho,\sigma)$; here $c(\rho,\zeta)=\prod_ic(\rho_i,\zeta_i)$. Every term of it serves every one of its outputs, and $c((\rho,\sigma),(\zeta,o''))=c(\rho,\zeta)c(\sigma,o'')$, so it is a box for any coefficients. It satisfies $G(0,\sigma,\emptyset)=\Pi(\sigma)$ and $G(j,\sigma,\zeta)=\sum_{r\in\Cov(\zeta_j)}c(r,\zeta_j)\,G(j-1,(r,\sigma),\zeta_{<j})$, with at most $\cB$ children, which again have order at least $\theta$. For an output $o$, a term $\tau\in\Cov(o)$ with $\mathrm{ord}(\tau)<\theta$ is a singleton of $o$; for $\mathrm{ord}(\tau)\ge\theta$ let $l$ be the largest index with $\mathrm{ord}(\tau_{\ge l})\ge\theta$, and let $\tau$ belong to the top box $G(l-1,\tau_{\ge l},o_{<l})$ of $o$, multiplied by $c(\tau_{\ge l},o_{\ge l})$. The index $l$ depends only on the suffix of $\tau$, and every term of that top box has the same suffix and hence the same $l$; so the top boxes and the singletons partition $\Cov(o)$.

\emph{Costs of one copy.} Let $N_h$ be the number of strings of length $h$ with order at least $\theta$; Markov's inequality with $\deg^s$ gives $N_h\le e^{-s\theta}S(s)^h$. There are at most $k^{2(n-h)}N_h$ generalised boxes with suffixes of length $h$, each with at most $\cB$ children, and $N_n$ term products are used. Since $S(s)\ge\RB\ge k^2$, $k^{2(n-h)}S(s)^h\le S(s)^n$, so the preprocessing cost is at most $(n\cB+1)e^{-s\theta}S(s)^n\le(\cB+1)nK^{2+\lambda(s)-su}$. For the query, put $V(z)=\sum_{\rho\in\Cov(z)}\deg(\rho)^{s-1}$ for a base output $z$. The singletons of $o$ number at most $\sum_{\tau\in\Cov(o)}(e^\theta/\deg\tau)^{1-s}=K^{u(1-s)}\prod_jV(o_j)$. A top box of $o$ is fixed by $l$, a letter of $\Cov(o_l)$ and a string of $\Cov(o_{>l})$ of order below $\theta$, so there are at most $\sum_l\cB K^{u(1-s)}\prod_{j>l}V(o_j)\le n\cB K^{u(1-s)}\prod_j\max(1,V(o_j))$ of them. So the query cost of $o$ is at most $(1+n\cB)K^{u(1-s)}\prod_j\max(1,V(o_j))$.

\emph{Worst-case outputs.} The $k^2$ cyclic shifts $(i,h)\mapsto(i+a,h+c)$ of the $k\times k$ output grid act simply transitively on it and are symmetries of $\mat kbk$, since $P_aX\cdot YQ_c=P_a(XY)Q_c$ for the cyclic permutation matrices $P_a$, $Q_c$; one shift $g$ applied in every coordinate is a symmetry of $\mat K{b^n}K$, and the copy $gB^{\otimes n}$ serves $o$ through the covering set of the shifted output $go$ in $B^{\otimes n}$. The mean of $V$ over the $k^2$ base outputs is $S(s)/k^2$, since each $\rho$ contributes $\deg(\rho)^{s-1}$ to $\deg(\rho)$ outputs. Lemma~\ref{cp:F22} with $f=\ln\max(1,V)$ and Jensen's inequality give, for every $o$, a shift $g$ with
\[
\prod_j\max(1,V(go_j))\le\Bigl(\frac1{k^2}\sum_z\max(1,V(z))\Bigr)^n\le\Bigl(1+\frac{S(s)}{k^2}\Bigr)^n\le\Bigl(\frac{2S(s)}{k^2}\Bigr)^n=K^{\lambda(s)+\log_k2}.
\]
Assign every output to such a shift, which is computed from $o$ alone, and keep in each copy only the $z$-forms of its assigned outputs. This gives one exact decomposition $\tilde B_n$ with $k^2\RB^n$ terms, and the union of the copies' box structures, restricted to the assigned outputs, is a box structure for it with the stated costs.

\emph{Units.} With $\ln D=\alB\ln K$ and $u=\alB u'$, the preprocessing is $K^2D^{-(su'-\lambda(s)/\alB)}$ and the query $D^{u'-(su'-\lambda(s)/\alB)+\log_k2/\alB}$, up to factors $\mathrm{poly}(n)=D^{o(1)}$. Choose $s\in[0,1]$ maximising $su'-\lambda(s)/\alB$; the maximum is $\Leg(u')>0$, so $s>0$, since the value at $s=0$ is $-\lambda(0)/\alB\le0$.
\end{proof}

The slack $\log_k2/\alB$ comes from the averaging step and does not vanish for a fixed base; along bases with $k=a^j$, as in Lemma~\ref{cp:T5:lemma4}, it tends to $0$. The lemma needs a group of symmetries that acts transitively on the outputs and coordinatewise on the powers; matrix products have one, but the decomposition of $F_q$ alone does not.

So for tensor powers the pairs (query exponent, saving) $(q,\gamma)=(u-\Leg(u),\Leg(u))$, $u>0$, are reached up to the slack. The next proposition records what such a frontier implies; it uses only Lemma~\ref{cp:F1:basic}.

\begin{proposition}[the frontier]\label{cp:F7}
Let $\Lambda$ be continuous and nondecreasing on $[0,1]$ with $\Lambda(0)\ge0$, $\Leg$ its Legendre function, and $q(u)=u-\Leg(u)$.
\begin{enumerate}[label=(\alph*)]
\item If $\Leg(u)>0$, then $q(u)<u$ and $\Leg(u)>\Leg(q(u))$: at query exponent $q$ the frontier saves more than $\Leg(q)$.
\item An application with $N^2/D^\beta$ queries pays $N^2D^{-\gamma}+N^2D^{-\beta+q}$ and saves $\min(\gamma,\beta-q)$. On the frontier the best saving is $\max_u\min(\Leg(u),\beta-u+\Leg(u))=\Leg(\beta)$, attained at $u=\beta$.
\item There is $u$ with $\Leg(u)>0$ and $q(u)<\beta$ if and only if $\Leg(\beta)>0$.
\end{enumerate}
\end{proposition}

\begin{proof}
(a) $q(u)<u$ is clear, and $\Leg(u)\ge\Leg(q(u))$ since $\Leg$ is nondecreasing. If equality held, $\Leg$ would be constant on $[q(u),u]$. Take $\beta$ strictly between and a maximising tilt $t_\beta$ at $\beta$ (it exists since $\Lambda$ is continuous). Then $\Leg(u)\ge t_\beta u-\Lambda(t_\beta)=\Leg(\beta)+t_\beta(u-\beta)=\Leg(u)+t_\beta(u-\beta)$, so $t_\beta=0$ and $\Leg(\beta)=-\Lambda(0)\le0$, contradicting $\Leg(\beta)=\Leg(u)>0$.
(b) For $u\le\beta$ the minimum is $\Leg(u)\le\Leg(\beta)$; for $u\ge\beta$ it is $\beta-u+\Leg(u)\le\Leg(\beta)$ by Lemma~\ref{cp:F1:basic}(a); at $u=\beta$ both terms equal $\Leg(\beta)$.
(c) If $\Leg(\beta)>0$, take $u=\beta$. Conversely, let $\Leg(u)>0$ and $u-\Leg(u)<\beta$. If $u\le\beta$, then $\Leg(\beta)\ge\Leg(u)>0$; if $u>\beta$, then $\Leg(\beta)\ge\Leg(u)-(u-\beta)>0$.
\end{proof}

For $\Hyp$ ($\beta=\frac12$) the balanced saving is $\Leg(\frac12)$. The trade-off of the data structure of~\cite{v2} has exactly this form:

\begin{remark}[the tile of~\cite{v2}]\label{cp:F5}
For the squared runs of~\cite{v2} at ratio $c$ read on their type class, the leaves per output and the log-degree have the rates $\varphi_c(x)=\int_0^x\ln\frac{4(1-y)}{c-1+y}\,dy$ and $\psi_c(x)=(c-1+x)\ln(c-1+x)-x\ln x-(c-1)\ln(c-1)$, and the leaf-degree profile is $\LamT c(t)=\frac2{\ln4}\max_{0\le x\le1}\bigl[\varphi_c(x)+t\psi_c(x)\bigr]$, by the leaf counts of~\cite{v2} and the method of types; let $\Leg^{\mathrm{tile}}_c$ be its Legendre function. Both rates are strictly concave and $\varphi_c+\psi_c=\Hb(x)+x\ln4$. For $c\ge2$ and $0<\theta<0.8$ with $\varphi_c'(\theta)<0$ the tilt $t=-\varphi_c'(\theta)/\psi_c'(\theta)$ lies in $(0,1)$, $\theta$ is the unique maximiser of $\varphi_c+t\psi_c$, and with $u=2\psi_c(\theta)/\ln4$ one gets $\Leg^{\mathrm{tile}}_c(u)=-2\varphi_c(\theta)/\ln4$ and $u-\Leg^{\mathrm{tile}}_c(u)=2(\varphi_c+\psi_c)(\theta)/\ln4$. These are the saving and the query exponent of the exact box count of~\cite[Corollary~\Vref{m:cor:exactboxes}]{v2}, whose hypothesis implies $\varphi_c'(\theta)<0$: every point of its trade-off lies on the frontier of Proposition~\ref{cp:F7} for the tile.
\end{remark}

\begin{remark}[ranges of the tile]\label{cp:F8}
A range is not a value of $\Leg$; it is set by the encoding cost at a zero of $\Leg$. The tile's rank excess $\LamT c(0)=\frac2{\ln4}\max\varphi_c$ is positive if and only if $c<5$, since $\varphi_c'(0)=\ln(4/(c-1))$. For $c\ge5$ the tile has $\cg=0$, like the merged level-1 family, so $\Leg^{\mathrm{tile}}_c(q)>0$ for every $q>0$; for $c<5$, $\Leg^{\mathrm{tile}}_c(q)<0$ for $q<\LamT c(0)$ by Lemma~\ref{cp:F1:basic}(b). The encoding cost $A(c)=2c\ln5-g_c$, with $g_c=c\Hb(1/c)+(c-1)\ln4$, sets the range $\ln4/(A(c)+\gamma\ln4)$ of the tile at ratio $c$ and saving $\gamma$~\cite{v2}, and it increases on $[5,\infty)$, so the best range with a saving at every query exponent is $\ln4/A(5)=\ln4/(5\ln5)=\es$: the range of~\cite{v2} is the encoding-limited range at the smallest ratio where the tile's rank excess vanishes. The ranges of~\cite[Theorem~\Vref{m:thm:fixedq}]{v2} at a fixed query exponent, $\cpv{hyp_half_v2}$ at $q=\frac12$ and $\cpv{hinted_omv_v2}$ at $q=1$, are $\ln4/A(c^*(q))$ with $c^*(q)$ the zero of $c\mapsto\Leg^{\mathrm{tile}}_c(q)$. For the composites the second parameter is the split $x$, with range $(1+\alpha x)/(x+1/\eps_A)$.
\end{remark}

\subsection{Mixtures and small budgets}\label{cp:sec:F:small}

\begin{proposition}[mixtures]\label{cp:F12}
Let $(\mathcal A_i,R_i)$, $i=1,2$, be decompositions with read sets and inner dimensions $D_i$, let $\mathcal A=\mathcal A_1\otimes\mathcal A_2$ be read on $R_1\times R_2$, $D=D_1D_2$ and $\theta_i=\ln D_i/\ln D$. With $\Lambda_i(t)=\ln(\frac1{M_i}\sum\deg^t)/\ln D_i$ and likewise $\Lambda$ for $\mathcal A$, at finite size,
\[
\Lambda=\theta_1\Lambda_1+\theta_2\Lambda_2,\qquad\Leg(\beta)\le\theta_1\Leg_1(\beta)+\theta_2\Leg_2(\beta)\le\max\bigl(\Leg_1(\beta),\Leg_2(\beta)\bigr).
\]
\end{proposition}

\begin{proof}
The degree of $(\tau_1,\tau_2)$ on $R_1\times R_2$ is $\deg\tau_1\cdot\deg\tau_2$, so $\frac1{M_1M_2}\sum\deg^t$ factorises and $\ln D=\ln D_1+\ln D_2$. A supremum of a sum is at most the sum of the suprema, and $\theta_1+\theta_2=1$.
\end{proof}

So mixing coordinates of two families never saves more at a budget than the better of them; for products of the level-1 top with tiles the corresponding statement on $\eps\gamma$ is Proposition~\ref{cp:G9}.

For the merged level-1 family the shape of the profile near $t=0$ is explicit. Recall from Lemma~\ref{cp:G1} that $\Lamcert(t)=\sum_q\lambda_q\Theta^*_q(t)/\ln a_1$, where $\Theta^*_q(t)$ is the maximum over the polytope $\Polq_q$ of $\Theta_q(t;d_B,d_O,y)=-\KL(x\,\|\,x^*)+t\,\sigma_q(d_B,d_O,y)$, with $\pi=1/(q+2)$, $x_A=q\pi-d_B+d_O$, $x_B=\pi+d_B$, $x_O=\pi-d_O$ and $\Hb$ the binary entropy. In the next proposition and its proof we abbreviate $u=d_B$, $v=d_O$, $a=x_A$, $b=x_B$, $o=x_O$ and $\KL=\KL(x\,\|\,x^*)$, so that
\begin{gather*}
\KL=a\ln\tfrac a{q\pi}+b\ln\tfrac b\pi+o\ln\tfrac o\pi,\qquad\sigma_q=b\Hb(\tfrac\pi b)+u\Hb(\tfrac yu)+(u-y)\ln q+a\Hb(\tfrac{v-y}a),\\
\Polq_q=\{0\le v\le\pi,\ 0\le y\le v,\ y\le u,\ u-y\le q\pi\}.
\end{gather*}

\begin{proposition}[small tilts, merged family]\label{cp:F17}
For every $q\ge2$, $\Theta^*_q(t)=\pi t^2\ln^2(1/t)\,(1+o(1))$ as $t\to0+$, with $\pi=1/(q+2)$. Hence, in units of $\ln D$,
\[
\frac{\Lamcert(t)}{\alpha(P)}=C_D\,t^2\ln^2(1/t)\,(1+o(1)),\qquad C_D=\frac{\sum_q\lambda_q/(q+2)}{\ln b_1},\quad\ln b_1=\alpha(P)\ln a_1,
\]
and $C_D=\cpv{F17_CD}\ldots$ at $\Pzero$ (certified; Section~\ref{cp:sec:certificates}). This is a statement about the upper bound $\Lamcert$ of Lemma~\ref{cp:G1}, not about \ThmD's profile $\Lam$, which is only bounded by it.
\end{proposition}

\begin{proof}
Write $\phi(r)=r\ln r-r+1$; since the deviations of $(a,b,o)$ from $(q\pi,\pi,\pi)$ sum to $0$, $\KL=q\pi\phi(\frac a{q\pi})+\pi\phi(\frac b\pi)+\pi\phi(\frac o\pi)$.
\emph{Lower bound.} Take $u=v=\bar x:=\pi t\ln(1/t)$ and $y=0$, a point of $\Polq_q$ for small $t$. Then $a=q\pi$ and, with $e=\bar x/\pi$, $\KL=\pi(\phi(1+e)+\phi(1-e))\le(\bar x^2/\pi)(1+e^2)$, since $\phi(1+e)+\phi(1-e)=\sum_{k\ge1}e^{2k}/(k(2k-1))$. Further $b\Hb(\pi/b)=\pi\ln(1+u/\pi)+u\ln((\pi+u)/u)\ge u\ln(\pi/u)$ and $a\Hb(v/a)\ge v\ln(a/v)$, so $\sigma_q\ge2\bar x\ln(1/\bar x)-O(\bar x)$. Hence $\Theta^*_q(t)\ge-\bar x^2(1+e^2)/\pi+2t\bar x\ln(1/\bar x)-O(t\bar x)=\pi t^2\ln^2(1/t)(1-o(1))$, because $\ln(1/\bar x)=\ln(1/t)(1-o(1))$.
\emph{Localisation.} At $u=v=y=0$ the objective is $0$, so $\Theta^*_q(t)\ge0$. By Pinsker's inequality $\KL\ge\frac12(\abs{v-u}+u+v)^2\ge(u+v)^2/2$, and $\sigma_q\le\sigma_{\max}$ on the compact set $\Polq_q$; so at every point with $u+v\ge\rho:=t^{1/3}$ the objective is at most $-t^{2/3}/2+t\sigma_{\max}<0$ for all small $t$, and the maximum is attained where $u+v\le\rho$.
\emph{Upper bound for $u+v\le\rho$.} From $\phi(1+e)\ge e^2/(2(1+\abs e/3))$, $\KL\ge Q(u,v)/(1+c_0\rho)$ for a constant $c_0$, with $Q=(v-u)^2/(2q\pi)+(u^2+v^2)/(2\pi)\ge(u+v)^2/(4\pi)$. Next $b\Hb(\pi/b)\le u\ln(1/u)+u$ (as $b\le1$), $\max_y[u\Hb(y/u)+(u-y)\ln q]=u\ln(1+q)$, and $a\Hb((v-y)/a)\le v\ln(ea/v)\le v\ln(1/v)+v$ (as $v\le a$ here); so $\sigma_q\le u\ln\frac1u+v\ln\frac1v+C_2(u+v)$ with $C_2=2+\ln(1+q)$, and $u\ln\frac1u+v\ln\frac1v\le s'\ln(2/s')$ for $s'=u+v$ by concavity. Hence
\[
\Theta^*_q(t)\le\max_{s'>0}\Bigl[-\frac{s'^2}{4\pi(1+c_0\rho)}+ts'\ln\frac{2e^{C_2}}{s'}\Bigr]=\pi(1+c_0\rho)\,t^2\ln^2(1/t)\,(1+o(1)),
\]
by an elementary maximisation (the maximiser is $s'=2\pi(1+c_0\rho)t\ln(1/t)(1+o(1))$). With $\rho\to0$ both bounds give the claim. The formula for $C_D$ follows from $\Lamcert/\alpha(P)=\sum_q\lambda_q\Theta^*_q/\ln b_1$.
\end{proof}

In particular $\Lamcert(t)/t^2\to\infty$: no finite variance describes the profile near $t=0$, and a Gaussian approximation of it is wrong already in the leading order.

\begin{remark}[small budgets]\label{cp:F18}
For a profile of the form $\Lambda(t)=Ct^2(W^2+2W)$ with $We^W=A/t$ (constants $C,A>0$) one has $\Lambda'(t)=2CtW(W+1)$, which is strictly increasing, and at the optimal tilt of budget $\beta$ (when it is at most $1$)
\[
\Leg(\beta)=Ct^2W^2=\frac{\beta^2}{4C(1+W^*)^2},
\]
$W^*$ the value of $W$ there. The merged family's $\Theta^*_q$ has this form only up to a relative error $O(tW)$, which is argued and not proved. As an approximation for $\Lamcert/\alpha(\Pzero)$, applied part by part with $A=q$ on the part $q$ and with $C_D$ of Proposition~\ref{cp:F17}, the formula $\beta^2/(4C_D(1+W^*)^2)$ is within $\cpv{F18_err_a}$ (relative) for $\beta\le\cpv{F18_beta_a}$, $\cpv{F18_err_b}$ at $\beta=\cpv{F18_beta_b}$ and $\cpv{F18_err_c}$ at $\beta=\frac12$ (computed). The leading asymptotic $\beta^2/(4C_D\ln^2(1/\beta))$ is correct only as $\beta\to0$, converging slowly, and is off by a factor $\cpv{F18_factor}$ at $\beta=\cpv{F18_beta_a}$ (computed); it should not be used for numbers.
\end{remark}

The quadratic-logarithmic shape is not a property of all profiles with $\cg=0$. It holds for the merged family (Proposition~\ref{cp:F17}), where the cost of the maximisation is quadratic at the optimum for $t=0$; for the tile at $c>5$ the cost is linear there, and the profile is much flatter (Proposition~\ref{cp:F19}).

\begin{proposition}[the tile at $c>5$]\label{cp:F19}
Let $c>5$ and $s=\ln((c-1)/4)>0$. Then $\LamT c(t)\le\frac{2c}{\ln4}\,te^{-s/t}$ for $t\in(0,1]$, and $\Leg^{\mathrm{tile}}_c(\beta)\ge\frac{s\beta}{\ln(1/\beta)}(1-o(1))$ as $\beta\to0+$. In particular $\LamT c(t)=o(t^p)$ for every $p>0$, and $\cg=0$.
\end{proposition}

\begin{proof}
$\varphi_c$ is concave with $\varphi_c(0)=0$ and $\varphi_c'(0)=\ln(4/(c-1))=-s$, so $\varphi_c(x)\le-sx$. Also $\psi_c(x)=x\ln\frac{c-1+x}x+(c-1)\ln(1+\frac x{c-1})\le x\ln\frac cx+x=x\ln\frac{ec}x$ for $0<x\le1$. Hence $\LamT c(t)\le\frac2{\ln4}\max_{x>0}x\bigl(-s+t\ln\frac{ec}x\bigr)$; the maximiser is $x=ce^{-s/t}$, with value $tce^{-s/t}$. For the lower bound take $t=s/\ln(A'/\beta)$ with $A'=\frac{2c}{\ln4}\ln\frac1\beta$, so that $\frac{2c}{\ln4}e^{-s/t}=\beta/\ln(1/\beta)$; then $\Leg^{\mathrm{tile}}_c(\beta)\ge t\beta-\frac{2c}{\ln4}te^{-s/t}=t\beta(1-1/\ln(1/\beta))$, and $\ln(A'/\beta)=\ln(1/\beta)(1+o(1))$.
\end{proof}

The \emph{rank exponent} of the tile is $\alpha^{\mathrm{tile}}_c=\ln4/g_c$, the ratio of $\ln D$ to the logarithm of the side of its output square. It is the tile's analogue of the dual exponent $\alpha$, not the range $\ln4/A(c)$ of~\cite{v2}, which is smaller for $c\ne5$, since $A(c)\ge g_c$ with equality only at $c=5$.

\begin{corollary}[no variance barrier below $\es$]\label{cp:F20}
No lower bound $\Lambda(t)\ge C(\alpha)t^2\ln^2(1/t)$ with $C(\alpha)>0$, and no lower bound by a fixed power of $t$, holds for all families with $\cg=0$ and rank exponent $\alpha<\es=\cpv{range_v1}$: the tiles at $c>5$ violate every such bound.
\end{corollary}

\begin{proof}
$g_c$ increases ($g_c'=\ln(4c/(c-1))>0$) from $g_5=5\ln5$ to $\infty$, so $\alpha^{\mathrm{tile}}_c$ runs over $(0,\es)$ as $c$ runs over $(5,\infty)$, and Proposition~\ref{cp:F19} bounds these profiles by $o(t^p)$ for every $p$.
\end{proof}

Above $\es$ no profile of this flat kind is known, and whether a bound $C_D\ge C_{\min}(\alpha)>0$ holds for all families with $\cg=0$ and rank exponent $\alpha\in(\es,\alpha(\Pzero)]$ is open. For powers of one fixed base there is a barrier, but it is a threshold and not a variance:

\begin{proposition}[fixed bases]\label{cp:F21}
\begin{enumerate}[label=(\alph*)]
\item Every exact decomposition of $\mat kbk$ with $b\ge2$ has more than $k^2$ terms of nonzero $z$-form.
\item Let $B$ be such a decomposition, read on all outputs, with $\lambda$ as above, $\delta=\lambda(0)$ and $\ell_B=\lambda'(0)$, the mean of $\log_k\deg$ over its terms. The powers of $B$ have, in units of $\ln D$, $\Leg(\beta)\le(\beta-\ell_B/\alB)_+-\delta/\alB$ for every $\beta\ge0$; so $\Leg(\beta)\le0$ at budgets $\beta\le(\delta+\ell_B)/\alB$, where neither the batch value nor the frontier of this section gives a saving, and $\delta>0$ by (a).
\end{enumerate}
\end{proposition}

\begin{proof}
(a) The $z$-slices $T_o$ of $\mat kbk$, the bilinear forms $(XY)_o$ written as matrices indexed by the entries of $X$ and $Y$, span a space of dimension $k^2$, and every nonzero element $\sum_{(i,h)}C_{ih}T_{ih}$ of it is block-diagonal with $b$ copies of the $k\times k$ matrix $C$, so it has rank $b\,\rk C\ge b\ge2$. If a decomposition had exactly $k^2$ terms of nonzero $z$-form, each $T_o=\sum_lw_l[o]\,u_l\otimes v_l$ would lie in the span of the $k^2$ matrices $u_l\otimes v_l$, so these spans would coincide and the $u_l\otimes v_l$ would be nonzero elements of rank one in the span of the $z$-slices, a contradiction. Fewer terms are impossible by Lemma~\ref{cp:T1:basic}(a).
(b) $\lambda$ is convex, so $\lambda(t)\ge\delta+t\ell_B$, and $\Leg(\beta)\le\sup_{0\le t\le1}t(\beta-\ell_B/\alB)-\delta/\alB$. By (a), $\delta=\log_k\RB-2>0$.
\end{proof}

Small budgets therefore need growing families in which $\delta$ and $\ell_B$ tend to $0$, as the bases of \ThmDk{} do (Corollary~\ref{cp:T5}).

\begin{remark}[a second certificate]\label{cp:F23}
The bound of Theorem~\ref{cp:G3} is certified a second time, independently of the tangent-plane certificate of Lemma~\ref{cp:G2}: weak Lagrange duality for the maximisation that defines $\Theta^*_q$ bounds $\Theta^*_q(t)$ from above at every choice of the dual potentials, and rational potentials give $\Theta^*_6(\cpv{G3_t})\le\cpv{F23_G6}$ and $\Theta^*_7(\cpv{G3_t})\le\cpv{F23_G7}$ in exact interval arithmetic, hence the same saving at $\Pzero$.
\end{remark}

\subsection{What is an identity and what is a bound}\label{cp:sec:F:wording}

For the tile of~\cite{v2} the profile $\LamT c$ is the exact leaf-degree log-moment, and Remarks~\ref{cp:F5} and~\ref{cp:F8} are identities. For tensor powers of one fixed exact base the profile $\lambda$ is exact at every power; Proposition~\ref{cp:F3} is an identity for the needed terms of a random wanted set, and Theorem~\ref{cp:F6} is a construction, whose optimality is not claimed. For the level-1 family, \ThmD's profile $\Lam$ is bounded above by $\Lamcert$ (Lemma~\ref{cp:G1}) and by $\Lampr$ (Lemma~\ref{cp:L8}), so every printed level-1 saving is a lower bound on what \ThmDk{} and \ThmB{} save, and whether these savings are exact is open. The statement ``the saving is a value of $\Leg$'' holds for batches (Section~\ref{cp:sec:F:batch}, with the tilt condition of Proposition~\ref{cp:F2}); data structures at a fixed query exponent follow the frontier of Proposition~\ref{cp:F7}, and ranges are not values of $\Leg$ (Remark~\ref{cp:F8}).

\section{Certificates and scripts}\label{cp:sec:certificates}

A number is \emph{certified} in this paper when a script obtains it in exact rational arithmetic, or in interval arithmetic with outward rounding, from exact rational inputs, and \emph{computed} when it is obtained in floating point. Only certified numbers appear in statements of the form ``for every $\eps<$'' or ``for every $\gamma<$''; computed numbers are called computed where they appear. Every printed threshold is rounded in the safe direction: lower bounds and thresholds down, upper bounds up.

\paragraph{Three kinds of arithmetic.} The scripts use three independent trust bases. (F) Exact \texttt{Fraction} arithmetic of Python, with the logarithm of a positive rational $x$ enclosed by $\ln x=k\ln2+2\operatorname{artanh}((y-1)/(y+1))$, $y=x/2^k\in[\frac23,\frac43)$, and the series of $\operatorname{artanh}$ truncated with an explicit bound on its tail; interval endpoints are rounded outward to a dyadic grid (for example $\cpv{cert_ln_tol}$). (M) Interval arithmetic of the \texttt{mpmath} library (\texttt{iv}, 60 digits, outward rounding). (D) Python's \texttt{decimal} module (libmpdec), whose $\ln$ is correctly rounded, with directed rounding on every operation and the logarithms widened by one unit in the last place. Floating-point searches (\texttt{numpy}, \texttt{scipy}, and \texttt{mpmath} without intervals) are used only to choose the rational points, tilts and splits at which a certificate is evaluated; validity never depends on them, since every condition is checked at the chosen rational point. The concave maximisations of Lemma~\ref{cp:G2} are certified by one interior rational point and an enclosure of the objective and its gradient there; that the point is interior is checked exactly.

\paragraph{The certificates.} Table~\ref{cp:tab:cert} lists, for every certified number of the paper, the script that certifies it, the lines of its output, the kind of arithmetic (F, M or D), and a second, independent certificate where there is one. Output lines are cited as \texttt{file:line}. Lines marked ``CHECK FAILS'' in the output of \texttt{r211b\_\allowbreak{}cert.\allowbreak{}py} refer to digits printed in earlier notes that were rounded the unsafe way; this paper prints the safe values next to them, which the same lines certify.

{\footnotesize
\renewcommand{\arraystretch}{1.12}\setlength{\tabcolsep}{3pt}
\begin{longtable}{@{}>{\raggedright}p{0.23\textwidth}>{\raggedright}p{0.15\textwidth}>{\raggedright}p{0.26\textwidth}>{\raggedright}p{0.05\textwidth}>{\raggedright\arraybackslash}p{0.21\textwidth}@{}}
\caption{The certified numbers and their scripts.}\label{cp:tab:cert}\\
printed values & statement & script, output lines & kind & second certificate\\\hline
\endfirsthead
printed values & statement & script, output lines & kind & second certificate\\\hline
\endhead
\hline\endfoot
$\cpv{alphaP0_lo}<\alpha(\Pzero)<\cpv{alphaP0_hi}$, $\slack(\Pzero)\in(\cpv{slackP0_lo},\cpv{slackP0_hi})$, $\co(\Pzero)\le\cpv{cbarP0_hi}$, $1/(\co+2)\ge\cpv{tcapP0_lo}$; hence $\cpv{thmD_range_P0}$ & Lemma~\ref{cp:L11}, Corollary~\ref{cp:A2} & \texttt{r211a\_\allowbreak{}p0\_\allowbreak{}certify.\allowbreak{}py}, \texttt{\_\allowbreak{}out.\allowbreak{}txt:1-15} & F & \texttt{r211c\_\allowbreak{}cert\_\allowbreak{}out.\allowbreak{}txt:1-8} (D); \texttt{r79\_\allowbreak{}certify\_\allowbreak{}out.\allowbreak{}txt:1-4} (M); \texttt{r211f\_\allowbreak{}certify\_\allowbreak{}out.\allowbreak{}txt:28-30} (F)\\
$2/(\co(\Pzero)+2)>\cpv{tpadP0_lo}$ & Lemma~\ref{cp:L11} & \texttt{r79b\_\allowbreak{}certify.\allowbreak{}py}, \texttt{\_\allowbreak{}out.\allowbreak{}txt:3} & M & \texttt{r211c\_\allowbreak{}cert\_\allowbreak{}out.\allowbreak{}txt:7} (D)\\
$\cpv{excess_crude_P0}$ at $t=\cpv{excess_crude_t}$, $x=\cpv{excess_crude_x}$ & Remark~\ref{cp:A3} & \texttt{r211a\_\allowbreak{}excess\_\allowbreak{}certify.\allowbreak{}py}, \texttt{\_\allowbreak{}out.\allowbreak{}txt:1-3} & F & --\\
$\Lamcert(\cpv{G3_t})\le\cpv{u7_lamcert_t0}$, $\gamma<\cpv{G3_gamma_exact}$, $\cpv{G3_eps_bind}$, $\cpv{u7_G3_margin}$; $\Hyp(\eps,\cpv{G3_gamma})$ for $\eps\le\cpv{G3_eps}$; $\eps\gamma\ge\cpv{G4_epsgamma}$ & Theorem~\ref{cp:G3}, Corollary~\ref{cp:G4}, Corollary~\ref{cp:B8} & \texttt{r211b\_\allowbreak{}cert.\allowbreak{}py} (with \texttt{r211b\_\allowbreak{}lib.\allowbreak{}py}), \texttt{\_\allowbreak{}out.\allowbreak{}txt:14-30} & F & \texttt{r74a\_\allowbreak{}certify\_\allowbreak{}out.\allowbreak{}txt:17-40} (M); the two bounds $\Theta^*_6$, $\Theta^*_7$ also by weak duality, \texttt{r109\_\allowbreak{}dual\_\allowbreak{}out.\allowbreak{}txt:31-36} (rational; Remark~\ref{cp:F23})\\
the five rows of Table~\ref{cp:tab:G6} and $\eps\gamma\ge\cpv{G6_epsgamma}$ & Theorem~\ref{cp:G6} & \texttt{r211b\_\allowbreak{}cert\_\allowbreak{}out.\allowbreak{}txt:41-84} & F & --\\
rows 1--12 of Table~\ref{cp:tab:G10} (products, point $\Pprod$) & Theorem~\ref{cp:G10} & \texttt{r211b\_\allowbreak{}curve.\allowbreak{}py}, \texttt{\_\allowbreak{}out.\allowbreak{}txt:1-2, 23-43} & F & --\\
rows 14--17 of Table~\ref{cp:tab:G10}, among them $\cpv{excess_sharp_P0}$ with $\gamma<\cpv{excess_sharp_gamma}$ & Theorem~\ref{cp:G10}, Remark~\ref{cp:A3} & \texttt{r211b\_\allowbreak{}curve\_\allowbreak{}out.\allowbreak{}txt:45-49} & F & --\\
the two crude rows at $\Ppr$ & Section~\ref{cp:G13}, after Remark~\ref{cp:G12} & \texttt{r211b\_\allowbreak{}cert\_\allowbreak{}out.\allowbreak{}txt:125-140} & F & \texttt{r117\_\allowbreak{}certify.\allowbreak{}py} (F), the same rows\\
$M_t\ge\Lamp(t)$ at the $\cpv{B2_grid}$ tilts of Table~\ref{cp:tab:B2} & Lemma~\ref{cp:B2} & \texttt{r79\_\allowbreak{}certify.\allowbreak{}py} ($t\le\cpv{u8_t_cap}$, \texttt{\_\allowbreak{}out.\allowbreak{}txt:6-30}) and \texttt{r79b\_\allowbreak{}certify.\allowbreak{}py} ($t\ge\cpv{u8_t_a}$, \texttt{\_\allowbreak{}out.\allowbreak{}txt:5-83}); exact bounds in \texttt{r79\_\allowbreak{}certify.\allowbreak{}json}, \texttt{r79b\_\allowbreak{}certify.\allowbreak{}json} & M & \texttt{r211c\_\allowbreak{}cert\_\allowbreak{}out.\allowbreak{}txt:10-24} (D), at 13 of the tilts\\
the values of Table~\ref{cp:tab:B3} & Theorem~\ref{cp:B3} & \texttt{r79\_\allowbreak{}curve.\allowbreak{}py} ($\kappa\le1.1$, \texttt{\_\allowbreak{}out.\allowbreak{}txt:12-40}) and \texttt{r79b\_\allowbreak{}curve.\allowbreak{}py} ($\kappa\in[1.1,1.8]$, \texttt{\_\allowbreak{}out.\allowbreak{}txt:3-144}) & M & \texttt{r211c\_\allowbreak{}cert\_\allowbreak{}out.\allowbreak{}txt:26-45} (D), at $\kappa=\frac12$, $1$, $1.02$, $1.05$ and at $1.1$ to $1.6$, with its own tilts; not at $\kappa=\cpv{u8_k_a},\cpv{u8_k_b},\cpv{u8_k_d}$\\
$\fcert(\cpv{u8_plateau_end})\ge\cpv{u8_plateau_f}$ & Corollary~\ref{cp:B7} & \texttt{r79\_\allowbreak{}curve\_\allowbreak{}out.\allowbreak{}txt:43-44} & M & --\\
$\gamma<\cpv{u8_gD_one}$ at $\kappa=1$; excesses $\cpv{u8_exc_half}$, $\cpv{u8_exc_one}$ & Corollary~\ref{cp:B8} & \texttt{r79\_\allowbreak{}curve\_\allowbreak{}out.\allowbreak{}txt:46-53} & M & excesses also from Table~\ref{cp:tab:B3} and Lemma~\ref{cp:L11}\\
the points $\Ptt$, $\Pthree$ and their intervals; Table~\ref{cp:tab:D4}; $\cpv{per_entry_P23}$, $\cpv{online_setdisj_R77}$, $\cpv{hinted_omv_R32}$ & Lemma~\ref{cp:D:points}, Theorems~\ref{cp:D1}, \ref{cp:D3}, \ref{cp:D5}, Proposition~\ref{cp:D4} & \texttt{r211e\_\allowbreak{}cert.\allowbreak{}py}, \texttt{\_\allowbreak{}out.\allowbreak{}txt:10-33}; \texttt{r211e\_\allowbreak{}cert2.\allowbreak{}py}, \texttt{\_\allowbreak{}out.\allowbreak{}txt:4-14} & F & --\\
$\cpv{R11_range}$ and, conditionally, $\cpv{D10_pre}$ & Remark~\ref{cp:D10} & \texttt{r211e\_\allowbreak{}cert\_\allowbreak{}out.\allowbreak{}txt:6-8} & F & --\\
$\cpv{D13_c}$ & Corollary~\ref{cp:D13} & \texttt{r215\_\allowbreak{}bcd\_\allowbreak{}checks.\allowbreak{}py}, \texttt{\_\allowbreak{}out.\allowbreak{}txt:31} & F & --\\
$\cpv{A11_vxxz_half}$, $\cpv{A11_tab1_half}$, $\cpv{A11_tab1_one}$, every $\eps<\AlphaVXXZ$ at $\kappa=1$, $\bar C\ge\cpv{A10_C}$ & Corollary~\ref{cp:A11}, Proposition~\ref{cp:E12} & \texttt{r215\_\allowbreak{}a\_\allowbreak{}certify.\allowbreak{}py}, \texttt{\_\allowbreak{}out.\allowbreak{}txt:1-72}, on the grid of \texttt{r79b\_\allowbreak{}certify.\allowbreak{}json} & F & --\\
$\cpv{C16_range}$ ($\theta^*$) & Section~\ref{cp:C16} (instance of Proposition~\ref{cp:C14}) & \texttt{r211d\_\allowbreak{}checks.\allowbreak{}py}, \texttt{\_\allowbreak{}out.\allowbreak{}txt:6} & F & \texttt{r217d\_\allowbreak{}checks\_\allowbreak{}out.\allowbreak{}txt:19-21}\\
$\cpv{C19_copy_lo}$, $\cpv{C21_cg}$, $\cpv{C13_cg}<\cpv{C13_half_alpha}$, $2^9\cdot7^7>3^{18}$, $\cpv{C19_window}$ & Lemma~\ref{cp:C19}, Remark~\ref{cp:C21}, Example~\ref{cp:C13} & \texttt{r217d\_\allowbreak{}checks.\allowbreak{}py}, \texttt{\_\allowbreak{}out.\allowbreak{}txt:1-37} & F, D & --\\
the seven points of Table~\ref{cp:E7:table} and their slacks & Lemma~\ref{cp:E7}, Theorem~\ref{cp:E8} & \texttt{r211f\_\allowbreak{}certify.\allowbreak{}py}, \texttt{\_\allowbreak{}out.\allowbreak{}txt:4-26} & F & --\\
the entries of Table~\ref{cp:E15:table} and the two extremes & Remark~\ref{cp:E15} & \texttt{r217b\_\allowbreak{}e15\_\allowbreak{}table.\allowbreak{}py}, \texttt{\_\allowbreak{}out.\allowbreak{}txt:1-14} & F & \texttt{r211f\_\allowbreak{}certify\_\allowbreak{}out.\allowbreak{}txt:41-48, 63-73}\\
$\cpv{E16_r_out}$; the triples of Proposition~\ref{cp:E16} & Proposition~\ref{cp:E16} & \texttt{r217b\_\allowbreak{}e15\_\allowbreak{}table\_\allowbreak{}out.\allowbreak{}txt:15}; \texttt{r215\_\allowbreak{}bcd\_\allowbreak{}checks\_\allowbreak{}out.\allowbreak{}txt:1-26} & F & \texttt{r211f\_\allowbreak{}certify\_\allowbreak{}out.\allowbreak{}txt:50-61}\\
$\cpv{E18_g}$ & Remark~\ref{cp:E18} & \texttt{r215\_\allowbreak{}bcd\_\allowbreak{}checks\_\allowbreak{}out.\allowbreak{}txt:28} & F & --\\
$C_D=\cpv{F17_CD}\ldots$ at $\Pzero$ & Proposition~\ref{cp:F17} & \texttt{r220e\_\allowbreak{}checks.\allowbreak{}py}, \texttt{\_\allowbreak{}out.\allowbreak{}txt:70-71} & F & --\\
\end{longtable}}

\paragraph{Computed numbers.} The numbers that are only computed are called so in the text: the supremum $\cpv{alphaC_computed}$ of $\alpha(P)$ (Remark~\ref{cp:A4}), the values of Remark~\ref{cp:G11}, the rates of Remark~\ref{cp:E9}, the typical support exponents $\cpv{C12_cg_q7}$ and $\cpv{C13_cg}$ of Section~\ref{cp:sec:transfer} (the inequality $\cpv{C13_cg}<\cpv{C13_half_alpha}$ is certified), the ratio $\cpv{R27_r}$ and the ceiling $\cpv{R27_ceiling}$ of Remark~\ref{cp:A5}, the level-1 value about $\cpv{G12_cw_epsgamma}$ of $\eps\gamma$ in Remark~\ref{cp:G12}, the strip $\kappa\in(1,\cpv{A10_strip})$ of Remark~\ref{cp:A11:rem}, and the numbers of Section~\ref{cp:sec:profile} that are labelled computed there.

\paragraph{Which scripts need \texttt{mpmath}.} The scripts of kind M, \texttt{r74a\_\allowbreak{}certify.\allowbreak{}py}, \texttt{r79\_\allowbreak{}certify.\allowbreak{}py}, \texttt{r79b\_\allowbreak{}certify.\allowbreak{}py}, \texttt{r79\_\allowbreak{}curve.\allowbreak{}py} and \texttt{r79b\_\allowbreak{}curve.\allowbreak{}py}, use the interval arithmetic of \texttt{mpmath} (version 1.3.0) for their certified values. \texttt{r211b\_\allowbreak{}lib.\allowbreak{}py}, and with it \texttt{r211b\_\allowbreak{}cert.\allowbreak{}py}, \texttt{r211b\_\allowbreak{}curve.\allowbreak{}py} and \texttt{r211b\_\allowbreak{}cross.\allowbreak{}py}, imports \texttt{mpmath} for its floating-point search of the tangent points and for self-tests against \texttt{mpmath}; its certified values are of kind F. The other scripts need no \texttt{mpmath}; \texttt{r117\_\allowbreak{}certify.\allowbreak{}py}, \texttt{r211c\_\allowbreak{}cert.\allowbreak{}py}, \texttt{r211e\_\allowbreak{}cert.\allowbreak{}py}, \texttt{r211e\_\allowbreak{}cert2.\allowbreak{}py} and \texttt{r109\_\allowbreak{}dual.\allowbreak{}py} use \texttt{numpy} or \texttt{scipy} for floating-point searches only. The certificate of Corollary~\ref{cp:A11} reads the grid of exact upper bounds that \texttt{r79b\_\allowbreak{}certify.\allowbreak{}py} wrote (\texttt{r79b\_\allowbreak{}certify.\allowbreak{}json}); so it rests on the M-certificate of that grid, which \texttt{r211c\_\allowbreak{}cert.\allowbreak{}py} confirms in arithmetic D at 13 of its tilts.
\paragraph{Ancillary files.} The scripts of Table~\ref{cp:tab:cert}, the modules they import and their data files are the ancillary files of this submission (folder \texttt{anc/}), each script with its recorded output \texttt{NAME\_\allowbreak{}out.\allowbreak{}txt}, whose lines the table cites; \texttt{anc/README.\allowbreak{}txt} gives the command lines, the order in which the scripts that read another script's output must run, and the running times. All of them were rerun on 19 November 2026 from a clean copy of the folder (Python 3.13, \texttt{numpy} 2.3, \texttt{scipy} 1.16, \texttt{mpmath} 1.3.0), and every recorded output was reproduced.

\appendix
\section{The structure of the degeneration of~\texorpdfstring{\cite{VXXZ24}}{VXXZ24}}\label{cp:app:vxxz}

This appendix lists, for Lemma~\ref{cp:C18}, every step of Algorithm~1 of~\cite{VXXZ24} and of the proofs it invokes, with the operation it performs on variables. Numbers of statements and pages refer to arXiv:2307.07970v2. The definitions used are those of~\cite[Sections~3.1 and~3.3]{VXXZ24}: their degeneration allows a parameter $\lambda$; a restriction is given by linear maps, and a zero-out sets some variables to zero; an isomorphism ``$\equiv$'' is a relabeling of the variables by permutations of the index sets; a direct sum is a sum after relabeling to disjoint supports. In the table, ``variable map'' means that every variable goes to $0$ or to one variable with coefficient~$1$, and ``letters kept'' means that the image of a variable is indexed by a permutation of the index string of the source; the second property is needed only for the global stage (rows~1--13 and the hole fixing inside it, rows~20--22), because the window of Lemma~\ref{cp:C19}(a) compares a variable of the input with its image after step~1 only.

{\small
\setlength{\tabcolsep}{3pt}
\begin{center}
\begin{tabular}{@{}r >{\raggedright\arraybackslash}p{0.36\textwidth} >{\raggedright\arraybackslash}p{0.17\textwidth} >{\raggedright\arraybackslash}p{0.22\textwidth} >{\raggedright\arraybackslash}p{0.09\textwidth} >{\raggedright\arraybackslash}p{0.07\textwidth}@{}}
\multicolumn{6}{@{}l}{\emph{Global stage} (step~1 of Algorithm~1)}\\
\# & step & place in~\cite{VXXZ24} & operation & variable map & letters kept\\\hline
1 & one input copy for each complete split distribution $\xi$ in $\tilde{\mathcal D}$, $\abs{\tilde{\mathcal D}}$ polynomial in $n$ & Thm.~5.3, proof, p.~24 & choice of copies & yes (direct sum) & yes\\
2 & regions: $CW^{\otimes2^{\ell-1}n}\equiv\bigotimes_rCW^{\otimes A_rn}$ & Prop.~5.1, Sect.~5.1, p.~25 & renaming of positions & yes & yes\\
3 & zero out the blocks not consistent with the marginals & Sect.~5.2, p.~25 & zero-out & yes & yes\\
4 & zero out the blocks whose hash value is not in the target set; cleanup of shared $X$- and $Y$-blocks, repeated; drop the $X_I$, $Y_J$ whose triple is not consistent with the distribution & Sect.~5.2, p.~26 & zero-out & yes & yes\\
5 & compatibility zero-out I & Sect.~5.3, p.~28 & zero-out & yes & yes\\
6 & zero out the $Z$-blocks compatible with several triples (holes) & Sect.~5.4, p.~29 & zero-out & yes & yes\\
7 & usefulness zero-out & Sect.~5.5, p.~29 & zero-out & yes & yes\\
8 & the hashed tensor on $X_IY_JZ_K$ is $\equiv$ a tensor power of its constituents (positions sorted by $(I_t,J_t,K_t)$) & Claim~5.11, p.~29 & renaming of positions & yes & yes\\
9 & the useful tensor is the direct sum of broken copies of the target & eq.~(8), p.~30 & equality, no map & -- & --\\
10 & keep the copies with at most $N/8$ holes & Sect.~5.6, p.~33 & zero-out of whole copies & yes & yes\\
11 & degenerate broken copies into unbroken ones by Cor.~4.2 & Sect.~5.6, p.~33 & hole fixing, rows~20--22 & yes & yes\\
12 & tensor product of the outputs of the three regions & Sect.~5.7, p.~33 & tensor product of maps & yes & yes\\
13 & equal numbers of copies for every $\xi$ (the minimum), then the sum over $\xi$: ``a direct sum of some tensors can be degenerated into the sum of these tensors'' & Thm.~5.3, proof, p.~25 & zero-out of surplus copies; identification $x^{(\xi)}\mapsto x$ & yes, coefficient~$1$ (the summands have disjoint supports: a level-1 triple has one split distribution per term) & yes\\
\end{tabular}
\end{center}
\begin{center}
\begin{tabular}{@{}r >{\raggedright\arraybackslash}p{0.36\textwidth} >{\raggedright\arraybackslash}p{0.17\textwidth} >{\raggedright\arraybackslash}p{0.22\textwidth} >{\raggedright\arraybackslash}p{0.09\textwidth} >{\raggedright\arraybackslash}p{0.07\textwidth}@{}}
\multicolumn{6}{@{}l}{\emph{Constituent stages and final steps} (steps~2--4 of Algorithm~1); hole fixing (rows~20--22)}\\
\# & step & place in~\cite{VXXZ24} & operation & variable map & letters kept\\\hline
14 & a constituent term with a zero index is $\equiv\mat1M1$ (or a rotation) & Alg.~1, step~2; Thm.~6.1, p.~34 & renaming (a bijection: every variable has a unique partner) & yes & not needed\\
15 & $\mat abc\otimes\mat def\equiv\mat{ad}{be}{cf}$; the split-off factors $\mat{a_\ell}{b_\ell}{c_\ell}$ are carried along: a direct sum of products of a matrix product with other tensors is that matrix product times their direct sum & Sect.~3.4, p.~14; Alg.~1, p.~50 & renaming; tensor product with the identity & yes & not needed\\
16 & for each $\xi$ an independent input copy; the $\eps$-interface tensor at $\xi$ is a subtensor of the $3\eps$-interface input; Prop.~6.2 on it; the sum over $\xi$ & Thm.~6.3, proof, p.~37 & choice of copies, zero-out, identification & yes, coefficient~$1$ & not needed\\
17 & zero out the level-1 blocks not $\eps$-consistent in each region; the result is ``isomorphic to'' the product over the regions & Prop.~6.2; Claim~6.4, p.~38 & zero-out, renaming of positions & yes & not needed\\
18 & zero out the level-$(\ell-1)$ blocks not consistent with the distribution; hashing and cleanup; compatibility I; unique triple (holes); usefulness & Sects.~6.2--6.5, pp.~38--41 & zero-outs & yes & not needed\\
19 & both kinds of holes fixed by Cor.~4.2 & Sect.~6.6, p.~44 & hole fixing, rows~20--22 & yes & not needed\\
20 & relabel the variables of the broken copy by permutations of the group of Property~7.1, which preserve the tensor & Thm.~7.2, proof, p.~48 & bijection of variables & yes & yes, for the group of row~22\\
21 & zero out all parts outside the chosen parts; the target as a sum of $8$ subtensors over disjoint part-triples, the other $7$ obtained recursively from other broken copies & Thm.~7.2, proof, p.~48 & zero-out; identification of the $8$ outputs (disjoint supports) & yes, coefficient~$1$ & yes\\
22 & the group: permutations of chunks within each term & Cor.~4.2, proof, p.~49 & permutation of positions (keeps the split distributions of the terms) & yes & yes\\
23 & the level-1 interface tensor is $\equiv\mat{a_1}{b_1}{c_1}$ (every level-1 constituent has a zero index) & Alg.~1, step~3; Thm.~6.1, p.~50 & renaming & yes & not needed\\
24 & $V$ copies of $\mat ABC\equiv\bigotimes_\ell\mat{a_\ell}{b_\ell}{c_\ell}$; zero-out to $\mat m{\lfloor m^\kappa\rfloor}m$, $m=\min(A,B^{1/\kappa},C)$, implicit in~(27) & Alg.~1, step~4, p.~50 & renaming; zero-out of index ranges & yes & not needed\\
\end{tabular}
\end{center}
}

The random choices of the construction (the hash weights, and the permutation of~\cite[Lemma~7.3]{VXXZ24}) are used only to show that a good choice exists~\cite[pp.~26 and~47]{VXXZ24}; each member of the family is one fixed restriction, which is all that the algorithms of Theorem~\ref{cp:C20} need, since they use one fixed member. The symbol $\lambda$ appears in~\cite{VXXZ24} outside Section~3.3 only as a name of entropy quantities and of Lagrange multipliers~\cite[for example pp.~22, 35 and~51]{VXXZ24}; no step uses a degeneration with a parameter. The roles of $X$, $Y$ and $Z$ are never exchanged: ``up to rotation of the three dimensions''~\cite[p.~21]{VXXZ24} says which dimension's blocks are shared in the hashing of each region, and the regions are combined by a tensor product with the roles kept (row~12), so the $z$-side of $CW_5$ becomes the $m\times m$ output side of $\mat m{m^\kappa}m$.

\section{Certificates: rational points and parameters}\label{cp:app:certificates}

This appendix collects the rational points at which the numbers of the paper are certified, with all their digits, and the parameters of the certificates that the text does not print. Section~\ref{cp:sec:certificates} names the scripts. A printed point must carry all its digits: the margins are small (for $\cpv{online_setdisj_R77}$ at $\Ptt$ about $4.5\cdot10^{-8}$), and a point rounded to a few digits can fail to be admissible.

\paragraph{Level-1 points of the mixture of $F_6$ and $F_7$.} In the coordinates $(\lambda_6,\,u_6(6+2),\,u_7(7+2))$ of Definition~\ref{cp:L3}:
\begin{center}\small
\renewcommand{\arraystretch}{1.15}
\begin{tabular}{@{}l l l l@{}}
point & coordinates & used for & printed in\\\hline
$\Pzero$ & $(\cpv{P0_lam},\,\cpv{P0_u6},\,\cpv{P0_u7})$ & \ThmD{} at level~1, $\gamma$ at $\Pzero$, batches, fast bottoms & Lemma~\ref{cp:L11}\\
$\Ptt$ & $(\cpv{P23_lam},\,\cpv{P23_u6},\,\cpv{P23_u7})$ & single entries, Online Set Disjointness & Lemma~\ref{cp:D:points}\\
$\Pthree$ & $(\cpv{P32_lam},\,\cpv{P32_u6},\,\cpv{P32_u7})$ & whole columns and rows & Lemma~\ref{cp:D:points}\\
$\Pprod$ & $(\cpv{u7_Pprod_lam},\,\cpv{u7_ubar},\,\cpv{u7_ubar})$ & rows 1--12 of Table~\ref{cp:tab:G10} & Table~\ref{cp:tab:G10}\\
$\Ppr$ & $(\cpv{u7_Ppr_lam},\,\cpv{u7_Ppr_u6},\,\cpv{u7_Ppr_u7})$ & the two crude rows & Section~\ref{cp:G13}\\\hline
\end{tabular}
\end{center}
The five points over other part sets are the rows of Table~\ref{cp:tab:G6}, with $u_q(q+2)=\cpv{u7_ubar}$ for every $q$; the seven single-$q$ points for unequal sides are the rows of Table~\ref{cp:E7:table}; and the single-$q$ point of Example~\ref{cp:C13} has $q=7$ and $u=w=\frac19$.

\paragraph{The tilts of the certified grid.} The grid $\Tcert$ of Lemma~\ref{cp:B2} has $\cpv{B2_grid}$ tilts: the $\cpv{u8_grid_lo}$ tilts $\cpv{B2_grid_lo_list}$; the $\cpv{u8_grid_mid}$ tilts $\cpv{u8_t_a},\cpv{u8_t_b},\dots,\cpv{u8_t_z}$ in steps of $0.01$ without $\cpv{u8_t_skip}$; and $t=1$, where $M_1$ is the certified upper bound on $\co(\Pzero)$. The exact upper bounds $M_t$ are binary interval endpoints written by the certificates as rationals; Table~\ref{cp:tab:B2} prints nine of them, rounded up.

\paragraph{The witnesses of Corollary~\ref{cp:A11}.} The certificate checks every cost line of the proof of Theorem~\ref{cp:A10} at $\eps=v$, the printed value, with $x=(1-v/e)/(v-a')$ and the tilt below:
\begin{center}\small
\renewcommand{\arraystretch}{1.15}
\begin{tabular}{@{}l l l l@{}}
instance & $(\kappa,e,g)$ & $v$ & tilt $t$\\\hline
(a) & $(\frac12,\AlphaVXXZ,0)$ & $\cpv{A11_v_a}=\cpv{A11_vxxz_half}$ & $\cpv{A11_t_half}$\\
(b) & $(\frac12,\cpv{A11_tab1_e_half},\cpv{A11_tab1_g_half})$ & $\cpv{A11_v_b}=\cpv{A11_tab1_half}$ & $\cpv{A11_t_tab_half}$\\
(c) & $(1,\cpv{A11_tab1_e_one},\cpv{A11_tab1_g_one})$ & $\cpv{A11_v_c}=\cpv{A11_tab1_one}$ & $\cpv{A11_t_one}$\\\hline
\end{tabular}
\end{center}
In each case $v+\cpv{cert_step}$ is certified at no tilt of the grid. At $\kappa=1$ with $g=0$ the certificate is the argument of Theorem~\ref{cp:A10} at one tilt with $a't>M'_t$; $\cpv{A11_grid_pos}$ of the $\cpv{A11_grid_n}$ tilts qualify, and $t=\cpv{A11_t_one}$ gives $\bar C\ge\cpv{A10_C}$.

\paragraph{Other parameters.} Remark~\ref{cp:A3}: $t=\cpv{excess_crude_t}$ and $x=\cpv{excess_crude_x}$ for the crude profile, $t=\cpv{G3_t}$ and $x=\cpv{excess_sharp_x}$ for the sharp one. Theorem~\ref{cp:G3}: $t=\cpv{G3_t}$; the tangent points of Lemma~\ref{cp:G2} at this tilt are rationals with 40 significant digits, listed in the output of \texttt{r74a\_\allowbreak{}certify.\allowbreak{}py}, and the dual potentials of the second certificate (Remark~\ref{cp:F23}) in the output of \texttt{r109\_\allowbreak{}dual.\allowbreak{}py}. Tables~\ref{cp:tab:G6} and~\ref{cp:tab:G10} print their tilts, tile ratios and shares. For Table~\ref{cp:tab:B3} the tilt, the bottom and the split $x$ of each value are listed in the outputs of \texttt{r79\_\allowbreak{}curve.\allowbreak{}py} and \texttt{r79b\_\allowbreak{}curve.\allowbreak{}py}; the split is a rational next to $\xmin$ or $x_{\max}$ chosen by the script, and the second certificate \texttt{r211c\_\allowbreak{}cert.\allowbreak{}py} uses its own tilts.

\begin{thebibliography}{VXXZ24}\setlength{\itemsep}{0pt}

\bibitem[ABFK24]{ABFK24} Amir Abboud, Karl Bringmann, Nick Fischer, and Marvin K\"unnemann. The time complexity of fully sparse matrix multiplication. In \emph{Proc.\ 35th ACM-SIAM Symposium on Discrete Algorithms (SODA 2024)}, pages 4670--4703. arXiv:2309.06317v1, whose numbering we cite.

\bibitem[AV26]{AV26} Josh Alman and Virginia Vassilevska Williams. Truly subquadratic 3SUM and truly subcubic APSP via triangles in sparse lopsided graphs. arXiv:2610.06783v1, 2026.

\bibitem[AvdBZ26]{AvdBZ26} Alimzhan Amanov, Maxim van den Berg, and Jeroen Zuiddam. Multiplying matrices via polynomials: exposition and barriers. arXiv:2610.30054, 2026.

\bibitem[Bin80]{Bin80} Dario Bini. Relations between exact and approximate bilinear algorithms. Applications. \emph{Calcolo}, 17(1):87--97, 1980.

\bibitem[vdBNS19]{BNS19} Jan van den Brand, Danupon Nanongkai, and Thatchaphol Saranurak. Dynamic matrix inverse: improved algorithms and matching conditional lower bounds. In \emph{Proc.\ 60th IEEE Symposium on Foundations of Computer Science (FOCS 2019)}, pages 456--480. arXiv:1905.05067.

\bibitem[vdBSZ24]{BSZ24} Jan van den Brand, Zhao Song, and Tianyi Zhou. Algorithm and hardness for dynamic attention maintenance in large language models. In \emph{Proc.\ 41st International Conference on Machine Learning (ICML 2024)}, PMLR 235, pages 49008--49028. arXiv:2304.02207.

\bibitem[Cop97]{Cop97} Don Coppersmith. Rectangular matrix multiplication revisited. \emph{Journal of Complexity}, 13(1):42--49, 1997.

\bibitem[CVZ18]{CVZ18} Matthias Christandl, P\'eter Vrana, and Jeroen Zuiddam. Universal points in the asymptotic spectrum of tensors. In \emph{Proc.\ 50th ACM Symposium on Theory of Computing (STOC 2018)}, pages 289--296; \emph{J.\ Amer.\ Math.\ Soc.}\ 36 (2023), 31--79. arXiv:1709.07851.

\bibitem[CW90]{CW90} Don Coppersmith and Shmuel Winograd. Matrix multiplication via arithmetic progressions. \emph{Journal of Symbolic Computation}, 9(3):251--280, 1990.

\bibitem[DMRU95]{DMRU95} Paul F. Dietz, Kurt Mehlhorn, Rajeev Raman, and Christian Uhrig. Lower bounds for set intersection queries. \emph{Algorithmica}, 14(2):154--168, 1995.

\bibitem[EZ26]{EZ26} Isaac Barouch Essayag and Aryeh Lev Zabokritskiy. Faster deterministic all-pairs shortest paths with real edge weights. arXiv:2611.17215v2, 2026.

\bibitem[GKLP17]{GKLP17} Isaac Goldstein, Tsvi Kopelowitz, Moshe Lewenstein, and Ely Porat. Conditional lower bounds for space/time tradeoffs. In \emph{Proc.\ 15th Algorithms and Data Structures Symposium (WADS 2017)}, LNCS 10389, pages 421--436. arXiv:1706.05847.

\bibitem[GLP19]{GLP19} Isaac Goldstein, Moshe Lewenstein, and Ely Porat. On the hardness of set disjointness and set intersection with bounded universe. In \emph{Proc.\ 30th International Symposium on Algorithms and Computation (ISAAC 2019)}, LIPIcs 149, pages 7:1--7:22. arXiv:1910.00831.

\bibitem[HLS24]{HLS24} Bernhard Haeupler, Yaowei Long, and Thatchaphol Saranurak. Dynamic deterministic constant-approximate distance oracles with $n^{\epsilon}$ worst-case update time. In \emph{Proc.\ 65th IEEE Symposium on Foundations of Computer Science (FOCS 2024)}, pages 2033--2044. arXiv:2402.18541.

\bibitem[Hoe63]{Hoe63} Wassily Hoeffding. Probability inequalities for sums of bounded random variables. \emph{Journal of the American Statistical Association}, 58(301):13--30, 1963.

\bibitem[KL26]{KL26} Sankeerth Rao Karingula and Shachar Lovett. A simplified proof of polynomial improvements for 3SUM and APSP. ECCC TR26-247, 2026.

\bibitem[LG12]{LG12} Fran\c{c}ois Le Gall. Faster algorithms for rectangular matrix multiplication. In \emph{Proc.\ 53rd IEEE Symposium on Foundations of Computer Science (FOCS 2012)}, pages 514--523. arXiv:1204.1111v2, whose numbering we cite.

\bibitem[LGU18]{LGU18} Fran\c{c}ois Le Gall and Florent Urrutia. Improved rectangular matrix multiplication using powers of the Coppersmith--Winograd tensor. In \emph{Proc.\ 29th ACM-SIAM Symposium on Discrete Algorithms (SODA 2018)}, pages 1029--1046. arXiv:1708.05622.

\bibitem[LR83]{LR83} Grazia Lotti and Francesco Romani. On the asymptotic complexity of rectangular matrix multiplication. \emph{Theoretical Computer Science}, 23:171--185, 1983.

\bibitem[Maj26b]{cop} Ruhollah Majdoddin. Thicker rectangular matrix multiplication with polylogarithmic overhead over fields of characteristic 2 and 3. 2026. \url{https://github.com/Majdoddin/thin-matrix-products-polylog}.

\bibitem[Maj26a]{v2} Ruhollah Majdoddin. Thin matrix products and lossless reductions from 3SUM, Exact Triangle and APSP, with a dichotomy for coordinatewise problems. arXiv:2610.17420, 2026 (version~1, 12 October 2026); version~2, November 2026, \url{https://github.com/Majdoddin/formal-math}. We cite version~2 by its statement numbers.%

\bibitem[Sch81]{Sch81} Arnold Sch\"onhage. Partial and total matrix multiplication. \emph{SIAM Journal on Computing}, 10(3):434--455, 1981.

\bibitem[Shr26]{Shr26} shreshth-s. Improving the displayed exponents in arXiv:2610.06783v1. Public working note, \url{https://github.com/shreshth-s/thin-matrix-results}, version of 24 October 2026 (commit 537b461).

\bibitem[Str73]{Str73} Volker Strassen. Vermeidung von Divisionen. \emph{Journal f\"ur die reine und angewandte Mathematik}, 264:184--202, 1973.

\bibitem[Tao16]{Tao16} Terence Tao. A symmetric formulation of the Croot--Lev--Pach--Ellenberg--Gijswijt capset bound. Blog post, 18 May 2016, \href{https://terrytao.wordpress.com/2016/05/18/a-symmetric-formulation-of-the-croot-lev-pach-ellenberg-gijswijt-capset-bound/}{\nolinkurl{terrytao.wordpress.com/2016/05/18}}.

\bibitem[Uzn26]{Uzn26} Przemys\l{}aw Uzna\'nski. Rectangular matrix multiplication from polynomial products. arXiv:2610.10494, 2026.

\bibitem[VXXZ24]{VXXZ24} Virginia Vassilevska Williams, Yinzhan Xu, Zixuan Xu, and Renfei Zhou. New bounds for matrix multiplication: from alpha to omega. In \emph{Proc.\ 35th ACM-SIAM Symposium on Discrete Algorithms (SODA 2024)}, pages 3792--3835. arXiv:2307.07970, whose numbering we cite.

\end{thebibliography}

\end{document}
